%% For double-blind review submission, w/o CCS and ACM Reference (max submission space)
%\documentclass[sigplan,10pt,review,anonymous]{acmart}
%\settopmatter{printfolios=false,printccs=false,printacmref=false}
%% For double-blind review submission, w/ CCS and ACM Reference
%\documentclass[sigplan,review,anonymous]{acmart}\settopmatter{printfolios=true}
%% For single-blind review submission, w/o CCS and ACM Reference (max submission space)
%\documentclass[sigplan,review]{acmart}\settopmatter{printfolios=true,printccs=false,printacmref=false}
%% For single-blind review submission, w/ CCS and ACM Reference
%\documentclass[sigplan,review]{acmart}\settopmatter{printfolios=true}
%% For final camera-ready submission, w/ required CCS and ACM Reference
%\documentclass[sigplan,nonacm]{acmart}
\documentclass[sigplan,review,acmsmall,nonacm]{acmart}\settopmatter{printfolios=false,printccs=false,printacmref=false}


%% Conference information
%% Supplied to authors by publisher for camera-ready submission;
%% use defaults for review submission.
%\acmConference[SPLASH'24]{ACM SIGPLAN conference on Systems, Programming, Languages, and Applications: Software for Humanity}{October 22-27, 2024}{Pasadena, California, United States}
%\acmConference{}{}{}
%\acmYear{2018}
%\acmISBN{} % \acmISBN{978-x-xxxx-xxxx-x/YY/MM}
%\acmDOI{} % \acmDOI{10.1145/nnnnnnn.nnnnnnn}
%\startPage{1}

%% Copyright information
%% Supplied to authors (based on authors' rights management selection;
%% see authors.acm.org) by publisher for camera-ready submission;
%% use 'none' for review submission.
\setcopyright{none}
%\setcopyright{acmcopyright}
%\setcopyright{acmlicensed}
%\setcopyright{rightsretained}
%\copyrightyear{2018}           %% If different from \acmYear

%% Bibliography style
\bibliographystyle{acmart}

\usepackage{dsfont}
\usepackage{stmaryrd}
\usepackage{colortbl}
\usepackage{hyperref}

\usepackage{amsmath}
\DeclareMathOperator*{\argmax}{argmax}
\DeclareMathOperator*{\argmin}{argmin}
\usepackage{amssymb}

\usepackage[dvipsnames, table]{xcolor}
\usepackage{textcomp}

% Packages
\usepackage[pdf]{graphviz}
\usepackage{mathrsfs}

\newcommand*\circled[1]{\tikz[baseline=-0.1cm]{
  \node[shape=circle,draw,inner sep=0.48pt] (char) {\fontsize{7}{12}\textsf{#1}};}}

\DeclareMathAlphabet{\mathcal}{OMS}{cmsy}{m}{n}
\usepackage{cancel}
\newcommand\ccancel[2][red]{\renewcommand\CancelColor{\color{#1}}\cancel{#2}}
\newcommand{\nDownarrow}{\ensuremath{\text{ }\cancel{\Downarrow}\text{ }}}
\usepackage{centernot}

\usepackage{pgfplots, pgfplotstable}
\pgfplotsset{compat=1.7}
\usepgfplotslibrary{fillbetween}
\usetikzlibrary{patterns}
\pgfmathdeclarefunction{gauss}{2}{\pgfmathparse{1/(#2*sqrt(2*pi))*exp(-((x-#1)^2)/(2*#2^2))}}
\pgfmathdeclarefunction{nil}{1}{\pgfmathparse{0.001}}

\usepackage{arydshln}
\usepackage{adjustbox}
\usepackage{enumerate}
\usepackage{enumitem}
\usepackage{tikz-cd}
\usetikzlibrary{calc}
\usepackage{amsfonts}
%\usepackage{prooftrees}
\usepackage{bussproofs}
\renewcommand{\sectionautorefname}{\S}
\renewcommand{\subsectionautorefname}{\S}
\usepackage{float}

\usepackage{tikz-3dplot}
\usetikzlibrary{3d}
\usetikzlibrary{calligraphy}
\newif\ifshowcellnumber
\showcellnumbertrue

\usepackage{algorithm}
\usepackage{algpseudocode}
\usepackage{algorithmicx}
\usepackage{sourcecodepro}
\usepackage{tikz-qtree}
\usepackage{amsthm}
\usepackage{bm}
\usetikzlibrary{bayesnet}
\usetikzlibrary{arrows}
\usepackage{subcaption}
\usetikzlibrary{backgrounds}
\usetikzlibrary{tikzmark}
\usetikzlibrary{hobby}

\usepackage{mwe}

\newcommand{\E}{\mathbb{E}}
\newcommand{\Var}{\mathrm{Var}}
\newcommand{\Cov}{\mathrm{Cov}}

\newcommand{\CompOrder}{\mathcal{O}}
\def\graphspace{\mathbf{G}}
\def\Uniform{\mbox{\rm Uniform}}
\def\Gaussian{\mbox{\rm Gaussian}}
\def\Bernoulli{\mbox{\rm Bernoulli}}
\def\Dirichlet{\mbox{\rm Dirichlet}}

\usepackage{mathtools}% superior to amsmath
\usepackage{tikz}
% Packages
\usepackage{listings}
\DeclareRobustCommand{\hlred}[1]{{\sethlcolor{pink}\hl{#1}}}
\usepackage{fontspec}

\setmonofont[Scale=0.8]{JetBrainsMono}[
  Contextuals={Alternate},
  Path=./font/,
  Extension = .ttf,
  UprightFont=*-Regular,
  BoldFont=*-Bold,
  ItalicFont=*-Italic,
  BoldItalicFont=*-BoldItalic
]

\usepackage[skins,breakable,listings]{tcolorbox}

\lstdefinelanguage{kotlin}{
  comment=[l]{//},
  commentstyle={\color{gray}\ttfamily},
  emph={delegate, filter, firstOrNull, forEach, it, lazy, mapNotNull, println, repeat, assert, with, head, tail, len, return@},
  numberstyle=\noncopyable,
  identifierstyle=\color{black},
  keywords={abstract, actual, as, as?, break, by, class, companion, continue, data, do, dynamic, else, enum, expect, false, final, for, fun, get, if, import, in, infix, interface, internal, is, null, object, open, operator, override, package, private, public, return, sealed, set, super, suspend, this, throw, true, try, catch, typealias, val, var, vararg, when, where, while, tailrec, reified},
  keywordstyle={\bfseries},
  morecomment=[s]{/*}{*/},
  morestring=[b]",
  morestring=[s]{"""*}{*"""},
  ndkeywords={@Deprecated, @JvmField, @JvmName, @JvmOverloads, @JvmStatic, @JvmSynthetic, Array, Byte, Double, Float, Boolean, Int, Integer, Iterable, Long, Runnable, Short, String, int},
  ndkeywordstyle={\bfseries},
  sensitive=true,
  stringstyle={\ttfamily},
  literate={`}{{\char0}}1,
  escapeinside={(*@}{@*)}
}
\lstdefinelanguage{tidy}{
  comment=[l]{//},
  commentstyle={\color{gray}\ttfamily},
  emph={|, ->, ---},
  emphstyle={\color{red}},
  identifierstyle=\color{black},
  keywords={\|, ->, ---},
  otherkeywords={|,->},
  morekeywords={|,->},
  keywordstyle={\color{blue}\bfseries},
  morecomment=[s]{/*}{*/},
  morestring=[b]",
  morestring=[s]{"""*}{*"""},
  ndkeywords={@Deprecated, @JvmField, @JvmName, @JvmOverloads, @JvmStatic, @JvmSynthetic, Array, Byte, Double, Float, Int, Integer, Iterable, Long, Runnable, Short, String},
  ndkeywordstyle={\color{orange}\bfseries},
  sensitive=true,
  stringstyle={\color{green}\ttfamily},
  literate={`}{{\char0}}1
}

%%%%%%%%%%%%%%%%%%%%%%%%%%%%%%%%%%%%%%%%%%%
%
% Color boxes
%
%%%%%%%%%%%%%%%%%%%%%%%%%%%%%%%%%%%%%%%%%%%

\tcbset{
  enhanced jigsaw,
  breakable,
  listing only,
%  boxsep=-1pt,
%  top=-1pt,
  bottom=0.1cm,
  right=0.5cm,
  overlay first={
    \node[black!50] (S) at (frame.south) {\Large\ding{34}};
    \draw[dashed,black!50] (frame.south west) -- (S) -- (frame.south east);
  },
  overlay middle={
    \node[black!50] (S) at (frame.south) {\Large\ding{34}};
    \draw[dashed,black!50] (frame.south west) -- (S) -- (frame.south east);
    \node[black!50] (S) at (frame.north) {\Large\ding{34}};
    \draw[dashed,black!50] (frame.north west) -- (S) -- (frame.north east);
  },
  overlay last={
    \node[black!50] (S) at (frame.north) {\Large\ding{34}};
    \draw[dashed,black!50] (frame.north west) -- (S) -- (frame.north east);
  },
  before={\par\vspace{5pt}},
  after={\par\vspace{\parskip}\noindent}
}

\definecolor{slightgray}{rgb}{0.90, 0.90, 0.90}

\usepackage{soul}
\makeatletter
\def\SOUL@hlpreamble{%
  \setul{}{3.0ex}%
  \let\SOUL@stcolor\SOUL@hlcolor%
  \SOUL@stpreamble%
}
\makeatother

\newcommand{\inline}[1]{%
  \begingroup%
  \sethlcolor{slightgray}%
  \hl{\ttfamily\footnotesize #1}%
  \endgroup
}

\newcommand{\tinline}[1]{%
  \begingroup%
  \sethlcolor{slightgray}%
  \hl{\ttfamily\tiny #1}%
  \endgroup
}

\newtcblisting{halftidyinput}[1][]{%
  left skip=0.7cm,
  left=0.35cm,
  width=6cm,
%  left=-0.01cm,
  top=-0.1cm,
  bottom=-0.35cm,
  listing options={
    language=tidy,
    basicstyle=\ttfamily\small,
%numberstyle=\footnotesize,
    showstringspaces=false,
    tabsize=2,
    breaklines=true,
    numbers=none,
    inputencoding=utf8,
    escapeinside={(*@}{@*)},
    #1
  },
  underlay unbroken and first={%
    \path[draw=none] (interior.north west) rectangle node[white]{\includegraphics[width=4mm]{../figures/tidyparse_logo.png}} ([xshift=-10mm,yshift=-7mm]interior.north west);
  }
}

\newtcblisting{wholetidyinput}[1][]{%
  left skip=0.7cm,
  left=0.35cm,
  top=0.1cm,
  middle=0mm,
  boxsep=0mm,
  listing options={
    language=tidy,
    basicstyle=\ttfamily\small,
%numberstyle=\footnotesize,
    showstringspaces=false,
    tabsize=2,
    breaklines=true,
    numbers=none,
    inputencoding=utf8,
    escapeinside={(*@}{@*)},
    #1
  },
  underlay unbroken and first={%
      \path[draw=none] (interior.north west) rectangle node[white]{\includegraphics[width=4mm]{../figures/tidyparse_logo.png}} ([xshift=-10mm,yshift=-9mm]interior.north west);
  }
}

\definecolor{A}{RGB}{6,150,104}
\definecolor{B}{RGB}{196,74,137}
\definecolor{C}{RGB}{117,237,133}
\definecolor{D}{RGB}{246,46,243}
\definecolor{E}{RGB}{89,162,12}
\definecolor{F}{RGB}{113,12,158}
\definecolor{G}{RGB}{191,205,142}
\definecolor{H}{RGB}{51,58,158}
\definecolor{I}{RGB}{244,212,3}
\definecolor{J}{RGB}{37,36,249}
\definecolor{K}{RGB}{253,165,71}
\definecolor{L}{RGB}{27,81,29}
\colorlet{LA}{A!30}
\colorlet{LB}{B!30}
\colorlet{LC}{C!30}
\colorlet{LD}{D!30}
\colorlet{LE}{E!30}
\colorlet{LF}{F!30}
\colorlet{LG}{G!30}
\colorlet{LH}{H!30}
\colorlet{LI}{I!30}
\colorlet{LJ}{J!30}
\colorlet{LK}{K!30}
\colorlet{LL}{L!30}
\newcommand{\hiliA}[1]{%
  \colorbox{LA}{$#1$}}
\newcommand{\hiliB}[1]{%
  \colorbox{LB}{$#1$}}
\newcommand{\hiliC}[1]{%
  \colorbox{LC}{$#1$}}
\newcommand{\hiliD}[1]{%
  \colorbox{LD}{$#1$}}
\newcommand{\hiliE}[1]{%
  \colorbox{LE}{$#1$}}
\newcommand{\hiliF}[1]{%
  \colorbox{LF}{$#1$}}
\newcommand{\hiliG}[1]{%
  \colorbox{LG}{$#1$}}
\newcommand{\hiliH}[1]{%
  \colorbox{LH}{$#1$}}
\newcommand{\hiliI}[1]{%
  \colorbox{LI}{$#1$}}
\newcommand{\hiliJ}[1]{%
  \colorbox{LJ}{$#1$}}
\newcommand{\hiliK}[1]{%
  \colorbox{LK}{$#1$}}
\newcommand{\hiliL}[1]{%
  \colorbox{LL}{$#1$}}
\newcommand{\highlight}[1]{%
  \colorbox{lgray}{$#1$}}
\colorlet{lred}{red!30}
\colorlet{lorange}{orange!30}
\colorlet{lgreen}{green!30}
\colorlet{lgray}{black!15}
\colorlet{dgray}{black!75}
\DeclareRobustCommand{\hlred}[1]{{\sethlcolor{lred}\hl{#1}}}
\DeclareRobustCommand{\hlorange}[1]{{\sethlcolor{lorange}\hl{#1}}}
\DeclareRobustCommand{\hlgreen}[1]{{\sethlcolor{lgreen}\hl{#1}}}
\DeclareRobustCommand{\hlgray}[1]{{\sethlcolor{lgray}\hl{#1}}}
\DeclareRobustCommand{\caret}[1]{{\sethlcolor{dgray}\textcolor{white}{\hl{#1}}}}

\usepackage{url}
\usepackage{qtree}

\usepackage{filecontents}
\usepackage{pstricks-add}
\usepackage{emoji}
\usepackage{alltt}
\usepackage{nicematrix}
\usepackage{graphicx}
\usepackage{ulem}
\usepackage{upquote}
\tikzstyle{every picture}+=[remember picture]
\usepackage{menukeys}
\pgfplotstableread[col sep=comma,]{timings_loc.csv}\loctimings
\pgfplotstableread[col sep=comma,]{timings_unloc.csv}\unloctimings

\makeatletter
\DeclareRobustCommand{\cev}[1]{%
  {\mathpalette\do@cev{#1}}%
}
\newcommand{\do@cev}[2]{%
  \vbox{\offinterlineskip
  \sbox\z@{$\m@th#1 x$}%
  \ialign{##\cr
  \hidewidth\reflectbox{$\m@th#1\vec{}\mkern4mu$}\hidewidth\cr
  \noalign{\kern-\ht\z@}
    $\m@th#1#2$\cr
  }%
  }%
}
\makeatother

\makeatletter
\DeclareRobustCommand{\pder}[1]{%
  \@ifnextchar\bgroup{\@pder{#1}}{\@pder{}{#1}}}
\newcommand{\@pder}[2]{\frac{\partial#1}{\partial#2}}
\makeatother

\newcommand{\shup}{\shortuparrow}
\newcommand{\shri}{\shortrightarrow}
\newcommand{\shur}{\shup\hspace{-5pt}\shri}

\makeatletter
\def\squigglyred{\bgroup \markoverwith{\textcolor{red}{\lower3\p@\hbox{\sixly \char58}}}\ULon}
\makeatother

\makeatletter
\def\squigglyblu{\bgroup \markoverwith{\textcolor{blue}{\lower3\p@\hbox{\sixly \char58}}}\ULon}
\makeatother

\makeatletter
\def\squigglyora{\bgroup \markoverwith{\textcolor{orange}{\lower3\p@\hbox{\sixly \char58}}}\ULon}
\makeatother

\newcommand{\err}[1]{\smash{\squigglyred{#1}{}}}
\newcommand{\erb}[1]{\smash{\squigglyblu{#1}{}}}
\newcommand{\ero}[1]{\smash{\squigglyora{#1}{}}}
\newcommand{\stirlingii}{\genfrac{\{}{\}}{0pt}{}}

%======== Arrows =========
\newcommand{\knightarrow}{
  \tikz{
    \fill (0pt,0pt) circle [radius = 1pt];
    \fill (0pt,6pt) circle [radius = 1pt];
    \fill (6pt,0pt) circle [radius = 1pt];
    \fill (6pt,6pt) circle [radius = 1pt];
    \fill (12pt,0pt) circle [radius = 1pt];
    \fill (12pt,6pt) circle [radius = 1pt];
    \fill (6pt,0pt) circle [radius = 1pt];
    \fill (12pt,0pt) circle [radius = 1pt];
    \draw [-to] (0pt,0pt) -- (12pt,6pt);
  }
}

\newcommand{\kingarrow}{
  \tikz{
    \fill (0pt,0pt) circle [radius = 1pt];
    \fill (6pt,0pt) circle [radius = 1pt];
    \fill (0pt,6pt) circle [radius = 1pt];
    \fill (6pt,6pt) circle [radius = 1pt];
    \draw [-to] (0pt,0pt) -- (6pt,6pt);
    \draw [-to] (0pt,0pt) -- (0pt,6pt);
    \draw [-to] (0pt,0pt) -- (6pt,0pt);
  }
}

\newcommand{\duparrow}{
  \tikz{
    \fill (0pt,0pt) circle [radius = 1pt];
    \fill (0pt,6pt) circle [radius = 1pt];
    \draw [-to] (0pt,0pt) -- (0pt,6pt);
  }
}

\newcommand{\drightarrow}{
  \tikz{
    \fill (0pt,0pt) circle [radius = 1pt];
    \fill (6pt,0pt) circle [radius = 1pt];
    \draw [-to] (0pt,0pt) -- (6pt,0pt);
  }
}

\newcommand{\ddiagarrow}{
  \tikz{
    \fill (0pt,0pt) circle [radius = 1pt];
    \fill (6pt,0pt) circle [radius = 1pt];
    \fill (0pt,6pt) circle [radius = 1pt];
    \fill (6pt,6pt) circle [radius = 1pt];
    \draw [-to] (0pt,0pt) -- (6pt,6pt);
  }
}

\newcommand{\knightkingarrow}{
  \tikz{
    \fill (0pt,0pt) circle [radius = 1pt];
    \fill (0pt,6pt) circle [radius = 1pt];
    \fill (6pt,0pt) circle [radius = 1pt];
    \fill (6pt,6pt) circle [radius = 1pt];
    \fill (12pt,0pt) circle [radius = 1pt];
    \fill (12pt,6pt) circle [radius = 1pt];
    \draw [-to] (0pt,0pt) -- (6pt,6pt);
    \draw [-to] (0pt,0pt) -- (0pt,6pt);
    \draw [-to] (0pt,0pt) -- (6pt,0pt);
    \draw [-to] (0pt,0pt) -- (12pt,6pt);
  }
}

%======== Arrows =========

\usetikzlibrary{decorations.pathreplacing,automata,calc,positioning,matrix,fit,decorations.pathmorphing}

\usepackage{wrapfig}

\newcommand{\mkTrellis}[1]{
  \begin{tikzpicture}
    \def\dx{20pt}
    \def\dy{30pt}
    \newcounter{i}
    \stepcounter{i}
    \node[circle, draw, fill=black!30] (\arabic{i}) at (0,0){};
    \foreach [count=\i] \x in {2,...,#1}{
      \pgfmathsetmacro{\lox}{\x-1}%
      \pgfmathsetmacro{\loxt}{\x-3}%
      \foreach [count=\j] \xx in {-\lox,-\loxt,...,\lox}{
        \pgfmathsetmacro{\jj}{\j-1}%
        \stepcounter{i}
        \pgfmathsetmacro{\kk}{\xx-2}%
        \pgfmathsetmacro{\lbl}{\lox!/(\jj!*(\lox-\jj)!)}
        \ifnum\x<\kk
        \pgfmath\node[circle, draw]  (\arabic{i}) at (\xx*\dx, -\lox*\dy) {};
        \else
        \pgfmath\node[circle, draw, fill=black!30]  (\arabic{i}) at (\xx*\dx, -\lox*\dy) {};
        \fi
      }
    }
    \newcounter{z}
    \newcounter{xn}
    \newcounter{xnn}
    \pgfmathsetmacro{\maxx}{#1 - 1}
    \foreach \x in {1,...,\maxx}{
      \foreach \xx in {1,...,\x}{
        \stepcounter{z}
        \setcounter{xn}{\arabic{z}}
        \addtocounter{xn}{\x}
        \setcounter{xnn}{\arabic{xn}}
        \stepcounter{xnn}
        \draw [<-] (\arabic{z}) -- (\arabic{xn});
        \draw [<-] (\arabic{z}) -- (\arabic{xnn});
      }
    }
  \end{tikzpicture}
}

\newcommand{\dx}{20pt}
\newcommand{\dy}{30pt}
\newcounter{i}
\newcounter{z}
\newcounter{xn}
\newcounter{xnn}
\newcommand{\mkTrellisAppend}[1]{
  \begin{tikzpicture}
    \setcounter{i}{0}
    \setcounter{z}{0}
    \setcounter{xn}{0}
    \setcounter{xnn}{0}
    \stepcounter{i}
    \node[circle, draw] (\arabic{i}) at (0,0){};
    \foreach [count=\i] \x in {2,...,#1}{
      \pgfmathsetmacro{\lox}{\x-1}%
      \pgfmathsetmacro{\loxt}{\x-3}%
      \foreach [count=\j] \xx in {-\lox,-\loxt,...,\lox}{
        \pgfmathsetmacro{\jj}{\j-1}%
        \stepcounter{i}
        \pgfmathsetmacro{\kk}{\xx+2}%
        \pgfmathsetmacro{\lbl}{\lox!/(\jj!*(\lox-\jj)!)}
        \ifnum\x>\kk
        \pgfmath\node[circle, draw, fill=black!30]  (\arabic{i}) at (\xx*\dx, -\lox*\dy) {};
        \else
        \pgfmath\node[circle, draw]  (\arabic{i}) at (\xx*\dx, -\lox*\dy) {};
        \fi
      }
    }
    \pgfmathsetmacro{\maxx}{#1 - 1}
    \foreach \x in {1,...,\maxx}{
      \foreach \xx in {1,...,\x}{
        \stepcounter{z}
        \setcounter{xn}{\arabic{z}}
        \addtocounter{xn}{\x}
        \setcounter{xnn}{\arabic{xn}}
        \stepcounter{xnn}
        \draw [<-] (\arabic{z}) -- (\arabic{xn});
        \draw [<-] (\arabic{z}) -- (\arabic{xnn});
      }
    }
  \end{tikzpicture}
}

\newcommand{\mkTrellisInsert}[1]{
  \begin{tikzpicture}
    \setcounter{i}{0}
    \setcounter{z}{0}
    \setcounter{xn}{0}
    \setcounter{xnn}{0}
    \stepcounter{i}
    \node[circle, draw] (\arabic{i}) at (0,0){};
    \foreach [count=\i] \x in {2,...,#1}{
      \pgfmathsetmacro{\lox}{\x-1}%
      \pgfmathsetmacro{\loxt}{\x-3}%
      \foreach [count=\j] \xx in {-\lox,-\loxt,...,\lox}{
        \pgfmathsetmacro{\jj}{\j-1}%
        \stepcounter{i}
        \pgfmathsetmacro{\mp}{\xx+#1}%
        \pgfmathsetmacro{\mq}{\xx+\x}%
        \pgfmathsetmacro{\lbl}{\lox!/(\jj!*(\lox-\jj)!)}
        \ifnum\x>\mp
        \pgfmath\node[circle, draw, fill=black!30]  (\arabic{i}) at (\xx*\dx, -\lox*\dy) {};
        \else
        \ifnum#1<\mq
        \pgfmath\node[circle, draw, fill=black!30]  (\arabic{i}) at (\xx*\dx, -\lox*\dy) {};
        \else
        \pgfmath\node[circle, draw]  (\arabic{i}) at (\xx*\dx, -\lox*\dy) {};
        \fi
        \fi

      }
    }
    \pgfmathsetmacro{\maxx}{#1 - 1}
    \foreach \x in {1,...,\maxx}{
      \foreach \xx in {1,...,\x}{
        \stepcounter{z}
        \setcounter{xn}{\arabic{z}}
        \addtocounter{xn}{\x}
        \setcounter{xnn}{\arabic{xn}}
        \stepcounter{xnn}
        \draw [<-] (\arabic{z}) -- (\arabic{xn});
        \draw [<-] (\arabic{z}) -- (\arabic{xnn});
      }
    }
  \end{tikzpicture}
}

\usetikzlibrary{automata, positioning, arrows}

\newcommand{\nobarfrac}{\genfrac{}{}{0pt}{}}
\pgfplotstableread[col sep=comma,]{timings_loc.csv}\loctimings
\pgfplotstableread[col sep=comma,]{timings_unloc.csv}\unloctimings
\pgfplotstableread[col sep=comma,]{natural_errors.csv}\naturalerrors
\pgfplotstableread[col sep=comma,]{synthetic_errors.csv}\syntheticerrors

\usepackage[all,pdf]{xy}

\newcommand{\bs}{\blacksquare}
\newcommand{\ws}{\square}
\begin{document}
%
  \title{Syntax Repair as Language Intersection}
  %
  \begin{abstract}
    We introduce a new technique for correcting syntax errors in arbitrary context-free languages. Our work stems from the observation that syntax errors with a small repair typically have very few unique small repairs, which can usually be enumerated up to a small edit distance then quickly reranked. We place a heavy emphasis on precision: the enumerated set must contain every possible repair within a few edits and no invalid repairs. To do so, we model error correction as a language intersection problem between a Levenshtein automaton and a context-free grammar. To extract the repairs, we sample trees from the intersection grammar, yielding valid repairs within a certain Levenshtein distance. Finally, we rank those repairs by n-gram likelihood.
    \keywords{Error correction \and CFL reachability \and Langauge games.}
  \end{abstract}

%\titlerunning{Abbreviated paper title}
% If the paper title is too long for the running head, you can set
% an abbreviated paper title here
%
  \author{Breandan Considine}
  \email{bre@ndan.co}
%
%  \authorrunning{Considine et al.}
% First names are abbreviated in the running head.
% If there are more than two authors, 'et al.' is used.
%
%  \institute{McGill University, Montr\'eal, QC H2R 2Z4, Canada\\
%  \email{{breandan.considine@mail, jguo@cs}.mcgill.ca}\and
%  University of Toronto, Toronto, ON, M5S 1A1 Canada\\
%  \email{six@utoronto.ca}}

  \author{Jin Guo}
  \email{jguo@cs.mcgill.ca}
  \author{Xujie Si}
  \email{six@utoronto.ca}

  \maketitle

  \section{Introduction}

  Syntax errors are a familiar nuisance for programmers, arising due to a variety of factors, from inexperience, typographic error, to cognitive load. Often the mistake itself is simple to fix, but manual correction can disrupt concentration, a developer's most precious and fickle resource. Syntax repair attempts to automate the correction process by modifying a syntactically invalid program so that it conforms to the grammar, saving time and attention.

  Early work on syntax repair by Irons~\cite{irons1963error} and Aho~\cite{aho1972minimum} use techniques from dynamic programming to find the nearest parse trees for an erroneous input. These methods guarantee correctness, but do not attempt to completely recover all nearby corrections. Instead they find just one or a small number of corrections, which are not necessarily the most likely or natural repairs. Nevertheless, these methods are appealing for their interpretability and well-understood algorithmic properties.

  More recently, probabilistic repair techniques have been introduced using neural language models to predict the most likely correction~\cite{allamanis2021self, yasunaga2021break, sakkas2022seq2parse}. While these techniques generate far more natural edits, they are often costly to train, prone to misgeneralization, and difficult to incorporate new constraints thereafter. Furthermore, the generated repairs are not necessarily sound without additional filtering, and we observe the released models often hallucinate false positive repairs.

  Recent work by Merrill et al.~\cite{merrill2022saturated} and Chiang et al.~\cite{chiang2023tighter} suggest that the issue may be more foundational: transformer-based language models, a popular class of neural language models used in probabilistic program repair, are fundamentally less expressive than context-free grammars, which formally describe the syntax of most programming languages. This suggests such models, despite their useful approximation properties, are ill-suited for the task of end-to-end syntax repair. Yet, they may still be useful for resolving ambiguity between valid repairs of differing likelihood.

  In this work, we consider the problem of ranked syntax repair under finite Levenshtein bounds. We demonstrate it is possible to attain a significant advantage over state-of-the-art neural repair techniques by exhaustively retrieving every valid Levenshtein edit in a certain distance and scoring it. Not only does this approach guarantee both soundness and completeness, we find it also improves precision when ranking by naturalness. Our proposed solution is straightforward:

  \begin{enumerate}
    \item We model syntax repair as a language intersection problem between the Levenshtein ball and a context-free language, then materialize the grammar using a specialized version of the Bar-Hillel construction to Levenshtein intersections. (\S~\ref{sec:lev_bh})
    \item We construct a data structure via idempotent matrix completion that compactly represents parse forests in context-free languages. This data structure is used to index syntax trees, significantly reducing the size of the intersection grammar. (\S~\ref{sec:matrix_completion},~\ref{sec:ptree})
    \item To extract the repairs, primarily we sample trees without replacement by constructing an explicit bijection between syntax trees and integers, sample integers uniformly without replacement from a finite range, then decode them as trees. (\S~\ref{sec:ptree})
    \item Finally, we rerank all repairs found before a fixed timeout by n-gram likelihood. (\S~\ref{sec:ranking})
  \end{enumerate}

  Our primary technical contributions are threefold: (1) the adaptation of the Levenshtein automaton and Bar-Hillel construction to syntax repair (2) a theoretical connection between idempotent matrix completion and CFL parsing with holes, and (3) an algebraic datatype and integer bijection for enumerating or sampling valid sentences in context-free languages. The efficacy of our technique owes to the fact it does not synthesize probable edits, but unique, fully formed repairs within a certain edit distance. This enables us to suggest correct and natural repairs with far less compute and data than would otherwise be required by a large language model to attain the same precision.

  \section{Example}

  Syntax errors are usually fixable with a small number of edits. If we assume the intended repair contains just a few edits, this imposes strong locality constraints on the space of possible edits. For example, let us consider the following Python snippet, which contains a small syntax error:\\

  \texttt{def prepend(i, k, L=[]) n and [prepend(i - 1, k, [b] + L) for b in range(k)]}\\

  We can fix it by inserting a colon after the function definition, yielding:\\

  \texttt{def prepend(i, k, L=[])\hlgreen{:} n and [prepend(i - 1, k, [b] + L) for b in range(k)]}\\

  A careful observer will note that there is only one way to repair this Python snippet by making a single edit. In fact, many programming languages share this curious property: syntax errors with a small repair have few uniquely small repairs. Valid sentences corrupted by a few small errors rarely have many small corrections. We call such sentences \textit{metastable}, since they are relatively stable to small perturbations, as likely to be incurred by a careless typist or novice programmer.
%  Consider the following Kotlin snippet:\\
%
%  \texttt{fun main() = try \{ fetch() \} except(e: Exception) \{ handle(e) \}}\\
%
%  \noindent Again, there are thousands of possible single-token edits, only one of which is a valid repair:\\
%
%  \texttt{fun main() = try \{ fetch() \} \hlorange{catch}(e: Exception) \{ handle(e) \}}\\

  Let us consider a slightly more ambiguous error: \texttt{v = df.iloc(5:, 2:)}. Assuming an alphabet of just a hundred lexical tokens, this tiny statement has millions of possible two-token edits, yet only six of those possibilities are accepted by the Python parser:

%\setlength{\columnsep}{-10pt}
%\setlength{\columnseprule}{-10pt}
%\noindent\begin{multicols}{3}
%  \begin{enumerate}
%    \item\texttt{v = df.iloc(5\hlred{:}, 2\hlorange{,})}\\
%    \item\texttt{v = df.iloc(5\hlgreen{[}:, 2:\hlgreen{]})}\\
%    \item\texttt{v = df.iloc\hlorange{[}5:, 2:\hlorange{]}}\\
%    \item\texttt{v = df.iloc(5\hlorange{)}, 2\hlorange{(})}\\
%    \item\texttt{v = df.iloc(5\hlred{:}, 2\hlred{:})}\\
%    \item\texttt{v = df.iloc(5\hlgreen{[}:, 2\hlorange{]})}
%  \end{enumerate}
%\end{multicols}
  \begin{figure}[h!]
    \noindent\begin{tabular}{@{}l@{\hspace{10pt}}l@{\hspace{10pt}}l@{}}
    (1) \texttt{v = df.iloc(5\hlred{:}, 2\hlorange{,})} & (3) \texttt{v = df.iloc(5\hlgreen{[}:, 2:\hlgreen{]})} & (5) \texttt{v = df.iloc\hlorange{[}5:, 2:\hlorange{]}} \\\\
    (2) \texttt{v = df.iloc(5\hlorange{)}, 2\hlorange{(})} & (4) \texttt{v = df.iloc(5\hlred{:}, 2\hlred{:})} & (6) \texttt{v = df.iloc(5\hlgreen{[}:, 2\hlorange{]})} \\
    \end{tabular}
  \end{figure}

  With some typing information, we could easily narrow the results, but even in the absence of semantic constraints, one can probably rule out (2, 3, 6) given that \texttt{5[} and \texttt{2(} are rare bigrams in Python, and knowing \texttt{df.iloc} is often followed by \texttt{[}, determine (5) is most natural. This is the key insight behind our approach: we can usually locate the intended fix by exhaustively searching small repairs. As the set of small repairs is itself often small, if only we had some procedure to distinguish valid from invalid patches, the resulting solutions could be simply ranked by naturalness.

  The trouble is that any such procedure must be highly sample-efficient. We cannot afford to sample the universe of possible $d$-token edits, then reject invalid samples -- assuming it takes just 10ms to generate and check each sample, (1-6) could take 24+ hours to find. The hardness of brute-force search grows superpolynomially with edit distance, sentence length and alphabet size. We will need a more efficient procedure for sampling all and only small valid repairs.

  \clearpage\section{Problem statement}

  Source code in a programming language can be treated as a string over a finite alphabet, $\Sigma$. We use a lexical alphabet for convenience. The language has a syntax, $\ell \subset \Sigma^*$, containing every acceptable program. A syntax error is an unacceptable string, $\err\sigma \notin \ell$. We can model syntax repair as a language intersection between a context-free language (CFL) and a regular language. Henceforth, $\err\sigma$ will always and only be used to denote a syntactically invalid string whose target language is known.

  \begin{definition}[Bounded Levenshtein-CFL reachability]\label{def:bcflr}
    Given a CFL, $\ell$, and an invalid string, $\err{\sigma}: \ell^\complement$, find every valid string reachable within $d$ edits of $\err{\sigma}$, i.e., letting $\Delta$ be the Levenshtein metric and $L(\err\sigma, d) = \{\sigma' \mid \Delta(\err{\sigma}, \sigma') \leq d\}$ be the Levenshtein $d$-ball, we seek to find $A = L(\err\sigma, d) \cap \ell$.
  \end{definition}

%  To solve this problem, it is convenient to first consider intersections with a finite-length string with holes, then turn our attention back to BCFLR.
%

  As the admissible set $A$ is typically under-constrained, we want a procedure which surfaces natural and valid repairs over unnatural but valid repairs:

  \begin{definition}[Ranked repair]\label{def:ranked-repair}
    Given a finite language $A = L(\err\sigma, d) \cap \ell$ and a probabilistic language model $\text{P}_\theta: \Sigma^* \rightarrow [0, 1] \subset \mathbb{R}$, the ranked repair problem is to find the top-$k$ maximum likelihood repairs under the language model. That is,
    \begin{equation}
      R(A, P_\theta) = \argmax_{\bm{\sigma} \subseteq A, |\bm{\sigma}| \leq k} \sum_{\sigma \in \bm{\sigma}}\text{P}_\theta(\sigma)
    \end{equation}
    % On average, across all $G, \sigma$ $\hat{R}$ should approximate $R$.
%    We want a procedure $\hat{R}$, minimizing $\mathbb{E}_{G, \sigma}\big[D_{\text{KL}}(\hat{R} \parallel R)\big]$ and wallclock runtime.
  \end{definition}

  A popular approach to ranked repair involves learning a distribution over strings, however this is highly sample-inefficient and generalizes poorly to new languages. Approximating a distribution over $\Sigma^*$ forces the model to jointly learn syntax and stylometry. Furthermore, even with an extremely efficient approximate sampler for $\sigma \sim \ell_\cap$, due to the size of $\ell$ and $L(\err\sigma, d)$, it would be intractable to sample either $\ell$ or $L(\err\sigma, d)$, reject duplicates, then reject invalid ($\sigma \notin \ell$) or unreachable ($\sigma \notin L(\err\sigma, d)$) edits, and completely out of the question to sample $\sigma \sim \Sigma^*$ as do many neural language models.

  As we will demonstrate, the ranked repair problem can be factorized into a bilevel objective: first maximal retrieval, then ranking. Instead of working with strings, we will explicitly construct a grammar which soundly and completely generates the set $\ell \cap L(\err\sigma, d)$, then retrieve repairs from its language. By ensuring retrieval is sufficiently precise and exhaustive, maximizing likelihood over the retrieved set can be achieved with a much simpler, syntax-oblivious language model.

  Assuming we have a grammar that recognizes the Levenshtein-CFL intersection, the question then becomes how to maximize the number of unique valid sentences in a given number of samples. Top-down incremental sampling with replacement eventually converges to the language, but does so superlinearly~\cite{flajolet1992birthday}. Due to practical considerations including latency, we require the sampler to converge linearly, ensuring with much higher probability that natural repairs are retrieved in a timely manner. This motivates the need for a specialized generating function. More precisely,

  \begin{definition}[Linear convergence]\label{def:linear-convergence}
    Given a finite CFL, $\ell$, we want a randomized generating function, $\bm{\varphi}: \mathbb{N}_{<|\ell|} \rightarrow 2^\ell$, whose rate of convergence is linear in expectation, i.e., $\mathbb{E}_{i \in [1, n]}|\bm{\varphi}(i)| \propto n$.
  \end{definition}

  To satisfy Def.~\ref{def:linear-convergence}, we construct a bijection from syntax trees to integers (\S~\ref{sec:ptree}), sample integers uniformly without replacement, then decode them as trees. This will produce a set of unique trees, and each tree, assuming grammatical unambiguity, will correspond to a unique sentence in the language. As long as $|\ell_\cap|$ is sufficiently small and enough samples are drawn, $\bm\varphi$ is sure to include the most natural repairs, and additionally, will terminate after exhausting all sentences.

  Finally, once we have a set of small and valid repairs, the problem of ranked repair reduces to sorting retrieved samples by likelihood, which can be approximated using an autoregressive language model or any suitable scoring function of the implementer's choice. In our case, we use a low-order Markov model for its inference speed, data efficiency, and simplicity.

  \clearpage\section{Method}

  The method we describe in this paper takes as input the invalid code fragment, and returns a set of plausible repairs. We assume to know the target syntax and a low-rank distribution of lexical n-grams to estimate the likelihood of candidate repairs. At a high level, our method can be decomposed into three main steps: (1) language intersection, (2) k-best decoding, and (3) reranking.

  \begin{figure}[h!]
  \begin{center}
    \resizebox{0.8\textwidth}{!}{
      \begin{tikzpicture}[node distance=5cm]
        \node (start) [io] {Broken code};
        \node (node1) [plain, right of=start] {\phantom{...}\textbf{Language intersection}\phantom{...}};
        \node (gram1) [io2, above of=node1, yshift=-3cm] {Syntax};
%        \node (node2) [plain, right of=node1] {\textbf{Repair extraction}};
%        \node (ptree) [io, above of=node2, yshift=-3cm] {$\mathbb{T}_2$};
        \node (node3) [plain, right of=node1] {\textbf{k-best decoding}};
        \node (ngram) [io2, above of=node3, yshift=-3cm] {PCFG};
        \node (node4) [io, right of=node3] {Repairs};
        \draw [arrow] (start) -- (node1);
        \draw [arrow] (gram1) -- (node1);
        \draw [arrow] (node1) -- (node3);
        \draw [arrow] (node3) -- (node4);
        \draw [arrow] (ngram) -- (node3);
%        \draw [arrow] (ptree) -- (node2);
      \end{tikzpicture}
    }
  \end{center}
%  \caption{Line chart of our proposed method.}\label{fig:linechart}
  \end{figure}

First, we generate a synthetic grammar representing the intersection between the syntax and the Levenshtein ball around the source code. During extraction, we retrieve as many repairs as possible from the intersection grammar via sampling or enumeration. Finally, in the reranking step, we rank all repairs by n-gram likelihood. This can be depicted in more detail as a flowchart (Fig.~\ref{fig:flowchart}).

  \begin{wrapfigure}{r}{0.4\textwidth}
%\begin{figure}[h!]
    \vspace{-0.4cm}
    \begin{center}
      \resizebox{0.39\textwidth}{!}{
        \begin{tikzpicture}[node distance=2cm]
          \node (start) [startstop, draw=none];
          \node (pro1) [process, below of=start, yshift=-0.3cm] {$G_\cap \leftarrow G\cap\Delta(\err\sigma, d)$};
          \node [above=0.07cm of pro1] {(\S~\ref{sec:lev_bh})};
          \node (pcfg) [io2, left of=pro1, xshift=-3cm] {[PHO]CFG};
          \node [above=0.07cm of pcfg, xshift=1cm] {(\S~\ref{sec:prelim})};
          \node (lnfa) [io, right of=pro1, xshift=3cm] {L-NFA};
          \node [above=0.07cm of lnfa, xshift=1cm] {(\S~\ref{sec:lev_nfa})};

          \node (code) [io, right of=start,xshift=3cm] {Code};
          \node (synt) [io2, left of=start,xshift=-3cm] {Syntax};

          \node (dec1) [decision, below of=pro1, yshift=-0.5cm] {$[G_{\cap} = \varnothing]$};

          \node (pro2b) [process, right of=dec1, xshift=3cm] {Increase radius, $d$};

          \node [below=0.7cm of pro2b, xshift=0.05cm] {\Large\textbf{Language intersection}};
          \draw[thick,dotted, rounded corners] ($(pcfg.north west)+(-1.9,0.8)$) rectangle ($(pro2b.south east)+(0.3,-1.5)$);

          \node (const) [process, below of=dec1, yshift=-1.8cm] {Construct $\mathbb{T}_2$ from $G_\cap'$};
          \node [above=0.07cm of const, xshift=1.5cm] {(\S~\ref{sec:matrix_completion})};

%          \node (dec2) [decision, below of=const, yshift=-0.5cm] {$|\mathcal{L}(G_\cap)|$};
%
%          \node (samp1) [process, left of=dec2, xshift=-3cm] {Enumerate $\sigma' \in \mathcal{L}(G_\cap)$};
%          \node [above=0.07cm of samp1] {(\S~\ref{sec:ptree})};
%          \node (samp2) [process, right of=dec2, xshift=3cm] {Sample $\sigma' \sim P(G_\cap)$};
%          \node [above=0.07cm of samp2] {(\S~\ref{sec:ptree})};

%          \draw[thick,dotted, rounded corners] ($(const.north west)+(-5.3,0.7)$) rectangle ($(samp2.south east)+(0.3,-0.6)$);

          \node (rank) [process, below of=const, yshift=-0.5cm] {Decode top-$k$ by $L_\theta(\sigma')$};
%          \node (vlmc) [io2, right of=rank, xshift=3cm] {Markov chain};
          \node [below=0.01cm of rank, xshift=5.5cm] {\Large\textbf{Repair decoding}};
          \node [above=0.1cm of rank, xshift=1.5cm] {(\S~\ref{sec:ranking})};
          \draw[thick,dotted, rounded corners] ($(rank.north west)+(-5.3,3.2)$) rectangle ($(rank.south east)+(5.3,-0.8)$);

%  \node (out1) [io, below of=pro2a] {Output};
          \node (stop) [startstop, below of=rank, yshift=-0.6cm];

%  \draw [arrow] (dec0) -- node[anchor=east] {no} (pro1);

%          \draw [->,thick] (-5, 1.3) -- (synt);
%          \draw [->,thick] (5, 1.3) -- (code);

%          \draw [arrow] (start) -- (code);
%          \draw [arrow] (start) -- (synt);
          \draw [arrow] (code) -- (lnfa);
          \draw [arrow] (synt) -- (pcfg);
          \draw [arrow] (lnfa) -- (pro1);
          \draw [arrow] (pcfg) -- (pro1);

%  \draw [arrow] (in1) -- (pro1);
          \draw [arrow] (pro1) -- (dec1);
          \draw [arrow] (dec1) -- node[anchor=south] {yes} (pro2b);
          \draw [arrow] (dec1) -- node[anchor=east,yshift=0.5cm] {no} (const);
          \draw [arrow] (const) -- (rank);
          \draw [arrow] (pro2b) -- (lnfa);
%          \draw [arrow] (dec2) -- node[anchor=south] {small} (samp1);
%          \draw [arrow] (dec2) -- node[anchor=south] {large} (samp2);

%          \draw [arrow] (vlmc) -- (rank);
          \draw [arrow] (pcfg) |- ([shift={(-1.3cm,0)}]rank.west)--(rank.west);
%          \draw [arrow] (samp2) |- ([shift={(0,1.3cm)}]rank.north)--(rank.north);
%  \draw [arrow] (pro2a) -- (out1);
          \draw [arrow] (rank) -- (stop);
%          \draw [arrow] (dec2) -- node[anchor=east] {1} (stop);

        \end{tikzpicture}
      }
    \end{center}
    \vspace{-0.7cm}
    \caption{Dataflow of our proposed method.}\label{fig:flowchart}
    \vspace{-0.5cm}
  \end{wrapfigure}

  Since the syntax of most programming languages is context-free, we first construct a context-free grammar (CFG), $G_\cap$, that represents the intersection between the programming language syntax ($G$) and an automaton recognizing the Levenshtein edit ball of a given radius, $L(\err\sigma, d)$. As the CFL family is closed under intersection with regular languages, the intersection language $\mathcal{L}(G_\cap)$ should contain every repair within a given Levenshtein distance and no invalid repairs. Either the grammar will be empty, in which case there are no repairs within the given radius, or it will be nonempty, in which case we can directly proceed to decode repairs.

%  \begin{enumerate}
%    \item $G_\cap$ is empty, in which case there is no repair within the given radius. In this case, we simply increase the radius and try again.
%    \item $\mathcal{L}(G_\cap)$ is small, in which case we enumerate all possible repairs. Complete enumeration is tractable for the majority of all syntax repairs.
%    \item $\mathcal{L}(G_\cap)$ is too large to completely enumerate, so we sample instead from $G_\cap$, top-down. Sampling is necessary for a minority of remaining cases.
%  \end{enumerate}

  To decode the repairs, we use a technique called cube pruning, which has been shown to be effective for decoding likely sentences from PCFGs~\cite{huang2005better}. If the language is sufficiently small, it will generate every possible repair, otherwise, it will draw a representative subset containing with high probability the top-k most likely repairs under the PCFG.

  We will first describe how to generate the intersection grammar (\S~\ref{sec:lev_nfa},~\ref{sec:lev_bh}), then, describe a data structure compactly representing its language, allowing us to efficiently decode all repairs contained within (\S~\ref{sec:ptree}). Optionally, we can choose to rerank the repairs by a more sophisticated language model, such as a neural network, to improve the naturalness of the top-k repairs (\S~\ref{sec:ranking}).

  \subsection{Preliminaries}\label{sec:prelim}

  Recall that a CFG, $\mathcal{G} = \langle \Sigma, V, P, S\rangle$, is a quadruple consisting of terminals $(\Sigma)$, nonterminals $(V)$, productions $(P\colon V \rightarrow (V \mid \Sigma)^*)$, and a start symbol, $(S)$. Every CFG is reducible to so-called \textit{Chomsky Normal Form}, $P'\colon V \rightarrow (V^2 \mid \Sigma)$, where every production is either (1) a binary production $w \rightarrow xz$, or (2) a unit production $w \rightarrow t$, where $w, x, z: V$ and $t: \Sigma$. For example:\vspace{-3pt}

  \begin{table}[H]
    \begin{tabular}{llll}
      $G = \big\{\;S \rightarrow S\:S \mid (\:S\:) \mid (\:)\;\big\} \Longrightarrow G' = \big\{\;S\rightarrow Q\:R \mid S\:S \mid L\:R,$ & $R \rightarrow\:),$ & $L \rightarrow (,$ & $Q\rightarrow L\:S\;\big\}$
    \end{tabular}
  \end{table}\vspace{-8pt}

  Likewise, a finite state automaton is a quintuple $\mathcal{A} = \langle Q, \Sigma, \delta, I, F\rangle$, where $Q$ is a finite set of states, $\Sigma$ is a finite alphabet, $\delta \subseteq Q \times \Sigma \times Q$ is the transition function, and $I, F \subseteq Q$ are the set of initial and final states, respectively. We will adhere to this notation in the following sections.

  \clearpage\subsection{Modeling lexical edits with the nominal Levenshtein automaton}\label{sec:lev_nfa}

  \begin{wrapfigure}{r}{0.5\textwidth}
    \vspace{-0.3cm}
    \begin{center}
        \resizebox{\textwidth}{!}{
  \begin{tikzpicture}[
%->, % makes the edges directed
  >=stealth',
  node distance=2.5cm, % specifies the minimum distance between two nodes. Change if necessary.
%  every state/.style={thick, fill=gray!10}, % sets the properties for each ’state’ node
  initial text=$ $, % sets the text that appears on the start arrow
  ]
  \node[state, initial]                (00) {$q_{0,0}$};
  \node[state, right of=00]            (10) {$q_{1,0}$};
  \node[accepting, state, right of=10] (20) {$q_{2,0}$};
  \node[accepting, state, right of=20] (30) {$q_{3,0}$};
  \node[accepting, state, right of=30] (40) {$q_{4,0}$};
  \node[accepting, state, right of=40] (50) {$q_{5,0}$};

  \node[state, above of=00, shift={(-1.5cm,-0.0cm)}] (01) {$q_{0,1}$};
  \node[state, right of=01]                          (11) {$q_{1,1}$};
  \node[state, right of=11]                          (21) {$q_{2,1}$};
  \node[accepting, state, right of=21]               (31) {$q_{3,1}$};
  \node[accepting, state, right of=31]               (41) {$q_{4,1}$};
  \node[accepting, state, right of=41]               (51) {$q_{5,1}$};

\node[state, above of=01, shift={(-1.5cm,-0.0cm)}] (0j) {$q_{0,2}$};
\node[state, right of=0j]                          (1j) {$q_{1,2}$};
\node[state, right of=1j]                          (2j) {$q_{2,2}$};
\node[state, right of=2j]                          (3j) {$q_{3,2}$};
\node[accepting, state, right of=3j]               (4j) {$q_{4,2}$};
\node[accepting, state, right of=4j]               (5j) {$q_{5,2}$};

\node[state, above of=0j, shift={(-1.5cm,0.0cm)}] (0k) {$q_{0,3}$};
\node[state, right of=0k]                         (1k) {$q_{1,3}$};
\node[state, right of=1k]                         (2k) {$q_{2,3}$};
\node[state, right of=2k]                         (3k) {$q_{3,3}$};
\node[state, right of=3k]                         (4k) {$q_{4,3}$};
\node[accepting, state, right of=4k]              (5k) {$q_{5,3}$};

\draw [->] (00) edge[below] node{$\sigma_1$} (10);
\draw [->] (10) edge[below] node{$\sigma_2$} (20);
\draw [->] (20) edge[below] node{$\sigma_3$} (30);
\draw [->] (30) edge[below] node{$\sigma_4$} (40);
\draw [->] (40) edge[below] node{$\sigma_5$} (50);

\draw [->] (01) edge[below] node{$\sigma_1$} (11);
\draw [->] (11) edge[below] node{$\sigma_2$} (21);
\draw [->] (21) edge[below] node{$\sigma_3$} (31);
\draw [->] (31) edge[below] node{$\sigma_4$} (41);
\draw [->] (41) edge[below] node{$\sigma_5$} (51);

\draw [->] (0j) edge[below] node{$\sigma_1$} (1j);
\draw [->] (1j) edge[below] node{$\sigma_2$} (2j);
\draw [->] (2j) edge[below] node{$\sigma_3$} (3j);
\draw [->] (3j) edge[below] node{$\sigma_4$} (4j);
\draw [->] (4j) edge[below] node{$\sigma_5$} (5j);

\draw [->] (0k) edge[below] node{$\sigma_1$} (1k);
\draw [->] (1k) edge[below] node{$\sigma_2$} (2k);
\draw [->] (2k) edge[below] node{$\sigma_3$} (3k);
\draw [->] (3k) edge[below] node{$\sigma_4$} (4k);
\draw [->] (4k) edge[below] node{$\sigma_5$} (5k);

\draw [->] (00) edge[left] node{$\phantom{\cdot}$} (11);
\draw [->] (10) edge[left] node{$\phantom{\cdot}$} (21);
\draw [->] (20) edge[left] node{$\phantom{\cdot}$} (31);
\draw [->] (30) edge[left] node{$\phantom{\cdot}$} (41);
\draw [->] (40) edge[left] node{$\phantom{\cdot}$} (51);

% Super-knight arcs
\draw [->, red] (00) edge[bend left=3] node[east, shift={(-0.3cm,-0.65cm)}]{$\color{red}\sigma_3$}         (3j);
\draw [->, red] (10) edge[bend left=3] node[east, shift={(-0.3cm,-0.65cm)}]{$\color{red}\sigma_4$}         (4j);
\draw [->, red] (20) edge[bend left=3] node[east, shift={(-0.3cm,-0.65cm)}]{$\color{red}\sigma_5$}         (5j);
\draw [->, red] (01) edge[bend left=3] node[east, shift={(-0.3cm,-0.65cm)}]{$\color{red}\sigma_3$}         (3k);
\draw [->, red] (11) edge[bend left=3] node[east, shift={(-0.3cm,-0.65cm)}]{$\color{red}\sigma_4$}         (4k);
\draw [->, red] (21) edge[bend left=3] node[east, shift={(-0.3cm,-0.65cm)}]{$\color{red}\sigma_5$}         (5k);

\draw [->, violet] (00) edge[bend right=1] node[east, shift={(0cm,0.35cm)}]{$\color{violet}\sigma_4$}  (4k);
\draw [->, violet] (10) edge[bend right=1] node[east, shift={(0cm,0.35cm)}]{$\color{violet}\sigma_5$}  (5k);

\draw [->] (01) edge[left] node{$\phantom{\cdot}$} (1j);
\draw [->] (11) edge[left] node{$\phantom{\cdot}$} (2j);
\draw [->] (21) edge[left] node{$\phantom{\cdot}$} (3j);
\draw [->] (31) edge[left] node{$\phantom{\cdot}$} (4j);
\draw [->] (41) edge[left] node{$\phantom{\cdot}$} (5j);

\draw [->] (0j) edge[left] node{$\phantom{\cdot}$} (1k);
\draw [->] (1j) edge[left] node{$\phantom{\cdot}$} (2k);
\draw [->] (2j) edge[left] node{$\phantom{\cdot}$} (3k);
\draw [->] (3j) edge[left] node{$\phantom{\cdot}$} (4k);
\draw [->] (4j) edge[left] node{$\phantom{\cdot}$} (5k);

\draw [->] (00) edge[bend left=10, left] node{$\phantom{\cdot}$} (01);
\draw [->] (10) edge[bend left=10, left] node{$\phantom{\cdot}$} (11);
\draw [->] (20) edge[bend left=10, left] node{$\phantom{\cdot}$} (21);
\draw [->] (30) edge[bend left=10, left] node{$\phantom{\cdot}$} (31);
\draw [->] (40) edge[bend left=10, left] node{$\phantom{\cdot}$} (41);
\draw [->] (50) edge[bend left=10, left] node{$\phantom{\cdot}$} (51);

\draw [->] (01) edge[bend left=10, left] node{$\phantom{\cdot}$} (0j);
\draw [->] (11) edge[bend left=10, left] node{$\phantom{\cdot}$} (1j);
\draw [->] (21) edge[bend left=10, left] node{$\phantom{\cdot}$} (2j);
\draw [->] (31) edge[bend left=10, left] node{$\phantom{\cdot}$} (3j);
\draw [->] (41) edge[bend left=10, left] node{$\phantom{\cdot}$} (4j);
\draw [->] (51) edge[bend left=10, left] node{$\phantom{\cdot}$} (5j);

\draw [->] (0j) edge[bend left=10, left] node{$\phantom{\cdot}$} (0k);
\draw [->] (1j) edge[bend left=10, left] node{$\phantom{\cdot}$} (1k);
\draw [->] (2j) edge[bend left=10, left] node{$\phantom{\cdot}$} (2k);
\draw [->] (3j) edge[bend left=10, left] node{$\phantom{\cdot}$} (3k);
\draw [->] (4j) edge[bend left=10, left] node{$\phantom{\cdot}$} (4k);
\draw [->] (5j) edge[bend left=10, left] node{$\phantom{\cdot}$} (5k);

\draw [->, blue] (00) edge[bend left=2,below] node[shift={(-0.7cm,-0.5cm)}]{$\color{blue}\sigma_2$}    (21);
\draw [->, blue] (10) edge[bend left=2,below] node[shift={(-0.7cm,-0.5cm)}]{$\color{blue}\sigma_3$}    (31);
\draw [->, blue] (20) edge[bend left=2,below] node[shift={(-0.7cm,-0.5cm)}]{$\color{blue}\sigma_4$}    (41);
\draw [->, blue] (30) edge[bend left=2,below] node[shift={(-0.7cm,-0.5cm)}]{$\color{blue}\sigma_5$}    (51);

\draw [->, blue] (01) edge[bend left=2,below] node[shift={(-0.7cm,-0.5cm)}]{$\color{blue}\sigma_2$}    (2j);
\draw [->, blue] (11) edge[bend left=2,below] node[shift={(-0.7cm,-0.5cm)}]{$\color{blue}\sigma_3$}    (3j);
\draw [->, blue] (21) edge[bend left=2,below] node[shift={(-0.7cm,-0.5cm)}]{$\color{blue}\sigma_4$}    (4j);
\draw [->, blue] (31) edge[bend left=2,below] node[shift={(-0.7cm,-0.5cm)}]{$\color{blue}\sigma_4$}    (5j);

\draw [->, blue] (0j) edge[bend left=2,below] node[shift={(-0.7cm,-0.5cm)}]{$\color{blue}\sigma_2$}    (2k);
\draw [->, blue] (1j) edge[bend left=2,below] node[shift={(-0.7cm,-0.5cm)}]{$\color{blue}\sigma_3$}    (3k);
\draw [->, blue] (2j) edge[bend left=2,below] node[shift={(-0.7cm,-0.5cm)}]{$\color{blue}\sigma_4$}    (4k);
\draw [->, blue] (3j) edge[bend left=2,below] node[shift={(-0.7cm,-0.5cm)}]{$\color{blue}\sigma_5$}    (5k);

%https://tex.stackexchange.com/a/20986/139648
\draw [decorate,decoration={brace,amplitude=10pt,raise=10pt,mirror}] (00.south west) -- (50.south east) node[midway,yshift=-3em]{\textbf{String length}};
\draw [decorate,decoration={brace,amplitude=10pt,raise=20pt}] (00.south west) -- (0k.north west) node[midway,xshift=-1.2cm,yshift=-0.6cm,rotate=-60]{\textbf{Edit distance}};
\end{tikzpicture}
}
    \end{center}
    \caption{NFA recognizing Levenshtein $L(\sigma: \Sigma^5, 3)$.}\label{fig:lev_nfa}
    \vspace{-0.5cm}
  \end{wrapfigure}

  Levenshtein edits are recognized by an automaton known as the Levenshtein automaton. As the original construction defined by Schultz and Mihov~\cite{schulz2002fast} contains cycles and $\varepsilon$-transitions, we propose a variant which is $\varepsilon$-free and acyclic. Furthermore, we adopt a nominal form which supports infinite alphabets and considerably simplifies the language intersection to follow. Illustrated in Fig.~\ref{fig:lev_nfa} is an example of a small Levenshtein automaton recognizing $L(\sigma: \Sigma^5, 3)$. Unlabeled arcs accept any terminal from the alphabet, $\Sigma$. Equivalently, this transition system can be viewed as a kind of proof system within an unlabeled lattice. The following construction is equivalent to Schultz and Mihov's original Levenshtein automaton, but is more amenable to our purposes as it does not any contain $\varepsilon$-arcs, and instead uses skip connections to recognize consecutive deletions of varying lengths.

  \begin{prooftree}
    \AxiomC{$s\in\Sigma \phantom{\land} i \in [0, n] \phantom{\land} j \in [1, d_{\max}]$}
    \RightLabel{$\duparrow$}
    \UnaryInfC{$(q_{i, j-1} \overset{s}{\rightarrow} q_{i,j}) \in \delta$}
    \DisplayProof
    \hskip 1.5em
    \AxiomC{$s\in\Sigma \phantom{\land} i \in [1, n] \phantom{\land} j \in [1, d_{\max}]$}
    \RightLabel{$\ddiagarrow$}
    \UnaryInfC{$(q_{i-1, j-1} \overset{s}{\rightarrow} q_{i,j}) \in \delta$}
  \end{prooftree}
  \begin{prooftree}
    \AxiomC{$i \in [1, n] \phantom{\land} j \in [0, d_{\max}]$}
    \RightLabel{$\drightarrow$}
    \UnaryInfC{$(q_{i-1, j} \overset{\sigma_i}{\rightarrow} q_{i,j}) \in \delta$}
    \DisplayProof
    \hskip 1.5em
    \AxiomC{$d \in [1, d_{\max}] \phantom{\land} i \in [d + 1, n] \phantom{\land} j \in [d, d_{\max}]$}
    \RightLabel{$\knightarrow$}
    \UnaryInfC{$(q_{i-d-1, j-d} \overset{\sigma_i}{\rightarrow} q_{i,j}) \in \delta$}
  \end{prooftree}
  \begin{prooftree}
    \AxiomC{$\vphantom{|}$}
    \RightLabel{$\textsc{Init}$}
    \UnaryInfC{$q_{0,0} \in I$}
    \DisplayProof
    \hskip 1.5em
    \AxiomC{$q_{i, j}$}
    \AxiomC{$|n-i+j| \leq d_{\max}$}
    \RightLabel{$\textsc{Done}$}
    \BinaryInfC{$q_{i, j}\in F$}
  \end{prooftree}

  \newcommand{\substitutionExample}{
    \tikz{
      \foreach \x in {0,8,16,24,32,40}{
        \fill (\x pt,0pt) circle [radius = 1pt];
        \fill (\x pt,8pt) circle [radius = 1pt];
      }
      \phantom{\fill (0pt,-8pt) circle [radius = 1pt];}
      \draw [-to] (0pt,0pt) -- (8pt,0pt);
      \draw [-to] (8pt,0pt) -- (16pt,0pt);
      \draw [-to] (16pt,0pt) -- (24pt,8pt);
      \draw [-to] (24pt,8pt) -- (32pt,8pt);
      \draw [-to] (32pt,8pt) -- (40pt,8pt);
    }
  }

  \newcommand{\insertionExample}{
    \tikz{
      \foreach \x in {0,8,16,24,32,40}{
        \fill (\x pt,0pt) circle [radius = 1pt];
        \fill (\x pt,8pt) circle [radius = 1pt];
      }
      \phantom{\fill (0pt,-8pt) circle [radius = 1pt];}
      \fill[white] (16pt,0pt) circle [radius = 1.2pt];
      \fill[white] (24pt,8pt) circle [radius = 1.2pt];
      \draw [-to] (0pt,0pt) -- (8pt,0pt);
      \draw [-to] (8pt,0pt) -- (24pt,0pt);
      \draw [-to] (24pt,0pt) -- (16pt,8pt);
      \draw [-to] (16pt,8pt) -- (32pt,8pt);
      \draw [-to] (32pt,8pt) -- (40pt,8pt);
    }
  }

  \newcommand{\deletionExample}{
    \tikz{
      \foreach \x in {0,8,16,24,32,40}{
        \fill (\x pt,0pt) circle [radius = 1pt];
        \fill (\x pt,8pt) circle [radius = 1pt];
      }
      \phantom{\fill (0pt,-8pt) circle [radius = 1pt];}
      \draw [-to] (0pt,0pt) -- (8pt,0pt);
      \draw [-to] (8pt,0pt) -- (16pt,0pt);
      \draw [-to] (16pt,0pt) -- (24pt,0pt);
      \draw [-to] (24pt,0pt) -- (40pt,8pt);
    }
  }

  \newcommand{\doubleDeletionExample}{
    \tikz{
      \foreach \x in {0,8,16,24,32,40}{
        \fill (\x pt,0pt) circle [radius = 1pt];
        \fill (\x pt,8pt) circle [radius = 1pt];
        \fill (\x pt,16pt) circle [radius = 1pt];
      }
      \draw [-to] (0pt,0pt) -- (24pt,16pt);
      \draw [-to] (24pt,16pt) -- (32pt,16pt);
      \draw [-to] (32pt,16pt) -- (40pt,16pt);
    }
  }

  \newcommand{\subDelExample}{
    \tikz{
      \foreach \x in {0,8,16,24,32,40}{
        \fill (\x pt,0pt) circle [radius = 1pt];
        \fill (\x pt,8pt) circle [radius = 1pt];
        \fill (\x pt,16pt) circle [radius = 1pt];
      }
      \draw [-to] (0pt,0pt) -- (8pt,0pt);
      \draw [-to] (8pt,0pt) -- (16pt,8pt);
      \draw [-to] (16pt,8pt) -- (32pt,16pt);
      \draw [-to] (32pt,16pt) -- (40pt,16pt);
    }
  }

  \newcommand{\subSubExample}{
    \tikz{
      \foreach \x in {0,8,16,24,32,40}{
        \fill (\x pt,0pt) circle [radius = 1pt];
        \fill (\x pt,8pt) circle [radius = 1pt];
        \fill (\x pt,16pt) circle [radius = 1pt];
      }
      \draw [-to] (0pt,0pt) -- (8pt,0pt);
      \draw [-to] (8pt,0pt) -- (16pt,8pt);
      \draw [-to] (16pt,8pt) -- (24pt,16pt);
      \draw [-to] (24pt,16pt) -- (32pt,16pt);
      \draw [-to] (32pt,16pt) -- (40pt,16pt);
    }
  }

  \newcommand{\insertDeleteExample}{
    \tikz{
      \foreach \x in {0,8,16,24,32,40,48}{
        \fill (\x pt,0pt) circle [radius = 1pt];
        \fill (\x pt,8pt) circle [radius = 1pt];
        \fill (\x pt,16pt) circle [radius = 1pt];
      }
      \fill[white] (16pt,16pt) circle [radius = 1.2pt];
      \fill[white] (8pt,0pt) circle [radius = 1.2pt];
      \fill[white] (16pt,8pt) circle [radius = 1.2pt];
      \draw [-to] (0pt,0pt) -- (16pt,0pt);
      \draw [-to] (16pt,0pt) -- (8pt,8pt);
      \draw [-to] (8pt,8pt) -- (24pt,8pt);
      \draw [-to] (24pt,8pt) -- (40pt,16pt);
      \draw [-to] (40pt,16pt) -- (48pt,16pt);
    }
  }

  Each arc plays a specific role. $\duparrow$ handles insertions, $\ddiagarrow$ handles substitutions and $\knightarrow$ handles deletions of one or more terminals. Let us consider some illustrative cases.

  \begin{table}[h!]
    \begin{tabular}{ccccccc}

      \texttt{f\hspace{3pt}.\hspace{3pt}\hlorange{[}\hspace{3pt}x\hspace{3pt})} &
      \texttt{f\hspace{3pt}.\hspace{3pt}\phantom{(}\hspace{3pt}x\hspace{3pt})} &
      \texttt{f\hspace{3pt}.\hspace{3pt}(\hspace{3pt}\hlred{x}\hspace{3pt})} &
      \texttt{\hlred{.}\hspace{3pt}\hlred{+}\hspace{3pt}(\hspace{3pt}x\hspace{3pt})} &
      \texttt{f\hspace{3pt}\hlorange{.}\hspace{3pt}\hlred{(}\hspace{3pt}x\hspace{3pt};} &
      \texttt{[\hspace{3pt}\hlorange{,}\hspace{3pt}\hlorange{x}\hspace{3pt}y\hspace{3pt}]} &
      \texttt{[\hspace{3pt}\phantom{,}\hspace{3pt},\hspace{3pt}\hlred{x}\hspace{3pt}y\hspace{3pt}]} \\

      \texttt{f\hspace{3pt}.\hspace{3pt}\hlorange{(}\hspace{3pt}x\hspace{3pt})} &
      \texttt{f\hspace{3pt}.\hspace{3pt}\hlgreen{(}\hspace{3pt}x\hspace{3pt})} &
      \texttt{f\hspace{3pt}.\hspace{3pt}(\hspace{3pt}\phantom{x}\hspace{3pt})} &
      \texttt{\phantom{f}\hspace{3pt}\phantom{.}\hspace{3pt}(\hspace{3pt}x\hspace{3pt})} &
      \texttt{f\hspace{3pt}\hlorange{*}\hspace{3pt}\phantom{(}\hspace{3pt}x\hspace{3pt};} &
      \texttt{[\hspace{3pt}\hlorange{x}\hspace{3pt}\hlorange{,}\hspace{3pt}y\hspace{3pt}]} &
      \texttt{[\hspace{3pt}\hlgreen{x}\hspace{3pt},\hspace{3pt}\phantom{x}\hspace{3pt}y\hspace{3pt}]} \\

      \substitutionExample & \insertionExample & \deletionExample & \doubleDeletionExample & \subDelExample & \subSubExample & \insertDeleteExample
    \end{tabular}
  \end{table}

  Note that the same patch can have multiple Levenshtein alignments. $\textsc{Done}$ constructs the final states, which are all states accepting strings $\sigma'$ whose Levenshtein distance $\Delta(\sigma, \sigma') \leq d_\max$.

  To avoid creating a parallel bundle of arcs for each insertion and substutition point, we instead decorate each arc with a nominal predicate, accepting or rejecting $\sigma_i$. To distinguish this nominal variant from the original construction, we highlight the modified rules in orange below.

  \begin{prooftree}
    \AxiomC{$i \in [0, n] \phantom{\land} j \in [1, k]$}
    \RightLabel{$\duparrow$}
    \UnaryInfC{$(q_{i, j-1} \overset{{\color{orange}[\neq \sigma_i]}}{\rightarrow} q_{i,j}) \in \delta$}
    \DisplayProof
    \hskip 1.5em
    \AxiomC{$i \in [1, n] \phantom{\land} j \in [1, k]$}
    \RightLabel{$\ddiagarrow$}
    \UnaryInfC{$(q_{i-1, j-1} \overset{{\color{orange}[\neq \sigma_i]}}{\rightarrow} q_{i,j}) \in \delta$}
  \end{prooftree}
  \begin{prooftree}
    \AxiomC{$i \in [1, n] \phantom{\land} j \in [0, k]$}
    \RightLabel{$\drightarrow$}
    \UnaryInfC{$(q_{i-1, j} \overset{{\color{orange}[=\sigma_i]}}{\rightarrow} q_{i,j}) \in \delta$}
    \DisplayProof
    \hskip 1.5em
    \AxiomC{$d \in [1, d_{\max}] \phantom{\land} i \in [d + 1, n] \phantom{\land} j \in [d, k]$}
    \RightLabel{$\knightarrow$}
    \UnaryInfC{$(q_{i-d-1, j-d} \overset{{\color{orange}[=\sigma_i]}}{\rightarrow} q_{i,j}) \in \delta$}
  \end{prooftree}

  Nominalizing the NFA eliminates the creation of $e=2(|\Sigma| - 1)\cdot|\sigma|\cdot d_\max$ unnecessary arcs over the entire Levenshtein automaton and drastically reduces the size of the construction to follow, but does not affect the underlying semantics. Thus, it is essential to first nominalize the automaton before proceeding to avoid a large blowup in the intermediate grammar.

  \subsection{Recognizing syntactically valid edits via language intersection}\label{sec:lev_bh}

  We now describe the Bar-Hillel construction, which generates a grammar recognizing the intersection between a regular and a context-free language, then specialize it to Levenshtein intersections.

  \begin{lemma}\label{lemma:bar-hillel}
  For any context-free language $\ell$ and finite state automaton $\alpha$, there exists a context-free grammar $G_\cap$ such that $\mathcal{L}(G_\cap) = \ell \cap \mathcal{L}(\alpha)$. See Bar-Hillel~\cite{bar1961formal}.
  \end{lemma}

  \noindent Although Bar-Hillel~\cite{bar1961formal} lacks an explicit construction, Beigel and Gasarch~\cite{beigelproof} construct $G_\cap$ like so:\vspace{-2pt}

  \begin{prooftree}
    \AxiomC{$q \in I \phantom{\land} r \in F\vphantom{\overset{a}{\rightarrow}}$}
    \RightLabel{$\downarrow$}
    \UnaryInfC{$\big(S\rightarrow q S r\big) \in P_\cap$}
    \DisplayProof
    \hskip 1em
    \AxiomC{$(A \rightarrow a) \in P$}
    \AxiomC{$(q\overset{a}{\rightarrow}r) \in \delta$}
    \RightLabel{$\uparrow$}
    \BinaryInfC{$\big(qAr\rightarrow a\big)\in P_\cap$}
    \DisplayProof
    \hskip 1em
%\end{prooftree}
%\begin{prooftree}
    \AxiomC{$(w \rightarrow xz) \in P\vphantom{\overset{a}{\rightarrow}}$}
    \AxiomC{$p,q,r \in Q$}
    \RightLabel{\Join}
    \BinaryInfC{$\big(pwr\rightarrow (pxq)(qzr)\big) \in P_\cap$}
  \end{prooftree}\vspace{2pt}

  This, now standard, Bar-Hillel construction applies to any CFL and REG language intersection, but generates a grammar whose cardinality is approximately $|P_\cap|=|I|\cdot|F| + |P|\cdot|\Sigma|\cdot|\sigma|\cdot2d_{\max} + |P|\cdot|Q|^3$. Applying the BH construction directly to practical languages and code snippets can generate hundreds of trillions of productions for even modestly-sized grammars and Levenshtein automata. Instead, we will describe a kind of reachability analysis that elides many superfluous productions in the case of Levenshtein intersection, greatly reducing the size of the intersection grammar, $G_\cap$.

%  \begin{wrapfigure}{r}{0.3\textwidth}
%    \vspace{-10pt}
%    \resizebox{0.3\textwidth}{!}{\def\secondcirclepath{(1.15,0) coordinate (e) circle (2cm)}

\begin{tikzpicture}[
  dot/.style = {circle, inner sep=0pt, minimum size=1mm, fill,
  node contents={}}
]
  \def\firstcircle{(-2.1,0) coordinate (a) circle (2.4cm)}
  \def\firstcirclea{(-2.1,0) coordinate (b) circle (0.6cm)}
  \def\firstcircleb{(-2.1,0) coordinate (c) circle (1.2cm)}
  \def\firstcirclec{(-2.1,0) coordinate (d) circle (1.8cm)}
  \def\secondcircle{(1.2,0) coordinate (e) circle (1.5cm)}

  \begin{scope}
    \clip[decorate, decoration={snake, amplitude=0.6mm, segment length=5.01mm}] \secondcirclepath;
    \fill[black!15] \firstcircle;
  \end{scope}

  \draw \firstcircle node[dot,label=$\err{\sigma}$](z0);
  \draw [dashed] \firstcirclea;
  \draw [dashed] \firstcircleb;
  \draw [dashed] \firstcirclec;
  \draw[-stealth] (-2.1,0) -- (-1.5, 0) node[midway,below]{$d_1$};
  \draw[-stealth] (-1.5,0) -- (-0.9, 0) node[midway,below]{$d_2$};
  \draw[-stealth] (-0.9,0) -- (-0.3, 0) node[midway,below]{$d^*$};
  \draw[-stealth] (-0.3,0) -- (0.3, 0) node[midway,above]{$\tilde{\sigma}$};
  \draw[-stealth] (-0.3,0) -- (0.3, 0) node[midway,below]{$d^+$};

  \draw[decorate, decoration={snake, amplitude=0.6mm, segment length=5.01mm}] \secondcirclepath;
  \node [above] at (current bounding box.north -| a) {$\mathcal{L}\bigl(G(\err\sigma, d^*)\bigr)$};
  \node [above,yshift=2.1cm] at (e) {$\mathcal{L}(\mathcal{G}')$};
\end{tikzpicture}}
%    \caption{$G_\cap$ is gradually expanded by increasing $d$ until $\mathcal{L}(G_\cap)\neq \varnothing$.}
%    \vspace{-10pt}
%  \end{wrapfigure}

  Consider $\Join$, the most expensive rule. What $\Join$ tells us is each nonterminal in the intersection grammar $\langle q, v, q'\rangle$ matches a substring simultaneously recognized by (1) a pair of states $q, q'$ in the original NFA and (2) a nonterminal, $v$, in the original CFG. A key observation is that $\Join$ generates the Cartesian product of every such triple, but this is a gross overapproximation for most NFAs and CFGs, as the vast majority of all state pairs and nonterminals recognize no strings in common.

  To identify these superfluous triples, we define an interval domain that soundly overapproximates the Parikh image, encoding the minimum and maximum number of terminals each nonterminal can generate. Since some intervals may be right-unbounded, we write $\mathbb{N}^*=\mathbb{N} \cup \{\infty\}$ to denote the upper bound, and $\Pi = \{[a, b] \in \mathbb{N} \times \mathbb{N}^* \mid a \leq b\}^{|\Sigma|}$ to denote the Parikh image of all terminals.

  \begin{definition}[Parikh mapping of a nonterminal]\label{def:parikh}
    Let $p: \Sigma^*\rightarrow\mathbb{N}^{|\Sigma|}$ be the Parikh operator~\cite{parikh1966context}, which counts the frequency of terminals in a string. We define the Parikh map, $\pi: V \rightarrow \Pi$, as a function returning the smallest interval such that $\forall \sigma: \Sigma^*, \forall v: V$, $v \Rightarrow^* \sigma \vdash p(\sigma) \in \pi(v)$.
  \end{definition}

  In other words, the Parikh mapping computes the greatest lower and least upper bound of the Parikh image over all strings in the language of a nonterminal. The infimum of a nonterminal's Parikh interval tells us how many of each terminal a nonterminal \textit{must} generate, and the supremum tells us how many it \textit{can} generate. Likewise, we define a similar relation over NFA state pairs:

  \begin{definition}[Parikh mapping of NFA states]
    We define $\pi: Q\times Q \rightarrow \Pi$ as returning the smallest interval such that $\forall \sigma: \Sigma^*, \forall q, q': Q$, $q \overset{\sigma}{\Longrightarrow} q' \vdash p(\sigma) \in \pi(q, q')$.
  \end{definition}

  Next, we will define a measure on Parikh intervals representing the minimum total edits required to transform a string in one Parikh interval to a string in another, across all such pairings.

  \begin{definition}[Parikh divergence]
    Given two Parikh intervals $\pi, \pi': \Pi$, we define the divergence between them as $\pi \parallel \pi' = \sum_{n=1}^{|\Sigma|} \min_{(i, i') \in \pi[n]\times \pi'[n]} |i - i'|$.
  \end{definition}

  Now, we know that if the Parikh divergence between two intervals exceeds the Levenshtein margin between two states in a Lev-NFA, those intervals must be incompatible as no two strings, one from each Parikh interval, can be transformed into the other with fewer than $\pi \parallel \pi'$ edits.

  \begin{definition}[Levenshtein-Parikh compatibility]
    Let $q = q_{h,i}, q'=q_{j,k}$ be two states in a Lev-NFA and V be a CFG nonterminal. We say that $\langle q, v, q'\rangle: Q\times V\times Q$ are compatible iff the Parikh divergence is bounded by the Levenshtein margin $k-i$, i.e., $v \lhd qq' \iff (\pi(v) \parallel \pi(q, q')) \leq k-i$.
  \end{definition}

  Finally, we define the modified Bar-Hillel construction for nominal Levenshtein automata as:\vspace{-2pt}

  \begin{prooftree}
    \hskip -0.9em
    \def\defaultHypSeparation{\hskip 0.14cm}
    \AxiomC{$(A \rightarrow a) \in P$}
    \AxiomC{$(q\overset{{\color{orange}[\cdot]}}{\rightarrow}r) \in \delta$}
    \AxiomC{$\color{orange}a[\cdot]$}
    \RightLabel{$\uparrow$}
    \TrinaryInfC{$\big(qAr\rightarrow a\big)\in P_\cap$}
    \DisplayProof
    \AxiomC{$\vphantom{\overset{[\cdot]}{\rightarrow}}\color{orange} w \lhd pr \phantom{\land} x \lhd pq \phantom{\land} z \lhd qr$}
    \AxiomC{$(w \rightarrow xz) \in P\vphantom{\overset{a}{\rightarrow}}$}
    \AxiomC{$p,q,r \in Q$}
    \RightLabel{\Join}
    \TrinaryInfC{$\big(pwr\rightarrow (pxq)(qzr)\big) \in P_\cap$}
  \end{prooftree}\vspace{2pt}

  \noindent Once constructed, we normalize $G_\cap$ by removing unreachable and non-generating productions~\cite{firsov2015certified} to obtain $G_\cap'$, which is a recognizer for the admissible set, i.e., $\mathcal{L}(G_\cap') = A$, satisfying Def.~\ref{def:bcflr}.

%Specifically, we compute Parikh intervals generated by every path though the Levenshtein automaton, then intersect the Parikh intervals for the candidate nonterminals in question. For example, suppose we have a $p, q, r: Q$ and $w \rightarrow x z$. Let's check...

%To generate edits from it, we can use the same procedure as before, but instead of interleaving $\err\sigma$ with $\varepsilon$ and introducing holes, we simply use $A\big((\_)^{|\err{\sigma}| + d}\big, G_\cap)$.

  Now that we have a language to recognize nearby repairs, we will need a method to generate the repairs themselves. We impose specific criteria on such a procedure: it must generate only valid repairs and all repairs in the language if possible, otherwise as many as can be sampled in an arbitrary but fixed timeout. In the following sections, we will describe a constructor (\S~\ref{sec:matrix_completion}) for a data structure (\S~\ref{sec:ptree}) representing parse forests in a length-bounded CFL. Among other features, this data structure provides an explicit way to construct the Parikh map for the Levenshtein Bar-Hillel (LBH) construction, and a method for sampling the language with or without replacement.


  \subsection{Code completion as idempotent matrix completion}\label{sec:matrix_completion}

  In this section, we will introduce the porous completion problem and show how it can be translated to a kind of idempotent matrix completion, whose roots are valid strings in a context-free language. This technique is convenient for its geometric interpretability, parallelizability, and generalizability to any CFG, regardless of finitude or ambiguity. We will see how, by redefining the algebraic operations $\oplus, \otimes$ over different carrier sets, one can obtain a recognizer, porous synthesizer, parser, generator, Parikh map and other convenient structures for CFL intersection and membership.

  Given a CFG, $G' : \mathcal{G}$ in Chomsky Normal Form (CNF), we can construct a recognizer $R: \mathcal{G} \rightarrow \Sigma^n \rightarrow \mathbb{B}$ for strings $\sigma: \Sigma^n$ as follows. Let $2^V$ be our domain, $0$ be $\varnothing$, $\oplus$ be $\cup$, and $\otimes$ be defined as:\vspace{-10pt}

  \begin{align}
    X \otimes Z = \big\{\;w \mid \langle x, z\rangle \in X \times Z, (w\rightarrow xz) \in P\;\big\}
  \end{align}

  \noindent If we define $\hat\sigma_r = \{w \mid (w \rightarrow \sigma_r) \in P\}$, then construct a matrix with nonterminals on the superdiagonal representing each token, $M_0[r+1=c](G', \sigma) = \;\hat\sigma_r$, the fixpoint $M_{i+1} = M_i + M_i^2$ is uniquely determined by the superdiagonal entries. The fixedpoint iteration proceeds as follows:\vspace{-10pt}

  \begin{align*}
    M_0=
    \begin{pNiceMatrix}[nullify-dots,xdots/line-style=loosely dotted]
      \varnothing & \hat\sigma_1 & \varnothing & \Cdots & \varnothing \\
      \Vdots      & \Ddots       & \Ddots      & \Ddots & \Vdots\\
                  &              &             &        & \varnothing\\
                  &              &             &        & \hat\sigma_n \\
      \varnothing & \Cdots       &             &        & \varnothing
    \end{pNiceMatrix} &\Rightarrow
    \begin{pNiceMatrix}[nullify-dots,xdots/line-style=loosely dotted]
      \varnothing & \hat\sigma_1 & \Lambda & \Cdots & \varnothing \\
      \Vdots      & \Ddots       & \Ddots  & \Ddots & \Vdots\\
                  &              &         &        & \Lambda\\
                  &              &         &        & \hat\sigma_n \\
      \varnothing & \Cdots       &         &        & \varnothing
    \end{pNiceMatrix} &\Rightarrow \ldots \Rightarrow M_\infty =
    \begin{pNiceMatrix}[nullify-dots,xdots/line-style=loosely dotted]
      \varnothing & \hat\sigma_1 & \Lambda & \Cdots & \Lambda^*_\sigma\\
      \Vdots      & \Ddots       & \Ddots  & \Ddots & \Vdots\\
                  &              &         &        & \Lambda\\
                  &              &         &        & \hat\sigma_n \\
      \varnothing & \Cdots       &         &        & \varnothing
    \end{pNiceMatrix}
  \end{align*}

  Once obtained, the proposition $[S \in \Lambda^*_\sigma]$ decides language membership, i.e., $[\sigma \in \mathcal{L}(\mathcal{G})]$~\footnote{Hereinafter, we use Iverson brackets to denote the indicator function of a predicate with free variables, i.e., $[P] \Leftrightarrow \mathds{1}(P)$.}. So far, this procedure is essentially the textbook CYK algorithm in a linear algebraic notation~\cite{goodman1999semiring}.

  This procedure can be lifted to the domain of strings containing free variables, which we call the \textit{porous completion problem}. In this case, the fixpoint is characterized by a system of language equations, whose solutions are the set of all sentences consistent with the template.

  \begin{definition}[Porous completion]
    Let $\underline\Sigma = \Sigma \cup \{\_\}$, where $\_$ denotes a hole. We denote $\sqsubseteq: \Sigma^n \times \underline\Sigma^n$ as the relation $\{\langle\sigma', \sigma\rangle \mid \sigma_i \in \Sigma \implies \sigma_i' = \sigma_i\}$ and the set of all inhabitants $\{\sigma': \Sigma^+ \mid \sigma' \sqsubseteq \sigma\}$ as $\text{H}(\sigma)$. Given a \textit{porous string}, $\sigma: \underline\Sigma^*$ we seek all syntactically valid inhabitants, i.e., $A(\sigma)=\text{H}(\sigma)\cap\ell$.
  \end{definition}

  Let us consider an example with two holes, $\sigma = 1$ \_ \_, and the context-free grammar being $G=\{S\rightarrow N O N, O \rightarrow + \mid \times, N \rightarrow 0 \mid 1\}$. This grammar will first be rewritten into CNF as $G'= \{S \rightarrow N L, N \rightarrow 0 \mid 1, O \rightarrow \times \mid +, L \rightarrow O N\}$. Using the powerset algebra we just defined, the matrix fixpoint $M' = M + M^2$ can be computed as follows, shown in the leftmost column below:

  \begin{small}
  {\renewcommand{\arraystretch}{1.2}
  \noindent\phantom{...}\begin{tabular}{|c|c|c|c|}
    \hline
    & $2^V$ & $\mathbb{Z}_2^{|V|}$ & $\mathbb{Z}_2^{|V|}\rightarrow\mathbb{Z}_2^{|V|}$\\\hline
    $M_0$ & \begin{pmatrix}
              \phantom{V} & \tiny{\{N\}} &         &             \\
              &              & \{N,O\} &             \\
              &              &         & \{N,O\} \\
              &              &         &
    \end{pmatrix} & \begin{pmatrix}
                      \phantom{V} & \ws\bs\ws\ws &              &              \\
                      &              & \ws\bs\bs\ws &              \\
                      &              &              & \ws\bs\bs\ws \\
                      &              &              &
    \end{pmatrix} & \begin{pmatrix}
                      \phantom{V} & V_{0, 1} &          &          \\
                      &          & V_{1, 2} &          \\
                      &          &          & V_{2, 3} \\
                      &          &          &
    \end{pmatrix} \\\hline
    $M_1$ & \begin{pmatrix}
              \phantom{V} & \tiny{\{N\}} & \varnothing &         \\
              &              & \{N,O\}     & \{L\}   \\
              &              &             & \{N,O\} \\
              &              &             &
    \end{pmatrix} & \begin{pmatrix}
                      \phantom{V} & \ws\bs\ws\ws & \ws\ws\ws\ws &              \\
                      &              & \ws\bs\bs\ws & \bs\ws\ws\ws \\
                      &              &              & \ws\bs\bs\ws \\
                      &              &              &
    \end{pmatrix} & \begin{pmatrix}
                      \phantom{V} & V_{0, 1} & V_{0, 2} &          \\
                      &          & V_{1, 2} & V_{1, 3} \\
                      &          &          & V_{2, 3} \\
                      &          &          &
    \end{pmatrix} \\\hline
    \begin{tabular}{@{}c@{}}$M_2$\\$=$\\$M_\infty$\end{tabular} & \begin{pmatrix}
                   \phantom{V} & \tiny{\{N\}} & \varnothing & \{S\}   \\
                   &              & \{N,O\}     & \{L\}   \\
                   &              &             & \{N,O\} \\
                   &              &             &
    \end{pmatrix} & \begin{pmatrix}
                      \phantom{V} & \ws\bs\ws\ws & \ws\ws\ws\ws & \ws\ws\ws\bs \\
                      &              & \ws\bs\bs\ws & \bs\ws\ws\ws \\
                      &              &              & \ws\bs\bs\ws \\
                      &              &              &
    \end{pmatrix} & \begin{pmatrix}
                      \phantom{V} & V_{0, 1} & V_{0, 2} & V_{0, 3} \\
                      &          & V_{1, 2} & V_{1, 3} \\
                      &          &          & V_{2, 3} \\
                      &          &          &
    \end{pmatrix}\\\hline
  \end{tabular}\\
  }
  \end{small}

  The same procedure can be translated, without loss of generality, into the bit domain ($\mathbb{Z}_2^{|V|}$) using a lexicographic ordering, however $M_\infty$ in both $2^V$ and $\mathbb{Z}_2^{|V|}$ represents a decision procedure, i.e., $[S\in V_{0, 3}]\Leftrightarrow [V_{0, 3, 3}=\bs] \Leftrightarrow [A(\sigma) \neq \varnothing]$. Since $V_{0, 3} = \{S\}$, we know there exists at least one solution $\sigma' \in A$, but $M_\infty$ does not explicitly reveal its identity.

%$\{\text{xor}, \land, \top\}$ is a functionally complete set is equivalent to $\mathbb{Z}_2$ $\top := 1, \land := \times, \text{xor} := +$. We can define $=$ as $(a = b) \Leftrightarrow (a \text{ xor } b) \text{ xor } \top \Leftrightarrow (a + b) + \top$.

  To extract the inhabitants, we can translate the bitwise procedure into an equation with free variables. Here, we can encode the idempotency constraint directly as $M = M^2$. We first define $X \boxtimes Z = [X_2 \land Z_1, \bot, \bot, X_1 \land Z_0]$ and $X \boxplus Z = [X_i \lor Z_i]_{i \in [0, |V|)}$, mirroring $\oplus, \otimes$ from the powerset domain, now over bitvectors. Since the unit nonterminals $O, N$ can only occur on the superdiagonal, they may be safely ignored by $\boxtimes$. To solve for $M_\infty$, we proceed by first computing $V_{0, 2}, V_{1, 3}$:

  \begin{small}
  \begin{align*}
    V_{0, 2} &= V_{0, j} \cdot V_{j, 2} = V_{0, 1} \boxtimes V_{1, 2}                         &  V_{1, 3} &= V_{1, j} \cdot V_{j, 3} = V_{1, 2} \boxtimes V_{2, 3}\\
    &= [L \in V_{0, 2}, \bot, \bot, S \in V_{0, 2}]                                           &  &= [L \in V_{1, 3}, \bot, \bot, S \in V_{1, 3}]\\
    &= [O \in V_{0, 1} \land N \in V_{1, 2}, \bot, \bot, N \in V_{0, 1} \land L \in V_{1, 2}] &  &= [O \in V_{1, 2} \land N \in V_{2, 3}, \bot, \bot, N \in V_{1, 2} \land L \in V_{2, 3}]\\
    &= [V_{0, 1, 2} \land V_{1, 2, 1}, \bot, \bot, V_{0, 1, 1} \land V_{1, 2, 0}]             &  &= [V_{1, 2, 2} \land V_{2, 3, 1}, \bot, \bot, V_{1, 2, 1} \land V_{2, 3, 0}]
  \end{align*}
  \end{small}

  \noindent Now we solve for the corner entry $V_{0, 3}$ by dotting the first row and last column, which yields:

  \begin{align*}
    V_{0, 3} &= V_{0, j} \cdot V_{j, 3} = (V_{0, 1} \boxtimes V_{1, 3}) \boxplus (V_{0, 2} \boxtimes V_{2, 3})\\
%  &= [V_{0, 1, 2} \land V_{1, 3, 1}, \bot, \bot, V_{0, 1, 1} \land V_{1, 3, 0}] + [V_{0, 2, 2} \land V_{2, 3, 1}, \bot, \bot, V_{0, 2, 1} \land V_{2, 3, 0}]\\
    &= [V_{0, 1, 2} \land V_{1, 3, 1} \lor V_{0, 2, 2} \land V_{2, 3, 1}, \bot, \bot, V_{0, 1, 1} \land V_{1, 3, 0} \lor V_{0, 2, 1} \land V_{2, 3, 0}]
  \end{align*}

  \noindent Since we only care about $V_{0, 3, 3} \Leftrightarrow [S \in V_{0, 3}]$, we can ignore the first three entries and solve for:

  \begin{align*}
    V_{0, 3, 3} &= V_{0, 1, 1} \land V_{1, 3, 0} \lor V_{0, 2, 1} \land V_{2, 3, 0}\\
    &= V_{0, 1, 1} \land (V_{1, 2, 2} \land V_{2, 3, 1}) \lor V_{0, 2, 1} \land \bot\\
    &= V_{0, 1, 1} \land V_{1, 2, 2} \land V_{2, 3, 1}\\
    &= [N \in V_{0, 1}] \land [O \in V_{1, 2}] \land [N \in V_{2, 3}]
  \end{align*}

  Now we know that $\sigma =$ 1 \underline{O} \underline{N} is a valid solution, and we can take the product $\{1\}\times \hat\sigma_2^{-1}(O) \times \hat\sigma_3^{-1}(N)$ to recover the inhabitants, yielding $A=\{1+0, 1+1, 1\times 0, 1\times 1\}$. In this case, since $G$ is unambiguous, there is only one parse tree satisfying $V_{0, |\sigma|, 3}$, but in general, there can be multiple valid parse trees.

  \subsection{An algebraic datatype for context-free parse forests}\label{sec:ptree}

  The procedure described in \S~\ref{sec:matrix_completion} generates solutions satisfying the matrix fixpoint, but forgets provenance. The question naturally arises, is there a way to solve for the parse trees directly? This would allow us to handle ambiguous grammars, whilst preserving the natural treelike structure.

  \begin{wrapfigure}{r}{0.47\textwidth}
    \vspace{-9pt}
    \resizebox{0.47\textwidth}{!}{
      \begin{tikzpicture}
      [
        grow                    = right,
        sibling distance        = 3em,
        level distance          = 5em,
        edge from parent/.style = {draw, -latex},
        every node/.style       = {font=\footnotesize},
        sloped,
        treenode/.style = {shape=rectangle, rounded corners,
        draw, align=center,
        top color=white, bottom color=blue!20},
        root/.style     = {treenode, font=\tiny, bottom color=red!30},
        env/.style      = {treenode, font=\tiny},
        dummy/.style    = {circle,draw}
      ]
        \node [root] {S}
        child { node [env] {BC}
        child { node [root] {B}
        child { node [env] {RD}
        child { node [root] {R} edge from parent node [above] }
        child { node [root] {D} edge from parent node [above] }
        edge from parent node [above] }
        edge from parent node [below] }
        child { node [root] {C}
        child { node [env] {\ldots\vphantom{BB}} edge from parent node [below] }
%  child { edge from parent node [above] {\ldots} }
        edge from parent node [below] }
        edge from parent node [above] }
        child { node [env] {\ldots\vphantom{BB}} edge from parent node [below] }
        child { node [env] {AB}
        child { node [root] {A}
        child {
          node [env] {QC}
          child { node [root] {Q} edge from parent node [above] }
          child { node [root] {C} edge from parent node [above] }
          edge from parent node [above]
        }
%    child { node [env] {ZQ} edge from parent node [above] }
        child { node [env] {\ldots\vphantom{BB}} edge from parent node [below] }
        edge from parent node [below] }
        child { node [root] {B}
        child { node [env] {RD}
        child { node [root] {R} edge from parent node [above] }
        child { node [root] {D} edge from parent node [above] }
        edge from parent node [above] }
        edge from parent node [below] }
        edge from parent node [above] };
      \end{tikzpicture}
    }
    \caption{A partial $\mathbb{T}_2$ corresponding to the grammar $\{S \rightarrow BC \mid \ldots \mid AB, B\rightarrow RD \mid \ldots, A\rightarrow QC \mid \ldots\}$.}
    \label{fig:ptree}
    \vspace{-12pt}
  \end{wrapfigure}

  We will now describe a datatype for compactly representing CFL parse forests, then redefine the matrix algebra over this domain. This datatype is particularly convenient for tracking provenance under ambiguity, constructing the Parikh map for a CFG, counting the size of a finite CFL, and sampling parse trees with or without replacement.

  We first define a datatype $\mathbb{T}_3 = (V \cup \Sigma) \rightharpoonup \mathbb{T}_2$ where $\mathbb{T}_2 = (V \cup \Sigma) \times (\mathbb{N} \rightharpoonup \mathbb{T}_2\times\mathbb{T}_2)$\footnote{Given a $T:\mathbb{T}_2$, we may also refer to $\pi_1(T), \pi_2(T)$ as $\texttt{root}(T)$ and $\texttt{children}(T)$ respectively.}. Morally, we can think of $\mathbb{T}_2$ as an implicit set of possible trees that can be generated by a CFG in CNF, consistent with a finite-length porous string. Structurally, we may interpret $\mathbb{T}_2$ as an algebraic data type corresponding to the fixpoints of the following recurrence, which tells us each $\mathbb{T}_2$ can be a terminal, nonterminal, or a nonterminal and a sequence of nonterminal pairs and their two children:\vspace{-10pt}

  \begin{equation}
    L(p) = 1 + p L(p) \phantom{addspace} P(a) = \Sigma + V + V L\big(V^2P(a)^2\big)
  \end{equation}

  Depicted in Fig.~\ref{fig:ptree} is a partial $\mathbb{T}_2$, where red nodes are \texttt{root}s and blue nodes are \texttt{children}. The shape of type $\mathbb{T}_2$ is congruent with an acyclic CFG in Chomsky Normal Form, i.e., $\mathbb{T}_2\cong\mathcal{G}'$, so assuming the CFG recognizes a finite language, as the case for $G_\cap'$, we can translate it directly. If the language is infinite, we slice the CFL, $\mathcal{L}(G)\cap \Sigma^n$, and compute the fixpoint for each slice.

  Given a porous string $\sigma: \underline\Sigma^n$ representing the slice, we can construct $\mathbb{T}_2$ from the bottom-up, and read off structures from the top-down. We construct the first upper diagonal $\hat\sigma_r = \Lambda(\sigma_r)$ as follows:

\vspace{-5pt}\begin{equation}
  \begin{footnotesize}
\Lambda(s: \underline\Sigma^n) \mapsto \begin{cases}
\bigoplus_{s'\in \Sigma} \Lambda(s') & \text{if $s$ is a hole,} \vspace{5pt}\\
\big\{\mathbb{T}_2\big(w, \big[\langle\mathbb{T}_2(s), \mathbb{T}_2(\varepsilon)\rangle\big]\big) \mid (w \rightarrow s)\in P\big\} & \text{otherwise.}
\end{cases}
  \end{footnotesize}
\end{equation}

\noindent This initializes the superdiagonal entries of $M_0$, enabling us to compute the fixpoint $M_\infty$ in the same manner described in \S~\ref{sec:matrix_completion} by redefining $\oplus, \otimes: \mathbb{T}_3 \times \mathbb{T}_3 \rightarrow \mathbb{T}_3$ as:

\vspace{-5pt}\begin{align}
  X \oplus Z &\mapsto \bigcup_{\mathclap{k\in \pi_1(X \cup Z)}}\Big\{\phantom{.}k \Rightarrow \mathbb{T}_2(k, x \cup z) \mid x \in \pi_2(X\circ k), z \in \pi_2(Z\circ k)\phantom{.}\Big\}\\
  X \otimes Z &\mapsto \bigoplus_{\mathclap{(w\rightarrow xz) \in P}}\Big\{\phantom{.}\mathbb{T}_2\Big(w, \big[\langle X\circ x, Z\circ z\rangle\big]\Big) \mid x \in \pi_1(X), z \in \pi_1(Z)\phantom{.}\Big\}
\end{align}

These operators group subtrees by their root nonterminal, then aggregate their children. Instead of tracking sets, each $\Lambda$ now becomes a dictionary of $\mathbb{T}_2$, indexed by their root nonterminals.

  $\mathbb{T}_2$ is a convenient datatype for many operations involving CFGs. We can use it to approximate the Parikh image, compute the size of a finite CFG, and sample parse trees with or without replacement. For example, to obtain the Parikh map of a CFG (Def.~\ref{def:parikh}), we may use the following recurrence,

\begin{equation}
  \pi(T: \mathbb{T}_2) \mapsto \begin{cases}
%    \big|\{s \mid \big(\texttt{root}(T) \rightarrow s\big) \in P^\cap\}\big| & \text{if $T$ is a leaf,} \\
  \big[[1, 1] \text{ if } \texttt{root}(T) = s \text{ else } [0, 0]\big]_{s\in \Sigma}  & \text{if $T$ is a leaf,} \\
  \bigoplus_{\langle T_1, T_2\rangle \in \texttt{children}(T)} \pi(T_1) \otimes \pi(T_2) & \text{otherwise.}
  \end{cases}
\end{equation}

  %infix fun IntRange.merge(other: IntRange) =
  %  minOf(start, other.first)..maxOf(last, other.last)
  %
  %operator fun ParikhBounds.plus(other: ParikhBounds) =
  %  (keys + other.keys).associateWith {
  %    (get(it) ?: 0..0) merge (other[it] ?: 0..0)
  %  }
  %
  %operator fun ParikhBounds.times(other: ParikhBounds) =
  %  (keys + other.keys).associateWith {
  %    (get(it) ?: 0..0) join (other[it] ?: 0..0)
  %  }
  %
  %infix fun IntRange.join(other: IntRange) =
  %  (first + other.first)..(last + other.last)

  \noindent where the operations over Parikh maps $\oplus, \otimes: \Pi \times \Pi \rightarrow \Pi$ are defined respectively as follows:

  \begin{align}
      X \oplus Z &\mapsto \big[[\min(X_s \cup Z_s), \max(X_s \cup Z_s)]\big]_{s \in \Sigma}\\ X \otimes Z &\mapsto \big[[\min(X_s) + \min(Z_s), \max(X_s) + \max(Z_s)]\big]_{s \in \Sigma}
  \end{align}

  To obtain the parameterized Parikh map of a length-bounded CFG, we abstractly parse the porous string and take the union of all intervals, which subsumes every repair in the Levenshtein ball. Given a specific programming language syntax, $G$, the following function can be precomputed and cached for all $v: V$, and small values of $n, d: \mathbb{N}$ for the sake of efficiency, then used to retrieve the Levenshtein-Parikh-$\langle v, n, d\rangle $ map for any invalid string $\err\sigma$ of length $n$ in constant time:

  \begin{equation}
  \pi(G: \mathcal{G}, v: V, n: \mathbb{N}, d: \mathbb{N}): \Pi = \bigoplus_{\mathclap{i\in [n-d, n+d]}}\pi\bigl(\Lambda^*(\{\_\}^i)\circ v\bigr)
  \end{equation}

  $\mathbb{T}_2$ also allows us to sample whole parse trees by constructing $(\Lambda^*(\sigma) \circ S): \mathbb{T}_2$ from the bottom-up and sampling top-down. Given a PCFG whose productions indexed by each nonterminal are decorated with a probability vector $\mathbf{p}$ (uniform in the non-probabilistic case), we define a tree sampler $\Gamma: (\mathbb{T}_2 \mid \mathbb{T}_2^2) \rightsquigarrow \mathbb{T}$ which recursively draws children according to a Multinoulli distribution:

\begin{equation}
  \Gamma(T) \mapsto \begin{cases}
        \texttt{BTree}\Big(\texttt{root}(T), \Gamma\big(\text{Multi}(\texttt{children}(T), \mathbf{p})\big)\Big) & \text{ if $T: \mathbb{T}_2$ } \\
        \big\langle \Gamma\big(\pi_1(T)\big), \Gamma\big(\pi_2(T)\big) \big\rangle & \text{ if $T: \mathbb{T}_2\times\mathbb{T}_2$ }
  \end{cases}
\end{equation}

\noindent This method is closely related to the generating function for the ordinary Boltzmann sampler,

\begin{equation}
  \Gamma C(x) \mapsto \begin{cases}
  \text{Bern} \left(\frac{A(x)}{A(x) + B(x)}\right) \rightarrow \Gamma A(x) \mid \Gamma B(x) & \text{ if } \mathcal{C}=\mathcal{A}+\mathcal{B} \\
  \big\langle \Gamma A(x), \Gamma B(x)\big\rangle & \text{ if } \mathcal{C}=\mathcal{A} \times \mathcal{B}
  \end{cases}
\end{equation}

\noindent from analytic combinatorics, however unlike Duchon et al.~\cite{duchon2004boltzmann}, our work does not depend on rejection to guarantee exact-size sampling, as all trees from $\mathbb{T}_2$ will necessarily be the same width.

The number of labeled binary trees inhabiting a single instance of $\mathbb{T}_2$ is sensititive to the number of nonterminals and rule expansions in the grammar. To obtain the total number of trees with breadth $n$, we abstractly parse the porous string, letting $T=\Lambda^*(\{\_\}^n) \circ S$, then use the recurrence below to compute the total number of unique trees in the language:

\begin{equation}
  |T: \mathbb{T}_2| \mapsto \begin{cases}
%    \big|\{s \mid \big(\texttt{root}(T) \rightarrow s\big) \in P^\cap\}\big| & \text{if $T$ is a leaf,} \\
    1 & \text{if $T$ is a leaf,} \\
    \sum_{\langle T_1, T_2\rangle \in \texttt{children}(T)} |T_1| \cdot |T_2| & \text{otherwise.}
  \end{cases}
\end{equation}

To sample all trees in a given $T: \mathbb{T}_2$ uniformly without replacement, we then construct a modular pairing function $\varphi: \mathbb{T}_2 \rightarrow \mathbb{Z}_{|T|} \rightarrow \texttt{BTree}$, that we define as follows:

\begin{equation}\label{eq:pairing}
  \varphi(T: \mathbb{T}_2, i: \mathbb{Z}_{|T|}) \mapsto \begin{cases}
  \Big\langle\texttt{BTree}\big(\texttt{root}(T)\big), i\Big\rangle & \text{if $T$ is a leaf,} \vspace{5pt}\\
  \text{Let } b = |\texttt{children}(T)|,\\
  \phantom{\text{Let }} q_1, r=\big\langle\lfloor\frac{i}{b}\rfloor, i \pmod{b}\big\rangle,\\
  \phantom{\text{Let }} lb, rb = \texttt{children}[r],\\
  \phantom{\text{Let }} T_1, q_2 = \varphi(lb, q_1),\\
  \phantom{\text{Let }} T_2, q_3 = \varphi(rb, q_2) \text{ in } \\
  \Big\langle\texttt{BTree}\big(\texttt{root}(T), T_1, T_2\big), q_3\Big\rangle & \text{otherwise.} \\
  \end{cases}
\end{equation}

  Then, instead of top-down incremental sampling, we can create a randomized $\varphi'$ from $\varphi$ by sampling integers uniformly without replacement from $\mathbb{Z}_{|T|}$ and decode them into whole parse trees. Assuming $\varphi(T, \cdot)$ is in fact bijective, then letting $\bm\varphi(i) = \bigcup_{j\in[1, i]} \{\varphi'(T, j)\}$ will satisfy Def.~\ref{def:linear-convergence} by construction. If the language being sampled is sufficiently small, we can enumerate every tree, otherwise, sample them uniformly without replacement, or with replacement using a PCFG. %This procedure is the basis for our enumerate sampler and the method we use to decode repairs from the intersection grammar.

  \subsection{Decoding repairs in order of maximal likelihood}\label{sec:decoding}

  The previous technique will enumerate parse trees in a given $\mathbb{T}_2$, but in an arbitrary order. While sufficient for the purpose of decoding small finite languages, it can be improved for large finite languages by decoding repairs in order of likelihood. To decode the top-k maximum likelihood results without extracting all repairs and reranking, one can instead define an automata algebra over $\mathcal{M}^{|V|}$, propagating a vector of NFAs, $A$ indexed by the nonterminal $v: V$:

  \begin{align}
    X \oplus Z &\mapsto \bigcup_{v \in V}\big\{\langle v, A_v \rangle \mid \mathcal{L}(A_v) = \mathcal{L}(X[v]) \cup \mathcal{L}(Z[v])\big\}\\
    X \otimes Z &\mapsto \big\{\langle v, A_v \rangle \mid \mathcal{L}(A_v) = \mathcal{L}(X[x]) \times \mathcal{L}(Z[z]), (v \rightarrow xz) \in P\big\}
  \end{align}

  \noindent where the unit nonterminals occupying the first upper diagonal are constructed as follows:

  \begin{equation}
    \begin{footnotesize}
      \Lambda(s: \underline\Sigma^n) \mapsto \Big\{\langle v, A_v\rangle\mid v: V, \mathcal{L}(A_v) = \{\varepsilon\}\Big\} \otimes \begin{cases}
        \big\{\langle v, A_v \rangle \mid \mathcal{L}(A_v) = \Sigma \big\} & \text{if $s$ is a hole,} \vspace{5pt}\\
        \big\{\langle v, A_v \rangle \mid \mathcal{L}(A_v) = \{s\}, (v \rightarrow s)\in P\big\} & \text{otherwise.}
      \end{cases}
    \end{footnotesize}
  \end{equation}

  Constructively, let $+, *: \mathcal{M}\times \mathcal{M} \rightarrow \mathcal{M}$ be the automata operators corresponding to language union and concatenation satisfying $\mathcal{L}(A_1 + A_2) = \mathcal{L}(A_1)\cup\mathcal{L}(A_2)$, and $\mathcal{L}(A_1 * A_2) = \mathcal{L}(A_1)\times\mathcal{L}(A_2)$. This can be implemented using the standard textbook construction, recalling that NFA are closed under these operations. $\Lambda^* \circ S$ would then yield an NFA recognizing $\Sigma^n\cap\mathcal{L}(G)$, which can be determinized and decoded using k-best paths to obtain the top-k maximum likelihood repairs.


  \definecolor{R}{RGB}{202,65,55}
  \definecolor{G}{RGB}{151,216,56}
  \definecolor{W}{RGB}{255,255,255}
  \definecolor{X}{RGB}{65,65,65}

  \newcommand{\TikZRubikFaceLeft}[9]{\def\myarrayL{#1,#2,#3,#4,#5,#6,#7,#8,#9}}
  \newcommand{\TikZRubikFaceRight}[9]{\def\myarrayR{#1,#2,#3,#4,#5,#6,#7,#8,#9}}
  \newcommand{\TikZRubikFaceTop}[9]{\def\myarrayT{#1,#2,#3,#4,#5,#6,#7,#8,#9}}
  \newcommand{\BuildArray}{\foreach \X [count=\Y] in \myarrayL%
  {\ifnum\Y=1%
  \xdef\myarray{"\X"}%
  \else%
  \xdef\myarray{\myarray,"\X"}%
  \fi}%
  \foreach \X in \myarrayR%
  {\xdef\myarray{\myarray,"\X"}}%
  \foreach \X in \myarrayT%
  {\xdef\myarray{\myarray,"\X"}}%
  \xdef\myarray{{\myarray}}%
  }
  \TikZRubikFaceLeft
  {LA}{W}{W}
  {W}{LB}{LC}
  {LD}{W}{W}
  \TikZRubikFaceRight
  {W}{LK}{W}
  {LC}{W}{LG}
  {W}{LH}{W}
  \TikZRubikFaceTop
  {LA}{W}{LI}
  {W}{W}{LJ}
  {W}{LK}{W}
  \BuildArray
  \pgfmathsetmacro\radius{0.1}
  \tdplotsetmaincoords{55}{135}

  \showcellnumberfalse

  \bgroup
  \newcommand\ddd{\Ddots}
  \newcommand\vdd{\Vdots}
  \newcommand\cdd{\Cdots}
  \newcommand\lds{\ldots}
  \newcommand\vno{\varnothing}
  \newcommand{\ts}[1]{\textsuperscript{#1}}
  \newcommand\non{1\ts{st}}
  \newcommand\ntw{2\ts{nd}}
  \newcommand\nth{3\ts{rd}}
  \newcommand\nfo{4\ts{th}}
  \newcommand\nfi{5\ts{th}}
  \newcommand\nsi{6\ts{th}}
  \newcommand\nse{7\ts{th}}
  \newcommand{\vs}[1]{\shup\sigma_{#1}}
  \newcommand\rcr{\rowcolor{black!15}}
  \newcommand\rcw{\rowcolor{white}}
  \newcommand\pcd{\cdot}
  \newcommand\pcp{\phantom\cdot}
  \newcommand\ppp{\phantom{\nse}}

  \begin{wrapfigure}{r}{0.32\textwidth}
    \vspace{-0.5cm}
      \begin{minipage}[r]{4cm}
      \begin{align*}
        o &\rightarrow \hiliD{so} \mid \hiliC{rs} \mid \hiliB{rr}\hspace{0.5pt} \mid \hiliA{oo}\\
        r &\rightarrow \hiliE{so} \mid \hiliH{ss}\hspace{0.4pt}\mid \hiliF{rr}\hspace{0.5pt} \mid \hiliK{os}\\
        s &\rightarrow \hiliL{so} \mid \hiliG{rs} \mid \hiliJ{or} \mid \hiliI{oo}
      \end{align*}
      \end{minipage}
%    \mathcal{J} = \begin{pmatrix}
%       \pder{o}{o} & \pder{o}{r} & \pder{o}{s}\\
%       \pder{r}{o} & \pder{r}{r} & \pder{r}{s}\\
%       \pder{s}{o} & \pder{s}{r} & \pder{s}{s}
%    \end{pmatrix}
    \hspace{-0.1cm}\begin{minipage}[r]{4cm}
       \scalebox{0.8}{\begin{tikzpicture}
          \clip (-3,-2.5) rectangle (3,2.5);
          \begin{scope}[tdplot_main_coords]
            \filldraw [canvas is yz plane at x=1.5] (-1.5,-1.5) rectangle (1.5,1.5);
            \filldraw [canvas is xz plane at y=1.5] (-1.5,-1.5) rectangle (1.5,1.5);
            \filldraw [canvas is yx plane at z=1.5] (-1.5,-1.5) rectangle (1.5,1.5);
            \foreach \X [count=\XX starting from 0] in {-1.5,-0.5,0.5}{
              \foreach \Y [count=\YY starting from 0] in {-1.5,-0.5,0.5}{
                \pgfmathtruncatemacro{\Z}{\XX+3*(2-\YY)}
                \pgfmathsetmacro{\mycolor}{\myarray[\Z]}
                \draw [thick,canvas is yz plane at x=1.5,shift={(\X,\Y)},fill=\mycolor] (0.5,0) -- ({1-\radius},0) arc (-90:0:\radius) -- (1,{1-\radius}) arc (0:90:\radius) -- (\radius,1) arc (90:180:\radius) -- (0,\radius) arc (180:270:\radius) -- cycle;
                \ifshowcellnumber
                \node[canvas is yz plane at x=1.5,shift={(\X+0.5,\Y+0.5)}] {\Z};
                \fi
                \pgfmathtruncatemacro{\Z}{2-\XX+3*(2-\YY)+9}
                \pgfmathsetmacro{\mycolor}{\myarray[\Z]}
                \draw [thick,canvas is xz plane at y=1.5,shift={(\X,\Y)},fill=\mycolor] (0.5,0) -- ({1-\radius},0) arc (-90:0:\radius) -- (1,{1-\radius}) arc (0:90:\radius) -- (\radius,1) arc (90:180:\radius) -- (0,\radius) arc (180:270:\radius) -- cycle;
                \ifshowcellnumber
                \node[canvas is xz plane at y=1.5,shift={(\X+0.5,\Y+0.5)},xscale=-1] {\Z};
                \fi
                \pgfmathtruncatemacro{\Z}{2-\YY+3*\XX+18}
                \pgfmathsetmacro{\mycolor}{\myarray[\Z]}
                \draw [thick,canvas is yx plane at z=1.5,shift={(\X,\Y)},fill=\mycolor] (0.5,0) -- ({1-\radius},0) arc (-90:0:\radius) -- (1,{1-\radius}) arc (0:90:\radius) -- (\radius,1) arc (90:180:\radius) -- (0,\radius) arc (180:270:\radius) -- cycle;
                \ifshowcellnumber
                \node[canvas is yx plane at z=1.5,shift={(\X+0.5,\Y+0.5)},xscale=-1,rotate=-90] {\Z};
                \fi
              }
            }
            \draw [decorate,decoration={calligraphic brace,amplitude=10pt,mirror},yshift=0pt, line width=1.25pt]
            (3,0) -- (3,3) node [black,midway,xshift=-8pt, yshift=-14pt] {\footnotesize $V_x$};
            \draw [decorate,decoration={calligraphic brace,amplitude=10pt},yshift=0pt, line width=1.25pt]
            (3,0) -- (0,-3) node [black,midway,xshift=-16pt, yshift=0pt] {\footnotesize $V_z$};
            \draw [decorate,decoration={calligraphic brace,amplitude=10pt},yshift=0pt, line width=1.25pt]
            (0,-3) -- (-3,-3) node [black,midway,xshift=-8pt, yshift=14pt] {\footnotesize $V_w$};
          \end{scope}
       \end{tikzpicture}}
    \end{minipage}
    \caption{CFGs are hypergraphs witnessed by a rank-3 tensor, whose nonempty inhabitants indicate CNF productions.\vspace{-10pt}}
  \end{wrapfigure}

  In the distributional setting, suppose we have a PCFG, $G$, such that $|\mathcal{L}(G)|<\infty$. This induces a distribution $F_G: \mathcal{L}(G) \rightarrow \mathbb{R}$ which can be bisimulated by a weighted automaton (WA), $A$, where $F_A: \mathcal{L}(A) \rightarrow \mathbb{R}$ such that $F_G \sim F_A$, i.e., sampling paths through $A$ yields members of $\mathcal{L}(G)$ equal probability as sampling the PCFG from the top, down. A key issue with this approach is that $|P_G| \ll |\delta_A|$ for bisimilar $G, A$, and renormalizing the transition weights can be computationally infeasible for large $\delta_A$.

   A more suitable decoding method is possible by considering a related technique called cube-pruning~\cite{huang2005better} that operates over hypergraphs, which are equiexpressive with both CFGs and ADTs. If we consider the PCFG in CNF, each production corresponds to a weighted hyperedge with at most two neighbors. While strictly more expressive than necessary to represent finite languages, cube pruning over context-free hypergraphs is capable of decoding top-k lists far more efficiently than k-longest paths over a WA. A further advantage of this representation is that we can decode likely subpaths from $\mathbb{T}_2$ directly, without translating it to a weighted automaton, minimizing, renormalizing transition weights, and then solving for the top-k best paths. The disadvantage however, is that CFGs are not as algebraically well-behaved, lacking the closure and decidability benefits of WAs under intersection.

  The question then becomes how to efficiently decode the top-k most likely repairs from a PCFG in CNF. Here, we propose a simple method that extracts the production from the intersection grammar $G_\cap$, scores subtrees, then returns the top-k most likely repairs.

 For every nonterminal in the intersection grammar $qvq' \in V_\cap$, we have $v \in V$. Since we do not use a weighted NFA, we may ignore the states $q, q'$ and treat every intersection nonterminal as the original nonterminal from the PCFG. By avoiding the NFA states, this simplifies the conditional probability table. Let's consider a few illustrative examples using graphical model notation. We use upper case to denote a generic nonterminal, and lower case to denote its concrete instantiation, $P$ to denote the normalized probability and $C$ to denote the coincidence count of a tuple. We can compute the probability of a production $W\rightarrow XZ$ with $W$ observed as follows:

  \begin{table}[h!]

    \begin{minipage}{3cm}\resizebox{.6\textwidth}{!}{\begin{tikzpicture}
        \node[obs] (z) {$W$};%
        \node[latent,below=of z,xshift=-1cm,fill] (x) {$X$}; %
        \node[latent,below=of z,xshift=1cm] (y) {$Z$}; %
        \edge {z} {x,y}
      \end{tikzpicture}}\end{minipage}\begin{minipage}{5cm}\begin{align*} P(W\rightarrow XZ) &= P\big((pwr) \rightarrow (pxq)(qzr)\big) \\ &= P(X=x, Z=z \mid W=w)\\ &= \frac{C(X=x, Z=z \mid W)}{\sum_{x', z'}C(X=x', Z=z' \mid W)} \end{align*}\end{minipage}
  \end{table}

Likewise, we can perform this same computation with higher-order dependencies, where $W$ stands for a generic grandparent, $X, Z$ are the generic parent and pibling, and all three are observed:

  \begin{table}[h!]
    \begin{minipage}{3cm}\resizebox{.6\textwidth}{!}{\begin{tikzpicture}
       \node[obs] (z) {$W$};%
       \node[obs,below=of z,xshift=-1cm,fill] (x) {$X$}; %
       \node[obs,below=of z,xshift=1cm] (y) {$Z$}; %
       \node[latent,below=of x,xshift=-1cm] (k) {$K$}; %
       \node[latent,below=of x,xshift=1cm] (l) {$L$}; %
       \edge {z} {x,y}
       \edge {x} {k,l}
    \end{tikzpicture}}\end{minipage}\begin{minipage}{5cm}\begin{align*} P(X\rightarrow KL \mid W\rightarrow XZ) &= P\big((\cdot x \cdot) \rightarrow (\cdot k \cdot)(\cdot l \cdot) \mid (\cdot w \cdot) \rightarrow (\cdot x \cdot)(\cdot z \cdot)\big) \\ &= P(K=k, L=l \mid W, X, Z)\\ &= \frac{C(K=k, L=l \mid W, X, Z)}{\sum_{k', l'}C(K=k', L=l' \mid W, X, Z)} \end{align*}\end{minipage}\end{table}

  Since the parent can be a left- or right-sibling, we also consider a symmetric case:

  \begin{table}[h!]
    \begin{minipage}{3cm}\resizebox{.6\textwidth}{!}{\begin{tikzpicture}
        \node[obs] (z) {$W$};%
        \node[obs,below=of z,xshift=-1cm,fill] (x) {$X$}; %
        \node[obs,below=of z,xshift=1cm] (y) {$Z$}; %
        \node[latent,below=of x,xshift=1cm] (m) {$M$}; %
        \node[latent,below=of x,xshift=3cm] (n) {$N$}; %
        \edge {z} {x,y}
        \edge {y} {m,n}
    \end{tikzpicture}}\end{minipage}\begin{minipage}{5cm}\begin{align*} P(Z\rightarrow MN \mid W\rightarrow XZ) &= P\big((\cdot z \cdot) \rightarrow (\cdot m \cdot)(\cdot n \cdot) \mid (\cdot w \cdot) \rightarrow (\cdot x \cdot)(\cdot z \cdot)\big) \\ &= P(M=m, N=n \mid W, X, Z)\\ &= \frac{C(M=m, N=n \mid W, X, Z)}{\sum_{m', n'}C(M=m', N=n' \mid W, X, Z)} \end{align*}\end{minipage}\\
  \end{table}

%  This method can be extended using a probabilistic higher-order grammar (PHOG). PHOGs are a generalization of PCFGs that allow for local dependencies between neighboring nonterminals, and increase the calibration of sampled trajectories.

  \clearpage\subsection{Ranked repair}\label{sec:ranking}

  Returning to the ranked repair problem (Def.~\ref{def:ranked-repair}), the above procedure returns a set of syntactically consistent repairs, and we need an ordering over them. We note that any metric is sufficient, such as the log-likelihood of the repair under a large language model or the probability under a PCFG. We implement the simplest solution: the likelihood of a low-order Markov chain. This solution is computationally fast, and as we will show, yields competitive results in practice.

  Specifically, given a string $\sigma: \Sigma^*$, we factorize the probability $P_\theta(\sigma)$ as a product of conditionals $\prod_{i=1}^{|\sigma|}P_\theta(\sigma_i \mid \sigma_{i-1}\ldots\sigma_{i-n})$, for some small $n \in \mathbb{N}$. To obtain the parameters $\theta$, we use the standard maximum likelihood estimator for Markov chains. We approximate the joint distribution $P(\Sigma^n)$ directly from data, then the conditionals by normalizing n-gram counts with Laplace smoothing.

  To score the repairs, we use the conventional length-normalized negative log likelihood:

  \begin{equation}
    \text{NLL}(\sigma) = -\frac{1}{|\sigma|}\sum_{i=1}^{|\sigma|}\log P_\theta(\sigma_i \mid \sigma_{i-1}\ldots\sigma_{i-n})
  \end{equation}

  Then, for each retrieved set $\hat{A} \subseteq A$ drawn by the sampler before a predetermined timeout elapses and each $\sigma \in \hat{A}$, we score the repair and return $\hat{A}$ in ascending order by score. If $\hat{A} = A$ and the Markov chain is itself the language model being maximized, then this procedure satisfies Def.~\ref{def:ranked-repair}. Otherwise it is a heuristic, and the quality of the ranking will depend on the quality of $\hat{A}$.

  To evaluate the quality of our ranking, we use the Precision@k statistic, which measures the frequency of repairs in the top-k results matching the true repair. Specifically, given a repair model, $R: \Sigma^* \rightarrow 2^{\Sigma^*}$ and a test set $\mathcal{D}_{\text{test}}$, we define Precision@k as:

  \begin{equation}
    \text{Precision@k}(R) = \frac{1}{|\mathcal{D}_{\text{test}}|}\sum_{\langle\sigma^\dagger, \sigma'\rangle \in \mathcal{D}_{\text{test}}} \mathds{1}\big[\sigma' \in \argmax_{\bm{\sigma} \subset R(\sigma^\dagger), |\bm{\sigma}| \leq k}\sum_{\sigma \in \bm{\sigma}}\text{NLL}(\sigma)\big]
  \end{equation}

  This is a variation on a standard metric used in information retrieval, and a common way to measure the quality of ranked results in machine translation and recommender systems.

  \section{Experimental Setup}

  We use syntax errors and fixes from the Python language to validate our approach.  Python source code fragments are abstracted as a sequence of lexical tokens using the official Python lexer, erasing numbers and identifiers, but retaining all other keywords. Precision is evaluated across a test set by checking for lexical equivalence with the ground-truth repair, following Sakkas et al. (2022)~\cite{sakkas2022seq2parse}.

  We compare our method against two separate baselines, Seq2Parse and Break-It-Fix-It (BIFI)~\cite{yasunaga2021break} on a single test set. This dataset~\cite{wong2019syntax} consists of 20k naturally-occurring pairs of Python errors and their corresponding human fixes from StackOverflow and is used to compare the precision of each method at blind recovery of the ground truth repair across varying edit distances, snippet lengths and latency cutoffs. We preprocess all source code by filtering for broken-fixed snippet pairs shorter than 80 tokens and fewer than five Levenshtein edits apart, whose broken and fixed form is accepted and rejected, respectively, by the Python 3.8.11 parser. We then balance the dataset by sampling an equal number of repairs from each length and Levenshtein edit distance.

%  In our synthetic experiments, we apply the pretrained BIFI breaker to synthetically corrupt Python snippets from the BIFI good code test set, using the clean source as the ground truth repair, and filter broken-fixed snippet pairs by the same criteria.

  The Seq2Parse and BIFI experiments were conducted on a single Nvidia V100 GPU with 32 GB of RAM. For Seq2Parse, we use the default pretrained model provided in commit \texttt{7ae0681}~\footnote{https://github.com/gsakkas/seq2parse/tree/7ae0681f1139cb873868727f035c1b7a369c3eb9}. Since it was unclear how to extract multiple repairs from their model, we only take a single repair prediction. For BIFI, we use the Round 2 breaker and fixer from commit \texttt{ee2a68c}\footnote{https://github.com/michiyasunaga/BIFI/tree/ee2a68cff8dbe88d2a2b2b5feabc7311d5f8338b}, the highest-performing model reported by the authors, with a variable-width beam search to control the number of predictions, and let the fixer model predict the top-k repairs, for $k=\{1, 5, 10, 20\times10^5\}$.

  The language intersection experiments were conducted on 40 Intel Skylake cores running at 2.4 GHz, with 150 GB of RAM, running bytecode compiled for JVM 17.0.2. To train our scoring function, we use an order-5 Markov chain trained on 55 million BIFI tokens. Training takes roughly 10 minutes, after which re-ranking is nearly instantaneous. Sequences are scored using NLL with Laplace smoothing and our evaluation measures the Precision@\{1, 5, 10, All\} for samples at varying latency cutoffs. We apply a 30-second latency cutoff for our sampler.

%  Pairwise naturally-occurring errors and human fixes are the most authentic source of real-world syntax repairs, but can be difficult to obtain due to the paucity of parallel syntax error corpi. In the absence of natural syntax repairs, one viable alternative is to collect a dataset of syntactically valid code, and synthetically corrupt it. The original source code becomes the ground truth repair for the synthetically generated typo, and the target for evaluating the precision of our repair procedure.

  %  We compute the Levenshtein alignment on the StackOverflow dataset, then approximate the conditional probability of each edit given the local context. During evaluation, we sample a corruption from the learned typo distribution, and measure the precision of our model at recovering the originally valid lexical sequence.

%  Suppose we have a dataset of Levenshtein edits and their local context. For simplicity, we shall assume a trigram language model, i.e., $P(\sigma_i' \mid \sigma_{i-1}, \sigma_i, \sigma_{i+1})$, however the approach can be generalized to higher-order Markov models. Given a string $\sigma$, we can sample error trajectories $q^1(\sigma), q^2(\sigma), \ldots, q^n(\sigma)$ by defining $q(\sigma)$ to sample a single edit from the set of all relevant edit actions $Q(\sigma)$, then recursively applying $q$ to the resulting string. More formally,
%
%  \begin{enumerate}
%    \item Given a string $\sigma$, compute $Q(\sigma)$, the set of all relevant edit actions for all possible edit locations by unioning the set of all possible edits at each location, i.e., $Q(\sigma) = \bigcup_{i=1}^{|\sigma| - 1} \big\{\sigma_i' \mid  0 < P(\sigma_i \mid \sigma_{i-1}, \sigma_i, \sigma_{i+1})\big\}$.
%    \item Renormalize the probabilities of each edit $P(q \mid \sigma)$ by $\sum_{q \in Q(\sigma)} P(q)$. This ensures the probability of sampling a particular edit is proportional to its relative probability under the language model and sums to 1.
%    \item Sample an edit $q(\sigma) \sim Q(\sigma)$, then repeat for $n$ steps where $n$ is sampled from a geometric distribution with mean $\mu$ matching the average edit distance of the dataset (this assumes the edit distance is independent of the edits).
%  \end{enumerate}
%
%  For example, suppose we have the following patch in our initial dataset:\\
%
%  \texttt{BOS \hlred{def} NAME ( NAME ) : NEWLINE \hlred{INDENT} return \hlorange{NAME} NEWLINE \hlgreen{INDENT} NEWLINE EOS}\\
%
%  From this patch, the following contextual typo probabilities will be incremented:
%
%  \begin{align*}
%    P(\texttt{BOS \hlred{def} NAME}) &\mathrel{+}= 1 &P(\texttt{NEWLINE \hlred{INDENT} return}) &\mathrel{+}= 1\\
%    P(\texttt{return \hlorange{NAME} NEWLINE}) &\mathrel{+}= 1 & P(\texttt{NEWLINE \hlgreen{INDENT} NEWLINE}) &\mathrel{+}= 1
%  \end{align*}
%
%  Later, these contextual probabilities will allow us to sample a synthetic corruption matching the distribution of typos in our dataset. We then measure the precision at recovering the originally valid string.


  \section{Evaluation}

  We call our method Tidyparse and consider the following research questions:

  \begin{itemize}
    \item \textbf{RQ 1}: What statistical properties do natural repairs exhibit? (e.g., length, edit distance)
    \item \textbf{RQ 2}: How performant is Tidyparse at fixing syntax errors? (i.e., vs. Seq2Parse and BIFI)
    \item \textbf{RQ 3}: Which design choices are most significant? (e.g., search vs. sampling, parallelism)
  \end{itemize}

  We address \textbf{RQ 1} in \S~\ref{sec:rq1} by analyzing the distribution of natural Python repair lengths and distances, \textbf{RQ 2} in \S~\ref{sec:rq2} by comparing Tidyparse against two existing syntax repair baselines, and \textbf{RQ 3} in \S~\ref{sec:rq3} by ablating various design choices and evaluating the impact on repair precision.

  \subsection{Dataset}\label{sec:rq1}

  In the following experiments, we use a dataset of Python snippets consisting of 20,500 pairwise-aligned human errors and fixes from StackOverflow~\cite{wong2019syntax}. We preprocess the dataset to lexicalize all code snippets, then filter by length and distance shorter than 80 lexical tokens and under five edits, i.e., where pairwise Levenshtein distance is under five lexical edits ($|\Sigma| = 50, |\err{\sigma}| < 80, \Delta(\err{\sigma}, \sigma') < 5$). We depict the length, edit distance, normalized edit locations and stability profile in Fig.~\ref{fig:patch_stats}.

  \begin{figure}[h!]
    \begin{tikzpicture}[scale=0.57]
      \begin{axis}[
      ybar,
      bar width=5pt,
      xlabel={$|\err\sigma|$},
      ylabel={Frequency},
      title={Cumulative length distribution},
      axis x line*=bottom,
      axis y line*=left,
      ymin=0,
      ymax=65,
      xtick=data,
      xticklabels={,<20,,<40,,<60,,<80,,<100},
      ymajorgrids=true,
      grid style=dashed,
      width=0.45\textwidth,
      height=0.3\textwidth
      ]

      \addplot[fill=black!30] table {
        X Y
        1 7.60
        2 14.52
        3 22.01
        4 30.54
        5 37.82
        6 44.30
        7 49.68
        8 55.21
        9 59.75
        10 63.59
      };
      \draw[red, dashed] (axis cs:8.5,0) -- (axis cs:8.5,65);
      \end{axis}
    \end{tikzpicture}
    \begin{tikzpicture}[scale=0.57]
      \begin{axis}[
        ybar,
        bar width=5pt,
        title={Human repair distance},
        xlabel={$\Delta(\err\sigma, \sigma')$},
        ylabel={Frequency},
        axis x line*=bottom,
        axis y line*=left,
        xtick=data,
        ymajorgrids=true,
        grid style=dashed,
        xticklabels={,\leq 2,,\leq 4,,\leq 6,,\leq 8,,\leq 10},
        ytick={0, 20, 40, 60, 80, 100},
        ymin=0,
        width=0.45\textwidth,
        height=0.3\textwidth
      ]
        \addplot[fill=black!30] table {
          X Y
          1  31.48
          2  47.52
          3  54.89
          4  60.44
          5  63.88
          6  66.38
          7  68.02
          8  70.04
          9  71.49
          10 72.22
        };
      \draw[red, dashed] (axis cs:4.5,0) -- (axis cs:4.5,80);
      \end{axis}
    \end{tikzpicture}
    \begin{tikzpicture}[scale=0.57]
      \begin{axis}[
      ybar,
      bar width=5pt,
      xlabel={Beginning $\longleftrightarrow$ End},
      ylabel={Frequency},
      title={Normalized edit locations},
      axis x line*=bottom,
      axis y line*=left,
      ymin=0,
      ymax=35,
      xtick=data,
      xticklabels={,20\%,,40\%,,60\%,,80\%,,100\%},
      ymajorgrids=true,
      grid style=dashed,
      width=0.45\textwidth,
      height=0.3\textwidth
      ]

      \addplot[fill=black!30] table {
        X Y
        10 11.6539
        20 5.7252
        30 6.2087
        40 5.9542
        50 5.5980
        60 7.9389
        70 7.0738
        80 6.9466
        90 12.4173
        100 30.4835
      };
      \end{axis}
    \end{tikzpicture}
%    \begin{tikzpicture}
%      \begin{axis}[
%        ybar,
%        bar width=5pt,
%        title={Intra-patch edit distance},
%        xlabel={Caret distance},
%        ylabel={Frequency},
%        axis x line*=bottom,
%        axis y line*=left,
%        xtick=data,
%        ymajorgrids=true,
%        grid style=dashed,
%        xticklabels={1,2,3,4,5,6,7,8,9,10+},
%        width=0.45\textwidth,
%        height=0.3\textwidth
%      ]
%
%        \addplot table {
%          X Y
%          1 40.66
%          2 15.00
%          3 5.80
%          4 4.86
%          5 4.26
%          6 2.98
%          7 2.05
%          8 2.73
%          9 1.62
%          10 13.64
%        };
%      \end{axis}
%    \end{tikzpicture}
    \begin{tikzpicture}[scale=0.57]
      \begin{axis}[
        legend cell align={left},
        legend style={fill opacity=1, draw opacity=1, text opacity=1, draw=lightgray204, legend columns=-1, legend pos=south east},
        xlabel={$|\err\sigma|$},
        ylabel={Stable region},
        title={Stability profile},
        ybar,
        axis lines*=left,
        xtick={0, 10, 20, 30, 40, 50, 60, 70},
        ytick={0, 0.2, 0.4, 0.6, 0.8, 1.0},
        xticklabels={, {[}10{,}20{)}, , {[}30{,}40{)}, , {[}50{,}60{)}, , {[}70{,}80{)}},
        yticklabels={0, 0.2, 0.4, 0.6, 0.8, 1.0},
        x tick label style={font=\scriptsize},
        ymax=1.0,
        ymin=0.0,
        bar width=3pt,
        grid style=dashed,
        ymajorgrids=true,
        width=0.45\textwidth,
        height=0.3\textwidth
      ]
      \addlegendimage{empty legend}
      \addlegendentry{$\Delta(\err\sigma, \sigma')=$}
      \addlegendimage{ybar,ybar legend,draw=none,green,fill=green!50}
      \addlegendentry{1,}
      \addlegendimage{ybar,ybar legend,draw=none,blue,fill=blue!50}
      \addlegendentry{2,}
      \addlegendimage{ybar,ybar legend,draw=none,orange,fill=orange!50}
      \addlegendentry{3}
      \addplot[green, fill=green!50] coordinates {(0, 0.80) (10, 0.91) (20, 0.96) (30, 0.97) (40, 0.99) (50, 0.99) (60, 0.99) (70, 0.99)};
      \addplot[blue, fill=blue!50] coordinates {(0, 0.35) (10, 0.59) (20, 0.69) (30, 0.73) (40, 0.79) (50, 0.82) (60, 0.84) (70, 0.86)};
      \addplot[orange, fill=orange!50] coordinates {(0, 0.23) (10, 0.45) (20, 0.58) (30, 0.66) (40, 0.70) (50, 0.77) (60, 0.78) (70, 0.86)};
      \end{axis}
    \end{tikzpicture}
    \caption{Repair statistics across the StackOverflow dataset, of which Tidyparse can handle about half in under $\sim$30s and $\sim$150 GB. Larger repairs and edit distances are possible, albeit requiring additional time and memory.}\label{fig:patch_stats}
  \end{figure}

  For the stability profile, we enumerate repairs for each syntax error and estimate the average fraction of all edit locations that were never altered by any repair in the $L\big(\err\sigma, \Delta(\err\sigma, \sigma')\big)$-ball. For example, on average roughly half of the string is stable for 3-edit syntax repairs in the $[10-20)$ token range, whereas 1-edit repairs of the same length could modify only $\sim 10\%$ of all locations. For a fixed edit distance, we observe an overall decrease in the number of degrees of caret freedom with increasing length, which intuitively makes sense, as the repairs are more heavily constrained by the surrounding context and their locations grow more concentrated relative to the entire string.

%  We also use a second training set of valid Python snippets from Yasunaga et al.~\cite{yasunaga2021break} on which BIFI was trained and we also use to train our scoring function. We lexicalize the dataset in the same way as StackOverflow, annotate each snippet with a \texttt{BOS} and \texttt{EOS} token, and train a 5-gram Markov model on a small subset. This model is used to score synthetic repairs generated by Tidyparse on repair instances in the StackOverflow dataset. No filtering criteria are applied to this dataset.

%  then synthetically corrupt them by introducing synthetic syntax errors and measure the Precision@k of our repair procedure at recovering the original, uncorrupted snippet. This dataset is effectively unlimited as we can generate as many synthetic errors across any range of edit distances and snippet lengths, however these repairs may be less representative of human syntax errors.

  \clearpage\subsection{StackOverflow evaluation}\label{sec:rq2}

  For our first experiment, we measure the precision of our repair procedure at various lengths and Levenshtein distances. We rebalance the StackOverflow dataset across each length interval and edit distance, sample uniformly from each category and compare Precision@1 of our method against Seq2Parse, vanilla BIFI and BIFI with a beam size and precision at 20k distinct samples.

  \begin{figure}[h!]
    \resizebox{.24\textwidth}{!}{\begin{tikzpicture}
  \begin{axis}[
    xlabel={$|\sigma|$},
    ylabel={Precision@1},
    title={Tidyparse Repair Precision@1},
    ybar,
    axis lines*=left,
    xtick={0, 10, 20, 30, 40, 50, 60, 70},
    ytick={0, 0.1, 0.2, 0.3, 0.4, 0.5, 0.6, 0.7, 0.8, 0.9, 1.0},
    xticklabels={{(}0{,}10{)}, {[}10{,}20{)}, {[}20{,}30{)}, {[}30{,}40{)}, {[}40{,}50{)}, {[}50{,}60{)}, {[}60{,}70{)}, {[}70{,}80{)}},
    x tick label style={font=\scriptsize},
    ymax=1.0,
    ymin=0.0,
    bar width=4pt,
  ]

  \addplot[green, fill=green] coordinates { (0, 1.0) (10, 1.0) (20, 0.975) (30, 0.98) (40, 1.0) (50, 1.0) (60, 0.95) (70, 0.9) };
  \addplot[blue, fill=blue] coordinates { (0, 0.51) (10, 0.36) (20, 0.24) (30, 0.26) (40, 0.24) (50, 0.23) (60, 0.12) (70, 0.1) };
  \addplot[orange, fill=orange] coordinates { (0, 0.38) (10, 0.26) (20, 0.37) (30, 0.25) (40, 0.22) (50, 0.16) (60, 0.14) (70, 0.08) };

%  \legend{Δ=1,Δ=2,Δ=3}
  \end{axis}
\end{tikzpicture}}
    \resizebox{.24\textwidth}{!}{\begin{tikzpicture}
  \begin{axis}[
  xlabel={$|\sigma|$},
  ylabel={Precision@20k},
  title={\textbf{BIFI Repair Precision@20k}},
  ybar,
  axis lines*=left,
  xtick={0, 10, 20, 30, 40, 50, 60, 70},
  ytick={0, 0.1, 0.2, 0.3, 0.4, 0.5, 0.6, 0.7, 0.8, 0.9, 1.0},
  xticklabels={{(}0{,}10{)}, {[}10{,}20{)}, {[}20{,}30{)}, {[}30{,}40{)}, {[}40{,}50{)}, {[}50{,}60{)}, {[}60{,}70{)}, {[}70{,}80{)}},
  x tick label style={font=\scriptsize},
  ymax=1.0,
  ymin=0.0,
  bar width=4pt,
  ]

  \addplot[green, fill=green!50] coordinates {(0, 0.65) (10, 0.67) (20, 0.71) (30, 0.63) (40, 0.6) (50, 0.62) (60, 0.59) (70, 0.64)};
  \addplot[blue, fill=blue!50] coordinates {(0, 0.52) (10, 0.41) (20, 0.35) (30, 0.31) (40, 0.27) (50, 0.27) (60, 0.21) (70, 0.22)};
  \addplot[orange, fill=orange!50] coordinates {(0, 0.25) (10, 0.08) (20, 0.08) (30, 0.17) (40, 0.11) (50, 0.17) (60, 0.08) (70, 0.08)};

  \end{axis}
\end{tikzpicture}
}
    \resizebox{.24\textwidth}{!}{\begin{tikzpicture}
  \begin{axis}[
    xlabel={$|\sigma|$},
    ylabel={Precision@1},
    title={Seq2Parse Repair Precision},
    ybar,
    axis lines*=left,
    xtick={0, 10, 20, 30, 40, 50, 60, 70},
    ytick={0, 0.1, 0.2, 0.3, 0.4},
    ymax=0.4,
    bar width=4pt,
  ]

    \addplot[green, fill=green] coordinates {(0, 0.368078) (10, 0.416931) (20, 0.393548) (30, 0.387001) (40, 0.3125) (50, 0.289926) (60, 0.258278) (70, 0.198157)};
    \addplot[blue, fill=blue] coordinates {(0, 0.127450) (10, 0.131627) (20, 0.133111) (30, 0.107287) (40, 0.099537) (50, 0.104729) (60, 0.100418) (70, 0.123762)};
    \addplot[orange, fill=orange] coordinates {(0, 0.037735) (10, 0.072) (20, 0.086093) (30, 0.08125) (40, 0.088235) (50, 0.022727) (60, 0.054945) (70, 0.061538)};

    \legend{Δ=1,Δ=2,Δ=3}
  \end{axis}
\end{tikzpicture}}
    \resizebox{.24\textwidth}{!}{\begin{tikzpicture}
  \begin{axis}[
    xlabel={$|\sigma|$},
    ylabel={Precision@1},
    title={BIFI Repair Precision@1},
    ybar,
    axis lines*=left,
    xtick={0, 10, 20, 30, 40, 50, 60, 70},
    ytick={0, 0.1, 0.2, 0.3, 0.4, 0.5, 0.6, 0.7, 0.8, 0.9, 1.0},
    xticklabels={{(}0{,}10{)}, {[}10{,}20{)}, {[}20{,}30{)}, {[}30{,}40{)}, {[}40{,}50{)}, {[}50{,}60{)}, {[}60{,}70{)}, {[}70{,}80{)}},
    x tick label style={font=\scriptsize},
    ymax=1.0,
    ymin=0.0,
    bar width=4pt,
  ]

    \addplot[green, fill=green] coordinates {(0, 0.196013) (10, 0.326401) (20, 0.318538) (30, 0.272843) (40, 0.213894) (50, 0.206651) (60, 0.247525) (70, 0.179245)};
    \addplot[blue, fill=blue] coordinates {(0, 0.174603) (10, 0.176651) (20, 0.209573) (30, 0.19195) (40, 0.18851) (50, 0.176166) (60, 0.110787) (70, 0.106383)};
    \addplot[orange, fill=orange] coordinates {(0, 0.015873) (10, 0.021858) (20, 0.030435) (30, 0.02439) (40, 0.032922) (50, 0.045) (60, 0.027397) (70, 0.017094)};

    \legend{Δ=1,Δ=2,Δ=3}
  \end{axis}
\end{tikzpicture}}
    \caption{Tidyparse, Seq2Parse and BIFI repair precision at various lengths and Levenshtein distances.}\label{fig:len_dist_prec}
  \end{figure}

  As we can see, Tidyparse has a highly competitive top-1 precision versus Seq2Parse and BIFI across all lengths and edit distances, and attains a significant advantage in the few-edit regime. The Precision@1 of our method is even competitive with BIFI's Precision@20k, whereas our Precision@All is Pareto-dominant across all lengths and edit distances, while requiring only a fraction of the data and compute. We report the raw data from these experiments in Appendix~\ref{sec:raw_prec_data}.

  Next, we measure the precision at various ranking cutoffs and wall-clock timeouts. Our method attains the same precision as Seq2Parse and BIFI for 1-edit repairs at comparable latency, however Tidyparse takes longer to attain the same precision for 2- and 3-edit repairs. BIFI and Seq2Parse both have subsecond single-shot latency but are neural models trained on a much larger dataset.

  \begin{figure}[h!]
%    \resizebox{.19\textwidth}{!}{% This file was created with tikzplotlib v0.10.1.
\begin{tikzpicture}

\definecolor{darkgray176}{RGB}{176,176,176}
\definecolor{darkviolet1270255}{RGB}{127,0,255}
\definecolor{deepskyblue3176236}{RGB}{3,176,236}
\definecolor{dodgerblue45123246}{RGB}{45,123,246}
\definecolor{royalblue8762253}{RGB}{87,62,253}

\begin{axis}[
tick align=outside,
tick pos=left,
title={\(\displaystyle \Delta\in[1,3]\) Repair Precision},
legend style={fill opacity=0.8, draw opacity=1, text opacity=1, legend columns=1, legend pos=north west},
x grid style={darkgray176},
xlabel={Seconds},
xmin=-0.5925, xmax=8.5925,
xtick style={color=black},
xtick={0,1,2,3,4,5,6,7,8},
xticklabels={1,2,3,4,5,6,7,8,9},
y grid style={darkgray176},
ylabel={Precision@k},
ymin=0, ymax=0.6993,
ytick style={color=black}
]
\draw[draw=none,fill=darkviolet1270255] (axis cs:-0.175,0) rectangle (axis cs:0.175,0.149);
\addlegendimage{ybar,ybar legend,draw=none,fill=darkviolet1270255}
\addlegendentry{P@All}

\draw[draw=none,fill=darkviolet1270255] (axis cs:0.825,0) rectangle (axis cs:1.175,0.303);
\draw[draw=none,fill=darkviolet1270255] (axis cs:1.825,0) rectangle (axis cs:2.175,0.412);
\draw[draw=none,fill=darkviolet1270255] (axis cs:2.825,0) rectangle (axis cs:3.175,0.501);
\draw[draw=none,fill=darkviolet1270255] (axis cs:3.825,0) rectangle (axis cs:4.175,0.569);
\draw[draw=none,fill=darkviolet1270255] (axis cs:4.825,0) rectangle (axis cs:5.175,0.605);
\draw[draw=none,fill=darkviolet1270255] (axis cs:5.825,0) rectangle (axis cs:6.175,0.63);
\draw[draw=none,fill=darkviolet1270255] (axis cs:6.825,0) rectangle (axis cs:7.175,0.645);
\draw[draw=none,fill=darkviolet1270255] (axis cs:7.825,0) rectangle (axis cs:8.175,0.666);
\draw[draw=none,fill=royalblue8762253] (axis cs:-0.175,0) rectangle (axis cs:0.175,0.128);
\addlegendimage{ybar,ybar legend,draw=none,fill=royalblue8762253}
\addlegendentry{P@10}

\draw[draw=none,fill=royalblue8762253] (axis cs:0.825,0) rectangle (axis cs:1.175,0.239);
\draw[draw=none,fill=royalblue8762253] (axis cs:1.825,0) rectangle (axis cs:2.175,0.323);
\draw[draw=none,fill=royalblue8762253] (axis cs:2.825,0) rectangle (axis cs:3.175,0.374);
\draw[draw=none,fill=royalblue8762253] (axis cs:3.825,0) rectangle (axis cs:4.175,0.421);
\draw[draw=none,fill=royalblue8762253] (axis cs:4.825,0) rectangle (axis cs:5.175,0.446);
\draw[draw=none,fill=royalblue8762253] (axis cs:5.825,0) rectangle (axis cs:6.175,0.459);
\draw[draw=none,fill=royalblue8762253] (axis cs:6.825,0) rectangle (axis cs:7.175,0.469);
\draw[draw=none,fill=royalblue8762253] (axis cs:7.825,0) rectangle (axis cs:8.175,0.477);
\draw[draw=none,fill=dodgerblue45123246] (axis cs:-0.175,0) rectangle (axis cs:0.175,0.11);
\addlegendimage{ybar,ybar legend,draw=none,fill=dodgerblue45123246}
\addlegendentry{P@5}

\draw[draw=none,fill=dodgerblue45123246] (axis cs:0.825,0) rectangle (axis cs:1.175,0.206);
\draw[draw=none,fill=dodgerblue45123246] (axis cs:1.825,0) rectangle (axis cs:2.175,0.282);
\draw[draw=none,fill=dodgerblue45123246] (axis cs:2.825,0) rectangle (axis cs:3.175,0.325);
\draw[draw=none,fill=dodgerblue45123246] (axis cs:3.825,0) rectangle (axis cs:4.175,0.369);
\draw[draw=none,fill=dodgerblue45123246] (axis cs:4.825,0) rectangle (axis cs:5.175,0.392);
\draw[draw=none,fill=dodgerblue45123246] (axis cs:5.825,0) rectangle (axis cs:6.175,0.405);
\draw[draw=none,fill=dodgerblue45123246] (axis cs:6.825,0) rectangle (axis cs:7.175,0.414);
\draw[draw=none,fill=dodgerblue45123246] (axis cs:7.825,0) rectangle (axis cs:8.175,0.422);
\draw[draw=none,fill=deepskyblue3176236] (axis cs:-0.175,0) rectangle (axis cs:0.175,0.092);
\addlegendimage{ybar,ybar legend,draw=none,fill=deepskyblue3176236}
\addlegendentry{P@1}

\draw[draw=none,fill=deepskyblue3176236] (axis cs:0.825,0) rectangle (axis cs:1.175,0.18);
\draw[draw=none,fill=deepskyblue3176236] (axis cs:1.825,0) rectangle (axis cs:2.175,0.244);
\draw[draw=none,fill=deepskyblue3176236] (axis cs:2.825,0) rectangle (axis cs:3.175,0.285);
\draw[draw=none,fill=deepskyblue3176236] (axis cs:3.825,0) rectangle (axis cs:4.175,0.324);
\draw[draw=none,fill=deepskyblue3176236] (axis cs:4.825,0) rectangle (axis cs:5.175,0.343);
\draw[draw=none,fill=deepskyblue3176236] (axis cs:5.825,0) rectangle (axis cs:6.175,0.356);
\draw[draw=none,fill=deepskyblue3176236] (axis cs:6.825,0) rectangle (axis cs:7.175,0.365);
\draw[draw=none,fill=deepskyblue3176236] (axis cs:7.825,0) rectangle (axis cs:8.175,0.37);
\addplot [red, thick] coordinates {(-0.8925,0.29) (14.8925,0.29)};
\end{axis}

\end{tikzpicture}
}
    \resizebox{.24\textwidth}{!}{% This file was created with tikzplotlib v0.10.1.
\begin{tikzpicture}
\begin{axis}[
legend cell align={left},
legend style={fill opacity=0.8, draw opacity=1, text opacity=1, draw=lightgray204, legend columns=-1, legend pos=north west},
tick align=outside,
tick pos=left,
axis lines*=left,
title={\(\displaystyle \Delta=1\) Repair Precision},
x grid style={darkgray176},
xlabel={Seconds},
xmin=-0.5925, xmax=8.5925,
xtick style={color=black},
xtick={0,1,2,3,4,5,6,7,8},
xticklabels={1,2,3,4,5,6,7,8,9},
y grid style={darkgray176},
ylabel={Precision@k},
ymin=0, ymax=1.0311,
ytick style={color=black}
]
\draw[draw=none,fill=darkviolet1270255] (axis cs:-0.175,0) rectangle (axis cs:0.175,0.546);
\addlegendimage{empty legend}
\addlegendentry{P@}
\addlegendimage{ybar,ybar legend,draw=none,fill=darkviolet1270255}
\addlegendentry{All}
\addlegendimage{ybar,ybar legend,draw=none,fill=royalblue8762253}
\addlegendentry{10}
\addlegendimage{ybar,ybar legend,draw=none,fill=dodgerblue45123246}
\addlegendentry{5}
\addlegendimage{ybar,ybar legend,draw=none,fill=deepskyblue3176236}
\addlegendentry{1}

\draw[draw=none,fill=darkviolet1270255] (axis cs:0.825,0) rectangle (axis cs:1.175,0.752);
\draw[draw=none,fill=darkviolet1270255] (axis cs:1.825,0) rectangle (axis cs:2.175,0.88);
\draw[draw=none,fill=darkviolet1270255] (axis cs:2.825,0) rectangle (axis cs:3.175,0.908);
\draw[draw=none,fill=darkviolet1270255] (axis cs:3.825,0) rectangle (axis cs:4.175,0.933);
\draw[draw=none,fill=darkviolet1270255] (axis cs:4.825,0) rectangle (axis cs:5.175,0.954);
\draw[draw=none,fill=darkviolet1270255] (axis cs:5.825,0) rectangle (axis cs:6.175,0.966);
\draw[draw=none,fill=darkviolet1270255] (axis cs:6.825,0) rectangle (axis cs:7.175,0.975);
\draw[draw=none,fill=darkviolet1270255] (axis cs:7.825,0) rectangle (axis cs:8.175,0.982);
\draw[draw=none,fill=royalblue8762253] (axis cs:-0.175,0) rectangle (axis cs:0.175,0.546);

\draw[draw=none,fill=royalblue8762253] (axis cs:0.825,0) rectangle (axis cs:1.175,0.752);
\draw[draw=none,fill=royalblue8762253] (axis cs:1.825,0) rectangle (axis cs:2.175,0.88);
\draw[draw=none,fill=royalblue8762253] (axis cs:2.825,0) rectangle (axis cs:3.175,0.908);
\draw[draw=none,fill=royalblue8762253] (axis cs:3.825,0) rectangle (axis cs:4.175,0.933);
\draw[draw=none,fill=royalblue8762253] (axis cs:4.825,0) rectangle (axis cs:5.175,0.954);
\draw[draw=none,fill=royalblue8762253] (axis cs:5.825,0) rectangle (axis cs:6.175,0.966);
\draw[draw=none,fill=royalblue8762253] (axis cs:6.825,0) rectangle (axis cs:7.175,0.975);
\draw[draw=none,fill=royalblue8762253] (axis cs:7.825,0) rectangle (axis cs:8.175,0.982);
\draw[draw=none,fill=dodgerblue45123246] (axis cs:-0.175,0) rectangle (axis cs:0.175,0.546);

\draw[draw=none,fill=dodgerblue45123246] (axis cs:0.825,0) rectangle (axis cs:1.175,0.752);
\draw[draw=none,fill=dodgerblue45123246] (axis cs:1.825,0) rectangle (axis cs:2.175,0.88);
\draw[draw=none,fill=dodgerblue45123246] (axis cs:2.825,0) rectangle (axis cs:3.175,0.908);
\draw[draw=none,fill=dodgerblue45123246] (axis cs:3.825,0) rectangle (axis cs:4.175,0.933);
\draw[draw=none,fill=dodgerblue45123246] (axis cs:4.825,0) rectangle (axis cs:5.175,0.954);
\draw[draw=none,fill=dodgerblue45123246] (axis cs:5.825,0) rectangle (axis cs:6.175,0.966);
\draw[draw=none,fill=dodgerblue45123246] (axis cs:6.825,0) rectangle (axis cs:7.175,0.975);
\draw[draw=none,fill=dodgerblue45123246] (axis cs:7.825,0) rectangle (axis cs:8.175,0.982);
\draw[draw=none,fill=deepskyblue3176236] (axis cs:-0.175,0) rectangle (axis cs:0.175,0.546);

\draw[draw=none,fill=deepskyblue3176236] (axis cs:0.825,0) rectangle (axis cs:1.175,0.752);
\draw[draw=none,fill=deepskyblue3176236] (axis cs:1.825,0) rectangle (axis cs:2.175,0.88);
\draw[draw=none,fill=deepskyblue3176236] (axis cs:2.825,0) rectangle (axis cs:3.175,0.908);
\draw[draw=none,fill=deepskyblue3176236] (axis cs:3.825,0) rectangle (axis cs:4.175,0.933);
\draw[draw=none,fill=deepskyblue3176236] (axis cs:4.825,0) rectangle (axis cs:5.175,0.954);
\draw[draw=none,fill=deepskyblue3176236] (axis cs:5.825,0) rectangle (axis cs:6.175,0.966);
\draw[draw=none,fill=deepskyblue3176236] (axis cs:6.825,0) rectangle (axis cs:7.175,0.975);
\draw[draw=none,fill=deepskyblue3176236] (axis cs:7.825,0) rectangle (axis cs:8.175,0.982);
\addplot [red, thick] coordinates {(-0.8925,0.39) (14.8925,0.39)};
\addplot [orange, thick] coordinates {(-0.8925,0.24) (14.8925,0.24)};
\end{axis}

\end{tikzpicture}
}
    \resizebox{.24\textwidth}{!}{% This file was created with tikzplotlib v0.10.1.
\begin{tikzpicture}
\begin{axis}[
tick align=outside,
tick pos=left,
axis lines=left,
title={\(\displaystyle \Delta=2\) Repair Precision},
x grid style={darkgray176},
xlabel={Seconds},
xmin=-0.5925, xmax=8.5925,
xtick style={color=black},
xtick={0,1,2,3,4,5,6,7,8},
xticklabels={10,20,30,40,50,60,70,80,90},
y grid style={darkgray176},
ylabel={Precision@k},
ymin=0, ymax=0.75285,
ytick style={color=black}
]

\draw[draw=none,fill=darkviolet1270255] (axis cs:-0.175,0) rectangle (axis cs:0.175,0.182);

\draw[draw=none,fill=darkviolet1270255] (axis cs:0.825,0) rectangle (axis cs:1.175,0.326);
\draw[draw=none,fill=darkviolet1270255] (axis cs:1.825,0) rectangle (axis cs:2.175,0.455);
\draw[draw=none,fill=darkviolet1270255] (axis cs:2.825,0) rectangle (axis cs:3.175,0.554);
\draw[draw=none,fill=darkviolet1270255] (axis cs:3.825,0) rectangle (axis cs:4.175,0.603);
\draw[draw=none,fill=darkviolet1270255] (axis cs:4.825,0) rectangle (axis cs:5.175,0.649);
\draw[draw=none,fill=darkviolet1270255] (axis cs:5.825,0) rectangle (axis cs:6.175,0.677);
\draw[draw=none,fill=darkviolet1270255] (axis cs:6.825,0) rectangle (axis cs:7.175,0.702);
\draw[draw=none,fill=darkviolet1270255] (axis cs:7.825,0) rectangle (axis cs:8.175,0.717);
\draw[draw=none,fill=royalblue8762253] (axis cs:-0.175,0) rectangle (axis cs:0.175,0.102);

\draw[draw=none,fill=royalblue8762253] (axis cs:0.825,0) rectangle (axis cs:1.175,0.178);
\draw[draw=none,fill=royalblue8762253] (axis cs:1.825,0) rectangle (axis cs:2.175,0.243);
\draw[draw=none,fill=royalblue8762253] (axis cs:2.825,0) rectangle (axis cs:3.175,0.289);
\draw[draw=none,fill=royalblue8762253] (axis cs:3.825,0) rectangle (axis cs:4.175,0.317);
\draw[draw=none,fill=royalblue8762253] (axis cs:4.825,0) rectangle (axis cs:5.175,0.335);
\draw[draw=none,fill=royalblue8762253] (axis cs:5.825,0) rectangle (axis cs:6.175,0.348);
\draw[draw=none,fill=royalblue8762253] (axis cs:6.825,0) rectangle (axis cs:7.175,0.357);
\draw[draw=none,fill=royalblue8762253] (axis cs:7.825,0) rectangle (axis cs:8.175,0.366);
\draw[draw=none,fill=dodgerblue45123246] (axis cs:-0.175,0) rectangle (axis cs:0.175,0.102);

\draw[draw=none,fill=dodgerblue45123246] (axis cs:0.825,0) rectangle (axis cs:1.175,0.178);
\draw[draw=none,fill=dodgerblue45123246] (axis cs:1.825,0) rectangle (axis cs:2.175,0.237);
\draw[draw=none,fill=dodgerblue45123246] (axis cs:2.825,0) rectangle (axis cs:3.175,0.283);
\draw[draw=none,fill=dodgerblue45123246] (axis cs:3.825,0) rectangle (axis cs:4.175,0.311);
\draw[draw=none,fill=dodgerblue45123246] (axis cs:4.825,0) rectangle (axis cs:5.175,0.329);
\draw[draw=none,fill=dodgerblue45123246] (axis cs:5.825,0) rectangle (axis cs:6.175,0.342);
\draw[draw=none,fill=dodgerblue45123246] (axis cs:6.825,0) rectangle (axis cs:7.175,0.351);
\draw[draw=none,fill=dodgerblue45123246] (axis cs:7.825,0) rectangle (axis cs:8.175,0.36);
\draw[draw=none,fill=deepskyblue3176236] (axis cs:-0.175,0) rectangle (axis cs:0.175,0.092);

\draw[draw=none,fill=deepskyblue3176236] (axis cs:0.825,0) rectangle (axis cs:1.175,0.16);
\draw[draw=none,fill=deepskyblue3176236] (axis cs:1.825,0) rectangle (axis cs:2.175,0.218);
\draw[draw=none,fill=deepskyblue3176236] (axis cs:2.825,0) rectangle (axis cs:3.175,0.258);
\draw[draw=none,fill=deepskyblue3176236] (axis cs:3.825,0) rectangle (axis cs:4.175,0.283);
\draw[draw=none,fill=deepskyblue3176236] (axis cs:4.825,0) rectangle (axis cs:5.175,0.302);
\draw[draw=none,fill=deepskyblue3176236] (axis cs:5.825,0) rectangle (axis cs:6.175,0.311);
\draw[draw=none,fill=deepskyblue3176236] (axis cs:6.825,0) rectangle (axis cs:7.175,0.32);
\draw[draw=none,fill=deepskyblue3176236] (axis cs:7.825,0) rectangle (axis cs:8.175,0.329);
\addplot [red, thick] coordinates {(-0.8925,0.15) (14.8925,0.15)};
\addplot [orange, thick] coordinates {(-0.8925,0.10) (14.8925,0.10)};
\end{axis}

\end{tikzpicture}
}
    \resizebox{.24\textwidth}{!}{% This file was created with tikzplotlib v0.10.1.
\begin{tikzpicture}

\definecolor{darkgray176}{RGB}{176,176,176}
\definecolor{darkviolet1270255}{RGB}{127,0,255}
\definecolor{deepskyblue3176236}{RGB}{3,176,236}
\definecolor{dodgerblue45123246}{RGB}{45,123,246}
\definecolor{royalblue8762253}{RGB}{87,62,253}

\begin{axis}[
tick align=outside,
tick pos=left,
axis lines=left,
title={\(\displaystyle \Delta=3\) Repair Precision},
legend style={fill opacity=0.8, draw opacity=1, text opacity=1, legend columns=1, legend pos=north west},
x grid style={darkgray176},
xlabel={Seconds},
xmin=-0.6425, xmax=9.6425,
xtick style={color=black},
xtick={0,1,2,3,4,5,6,7,8,9},
xticklabels={6,12,18,24,30,36,42,48,54,60},
y grid style={darkgray176},
ylabel={Precision@k},
ymin=0, ymax=0.5754,
ytick style={color=black}
]
\draw[draw=none,fill=darkviolet1270255] (axis cs:-0.175,0) rectangle (axis cs:0.175,0.014);
\addlegendimage{ybar,ybar legend,draw=none,fill=darkviolet1270255}
\addlegendentry{P@All}

\draw[draw=none,fill=darkviolet1270255] (axis cs:0.825,0) rectangle (axis cs:1.175,0.068);
\draw[draw=none,fill=darkviolet1270255] (axis cs:1.825,0) rectangle (axis cs:2.175,0.164);
\draw[draw=none,fill=darkviolet1270255] (axis cs:2.825,0) rectangle (axis cs:3.175,0.233);
\draw[draw=none,fill=darkviolet1270255] (axis cs:3.825,0) rectangle (axis cs:4.175,0.342);
\draw[draw=none,fill=darkviolet1270255] (axis cs:4.825,0) rectangle (axis cs:5.175,0.411);
\draw[draw=none,fill=darkviolet1270255] (axis cs:5.825,0) rectangle (axis cs:6.175,0.438);
\draw[draw=none,fill=darkviolet1270255] (axis cs:6.825,0) rectangle (axis cs:7.175,0.493);
\draw[draw=none,fill=darkviolet1270255] (axis cs:7.825,0) rectangle (axis cs:8.175,0.507);
\draw[draw=none,fill=darkviolet1270255] (axis cs:8.825,0) rectangle (axis cs:9.175,0.548);
\draw[draw=none,fill=royalblue8762253] (axis cs:-0.175,0) rectangle (axis cs:0.175,0);
\addlegendimage{ybar,ybar legend,draw=none,fill=royalblue8762253}
\addlegendentry{P@10}

\draw[draw=none,fill=royalblue8762253] (axis cs:0.825,0) rectangle (axis cs:1.175,0.014);
\draw[draw=none,fill=royalblue8762253] (axis cs:1.825,0) rectangle (axis cs:2.175,0.055);
\draw[draw=none,fill=royalblue8762253] (axis cs:2.825,0) rectangle (axis cs:3.175,0.068);
\draw[draw=none,fill=royalblue8762253] (axis cs:3.825,0) rectangle (axis cs:4.175,0.11);
\draw[draw=none,fill=royalblue8762253] (axis cs:4.825,0) rectangle (axis cs:5.175,0.137);
\draw[draw=none,fill=royalblue8762253] (axis cs:5.825,0) rectangle (axis cs:6.175,0.151);
\draw[draw=none,fill=royalblue8762253] (axis cs:6.825,0) rectangle (axis cs:7.175,0.164);
\draw[draw=none,fill=royalblue8762253] (axis cs:7.825,0) rectangle (axis cs:8.175,0.178);
\draw[draw=none,fill=royalblue8762253] (axis cs:8.825,0) rectangle (axis cs:9.175,0.178);
\draw[draw=none,fill=dodgerblue45123246] (axis cs:-0.175,0) rectangle (axis cs:0.175,0);
\addlegendimage{ybar,ybar legend,draw=none,fill=dodgerblue45123246}
\addlegendentry{P@5}

\draw[draw=none,fill=dodgerblue45123246] (axis cs:0.825,0) rectangle (axis cs:1.175,0.014);
\draw[draw=none,fill=dodgerblue45123246] (axis cs:1.825,0) rectangle (axis cs:2.175,0.055);
\draw[draw=none,fill=dodgerblue45123246] (axis cs:2.825,0) rectangle (axis cs:3.175,0.068);
\draw[draw=none,fill=dodgerblue45123246] (axis cs:3.825,0) rectangle (axis cs:4.175,0.11);
\draw[draw=none,fill=dodgerblue45123246] (axis cs:4.825,0) rectangle (axis cs:5.175,0.123);
\draw[draw=none,fill=dodgerblue45123246] (axis cs:5.825,0) rectangle (axis cs:6.175,0.137);
\draw[draw=none,fill=dodgerblue45123246] (axis cs:6.825,0) rectangle (axis cs:7.175,0.151);
\draw[draw=none,fill=dodgerblue45123246] (axis cs:7.825,0) rectangle (axis cs:8.175,0.164);
\draw[draw=none,fill=dodgerblue45123246] (axis cs:8.825,0) rectangle (axis cs:9.175,0.164);
\draw[draw=none,fill=deepskyblue3176236] (axis cs:-0.175,0) rectangle (axis cs:0.175,0);
\addlegendimage{ybar,ybar legend,draw=none,fill=deepskyblue3176236}
\addlegendentry{P@1}

\draw[draw=none,fill=deepskyblue3176236] (axis cs:0.825,0) rectangle (axis cs:1.175,0.014);
\draw[draw=none,fill=deepskyblue3176236] (axis cs:1.825,0) rectangle (axis cs:2.175,0.055);
\draw[draw=none,fill=deepskyblue3176236] (axis cs:2.825,0) rectangle (axis cs:3.175,0.068);
\draw[draw=none,fill=deepskyblue3176236] (axis cs:3.825,0) rectangle (axis cs:4.175,0.096);
\draw[draw=none,fill=deepskyblue3176236] (axis cs:4.825,0) rectangle (axis cs:5.175,0.11);
\draw[draw=none,fill=deepskyblue3176236] (axis cs:5.825,0) rectangle (axis cs:6.175,0.123);
\draw[draw=none,fill=deepskyblue3176236] (axis cs:6.825,0) rectangle (axis cs:7.175,0.137);
\draw[draw=none,fill=deepskyblue3176236] (axis cs:7.825,0) rectangle (axis cs:8.175,0.151);
\draw[draw=none,fill=deepskyblue3176236] (axis cs:8.825,0) rectangle (axis cs:9.175,0.151);
\addplot [red, thick] coordinates {(-0.8925,0.07) (14.8925,0.07)};
\end{axis}

\end{tikzpicture}
}
    \resizebox{.24\textwidth}{!}{% This file was created with tikzplotlib v0.10.1.
\begin{tikzpicture}

\definecolor{darkgray176}{RGB}{176,176,176}
\definecolor{darkviolet1270255}{RGB}{127,0,255}
\definecolor{deepskyblue3176236}{RGB}{3,176,236}
\definecolor{dodgerblue45123246}{RGB}{45,123,246}
\definecolor{royalblue8762253}{RGB}{87,62,253}

\begin{axis}[
tick align=outside,
tick pos=left,
axis lines=left,
title={\(\displaystyle \Delta=4\) Repair Precision},
legend style={fill opacity=0.8, draw opacity=1, text opacity=1, legend columns=1, legend pos=north west},
x grid style={darkgray176},
xlabel={Seconds},
xmin=-0.5925, xmax=8.5925,
xtick style={color=black},
xtick={0,1,2,3,4,5,6,7,8},
xticklabels={10,20,30,40,50,60,70,80,90},
y grid style={darkgray176},
ylabel={Precision@k},
ymin=0, ymax=0.54915,
ytick style={color=black}
]
\draw[draw=none,fill=darkviolet1270255] (axis cs:-0.175,0) rectangle (axis cs:0.175,0.123);
\addlegendimage{ybar,ybar legend,draw=none,fill=darkviolet1270255}
\addlegendentry{P@All}

\draw[draw=none,fill=darkviolet1270255] (axis cs:0.825,0) rectangle (axis cs:1.175,0.323);
\draw[draw=none,fill=darkviolet1270255] (axis cs:1.825,0) rectangle (axis cs:2.175,0.415);
\draw[draw=none,fill=darkviolet1270255] (axis cs:2.825,0) rectangle (axis cs:3.175,0.431);
\draw[draw=none,fill=darkviolet1270255] (axis cs:3.825,0) rectangle (axis cs:4.175,0.477);
\draw[draw=none,fill=darkviolet1270255] (axis cs:4.825,0) rectangle (axis cs:5.175,0.508);
\draw[draw=none,fill=darkviolet1270255] (axis cs:5.825,0) rectangle (axis cs:6.175,0.523);
\draw[draw=none,fill=darkviolet1270255] (axis cs:6.825,0) rectangle (axis cs:7.175,0.523);
\draw[draw=none,fill=darkviolet1270255] (axis cs:7.825,0) rectangle (axis cs:8.175,0.523);
\draw[draw=none,fill=royalblue8762253] (axis cs:-0.175,0) rectangle (axis cs:0.175,0.092);
\addlegendimage{ybar,ybar legend,draw=none,fill=royalblue8762253}
\addlegendentry{P@10}

\draw[draw=none,fill=royalblue8762253] (axis cs:0.825,0) rectangle (axis cs:1.175,0.169);
\draw[draw=none,fill=royalblue8762253] (axis cs:1.825,0) rectangle (axis cs:2.175,0.215);
\draw[draw=none,fill=royalblue8762253] (axis cs:2.825,0) rectangle (axis cs:3.175,0.231);
\draw[draw=none,fill=royalblue8762253] (axis cs:3.825,0) rectangle (axis cs:4.175,0.277);
\draw[draw=none,fill=royalblue8762253] (axis cs:4.825,0) rectangle (axis cs:5.175,0.308);
\draw[draw=none,fill=royalblue8762253] (axis cs:5.825,0) rectangle (axis cs:6.175,0.323);
\draw[draw=none,fill=royalblue8762253] (axis cs:6.825,0) rectangle (axis cs:7.175,0.323);
\draw[draw=none,fill=royalblue8762253] (axis cs:7.825,0) rectangle (axis cs:8.175,0.323);
\draw[draw=none,fill=dodgerblue45123246] (axis cs:-0.175,0) rectangle (axis cs:0.175,0.092);
\addlegendimage{ybar,ybar legend,draw=none,fill=dodgerblue45123246}
\addlegendentry{P@5}

\draw[draw=none,fill=dodgerblue45123246] (axis cs:0.825,0) rectangle (axis cs:1.175,0.169);
\draw[draw=none,fill=dodgerblue45123246] (axis cs:1.825,0) rectangle (axis cs:2.175,0.215);
\draw[draw=none,fill=dodgerblue45123246] (axis cs:2.825,0) rectangle (axis cs:3.175,0.231);
\draw[draw=none,fill=dodgerblue45123246] (axis cs:3.825,0) rectangle (axis cs:4.175,0.277);
\draw[draw=none,fill=dodgerblue45123246] (axis cs:4.825,0) rectangle (axis cs:5.175,0.308);
\draw[draw=none,fill=dodgerblue45123246] (axis cs:5.825,0) rectangle (axis cs:6.175,0.323);
\draw[draw=none,fill=dodgerblue45123246] (axis cs:6.825,0) rectangle (axis cs:7.175,0.323);
\draw[draw=none,fill=dodgerblue45123246] (axis cs:7.825,0) rectangle (axis cs:8.175,0.323);
\draw[draw=none,fill=deepskyblue3176236] (axis cs:-0.175,0) rectangle (axis cs:0.175,0.077);
\addlegendimage{ybar,ybar legend,draw=none,fill=deepskyblue3176236}
\addlegendentry{P@1}

\draw[draw=none,fill=deepskyblue3176236] (axis cs:0.825,0) rectangle (axis cs:1.175,0.138);
\draw[draw=none,fill=deepskyblue3176236] (axis cs:1.825,0) rectangle (axis cs:2.175,0.185);
\draw[draw=none,fill=deepskyblue3176236] (axis cs:2.825,0) rectangle (axis cs:3.175,0.2);
\draw[draw=none,fill=deepskyblue3176236] (axis cs:3.825,0) rectangle (axis cs:4.175,0.246);
\draw[draw=none,fill=deepskyblue3176236] (axis cs:4.825,0) rectangle (axis cs:5.175,0.277);
\draw[draw=none,fill=deepskyblue3176236] (axis cs:5.825,0) rectangle (axis cs:6.175,0.292);
\draw[draw=none,fill=deepskyblue3176236] (axis cs:6.825,0) rectangle (axis cs:7.175,0.292);
\draw[draw=none,fill=deepskyblue3176236] (axis cs:7.825,0) rectangle (axis cs:8.175,0.292);
\addplot [red, thick] coordinates {(-0.8925,0.10) (14.8925,0.10)};
\end{axis}

\end{tikzpicture}
}
%    \resizebox{.24\textwidth}{!}{% This file was created with tikzplotlib v0.10.1.
\begin{tikzpicture}

\definecolor{darkgray176}{RGB}{176,176,176}
\definecolor{darkviolet1270255}{RGB}{127,0,255}
\definecolor{deepskyblue3176236}{RGB}{3,176,236}
\definecolor{dodgerblue45123246}{RGB}{45,123,246}
\definecolor{royalblue8762253}{RGB}{87,62,253}

\begin{axis}[
tick align=outside,
tick pos=left,
axis lines=left,
title={\(\displaystyle \Delta=4\) Repair Precision},
legend style={fill opacity=0.8, draw opacity=1, text opacity=1, legend columns=1, legend pos=north west},
x grid style={darkgray176},
xlabel={Seconds},
xmin=-0.5925, xmax=8.5925,
xtick={0,1,2,3,4,5,6,7,8},
xticklabels={10,20,30,40,50,60,70,80,90},
y grid style={darkgray176},
ylabel={Precision@k},
ymin=0, ymax=0.54915,
]
\draw[draw=none,fill=darkviolet1270255] (axis cs:-0.175,0) rectangle (axis cs:0.175,0.123);
\addlegendimage{ybar,ybar legend,draw=none,fill=darkviolet1270255}
\addlegendentry{P@All}

\draw[draw=none,fill=darkviolet1270255] (axis cs:0.825,0) rectangle (axis cs:1.175,0.323);
\draw[draw=none,fill=darkviolet1270255] (axis cs:1.825,0) rectangle (axis cs:2.175,0.415);
\draw[draw=none,fill=darkviolet1270255] (axis cs:2.825,0) rectangle (axis cs:3.175,0.431);
\draw[draw=none,fill=darkviolet1270255] (axis cs:3.825,0) rectangle (axis cs:4.175,0.477);
\draw[draw=none,fill=darkviolet1270255] (axis cs:4.825,0) rectangle (axis cs:5.175,0.508);
\draw[draw=none,fill=darkviolet1270255] (axis cs:5.825,0) rectangle (axis cs:6.175,0.523);
\draw[draw=none,fill=darkviolet1270255] (axis cs:6.825,0) rectangle (axis cs:7.175,0.523);
\draw[draw=none,fill=darkviolet1270255] (axis cs:7.825,0) rectangle (axis cs:8.175,0.523);
\draw[draw=none,fill=royalblue8762253] (axis cs:-0.175,0) rectangle (axis cs:0.175,0.092);
\addlegendimage{ybar,ybar legend,draw=none,fill=royalblue8762253}
\addlegendentry{P@10}

\draw[draw=none,fill=royalblue8762253] (axis cs:0.825,0) rectangle (axis cs:1.175,0.169);
\draw[draw=none,fill=royalblue8762253] (axis cs:1.825,0) rectangle (axis cs:2.175,0.215);
\draw[draw=none,fill=royalblue8762253] (axis cs:2.825,0) rectangle (axis cs:3.175,0.231);
\draw[draw=none,fill=royalblue8762253] (axis cs:3.825,0) rectangle (axis cs:4.175,0.277);
\draw[draw=none,fill=royalblue8762253] (axis cs:4.825,0) rectangle (axis cs:5.175,0.308);
\draw[draw=none,fill=royalblue8762253] (axis cs:5.825,0) rectangle (axis cs:6.175,0.323);
\draw[draw=none,fill=royalblue8762253] (axis cs:6.825,0) rectangle (axis cs:7.175,0.323);
\draw[draw=none,fill=royalblue8762253] (axis cs:7.825,0) rectangle (axis cs:8.175,0.323);
\draw[draw=none,fill=dodgerblue45123246] (axis cs:-0.175,0) rectangle (axis cs:0.175,0.092);
\addlegendimage{ybar,ybar legend,draw=none,fill=dodgerblue45123246}
\addlegendentry{P@5}

\draw[draw=none,fill=dodgerblue45123246] (axis cs:0.825,0) rectangle (axis cs:1.175,0.169);
\draw[draw=none,fill=dodgerblue45123246] (axis cs:1.825,0) rectangle (axis cs:2.175,0.215);
\draw[draw=none,fill=dodgerblue45123246] (axis cs:2.825,0) rectangle (axis cs:3.175,0.231);
\draw[draw=none,fill=dodgerblue45123246] (axis cs:3.825,0) rectangle (axis cs:4.175,0.277);
\draw[draw=none,fill=dodgerblue45123246] (axis cs:4.825,0) rectangle (axis cs:5.175,0.308);
\draw[draw=none,fill=dodgerblue45123246] (axis cs:5.825,0) rectangle (axis cs:6.175,0.323);
\draw[draw=none,fill=dodgerblue45123246] (axis cs:6.825,0) rectangle (axis cs:7.175,0.323);
\draw[draw=none,fill=dodgerblue45123246] (axis cs:7.825,0) rectangle (axis cs:8.175,0.323);
\draw[draw=none,fill=deepskyblue3176236] (axis cs:-0.175,0) rectangle (axis cs:0.175,0.077);
\addlegendimage{ybar,ybar legend,draw=none,fill=deepskyblue3176236}
\addlegendentry{P@1}

\draw[draw=none,fill=deepskyblue3176236] (axis cs:0.825,0) rectangle (axis cs:1.175,0.138);
\draw[draw=none,fill=deepskyblue3176236] (axis cs:1.825,0) rectangle (axis cs:2.175,0.185);
\draw[draw=none,fill=deepskyblue3176236] (axis cs:2.825,0) rectangle (axis cs:3.175,0.2);
\draw[draw=none,fill=deepskyblue3176236] (axis cs:3.825,0) rectangle (axis cs:4.175,0.246);
\draw[draw=none,fill=deepskyblue3176236] (axis cs:4.825,0) rectangle (axis cs:5.175,0.277);
\draw[draw=none,fill=deepskyblue3176236] (axis cs:5.825,0) rectangle (axis cs:6.175,0.292);
\draw[draw=none,fill=deepskyblue3176236] (axis cs:6.825,0) rectangle (axis cs:7.175,0.292);
\draw[draw=none,fill=deepskyblue3176236] (axis cs:7.825,0) rectangle (axis cs:8.175,0.292);
\addplot [red, thick] coordinates {(-0.8925,0.10) (14.8925,0.10)};
\end{axis}

\end{tikzpicture}
}
%\resizebox{.3\textwidth}{!}{% This file was created with tikzplotlib v0.10.1.
\begin{tikzpicture}

\definecolor{darkgray176}{RGB}{176,176,176}
\definecolor{darkviolet1270255}{RGB}{127,0,255}
\definecolor{deepskyblue3176236}{RGB}{3,176,236}
\definecolor{dodgerblue45123246}{RGB}{45,123,246}
\definecolor{lightgray204}{RGB}{204,204,204}
\definecolor{royalblue8762253}{RGB}{87,62,253}

\begin{axis}[
legend cell align={left},
legend style={fill opacity=0.8, draw opacity=1, text opacity=1, draw=lightgray204, legend columns=-1, legend pos=north west},
tick align=outside,
tick pos=left,
title={$\Delta=1$ Repair Precision},
x grid style={darkgray176},
xlabel={Seconds},
xmin=-0.4925, xmax=6.4925,
xtick style={color=black},
xtick={0,1,2,3,4,5,6},
xticklabels={20,60,100,140,180,220,260},
y grid style={darkgray176},
ylabel={\phantom{Precision@k}},
ymin=0, ymax=0.77595,
ytick style={color=black}
]
\addlegendimage{empty legend}
\addlegendentry{P@}
\draw[draw=none,fill=darkviolet1270255] (axis cs:-0.175,0) rectangle (axis cs:0.175,0.145);
\addlegendimage{ybar,ybar legend,draw=none,fill=darkviolet1270255}
\addlegendentry{All}

\draw[draw=none,fill=darkviolet1270255] (axis cs:0.825,0) rectangle (axis cs:1.175,0.304);
\draw[draw=none,fill=darkviolet1270255] (axis cs:1.825,0) rectangle (axis cs:2.175,0.42);
\draw[draw=none,fill=darkviolet1270255] (axis cs:2.825,0) rectangle (axis cs:3.175,0.507);
\draw[draw=none,fill=darkviolet1270255] (axis cs:3.825,0) rectangle (axis cs:4.175,0.594);
\draw[draw=none,fill=darkviolet1270255] (axis cs:4.825,0) rectangle (axis cs:5.175,0.71);
\draw[draw=none,fill=darkviolet1270255] (axis cs:5.825,0) rectangle (axis cs:6.175,0.739);
\draw[draw=none,fill=royalblue8762253] (axis cs:-0.175,0) rectangle (axis cs:0.175,0.145);
\addlegendimage{ybar,ybar legend,draw=none,fill=royalblue8762253}
\addlegendentry{10}

\draw[draw=none,fill=royalblue8762253] (axis cs:0.825,0) rectangle (axis cs:1.175,0.304);
\draw[draw=none,fill=royalblue8762253] (axis cs:1.825,0) rectangle (axis cs:2.175,0.42);
\draw[draw=none,fill=royalblue8762253] (axis cs:2.825,0) rectangle (axis cs:3.175,0.507);
\draw[draw=none,fill=royalblue8762253] (axis cs:3.825,0) rectangle (axis cs:4.175,0.594);
\draw[draw=none,fill=royalblue8762253] (axis cs:4.825,0) rectangle (axis cs:5.175,0.703);
\draw[draw=none,fill=royalblue8762253] (axis cs:5.825,0) rectangle (axis cs:6.175,0.732);
\draw[draw=none,fill=dodgerblue45123246] (axis cs:-0.175,0) rectangle (axis cs:0.175,0.145);
\addlegendimage{ybar,ybar legend,draw=none,fill=dodgerblue45123246}
\addlegendentry{5}

\draw[draw=none,fill=dodgerblue45123246] (axis cs:0.825,0) rectangle (axis cs:1.175,0.304);
\draw[draw=none,fill=dodgerblue45123246] (axis cs:1.825,0) rectangle (axis cs:2.175,0.42);
\draw[draw=none,fill=dodgerblue45123246] (axis cs:2.825,0) rectangle (axis cs:3.175,0.493);
\draw[draw=none,fill=dodgerblue45123246] (axis cs:3.825,0) rectangle (axis cs:4.175,0.565);
\draw[draw=none,fill=dodgerblue45123246] (axis cs:4.825,0) rectangle (axis cs:5.175,0.638);
\draw[draw=none,fill=dodgerblue45123246] (axis cs:5.825,0) rectangle (axis cs:6.175,0.667);
\draw[draw=none,fill=deepskyblue3176236] (axis cs:-0.175,0) rectangle (axis cs:0.175,0.116);
\addlegendimage{ybar,ybar legend,draw=none,fill=deepskyblue3176236}
\addlegendentry{1}

\draw[draw=none,fill=deepskyblue3176236] (axis cs:0.825,0) rectangle (axis cs:1.175,0.203);
\draw[draw=none,fill=deepskyblue3176236] (axis cs:1.825,0) rectangle (axis cs:2.175,0.268);
\draw[draw=none,fill=deepskyblue3176236] (axis cs:2.825,0) rectangle (axis cs:3.175,0.312);
\draw[draw=none,fill=deepskyblue3176236] (axis cs:3.825,0) rectangle (axis cs:4.175,0.355);
\draw[draw=none,fill=deepskyblue3176236] (axis cs:4.825,0) rectangle (axis cs:5.175,0.428);
\draw[draw=none,fill=deepskyblue3176236] (axis cs:5.825,0) rectangle (axis cs:6.175,0.442);
\end{axis}

\end{tikzpicture}
}
%\resizebox{.307\textwidth}{!}{% This file was created with tikzplotlib v0.10.1.
\begin{tikzpicture}

\definecolor{darkgray176}{RGB}{176,176,176}
\definecolor{darkviolet1270255}{RGB}{127,0,255}
\definecolor{deepskyblue3176236}{RGB}{3,176,236}
\definecolor{dodgerblue45123246}{RGB}{45,123,246}
\definecolor{lightgray204}{RGB}{204,204,204}
\definecolor{royalblue8762253}{RGB}{87,62,253}

\begin{axis}[
legend cell align={left},
legend style={fill opacity=0.8, draw opacity=1, text opacity=1, draw=lightgray204, legend columns=-1, legend pos=north west},
tick align=outside,
tick pos=left,
title={$\Delta=2$ Repair Precision},
x grid style={darkgray176},
xlabel={Seconds},
xmin=-0.4925, xmax=6.4925,
xtick style={color=black},
xtick={0,1,2,3,4,5,6},
xticklabels={20,60,100,140,180,220,260},
y grid style={darkgray176},
ylabel={\phantom{Precision@k}},
ymin=0, ymax=0.40635,
ytick style={color=black}
]
\addlegendimage{empty legend}
\addlegendentry{P@}
\draw[draw=none,fill=darkviolet1270255] (axis cs:-0.175,0) rectangle (axis cs:0.175,0.048);
\addlegendimage{ybar,ybar legend,draw=none,fill=darkviolet1270255}
\addlegendentry{All}

\draw[draw=none,fill=darkviolet1270255] (axis cs:0.825,0) rectangle (axis cs:1.175,0.081);
\draw[draw=none,fill=darkviolet1270255] (axis cs:1.825,0) rectangle (axis cs:2.175,0.177);
\draw[draw=none,fill=darkviolet1270255] (axis cs:2.825,0) rectangle (axis cs:3.175,0.274);
\draw[draw=none,fill=darkviolet1270255] (axis cs:3.825,0) rectangle (axis cs:4.175,0.306);
\draw[draw=none,fill=darkviolet1270255] (axis cs:4.825,0) rectangle (axis cs:5.175,0.355);
\draw[draw=none,fill=darkviolet1270255] (axis cs:5.825,0) rectangle (axis cs:6.175,0.387);
\draw[draw=none,fill=royalblue8762253] (axis cs:-0.175,0) rectangle (axis cs:0.175,0.048);
\addlegendimage{ybar,ybar legend,draw=none,fill=royalblue8762253}
\addlegendentry{10}

\draw[draw=none,fill=royalblue8762253] (axis cs:0.825,0) rectangle (axis cs:1.175,0.081);
\draw[draw=none,fill=royalblue8762253] (axis cs:1.825,0) rectangle (axis cs:2.175,0.113);
\draw[draw=none,fill=royalblue8762253] (axis cs:2.825,0) rectangle (axis cs:3.175,0.21);
\draw[draw=none,fill=royalblue8762253] (axis cs:3.825,0) rectangle (axis cs:4.175,0.226);
\draw[draw=none,fill=royalblue8762253] (axis cs:4.825,0) rectangle (axis cs:5.175,0.258);
\draw[draw=none,fill=royalblue8762253] (axis cs:5.825,0) rectangle (axis cs:6.175,0.242);
\draw[draw=none,fill=dodgerblue45123246] (axis cs:-0.175,0) rectangle (axis cs:0.175,0.048);
\addlegendimage{ybar,ybar legend,draw=none,fill=dodgerblue45123246}
\addlegendentry{5}

\draw[draw=none,fill=dodgerblue45123246] (axis cs:0.825,0) rectangle (axis cs:1.175,0.032);
\draw[draw=none,fill=dodgerblue45123246] (axis cs:1.825,0) rectangle (axis cs:2.175,0.065);
\draw[draw=none,fill=dodgerblue45123246] (axis cs:2.825,0) rectangle (axis cs:3.175,0.129);
\draw[draw=none,fill=dodgerblue45123246] (axis cs:3.825,0) rectangle (axis cs:4.175,0.113);
\draw[draw=none,fill=dodgerblue45123246] (axis cs:4.825,0) rectangle (axis cs:5.175,0.129);
\draw[draw=none,fill=dodgerblue45123246] (axis cs:5.825,0) rectangle (axis cs:6.175,0.113);
\draw[draw=none,fill=deepskyblue3176236] (axis cs:-0.175,0) rectangle (axis cs:0.175,0.016);
\addlegendimage{ybar,ybar legend,draw=none,fill=deepskyblue3176236}
\addlegendentry{1}

\draw[draw=none,fill=deepskyblue3176236] (axis cs:0.825,0) rectangle (axis cs:1.175,0);
\draw[draw=none,fill=deepskyblue3176236] (axis cs:1.825,0) rectangle (axis cs:2.175,0.016);
\draw[draw=none,fill=deepskyblue3176236] (axis cs:2.825,0) rectangle (axis cs:3.175,0.065);
\draw[draw=none,fill=deepskyblue3176236] (axis cs:3.825,0) rectangle (axis cs:4.175,0.048);
\draw[draw=none,fill=deepskyblue3176236] (axis cs:4.825,0) rectangle (axis cs:5.175,0.032);
\draw[draw=none,fill=deepskyblue3176236] (axis cs:5.825,0) rectangle (axis cs:6.175,0.032);
\end{axis}

\end{tikzpicture}
}
    \caption{Human repair benchmarks. Note the y-axis across different edit distance plots has varying ranges. The red line indicates Seq2Parse and the orange line indicates BIFI's Precision@1 on the same repairs.}\label{fig:human}
  \end{figure}

  \begin{wrapfigure}{r}{0.4\textwidth}
    \vspace{-1.5cm}
    \hspace{-0.8cm}
    \resizebox{.53\textwidth}{!}{% This file was created with tikzplotlib v0.10.1.
\begin{tikzpicture}

\definecolor{darkgray176}{RGB}{176,176,176}
\definecolor{teal9147104}{RGB}{9,147,104}

\begin{axis}[
hide x axis,
hide y axis,
tick align=outside,
tick pos=left,
x grid style={darkgray176},
xmin=-0.65, xmax=2.3168,
xtick style={color=black},
y grid style={darkgray176},
ymin=-1.09996565315951, ymax=0.5934,
ytick style={color=black}
]
\path [draw=teal9147104, fill=teal9147104]
(axis cs:-0.25,0.1934)
--(axis cs:-0.125,0.1934)
.. controls (axis cs:-0.03125,0.1934) and (axis cs:0.03125,0.1934) .. (axis cs:0.125,0.1934)
--(axis cs:1.5928,0.1934)
.. controls (axis cs:1.670743743262,0.1934) and (axis cs:1.745574844138,0.16240394309) .. (axis cs:1.800689393614,0.107289393614)
.. controls (axis cs:1.85580394309,0.052174844138) and (axis cs:1.8868,-0.022656256738) .. (axis cs:1.8868,-0.1006)
--(axis cs:1.8868,-0.4434)
--(axis cs:1.9168,-0.4434)
--(axis cs:1.7898,-0.549965653159515)
--(axis cs:1.6628,-0.4434)
--(axis cs:1.6928,-0.4434)
--(axis cs:1.6928,-0.1006)
.. controls (axis cs:1.6928,-0.0740885227) and (axis cs:1.6822571235,-0.0486357673) .. (axis cs:1.6635106781,-0.0298893219)
.. controls (axis cs:1.6447642327,-0.0111428765) and (axis cs:1.6193114773,-0.000600000000000017) .. (axis cs:1.5928,-0.000600000000000017)
--(axis cs:1.4428,-0.000600000000000017)
--(axis cs:1.2476,-0.000600000000000017)
.. controls (axis cs:1.2993504036896,-0.000600000000000017) and (axis cs:1.3490341822304,-0.021179694928) .. (axis cs:1.3856272436512,-0.0577727563488)
.. controls (axis cs:1.422220305072,-0.0943658177696) and (axis cs:1.4428,-0.1440495963104) .. (axis cs:1.4428,-0.1958)
--(axis cs:1.4428,-0.4434)
--(axis cs:1.4728,-0.4434)
--(axis cs:1.3952,-0.508514131379357)
--(axis cs:1.3176,-0.4434)
--(axis cs:1.3476,-0.4434)
--(axis cs:1.3476,-0.1958)
.. controls (axis cs:1.3476,-0.1692885227) and (axis cs:1.3370571235,-0.1438357673) .. (axis cs:1.3183106781,-0.1250893219)
.. controls (axis cs:1.2995642327,-0.1063428765) and (axis cs:1.2741114773,-0.0958) .. (axis cs:1.2476,-0.0958)
--(axis cs:1.0976,-0.0958)
--(axis cs:0.9364,-0.0958)
.. controls (axis cs:0.9791365014076,-0.0958) and (axis cs:1.0201663431124,-0.112795116918) .. (axis cs:1.0503856130972,-0.1430143869028)
.. controls (axis cs:1.080604883082,-0.1732336568876) and (axis cs:1.0976,-0.2142634985924) .. (axis cs:1.0976,-0.257)
--(axis cs:1.0976,-0.4434)
--(axis cs:1.1276,-0.4434)
--(axis cs:1.067,-0.494249437649343)
--(axis cs:1.0064,-0.4434)
--(axis cs:1.0364,-0.4434)
--(axis cs:1.0364,-0.257)
.. controls (axis cs:1.0364,-0.2304885227) and (axis cs:1.0258571235,-0.2050357673) .. (axis cs:1.0071106781,-0.1862893219)
.. controls (axis cs:0.9883642327,-0.1675428765) and (axis cs:0.9629114773,-0.157) .. (axis cs:0.9364,-0.157)
--(axis cs:0.7864,-0.157)
--(axis cs:0.6848,-0.157)
.. controls (axis cs:0.7117356609368,-0.157) and (axis cs:0.7375956604232,-0.167711562524) .. (axis cs:0.7566420489496,-0.1867579510504)
.. controls (axis cs:0.775688437476,-0.2058043395768) and (axis cs:0.7864,-0.2316643390632) .. (axis cs:0.7864,-0.2586)
--(axis cs:0.7864,-0.4434)
--(axis cs:0.8164,-0.4434)
--(axis cs:0.7856,-0.46924426864026)
--(axis cs:0.7548,-0.4434)
--(axis cs:0.7848,-0.4434)
--(axis cs:0.7848,-0.2586)
.. controls (axis cs:0.7848,-0.2320885227) and (axis cs:0.7742571235,-0.2066357673) .. (axis cs:0.7555106781,-0.1878893219)
.. controls (axis cs:0.7367642327,-0.1691428765) and (axis cs:0.7113114773,-0.1586) .. (axis cs:0.6848,-0.1586)
--(axis cs:0.5348,-0.1586)
--(axis cs:0.4,-0.1586)
.. controls (axis cs:0.4357374714004,-0.1586) and (axis cs:0.4700477856796,-0.172811797522) .. (axis cs:0.4953179940788,-0.1980820059212)
.. controls (axis cs:0.520588202478,-0.2233522143204) and (axis cs:0.5348,-0.2576625285996) .. (axis cs:0.5348,-0.2934)
--(axis cs:0.5348,-0.4434)
--(axis cs:0.5648,-0.4434)
--(axis cs:0.5174,-0.483173322517803)
--(axis cs:0.47,-0.4434)
--(axis cs:0.5,-0.4434)
--(axis cs:0.5,-0.2934)
.. controls (axis cs:0.5,-0.2668885227) and (axis cs:0.4894571235,-0.2414357673) .. (axis cs:0.4707106781,-0.2226893219)
.. controls (axis cs:0.4519642327,-0.2039428765) and (axis cs:0.4265114773,-0.1934) .. (axis cs:0.4,-0.1934)
--(axis cs:0.25,-0.1934)
--(axis cs:0.25,-0.1934)
--(axis cs:0.125,-0.1934)
.. controls (axis cs:0.03125,-0.1934) and (axis cs:-0.03125,-0.1934) .. (axis cs:-0.125,-0.1934)
--(axis cs:-0.25,-0.1934)
--(axis cs:-0.25,-0.1934)
--(axis cs:-0.0877181313303141,0)
--(axis cs:-0.25,0.1934)
--(axis cs:0,0.1934)
--(axis cs:-0.25,0.1934)
--cycle;
\draw (axis cs:-0.237718131330314,0) node[
  scale=0.5,
  text=black,
  rotate=0.0,
  align=center
]{Total\\967};
\draw (axis cs:1.7898,-0.699965653159515) node[
  scale=0.5,
  text=black,
  rotate=0.0,
  align=center
]{Top-1\\485};
\draw (axis cs:1.3952,-0.658514131379357) node[
  scale=0.5,
  text=black,
  rotate=0.0,
  align=center
]{Top-2+\\238};
\draw (axis cs:1.067,-0.644249437649343) node[
  scale=0.5,
  text=black,
  rotate=0.0,
  align=center
]{Timeout\\153};
\draw (axis cs:0.7856,-0.61924426864026) node[
  scale=0.5,
  text=black,
  rotate=0.0,
  align=center
]{$\sigma'\notin\mathcal{L}_\cap$\\4};
\draw (axis cs:0.5174,-0.633173322517803) node[
  scale=0.5,
  text=black,
  rotate=0.0,
  align=center
]{OOM\\87};
\end{axis}

\end{tikzpicture}
}
    \vspace{-1.8cm}
    \caption{Outcomes in the repair pipeline.}
    \label{fig:sankey}
  \end{wrapfigure}

  %We believe that rewriting the sampler in CUDA or using a more informed prior could significantly improve the latency-precision tradeoff.

  We present a Sankey diagram of our repair pipeline in Fig.~\ref{fig:sankey}. We drew 967 total repairs from the StackOverflow dataset balanced evenly across lengths and edit distances ($\lfloor|\err\sigma| / 10\rfloor \in [0, 8], \Delta(\err\sigma, \sigma') < 4$) with a timeout of 30s and tracked individual outcomes. In 87 cases, the intersection grammar was too large to construct and threw an out-of-memory (OOM) error, in 4 cases the human repair was not recognized, in 153 cases the sampler timed out before drawing the human repair, in 238 cases the human repair was drawn but not ranked first, and in the remaining 485 cases the first prediction matched the human repair.

  \clearpage The remaining experiments in this section were run on a 10-core ARM64 M1 with 16 GB of memory. We balance the StackOverflow dataset across Levenshtein distances, then measure the number of samples required to draw the exact human repair across varying Levenshtein radii. This tells us of how many samples are required on average to saturate the admissible set.

  \begin{figure}[h!]
    % This file was created with tikzplotlib v0.10.1.
\begin{tikzpicture}

\definecolor{darkgray176}{RGB}{176,176,176}
\definecolor{darkorange25512714}{RGB}{255,127,14}
\definecolor{forestgreen4416044}{RGB}{44,160,44}
\definecolor{steelblue31119180}{RGB}{31,119,180}

\begin{axis}[
width=\textwidth,
height=.4\textwidth,
log basis x={10},
tick align=outside,
tick pos=left,
axis lines*=left,
title={Sample Efficiency},
legend style={fill opacity=0.8, draw opacity=1, text opacity=1, legend columns=1, legend pos=north west, at={(0.055, 0.96)}},
x grid style={darkgray176},
xlabel={Samples drawn (log scale)},
xmin=0.34892211545542, xmax=1001,
xmode=log,
%xtick style={color=black},
xtick={0.01,0.1,1,10,100,1000,10000},
xticklabels={
  \(\displaystyle 10^-2\),
  \(\displaystyle 10^-1\),
  \(\displaystyle 10^2\),
  \(\displaystyle 10^3\),
  \(\displaystyle 10^4\),
  \(\displaystyle 10^5\),
  \(\displaystyle 10^6\)
},
y grid style={darkgray176},
ylabel={Precision@All},
ymin=0, ymax=101,
%ytick style={color=black}
ytick={0,20,40,60,80,100}
]
\draw[draw=none,fill=steelblue31119180] (axis cs:-0.5,0) rectangle (axis cs:0.5,100);
\addlegendimage{ybar,ybar legend,draw=none,fill=steelblue31119180}
\addlegendentry{$\Delta_L=1$}

\draw[draw=none,fill=steelblue31119180] (axis cs:0.5,0) rectangle (axis cs:1.5,100);
\draw[draw=none,fill=steelblue31119180] (axis cs:1.5,0) rectangle (axis cs:2.5,100);
\draw[draw=none,fill=steelblue31119180] (axis cs:2.5,0) rectangle (axis cs:3.5,100);
\draw[draw=none,fill=steelblue31119180] (axis cs:3.5,0) rectangle (axis cs:4.5,100);
\draw[draw=none,fill=steelblue31119180] (axis cs:4.5,0) rectangle (axis cs:5.5,100);
\draw[draw=none,fill=steelblue31119180] (axis cs:5.5,0) rectangle (axis cs:6.5,100);
\draw[draw=none,fill=steelblue31119180] (axis cs:6.5,0) rectangle (axis cs:7.5,100);
\draw[draw=none,fill=steelblue31119180] (axis cs:7.5,0) rectangle (axis cs:8.5,100);
\draw[draw=none,fill=steelblue31119180] (axis cs:8.5,0) rectangle (axis cs:9.5,100);
\draw[draw=none,fill=steelblue31119180] (axis cs:9.5,0) rectangle (axis cs:10.5,100);
\draw[draw=none,fill=steelblue31119180] (axis cs:10.5,0) rectangle (axis cs:11.5,100);
\draw[draw=none,fill=steelblue31119180] (axis cs:11.5,0) rectangle (axis cs:12.5,100);
\draw[draw=none,fill=steelblue31119180] (axis cs:12.5,0) rectangle (axis cs:13.5,100);
\draw[draw=none,fill=steelblue31119180] (axis cs:13.5,0) rectangle (axis cs:14.5,100);
\draw[draw=none,fill=steelblue31119180] (axis cs:14.5,0) rectangle (axis cs:15.5,100);
\draw[draw=none,fill=steelblue31119180] (axis cs:15.5,0) rectangle (axis cs:16.5,100);
\draw[draw=none,fill=steelblue31119180] (axis cs:16.5,0) rectangle (axis cs:17.5,100);
\draw[draw=none,fill=steelblue31119180] (axis cs:17.5,0) rectangle (axis cs:18.5,100);
\draw[draw=none,fill=steelblue31119180] (axis cs:18.5,0) rectangle (axis cs:19.5,100);
\draw[draw=none,fill=steelblue31119180] (axis cs:19.5,0) rectangle (axis cs:20.5,100);
\draw[draw=none,fill=steelblue31119180] (axis cs:20.5,0) rectangle (axis cs:21.5,100);
\draw[draw=none,fill=steelblue31119180] (axis cs:21.5,0) rectangle (axis cs:22.5,100);
\draw[draw=none,fill=steelblue31119180] (axis cs:22.5,0) rectangle (axis cs:23.5,100);
\draw[draw=none,fill=steelblue31119180] (axis cs:23.5,0) rectangle (axis cs:24.5,100);
\draw[draw=none,fill=steelblue31119180] (axis cs:24.5,0) rectangle (axis cs:25.5,100);
\draw[draw=none,fill=steelblue31119180] (axis cs:25.5,0) rectangle (axis cs:26.5,100);
\draw[draw=none,fill=steelblue31119180] (axis cs:26.5,0) rectangle (axis cs:27.5,100);
\draw[draw=none,fill=steelblue31119180] (axis cs:27.5,0) rectangle (axis cs:28.5,100);
\draw[draw=none,fill=steelblue31119180] (axis cs:28.5,0) rectangle (axis cs:29.5,100);
\draw[draw=none,fill=steelblue31119180] (axis cs:29.5,0) rectangle (axis cs:30.5,100);
\draw[draw=none,fill=steelblue31119180] (axis cs:30.5,0) rectangle (axis cs:31.5,100);
\draw[draw=none,fill=steelblue31119180] (axis cs:31.5,0) rectangle (axis cs:32.5,100);
\draw[draw=none,fill=steelblue31119180] (axis cs:32.5,0) rectangle (axis cs:33.5,100);
\draw[draw=none,fill=steelblue31119180] (axis cs:33.5,0) rectangle (axis cs:34.5,100);
\draw[draw=none,fill=steelblue31119180] (axis cs:34.5,0) rectangle (axis cs:35.5,100);
\draw[draw=none,fill=steelblue31119180] (axis cs:35.5,0) rectangle (axis cs:36.5,100);
\draw[draw=none,fill=steelblue31119180] (axis cs:36.5,0) rectangle (axis cs:37.5,100);
\draw[draw=none,fill=steelblue31119180] (axis cs:37.5,0) rectangle (axis cs:38.5,100);
\draw[draw=none,fill=steelblue31119180] (axis cs:38.5,0) rectangle (axis cs:39.5,100);
\draw[draw=none,fill=steelblue31119180] (axis cs:39.5,0) rectangle (axis cs:40.5,100);
\draw[draw=none,fill=steelblue31119180] (axis cs:40.5,0) rectangle (axis cs:41.5,100);
\draw[draw=none,fill=steelblue31119180] (axis cs:41.5,0) rectangle (axis cs:42.5,100);
\draw[draw=none,fill=steelblue31119180] (axis cs:42.5,0) rectangle (axis cs:43.5,100);
\draw[draw=none,fill=steelblue31119180] (axis cs:43.5,0) rectangle (axis cs:44.5,100);
\draw[draw=none,fill=steelblue31119180] (axis cs:44.5,0) rectangle (axis cs:45.5,100);
\draw[draw=none,fill=steelblue31119180] (axis cs:45.5,0) rectangle (axis cs:46.5,100);
\draw[draw=none,fill=steelblue31119180] (axis cs:46.5,0) rectangle (axis cs:47.5,100);
\draw[draw=none,fill=steelblue31119180] (axis cs:47.5,0) rectangle (axis cs:48.5,100);
\draw[draw=none,fill=steelblue31119180] (axis cs:48.5,0) rectangle (axis cs:49.5,100);
\draw[draw=none,fill=steelblue31119180] (axis cs:49.5,0) rectangle (axis cs:50.5,100);
\draw[draw=none,fill=steelblue31119180] (axis cs:50.5,0) rectangle (axis cs:51.5,100);
\draw[draw=none,fill=steelblue31119180] (axis cs:51.5,0) rectangle (axis cs:52.5,100);
\draw[draw=none,fill=steelblue31119180] (axis cs:52.5,0) rectangle (axis cs:53.5,100);
\draw[draw=none,fill=steelblue31119180] (axis cs:53.5,0) rectangle (axis cs:54.5,100);
\draw[draw=none,fill=steelblue31119180] (axis cs:54.5,0) rectangle (axis cs:55.5,100);
\draw[draw=none,fill=steelblue31119180] (axis cs:55.5,0) rectangle (axis cs:56.5,100);
\draw[draw=none,fill=steelblue31119180] (axis cs:56.5,0) rectangle (axis cs:57.5,100);
\draw[draw=none,fill=steelblue31119180] (axis cs:57.5,0) rectangle (axis cs:58.5,100);
\draw[draw=none,fill=steelblue31119180] (axis cs:58.5,0) rectangle (axis cs:59.5,100);
\draw[draw=none,fill=steelblue31119180] (axis cs:59.5,0) rectangle (axis cs:60.5,100);
\draw[draw=none,fill=steelblue31119180] (axis cs:60.5,0) rectangle (axis cs:61.5,100);
\draw[draw=none,fill=steelblue31119180] (axis cs:61.5,0) rectangle (axis cs:62.5,100);
\draw[draw=none,fill=steelblue31119180] (axis cs:62.5,0) rectangle (axis cs:63.5,100);
\draw[draw=none,fill=steelblue31119180] (axis cs:63.5,0) rectangle (axis cs:64.5,100);
\draw[draw=none,fill=steelblue31119180] (axis cs:64.5,0) rectangle (axis cs:65.5,100);
\draw[draw=none,fill=steelblue31119180] (axis cs:65.5,0) rectangle (axis cs:66.5,100);
\draw[draw=none,fill=steelblue31119180] (axis cs:66.5,0) rectangle (axis cs:67.5,100);
\draw[draw=none,fill=steelblue31119180] (axis cs:67.5,0) rectangle (axis cs:68.5,100);
\draw[draw=none,fill=steelblue31119180] (axis cs:68.5,0) rectangle (axis cs:69.5,100);
\draw[draw=none,fill=steelblue31119180] (axis cs:69.5,0) rectangle (axis cs:70.5,100);
\draw[draw=none,fill=steelblue31119180] (axis cs:70.5,0) rectangle (axis cs:71.5,100);
\draw[draw=none,fill=steelblue31119180] (axis cs:71.5,0) rectangle (axis cs:72.5,100);
\draw[draw=none,fill=steelblue31119180] (axis cs:72.5,0) rectangle (axis cs:73.5,100);
\draw[draw=none,fill=steelblue31119180] (axis cs:73.5,0) rectangle (axis cs:74.5,100);
\draw[draw=none,fill=steelblue31119180] (axis cs:74.5,0) rectangle (axis cs:75.5,100);
\draw[draw=none,fill=steelblue31119180] (axis cs:75.5,0) rectangle (axis cs:76.5,100);
\draw[draw=none,fill=steelblue31119180] (axis cs:76.5,0) rectangle (axis cs:77.5,100);
\draw[draw=none,fill=steelblue31119180] (axis cs:77.5,0) rectangle (axis cs:78.5,100);
\draw[draw=none,fill=steelblue31119180] (axis cs:78.5,0) rectangle (axis cs:79.5,100);
\draw[draw=none,fill=steelblue31119180] (axis cs:79.5,0) rectangle (axis cs:80.5,100);
\draw[draw=none,fill=steelblue31119180] (axis cs:80.5,0) rectangle (axis cs:81.5,100);
\draw[draw=none,fill=steelblue31119180] (axis cs:81.5,0) rectangle (axis cs:82.5,100);
\draw[draw=none,fill=steelblue31119180] (axis cs:82.5,0) rectangle (axis cs:83.5,100);
\draw[draw=none,fill=steelblue31119180] (axis cs:83.5,0) rectangle (axis cs:84.5,100);
\draw[draw=none,fill=steelblue31119180] (axis cs:84.5,0) rectangle (axis cs:85.5,100);
\draw[draw=none,fill=steelblue31119180] (axis cs:85.5,0) rectangle (axis cs:86.5,100);
\draw[draw=none,fill=steelblue31119180] (axis cs:86.5,0) rectangle (axis cs:87.5,100);
\draw[draw=none,fill=steelblue31119180] (axis cs:87.5,0) rectangle (axis cs:88.5,100);
\draw[draw=none,fill=steelblue31119180] (axis cs:88.5,0) rectangle (axis cs:89.5,100);
\draw[draw=none,fill=steelblue31119180] (axis cs:89.5,0) rectangle (axis cs:90.5,100);
\draw[draw=none,fill=steelblue31119180] (axis cs:90.5,0) rectangle (axis cs:91.5,100);
\draw[draw=none,fill=steelblue31119180] (axis cs:91.5,0) rectangle (axis cs:92.5,100);
\draw[draw=none,fill=steelblue31119180] (axis cs:92.5,0) rectangle (axis cs:93.5,100);
\draw[draw=none,fill=steelblue31119180] (axis cs:93.5,0) rectangle (axis cs:94.5,100);
\draw[draw=none,fill=steelblue31119180] (axis cs:94.5,0) rectangle (axis cs:95.5,100);
\draw[draw=none,fill=steelblue31119180] (axis cs:95.5,0) rectangle (axis cs:96.5,100);
\draw[draw=none,fill=steelblue31119180] (axis cs:96.5,0) rectangle (axis cs:97.5,100);
\draw[draw=none,fill=steelblue31119180] (axis cs:97.5,0) rectangle (axis cs:98.5,100);
\draw[draw=none,fill=steelblue31119180] (axis cs:98.5,0) rectangle (axis cs:99.5,100);
\draw[draw=none,fill=steelblue31119180] (axis cs:99.5,0) rectangle (axis cs:100.5,100);
\draw[draw=none,fill=steelblue31119180] (axis cs:100.5,0) rectangle (axis cs:101.5,100);
\draw[draw=none,fill=steelblue31119180] (axis cs:101.5,0) rectangle (axis cs:102.5,100);
\draw[draw=none,fill=steelblue31119180] (axis cs:102.5,0) rectangle (axis cs:103.5,100);
\draw[draw=none,fill=steelblue31119180] (axis cs:103.5,0) rectangle (axis cs:104.5,100);
\draw[draw=none,fill=steelblue31119180] (axis cs:104.5,0) rectangle (axis cs:105.5,100);
\draw[draw=none,fill=steelblue31119180] (axis cs:105.5,0) rectangle (axis cs:106.5,100);
\draw[draw=none,fill=steelblue31119180] (axis cs:106.5,0) rectangle (axis cs:107.5,100);
\draw[draw=none,fill=steelblue31119180] (axis cs:107.5,0) rectangle (axis cs:108.5,100);
\draw[draw=none,fill=steelblue31119180] (axis cs:108.5,0) rectangle (axis cs:109.5,100);
\draw[draw=none,fill=steelblue31119180] (axis cs:109.5,0) rectangle (axis cs:110.5,100);
\draw[draw=none,fill=steelblue31119180] (axis cs:110.5,0) rectangle (axis cs:111.5,100);
\draw[draw=none,fill=steelblue31119180] (axis cs:111.5,0) rectangle (axis cs:112.5,100);
\draw[draw=none,fill=steelblue31119180] (axis cs:112.5,0) rectangle (axis cs:113.5,100);
\draw[draw=none,fill=steelblue31119180] (axis cs:113.5,0) rectangle (axis cs:114.5,100);
\draw[draw=none,fill=steelblue31119180] (axis cs:114.5,0) rectangle (axis cs:115.5,100);
\draw[draw=none,fill=steelblue31119180] (axis cs:115.5,0) rectangle (axis cs:116.5,100);
\draw[draw=none,fill=steelblue31119180] (axis cs:116.5,0) rectangle (axis cs:117.5,100);
\draw[draw=none,fill=steelblue31119180] (axis cs:117.5,0) rectangle (axis cs:118.5,100);
\draw[draw=none,fill=steelblue31119180] (axis cs:118.5,0) rectangle (axis cs:119.5,100);
\draw[draw=none,fill=steelblue31119180] (axis cs:119.5,0) rectangle (axis cs:120.5,100);
\draw[draw=none,fill=steelblue31119180] (axis cs:120.5,0) rectangle (axis cs:121.5,100);
\draw[draw=none,fill=steelblue31119180] (axis cs:121.5,0) rectangle (axis cs:122.5,100);
\draw[draw=none,fill=steelblue31119180] (axis cs:122.5,0) rectangle (axis cs:123.5,100);
\draw[draw=none,fill=steelblue31119180] (axis cs:123.5,0) rectangle (axis cs:124.5,100);
\draw[draw=none,fill=steelblue31119180] (axis cs:124.5,0) rectangle (axis cs:125.5,100);
\draw[draw=none,fill=steelblue31119180] (axis cs:125.5,0) rectangle (axis cs:126.5,100);
\draw[draw=none,fill=steelblue31119180] (axis cs:126.5,0) rectangle (axis cs:127.5,100);
\draw[draw=none,fill=steelblue31119180] (axis cs:127.5,0) rectangle (axis cs:128.5,100);
\draw[draw=none,fill=steelblue31119180] (axis cs:128.5,0) rectangle (axis cs:129.5,100);
\draw[draw=none,fill=steelblue31119180] (axis cs:129.5,0) rectangle (axis cs:130.5,100);
\draw[draw=none,fill=steelblue31119180] (axis cs:130.5,0) rectangle (axis cs:131.5,100);
\draw[draw=none,fill=steelblue31119180] (axis cs:131.5,0) rectangle (axis cs:132.5,100);
\draw[draw=none,fill=steelblue31119180] (axis cs:132.5,0) rectangle (axis cs:133.5,100);
\draw[draw=none,fill=steelblue31119180] (axis cs:133.5,0) rectangle (axis cs:134.5,100);
\draw[draw=none,fill=steelblue31119180] (axis cs:134.5,0) rectangle (axis cs:135.5,100);
\draw[draw=none,fill=steelblue31119180] (axis cs:135.5,0) rectangle (axis cs:136.5,100);
\draw[draw=none,fill=steelblue31119180] (axis cs:136.5,0) rectangle (axis cs:137.5,100);
\draw[draw=none,fill=steelblue31119180] (axis cs:137.5,0) rectangle (axis cs:138.5,100);
\draw[draw=none,fill=steelblue31119180] (axis cs:138.5,0) rectangle (axis cs:139.5,100);
\draw[draw=none,fill=steelblue31119180] (axis cs:139.5,0) rectangle (axis cs:140.5,100);
\draw[draw=none,fill=steelblue31119180] (axis cs:140.5,0) rectangle (axis cs:141.5,100);
\draw[draw=none,fill=steelblue31119180] (axis cs:141.5,0) rectangle (axis cs:142.5,100);
\draw[draw=none,fill=steelblue31119180] (axis cs:142.5,0) rectangle (axis cs:143.5,100);
\draw[draw=none,fill=steelblue31119180] (axis cs:143.5,0) rectangle (axis cs:144.5,100);
\draw[draw=none,fill=steelblue31119180] (axis cs:144.5,0) rectangle (axis cs:145.5,100);
\draw[draw=none,fill=steelblue31119180] (axis cs:145.5,0) rectangle (axis cs:146.5,100);
\draw[draw=none,fill=steelblue31119180] (axis cs:146.5,0) rectangle (axis cs:147.5,100);
\draw[draw=none,fill=steelblue31119180] (axis cs:147.5,0) rectangle (axis cs:148.5,100);
\draw[draw=none,fill=steelblue31119180] (axis cs:148.5,0) rectangle (axis cs:149.5,100);
\draw[draw=none,fill=steelblue31119180] (axis cs:149.5,0) rectangle (axis cs:150.5,100);
\draw[draw=none,fill=steelblue31119180] (axis cs:150.5,0) rectangle (axis cs:151.5,100);
\draw[draw=none,fill=steelblue31119180] (axis cs:151.5,0) rectangle (axis cs:152.5,100);
\draw[draw=none,fill=steelblue31119180] (axis cs:152.5,0) rectangle (axis cs:153.5,100);
\draw[draw=none,fill=steelblue31119180] (axis cs:153.5,0) rectangle (axis cs:154.5,100);
\draw[draw=none,fill=steelblue31119180] (axis cs:154.5,0) rectangle (axis cs:155.5,100);
\draw[draw=none,fill=steelblue31119180] (axis cs:155.5,0) rectangle (axis cs:156.5,100);
\draw[draw=none,fill=steelblue31119180] (axis cs:156.5,0) rectangle (axis cs:157.5,100);
\draw[draw=none,fill=steelblue31119180] (axis cs:157.5,0) rectangle (axis cs:158.5,100);
\draw[draw=none,fill=steelblue31119180] (axis cs:158.5,0) rectangle (axis cs:159.5,100);
\draw[draw=none,fill=steelblue31119180] (axis cs:159.5,0) rectangle (axis cs:160.5,100);
\draw[draw=none,fill=steelblue31119180] (axis cs:160.5,0) rectangle (axis cs:161.5,100);
\draw[draw=none,fill=steelblue31119180] (axis cs:161.5,0) rectangle (axis cs:162.5,100);
\draw[draw=none,fill=steelblue31119180] (axis cs:162.5,0) rectangle (axis cs:163.5,100);
\draw[draw=none,fill=steelblue31119180] (axis cs:163.5,0) rectangle (axis cs:164.5,100);
\draw[draw=none,fill=steelblue31119180] (axis cs:164.5,0) rectangle (axis cs:165.5,100);
\draw[draw=none,fill=steelblue31119180] (axis cs:165.5,0) rectangle (axis cs:166.5,100);
\draw[draw=none,fill=steelblue31119180] (axis cs:166.5,0) rectangle (axis cs:167.5,100);
\draw[draw=none,fill=steelblue31119180] (axis cs:167.5,0) rectangle (axis cs:168.5,100);
\draw[draw=none,fill=steelblue31119180] (axis cs:168.5,0) rectangle (axis cs:169.5,100);
\draw[draw=none,fill=steelblue31119180] (axis cs:169.5,0) rectangle (axis cs:170.5,100);
\draw[draw=none,fill=steelblue31119180] (axis cs:170.5,0) rectangle (axis cs:171.5,100);
\draw[draw=none,fill=steelblue31119180] (axis cs:171.5,0) rectangle (axis cs:172.5,100);
\draw[draw=none,fill=steelblue31119180] (axis cs:172.5,0) rectangle (axis cs:173.5,100);
\draw[draw=none,fill=steelblue31119180] (axis cs:173.5,0) rectangle (axis cs:174.5,100);
\draw[draw=none,fill=steelblue31119180] (axis cs:174.5,0) rectangle (axis cs:175.5,100);
\draw[draw=none,fill=steelblue31119180] (axis cs:175.5,0) rectangle (axis cs:176.5,100);
\draw[draw=none,fill=steelblue31119180] (axis cs:176.5,0) rectangle (axis cs:177.5,100);
\draw[draw=none,fill=steelblue31119180] (axis cs:177.5,0) rectangle (axis cs:178.5,100);
\draw[draw=none,fill=steelblue31119180] (axis cs:178.5,0) rectangle (axis cs:179.5,100);
\draw[draw=none,fill=steelblue31119180] (axis cs:179.5,0) rectangle (axis cs:180.5,100);
\draw[draw=none,fill=steelblue31119180] (axis cs:180.5,0) rectangle (axis cs:181.5,100);
\draw[draw=none,fill=steelblue31119180] (axis cs:181.5,0) rectangle (axis cs:182.5,100);
\draw[draw=none,fill=steelblue31119180] (axis cs:182.5,0) rectangle (axis cs:183.5,100);
\draw[draw=none,fill=steelblue31119180] (axis cs:183.5,0) rectangle (axis cs:184.5,100);
\draw[draw=none,fill=steelblue31119180] (axis cs:184.5,0) rectangle (axis cs:185.5,100);
\draw[draw=none,fill=steelblue31119180] (axis cs:185.5,0) rectangle (axis cs:186.5,100);
\draw[draw=none,fill=steelblue31119180] (axis cs:186.5,0) rectangle (axis cs:187.5,100);
\draw[draw=none,fill=steelblue31119180] (axis cs:187.5,0) rectangle (axis cs:188.5,100);
\draw[draw=none,fill=steelblue31119180] (axis cs:188.5,0) rectangle (axis cs:189.5,100);
\draw[draw=none,fill=steelblue31119180] (axis cs:189.5,0) rectangle (axis cs:190.5,100);
\draw[draw=none,fill=steelblue31119180] (axis cs:190.5,0) rectangle (axis cs:191.5,100);
\draw[draw=none,fill=steelblue31119180] (axis cs:191.5,0) rectangle (axis cs:192.5,100);
\draw[draw=none,fill=steelblue31119180] (axis cs:192.5,0) rectangle (axis cs:193.5,100);
\draw[draw=none,fill=steelblue31119180] (axis cs:193.5,0) rectangle (axis cs:194.5,100);
\draw[draw=none,fill=steelblue31119180] (axis cs:194.5,0) rectangle (axis cs:195.5,100);
\draw[draw=none,fill=steelblue31119180] (axis cs:195.5,0) rectangle (axis cs:196.5,100);
\draw[draw=none,fill=steelblue31119180] (axis cs:196.5,0) rectangle (axis cs:197.5,100);
\draw[draw=none,fill=steelblue31119180] (axis cs:197.5,0) rectangle (axis cs:198.5,100);
\draw[draw=none,fill=steelblue31119180] (axis cs:198.5,0) rectangle (axis cs:199.5,100);
\draw[draw=none,fill=steelblue31119180] (axis cs:199.5,0) rectangle (axis cs:200.5,100);
\draw[draw=none,fill=steelblue31119180] (axis cs:200.5,0) rectangle (axis cs:201.5,100);
\draw[draw=none,fill=steelblue31119180] (axis cs:201.5,0) rectangle (axis cs:202.5,100);
\draw[draw=none,fill=steelblue31119180] (axis cs:202.5,0) rectangle (axis cs:203.5,100);
\draw[draw=none,fill=steelblue31119180] (axis cs:203.5,0) rectangle (axis cs:204.5,100);
\draw[draw=none,fill=steelblue31119180] (axis cs:204.5,0) rectangle (axis cs:205.5,100);
\draw[draw=none,fill=steelblue31119180] (axis cs:205.5,0) rectangle (axis cs:206.5,100);
\draw[draw=none,fill=steelblue31119180] (axis cs:206.5,0) rectangle (axis cs:207.5,100);
\draw[draw=none,fill=steelblue31119180] (axis cs:207.5,0) rectangle (axis cs:208.5,100);
\draw[draw=none,fill=steelblue31119180] (axis cs:208.5,0) rectangle (axis cs:209.5,100);
\draw[draw=none,fill=steelblue31119180] (axis cs:209.5,0) rectangle (axis cs:210.5,100);
\draw[draw=none,fill=steelblue31119180] (axis cs:210.5,0) rectangle (axis cs:211.5,100);
\draw[draw=none,fill=steelblue31119180] (axis cs:211.5,0) rectangle (axis cs:212.5,100);
\draw[draw=none,fill=steelblue31119180] (axis cs:212.5,0) rectangle (axis cs:213.5,100);
\draw[draw=none,fill=steelblue31119180] (axis cs:213.5,0) rectangle (axis cs:214.5,100);
\draw[draw=none,fill=steelblue31119180] (axis cs:214.5,0) rectangle (axis cs:215.5,100);
\draw[draw=none,fill=steelblue31119180] (axis cs:215.5,0) rectangle (axis cs:216.5,100);
\draw[draw=none,fill=steelblue31119180] (axis cs:216.5,0) rectangle (axis cs:217.5,100);
\draw[draw=none,fill=steelblue31119180] (axis cs:217.5,0) rectangle (axis cs:218.5,100);
\draw[draw=none,fill=steelblue31119180] (axis cs:218.5,0) rectangle (axis cs:219.5,100);
\draw[draw=none,fill=steelblue31119180] (axis cs:219.5,0) rectangle (axis cs:220.5,100);
\draw[draw=none,fill=steelblue31119180] (axis cs:220.5,0) rectangle (axis cs:221.5,100);
\draw[draw=none,fill=steelblue31119180] (axis cs:221.5,0) rectangle (axis cs:222.5,100);
\draw[draw=none,fill=steelblue31119180] (axis cs:222.5,0) rectangle (axis cs:223.5,100);
\draw[draw=none,fill=steelblue31119180] (axis cs:223.5,0) rectangle (axis cs:224.5,100);
\draw[draw=none,fill=steelblue31119180] (axis cs:224.5,0) rectangle (axis cs:225.5,100);
\draw[draw=none,fill=steelblue31119180] (axis cs:225.5,0) rectangle (axis cs:226.5,100);
\draw[draw=none,fill=steelblue31119180] (axis cs:226.5,0) rectangle (axis cs:227.5,100);
\draw[draw=none,fill=steelblue31119180] (axis cs:227.5,0) rectangle (axis cs:228.5,100);
\draw[draw=none,fill=steelblue31119180] (axis cs:228.5,0) rectangle (axis cs:229.5,100);
\draw[draw=none,fill=steelblue31119180] (axis cs:229.5,0) rectangle (axis cs:230.5,100);
\draw[draw=none,fill=steelblue31119180] (axis cs:230.5,0) rectangle (axis cs:231.5,100);
\draw[draw=none,fill=steelblue31119180] (axis cs:231.5,0) rectangle (axis cs:232.5,100);
\draw[draw=none,fill=steelblue31119180] (axis cs:232.5,0) rectangle (axis cs:233.5,100);
\draw[draw=none,fill=steelblue31119180] (axis cs:233.5,0) rectangle (axis cs:234.5,100);
\draw[draw=none,fill=steelblue31119180] (axis cs:234.5,0) rectangle (axis cs:235.5,100);
\draw[draw=none,fill=steelblue31119180] (axis cs:235.5,0) rectangle (axis cs:236.5,100);
\draw[draw=none,fill=steelblue31119180] (axis cs:236.5,0) rectangle (axis cs:237.5,100);
\draw[draw=none,fill=steelblue31119180] (axis cs:237.5,0) rectangle (axis cs:238.5,100);
\draw[draw=none,fill=steelblue31119180] (axis cs:238.5,0) rectangle (axis cs:239.5,100);
\draw[draw=none,fill=steelblue31119180] (axis cs:239.5,0) rectangle (axis cs:240.5,100);
\draw[draw=none,fill=steelblue31119180] (axis cs:240.5,0) rectangle (axis cs:241.5,100);
\draw[draw=none,fill=steelblue31119180] (axis cs:241.5,0) rectangle (axis cs:242.5,100);
\draw[draw=none,fill=steelblue31119180] (axis cs:242.5,0) rectangle (axis cs:243.5,100);
\draw[draw=none,fill=steelblue31119180] (axis cs:243.5,0) rectangle (axis cs:244.5,100);
\draw[draw=none,fill=steelblue31119180] (axis cs:244.5,0) rectangle (axis cs:245.5,100);
\draw[draw=none,fill=steelblue31119180] (axis cs:245.5,0) rectangle (axis cs:246.5,100);
\draw[draw=none,fill=steelblue31119180] (axis cs:246.5,0) rectangle (axis cs:247.5,100);
\draw[draw=none,fill=steelblue31119180] (axis cs:247.5,0) rectangle (axis cs:248.5,100);
\draw[draw=none,fill=steelblue31119180] (axis cs:248.5,0) rectangle (axis cs:249.5,100);
\draw[draw=none,fill=steelblue31119180] (axis cs:249.5,0) rectangle (axis cs:250.5,100);
\draw[draw=none,fill=steelblue31119180] (axis cs:250.5,0) rectangle (axis cs:251.5,100);
\draw[draw=none,fill=steelblue31119180] (axis cs:251.5,0) rectangle (axis cs:252.5,100);
\draw[draw=none,fill=steelblue31119180] (axis cs:252.5,0) rectangle (axis cs:253.5,100);
\draw[draw=none,fill=steelblue31119180] (axis cs:253.5,0) rectangle (axis cs:254.5,100);
\draw[draw=none,fill=steelblue31119180] (axis cs:254.5,0) rectangle (axis cs:255.5,100);
\draw[draw=none,fill=steelblue31119180] (axis cs:255.5,0) rectangle (axis cs:256.5,100);
\draw[draw=none,fill=steelblue31119180] (axis cs:256.5,0) rectangle (axis cs:257.5,100);
\draw[draw=none,fill=steelblue31119180] (axis cs:257.5,0) rectangle (axis cs:258.5,100);
\draw[draw=none,fill=steelblue31119180] (axis cs:258.5,0) rectangle (axis cs:259.5,100);
\draw[draw=none,fill=steelblue31119180] (axis cs:259.5,0) rectangle (axis cs:260.5,100);
\draw[draw=none,fill=steelblue31119180] (axis cs:260.5,0) rectangle (axis cs:261.5,100);
\draw[draw=none,fill=steelblue31119180] (axis cs:261.5,0) rectangle (axis cs:262.5,100);
\draw[draw=none,fill=steelblue31119180] (axis cs:262.5,0) rectangle (axis cs:263.5,100);
\draw[draw=none,fill=steelblue31119180] (axis cs:263.5,0) rectangle (axis cs:264.5,100);
\draw[draw=none,fill=steelblue31119180] (axis cs:264.5,0) rectangle (axis cs:265.5,100);
\draw[draw=none,fill=steelblue31119180] (axis cs:265.5,0) rectangle (axis cs:266.5,100);
\draw[draw=none,fill=steelblue31119180] (axis cs:266.5,0) rectangle (axis cs:267.5,100);
\draw[draw=none,fill=steelblue31119180] (axis cs:267.5,0) rectangle (axis cs:268.5,100);
\draw[draw=none,fill=steelblue31119180] (axis cs:268.5,0) rectangle (axis cs:269.5,100);
\draw[draw=none,fill=steelblue31119180] (axis cs:269.5,0) rectangle (axis cs:270.5,100);
\draw[draw=none,fill=steelblue31119180] (axis cs:270.5,0) rectangle (axis cs:271.5,100);
\draw[draw=none,fill=steelblue31119180] (axis cs:271.5,0) rectangle (axis cs:272.5,100);
\draw[draw=none,fill=steelblue31119180] (axis cs:272.5,0) rectangle (axis cs:273.5,100);
\draw[draw=none,fill=steelblue31119180] (axis cs:273.5,0) rectangle (axis cs:274.5,100);
\draw[draw=none,fill=steelblue31119180] (axis cs:274.5,0) rectangle (axis cs:275.5,100);
\draw[draw=none,fill=steelblue31119180] (axis cs:275.5,0) rectangle (axis cs:276.5,100);
\draw[draw=none,fill=steelblue31119180] (axis cs:276.5,0) rectangle (axis cs:277.5,100);
\draw[draw=none,fill=steelblue31119180] (axis cs:277.5,0) rectangle (axis cs:278.5,100);
\draw[draw=none,fill=steelblue31119180] (axis cs:278.5,0) rectangle (axis cs:279.5,100);
\draw[draw=none,fill=steelblue31119180] (axis cs:279.5,0) rectangle (axis cs:280.5,100);
\draw[draw=none,fill=steelblue31119180] (axis cs:280.5,0) rectangle (axis cs:281.5,100);
\draw[draw=none,fill=steelblue31119180] (axis cs:281.5,0) rectangle (axis cs:282.5,100);
\draw[draw=none,fill=steelblue31119180] (axis cs:282.5,0) rectangle (axis cs:283.5,100);
\draw[draw=none,fill=steelblue31119180] (axis cs:283.5,0) rectangle (axis cs:284.5,100);
\draw[draw=none,fill=steelblue31119180] (axis cs:284.5,0) rectangle (axis cs:285.5,100);
\draw[draw=none,fill=steelblue31119180] (axis cs:285.5,0) rectangle (axis cs:286.5,100);
\draw[draw=none,fill=steelblue31119180] (axis cs:286.5,0) rectangle (axis cs:287.5,100);
\draw[draw=none,fill=steelblue31119180] (axis cs:287.5,0) rectangle (axis cs:288.5,100);
\draw[draw=none,fill=steelblue31119180] (axis cs:288.5,0) rectangle (axis cs:289.5,100);
\draw[draw=none,fill=steelblue31119180] (axis cs:289.5,0) rectangle (axis cs:290.5,100);
\draw[draw=none,fill=steelblue31119180] (axis cs:290.5,0) rectangle (axis cs:291.5,100);
\draw[draw=none,fill=steelblue31119180] (axis cs:291.5,0) rectangle (axis cs:292.5,100);
\draw[draw=none,fill=steelblue31119180] (axis cs:292.5,0) rectangle (axis cs:293.5,100);
\draw[draw=none,fill=steelblue31119180] (axis cs:293.5,0) rectangle (axis cs:294.5,100);
\draw[draw=none,fill=steelblue31119180] (axis cs:294.5,0) rectangle (axis cs:295.5,100);
\draw[draw=none,fill=steelblue31119180] (axis cs:295.5,0) rectangle (axis cs:296.5,100);
\draw[draw=none,fill=steelblue31119180] (axis cs:296.5,0) rectangle (axis cs:297.5,100);
\draw[draw=none,fill=steelblue31119180] (axis cs:297.5,0) rectangle (axis cs:298.5,100);
\draw[draw=none,fill=steelblue31119180] (axis cs:298.5,0) rectangle (axis cs:299.5,100);
\draw[draw=none,fill=steelblue31119180] (axis cs:299.5,0) rectangle (axis cs:300.5,100);
\draw[draw=none,fill=steelblue31119180] (axis cs:300.5,0) rectangle (axis cs:301.5,100);
\draw[draw=none,fill=steelblue31119180] (axis cs:301.5,0) rectangle (axis cs:302.5,100);
\draw[draw=none,fill=steelblue31119180] (axis cs:302.5,0) rectangle (axis cs:303.5,100);
\draw[draw=none,fill=steelblue31119180] (axis cs:303.5,0) rectangle (axis cs:304.5,100);
\draw[draw=none,fill=steelblue31119180] (axis cs:304.5,0) rectangle (axis cs:305.5,100);
\draw[draw=none,fill=steelblue31119180] (axis cs:305.5,0) rectangle (axis cs:306.5,100);
\draw[draw=none,fill=steelblue31119180] (axis cs:306.5,0) rectangle (axis cs:307.5,100);
\draw[draw=none,fill=steelblue31119180] (axis cs:307.5,0) rectangle (axis cs:308.5,100);
\draw[draw=none,fill=steelblue31119180] (axis cs:308.5,0) rectangle (axis cs:309.5,100);
\draw[draw=none,fill=steelblue31119180] (axis cs:309.5,0) rectangle (axis cs:310.5,100);
\draw[draw=none,fill=steelblue31119180] (axis cs:310.5,0) rectangle (axis cs:311.5,100);
\draw[draw=none,fill=steelblue31119180] (axis cs:311.5,0) rectangle (axis cs:312.5,100);
\draw[draw=none,fill=steelblue31119180] (axis cs:312.5,0) rectangle (axis cs:313.5,100);
\draw[draw=none,fill=steelblue31119180] (axis cs:313.5,0) rectangle (axis cs:314.5,100);
\draw[draw=none,fill=steelblue31119180] (axis cs:314.5,0) rectangle (axis cs:315.5,100);
\draw[draw=none,fill=steelblue31119180] (axis cs:315.5,0) rectangle (axis cs:316.5,100);
\draw[draw=none,fill=steelblue31119180] (axis cs:316.5,0) rectangle (axis cs:317.5,100);
\draw[draw=none,fill=steelblue31119180] (axis cs:317.5,0) rectangle (axis cs:318.5,100);
\draw[draw=none,fill=steelblue31119180] (axis cs:318.5,0) rectangle (axis cs:319.5,100);
\draw[draw=none,fill=steelblue31119180] (axis cs:319.5,0) rectangle (axis cs:320.5,100);
\draw[draw=none,fill=steelblue31119180] (axis cs:320.5,0) rectangle (axis cs:321.5,100);
\draw[draw=none,fill=steelblue31119180] (axis cs:321.5,0) rectangle (axis cs:322.5,100);
\draw[draw=none,fill=steelblue31119180] (axis cs:322.5,0) rectangle (axis cs:323.5,100);
\draw[draw=none,fill=steelblue31119180] (axis cs:323.5,0) rectangle (axis cs:324.5,100);
\draw[draw=none,fill=steelblue31119180] (axis cs:324.5,0) rectangle (axis cs:325.5,100);
\draw[draw=none,fill=steelblue31119180] (axis cs:325.5,0) rectangle (axis cs:326.5,100);
\draw[draw=none,fill=steelblue31119180] (axis cs:326.5,0) rectangle (axis cs:327.5,100);
\draw[draw=none,fill=steelblue31119180] (axis cs:327.5,0) rectangle (axis cs:328.5,100);
\draw[draw=none,fill=steelblue31119180] (axis cs:328.5,0) rectangle (axis cs:329.5,100);
\draw[draw=none,fill=steelblue31119180] (axis cs:329.5,0) rectangle (axis cs:330.5,100);
\draw[draw=none,fill=steelblue31119180] (axis cs:330.5,0) rectangle (axis cs:331.5,100);
\draw[draw=none,fill=steelblue31119180] (axis cs:331.5,0) rectangle (axis cs:332.5,100);
\draw[draw=none,fill=steelblue31119180] (axis cs:332.5,0) rectangle (axis cs:333.5,100);
\draw[draw=none,fill=steelblue31119180] (axis cs:333.5,0) rectangle (axis cs:334.5,100);
\draw[draw=none,fill=steelblue31119180] (axis cs:334.5,0) rectangle (axis cs:335.5,100);
\draw[draw=none,fill=steelblue31119180] (axis cs:335.5,0) rectangle (axis cs:336.5,100);
\draw[draw=none,fill=steelblue31119180] (axis cs:336.5,0) rectangle (axis cs:337.5,100);
\draw[draw=none,fill=steelblue31119180] (axis cs:337.5,0) rectangle (axis cs:338.5,100);
\draw[draw=none,fill=steelblue31119180] (axis cs:338.5,0) rectangle (axis cs:339.5,100);
\draw[draw=none,fill=steelblue31119180] (axis cs:339.5,0) rectangle (axis cs:340.5,100);
\draw[draw=none,fill=steelblue31119180] (axis cs:340.5,0) rectangle (axis cs:341.5,100);
\draw[draw=none,fill=steelblue31119180] (axis cs:341.5,0) rectangle (axis cs:342.5,100);
\draw[draw=none,fill=steelblue31119180] (axis cs:342.5,0) rectangle (axis cs:343.5,100);
\draw[draw=none,fill=steelblue31119180] (axis cs:343.5,0) rectangle (axis cs:344.5,100);
\draw[draw=none,fill=steelblue31119180] (axis cs:344.5,0) rectangle (axis cs:345.5,100);
\draw[draw=none,fill=steelblue31119180] (axis cs:345.5,0) rectangle (axis cs:346.5,100);
\draw[draw=none,fill=steelblue31119180] (axis cs:346.5,0) rectangle (axis cs:347.5,100);
\draw[draw=none,fill=steelblue31119180] (axis cs:347.5,0) rectangle (axis cs:348.5,100);
\draw[draw=none,fill=steelblue31119180] (axis cs:348.5,0) rectangle (axis cs:349.5,100);
\draw[draw=none,fill=steelblue31119180] (axis cs:349.5,0) rectangle (axis cs:350.5,100);
\draw[draw=none,fill=steelblue31119180] (axis cs:350.5,0) rectangle (axis cs:351.5,100);
\draw[draw=none,fill=steelblue31119180] (axis cs:351.5,0) rectangle (axis cs:352.5,100);
\draw[draw=none,fill=steelblue31119180] (axis cs:352.5,0) rectangle (axis cs:353.5,100);
\draw[draw=none,fill=steelblue31119180] (axis cs:353.5,0) rectangle (axis cs:354.5,100);
\draw[draw=none,fill=steelblue31119180] (axis cs:354.5,0) rectangle (axis cs:355.5,100);
\draw[draw=none,fill=steelblue31119180] (axis cs:355.5,0) rectangle (axis cs:356.5,100);
\draw[draw=none,fill=steelblue31119180] (axis cs:356.5,0) rectangle (axis cs:357.5,100);
\draw[draw=none,fill=steelblue31119180] (axis cs:357.5,0) rectangle (axis cs:358.5,100);
\draw[draw=none,fill=steelblue31119180] (axis cs:358.5,0) rectangle (axis cs:359.5,100);
\draw[draw=none,fill=steelblue31119180] (axis cs:359.5,0) rectangle (axis cs:360.5,100);
\draw[draw=none,fill=steelblue31119180] (axis cs:360.5,0) rectangle (axis cs:361.5,100);
\draw[draw=none,fill=steelblue31119180] (axis cs:361.5,0) rectangle (axis cs:362.5,100);
\draw[draw=none,fill=steelblue31119180] (axis cs:362.5,0) rectangle (axis cs:363.5,100);
\draw[draw=none,fill=steelblue31119180] (axis cs:363.5,0) rectangle (axis cs:364.5,100);
\draw[draw=none,fill=steelblue31119180] (axis cs:364.5,0) rectangle (axis cs:365.5,100);
\draw[draw=none,fill=steelblue31119180] (axis cs:365.5,0) rectangle (axis cs:366.5,100);
\draw[draw=none,fill=steelblue31119180] (axis cs:366.5,0) rectangle (axis cs:367.5,100);
\draw[draw=none,fill=steelblue31119180] (axis cs:367.5,0) rectangle (axis cs:368.5,100);
\draw[draw=none,fill=steelblue31119180] (axis cs:368.5,0) rectangle (axis cs:369.5,100);
\draw[draw=none,fill=steelblue31119180] (axis cs:369.5,0) rectangle (axis cs:370.5,100);
\draw[draw=none,fill=steelblue31119180] (axis cs:370.5,0) rectangle (axis cs:371.5,100);
\draw[draw=none,fill=steelblue31119180] (axis cs:371.5,0) rectangle (axis cs:372.5,100);
\draw[draw=none,fill=steelblue31119180] (axis cs:372.5,0) rectangle (axis cs:373.5,100);
\draw[draw=none,fill=steelblue31119180] (axis cs:373.5,0) rectangle (axis cs:374.5,100);
\draw[draw=none,fill=steelblue31119180] (axis cs:374.5,0) rectangle (axis cs:375.5,100);
\draw[draw=none,fill=steelblue31119180] (axis cs:375.5,0) rectangle (axis cs:376.5,100);
\draw[draw=none,fill=steelblue31119180] (axis cs:376.5,0) rectangle (axis cs:377.5,100);
\draw[draw=none,fill=steelblue31119180] (axis cs:377.5,0) rectangle (axis cs:378.5,100);
\draw[draw=none,fill=steelblue31119180] (axis cs:378.5,0) rectangle (axis cs:379.5,100);
\draw[draw=none,fill=steelblue31119180] (axis cs:379.5,0) rectangle (axis cs:380.5,100);
\draw[draw=none,fill=steelblue31119180] (axis cs:380.5,0) rectangle (axis cs:381.5,100);
\draw[draw=none,fill=steelblue31119180] (axis cs:381.5,0) rectangle (axis cs:382.5,100);
\draw[draw=none,fill=steelblue31119180] (axis cs:382.5,0) rectangle (axis cs:383.5,100);
\draw[draw=none,fill=steelblue31119180] (axis cs:383.5,0) rectangle (axis cs:384.5,100);
\draw[draw=none,fill=steelblue31119180] (axis cs:384.5,0) rectangle (axis cs:385.5,100);
\draw[draw=none,fill=steelblue31119180] (axis cs:385.5,0) rectangle (axis cs:386.5,100);
\draw[draw=none,fill=steelblue31119180] (axis cs:386.5,0) rectangle (axis cs:387.5,100);
\draw[draw=none,fill=steelblue31119180] (axis cs:387.5,0) rectangle (axis cs:388.5,100);
\draw[draw=none,fill=steelblue31119180] (axis cs:388.5,0) rectangle (axis cs:389.5,100);
\draw[draw=none,fill=steelblue31119180] (axis cs:389.5,0) rectangle (axis cs:390.5,100);
\draw[draw=none,fill=steelblue31119180] (axis cs:390.5,0) rectangle (axis cs:391.5,100);
\draw[draw=none,fill=steelblue31119180] (axis cs:391.5,0) rectangle (axis cs:392.5,100);
\draw[draw=none,fill=steelblue31119180] (axis cs:392.5,0) rectangle (axis cs:393.5,100);
\draw[draw=none,fill=steelblue31119180] (axis cs:393.5,0) rectangle (axis cs:394.5,100);
\draw[draw=none,fill=steelblue31119180] (axis cs:394.5,0) rectangle (axis cs:395.5,100);
\draw[draw=none,fill=steelblue31119180] (axis cs:395.5,0) rectangle (axis cs:396.5,100);
\draw[draw=none,fill=steelblue31119180] (axis cs:396.5,0) rectangle (axis cs:397.5,100);
\draw[draw=none,fill=steelblue31119180] (axis cs:397.5,0) rectangle (axis cs:398.5,100);
\draw[draw=none,fill=steelblue31119180] (axis cs:398.5,0) rectangle (axis cs:399.5,100);
\draw[draw=none,fill=steelblue31119180] (axis cs:399.5,0) rectangle (axis cs:400.5,100);
\draw[draw=none,fill=steelblue31119180] (axis cs:400.5,0) rectangle (axis cs:401.5,100);
\draw[draw=none,fill=steelblue31119180] (axis cs:401.5,0) rectangle (axis cs:402.5,100);
\draw[draw=none,fill=steelblue31119180] (axis cs:402.5,0) rectangle (axis cs:403.5,100);
\draw[draw=none,fill=steelblue31119180] (axis cs:403.5,0) rectangle (axis cs:404.5,100);
\draw[draw=none,fill=steelblue31119180] (axis cs:404.5,0) rectangle (axis cs:405.5,100);
\draw[draw=none,fill=steelblue31119180] (axis cs:405.5,0) rectangle (axis cs:406.5,100);
\draw[draw=none,fill=steelblue31119180] (axis cs:406.5,0) rectangle (axis cs:407.5,100);
\draw[draw=none,fill=steelblue31119180] (axis cs:407.5,0) rectangle (axis cs:408.5,100);
\draw[draw=none,fill=steelblue31119180] (axis cs:408.5,0) rectangle (axis cs:409.5,100);
\draw[draw=none,fill=steelblue31119180] (axis cs:409.5,0) rectangle (axis cs:410.5,100);
\draw[draw=none,fill=steelblue31119180] (axis cs:410.5,0) rectangle (axis cs:411.5,100);
\draw[draw=none,fill=steelblue31119180] (axis cs:411.5,0) rectangle (axis cs:412.5,100);
\draw[draw=none,fill=steelblue31119180] (axis cs:412.5,0) rectangle (axis cs:413.5,100);
\draw[draw=none,fill=steelblue31119180] (axis cs:413.5,0) rectangle (axis cs:414.5,100);
\draw[draw=none,fill=steelblue31119180] (axis cs:414.5,0) rectangle (axis cs:415.5,100);
\draw[draw=none,fill=steelblue31119180] (axis cs:415.5,0) rectangle (axis cs:416.5,100);
\draw[draw=none,fill=steelblue31119180] (axis cs:416.5,0) rectangle (axis cs:417.5,100);
\draw[draw=none,fill=steelblue31119180] (axis cs:417.5,0) rectangle (axis cs:418.5,100);
\draw[draw=none,fill=steelblue31119180] (axis cs:418.5,0) rectangle (axis cs:419.5,100);
\draw[draw=none,fill=steelblue31119180] (axis cs:419.5,0) rectangle (axis cs:420.5,100);
\draw[draw=none,fill=steelblue31119180] (axis cs:420.5,0) rectangle (axis cs:421.5,100);
\draw[draw=none,fill=steelblue31119180] (axis cs:421.5,0) rectangle (axis cs:422.5,100);
\draw[draw=none,fill=steelblue31119180] (axis cs:422.5,0) rectangle (axis cs:423.5,100);
\draw[draw=none,fill=steelblue31119180] (axis cs:423.5,0) rectangle (axis cs:424.5,100);
\draw[draw=none,fill=steelblue31119180] (axis cs:424.5,0) rectangle (axis cs:425.5,100);
\draw[draw=none,fill=steelblue31119180] (axis cs:425.5,0) rectangle (axis cs:426.5,100);
\draw[draw=none,fill=steelblue31119180] (axis cs:426.5,0) rectangle (axis cs:427.5,100);
\draw[draw=none,fill=steelblue31119180] (axis cs:427.5,0) rectangle (axis cs:428.5,100);
\draw[draw=none,fill=steelblue31119180] (axis cs:428.5,0) rectangle (axis cs:429.5,100);
\draw[draw=none,fill=steelblue31119180] (axis cs:429.5,0) rectangle (axis cs:430.5,100);
\draw[draw=none,fill=steelblue31119180] (axis cs:430.5,0) rectangle (axis cs:431.5,100);
\draw[draw=none,fill=steelblue31119180] (axis cs:431.5,0) rectangle (axis cs:432.5,100);
\draw[draw=none,fill=steelblue31119180] (axis cs:432.5,0) rectangle (axis cs:433.5,100);
\draw[draw=none,fill=steelblue31119180] (axis cs:433.5,0) rectangle (axis cs:434.5,100);
\draw[draw=none,fill=steelblue31119180] (axis cs:434.5,0) rectangle (axis cs:435.5,100);
\draw[draw=none,fill=steelblue31119180] (axis cs:435.5,0) rectangle (axis cs:436.5,100);
\draw[draw=none,fill=steelblue31119180] (axis cs:436.5,0) rectangle (axis cs:437.5,100);
\draw[draw=none,fill=steelblue31119180] (axis cs:437.5,0) rectangle (axis cs:438.5,100);
\draw[draw=none,fill=steelblue31119180] (axis cs:438.5,0) rectangle (axis cs:439.5,100);
\draw[draw=none,fill=steelblue31119180] (axis cs:439.5,0) rectangle (axis cs:440.5,100);
\draw[draw=none,fill=steelblue31119180] (axis cs:440.5,0) rectangle (axis cs:441.5,100);
\draw[draw=none,fill=steelblue31119180] (axis cs:441.5,0) rectangle (axis cs:442.5,100);
\draw[draw=none,fill=steelblue31119180] (axis cs:442.5,0) rectangle (axis cs:443.5,100);
\draw[draw=none,fill=steelblue31119180] (axis cs:443.5,0) rectangle (axis cs:444.5,100);
\draw[draw=none,fill=steelblue31119180] (axis cs:444.5,0) rectangle (axis cs:445.5,100);
\draw[draw=none,fill=steelblue31119180] (axis cs:445.5,0) rectangle (axis cs:446.5,100);
\draw[draw=none,fill=steelblue31119180] (axis cs:446.5,0) rectangle (axis cs:447.5,100);
\draw[draw=none,fill=steelblue31119180] (axis cs:447.5,0) rectangle (axis cs:448.5,100);
\draw[draw=none,fill=steelblue31119180] (axis cs:448.5,0) rectangle (axis cs:449.5,100);
\draw[draw=none,fill=steelblue31119180] (axis cs:449.5,0) rectangle (axis cs:450.5,100);
\draw[draw=none,fill=steelblue31119180] (axis cs:450.5,0) rectangle (axis cs:451.5,100);
\draw[draw=none,fill=steelblue31119180] (axis cs:451.5,0) rectangle (axis cs:452.5,100);
\draw[draw=none,fill=steelblue31119180] (axis cs:452.5,0) rectangle (axis cs:453.5,100);
\draw[draw=none,fill=steelblue31119180] (axis cs:453.5,0) rectangle (axis cs:454.5,100);
\draw[draw=none,fill=steelblue31119180] (axis cs:454.5,0) rectangle (axis cs:455.5,100);
\draw[draw=none,fill=steelblue31119180] (axis cs:455.5,0) rectangle (axis cs:456.5,100);
\draw[draw=none,fill=steelblue31119180] (axis cs:456.5,0) rectangle (axis cs:457.5,100);
\draw[draw=none,fill=steelblue31119180] (axis cs:457.5,0) rectangle (axis cs:458.5,100);
\draw[draw=none,fill=steelblue31119180] (axis cs:458.5,0) rectangle (axis cs:459.5,100);
\draw[draw=none,fill=steelblue31119180] (axis cs:459.5,0) rectangle (axis cs:460.5,100);
\draw[draw=none,fill=steelblue31119180] (axis cs:460.5,0) rectangle (axis cs:461.5,100);
\draw[draw=none,fill=steelblue31119180] (axis cs:461.5,0) rectangle (axis cs:462.5,100);
\draw[draw=none,fill=steelblue31119180] (axis cs:462.5,0) rectangle (axis cs:463.5,100);
\draw[draw=none,fill=steelblue31119180] (axis cs:463.5,0) rectangle (axis cs:464.5,100);
\draw[draw=none,fill=steelblue31119180] (axis cs:464.5,0) rectangle (axis cs:465.5,100);
\draw[draw=none,fill=steelblue31119180] (axis cs:465.5,0) rectangle (axis cs:466.5,100);
\draw[draw=none,fill=steelblue31119180] (axis cs:466.5,0) rectangle (axis cs:467.5,100);
\draw[draw=none,fill=steelblue31119180] (axis cs:467.5,0) rectangle (axis cs:468.5,100);
\draw[draw=none,fill=steelblue31119180] (axis cs:468.5,0) rectangle (axis cs:469.5,100);
\draw[draw=none,fill=steelblue31119180] (axis cs:469.5,0) rectangle (axis cs:470.5,100);
\draw[draw=none,fill=steelblue31119180] (axis cs:470.5,0) rectangle (axis cs:471.5,100);
\draw[draw=none,fill=steelblue31119180] (axis cs:471.5,0) rectangle (axis cs:472.5,100);
\draw[draw=none,fill=steelblue31119180] (axis cs:472.5,0) rectangle (axis cs:473.5,100);
\draw[draw=none,fill=steelblue31119180] (axis cs:473.5,0) rectangle (axis cs:474.5,100);
\draw[draw=none,fill=steelblue31119180] (axis cs:474.5,0) rectangle (axis cs:475.5,100);
\draw[draw=none,fill=steelblue31119180] (axis cs:475.5,0) rectangle (axis cs:476.5,100);
\draw[draw=none,fill=steelblue31119180] (axis cs:476.5,0) rectangle (axis cs:477.5,100);
\draw[draw=none,fill=steelblue31119180] (axis cs:477.5,0) rectangle (axis cs:478.5,100);
\draw[draw=none,fill=steelblue31119180] (axis cs:478.5,0) rectangle (axis cs:479.5,100);
\draw[draw=none,fill=steelblue31119180] (axis cs:479.5,0) rectangle (axis cs:480.5,100);
\draw[draw=none,fill=steelblue31119180] (axis cs:480.5,0) rectangle (axis cs:481.5,100);
\draw[draw=none,fill=steelblue31119180] (axis cs:481.5,0) rectangle (axis cs:482.5,100);
\draw[draw=none,fill=steelblue31119180] (axis cs:482.5,0) rectangle (axis cs:483.5,100);
\draw[draw=none,fill=steelblue31119180] (axis cs:483.5,0) rectangle (axis cs:484.5,100);
\draw[draw=none,fill=steelblue31119180] (axis cs:484.5,0) rectangle (axis cs:485.5,100);
\draw[draw=none,fill=steelblue31119180] (axis cs:485.5,0) rectangle (axis cs:486.5,100);
\draw[draw=none,fill=steelblue31119180] (axis cs:486.5,0) rectangle (axis cs:487.5,100);
\draw[draw=none,fill=steelblue31119180] (axis cs:487.5,0) rectangle (axis cs:488.5,100);
\draw[draw=none,fill=steelblue31119180] (axis cs:488.5,0) rectangle (axis cs:489.5,100);
\draw[draw=none,fill=steelblue31119180] (axis cs:489.5,0) rectangle (axis cs:490.5,100);
\draw[draw=none,fill=steelblue31119180] (axis cs:490.5,0) rectangle (axis cs:491.5,100);
\draw[draw=none,fill=steelblue31119180] (axis cs:491.5,0) rectangle (axis cs:492.5,100);
\draw[draw=none,fill=steelblue31119180] (axis cs:492.5,0) rectangle (axis cs:493.5,100);
\draw[draw=none,fill=steelblue31119180] (axis cs:493.5,0) rectangle (axis cs:494.5,100);
\draw[draw=none,fill=steelblue31119180] (axis cs:494.5,0) rectangle (axis cs:495.5,100);
\draw[draw=none,fill=steelblue31119180] (axis cs:495.5,0) rectangle (axis cs:496.5,100);
\draw[draw=none,fill=steelblue31119180] (axis cs:496.5,0) rectangle (axis cs:497.5,100);
\draw[draw=none,fill=steelblue31119180] (axis cs:497.5,0) rectangle (axis cs:498.5,100);
\draw[draw=none,fill=steelblue31119180] (axis cs:498.5,0) rectangle (axis cs:499.5,100);
\draw[draw=none,fill=steelblue31119180] (axis cs:499.5,0) rectangle (axis cs:500.5,100);
\draw[draw=none,fill=steelblue31119180] (axis cs:500.5,0) rectangle (axis cs:501.5,100);
\draw[draw=none,fill=steelblue31119180] (axis cs:501.5,0) rectangle (axis cs:502.5,100);
\draw[draw=none,fill=steelblue31119180] (axis cs:502.5,0) rectangle (axis cs:503.5,100);
\draw[draw=none,fill=steelblue31119180] (axis cs:503.5,0) rectangle (axis cs:504.5,100);
\draw[draw=none,fill=steelblue31119180] (axis cs:504.5,0) rectangle (axis cs:505.5,100);
\draw[draw=none,fill=steelblue31119180] (axis cs:505.5,0) rectangle (axis cs:506.5,100);
\draw[draw=none,fill=steelblue31119180] (axis cs:506.5,0) rectangle (axis cs:507.5,100);
\draw[draw=none,fill=steelblue31119180] (axis cs:507.5,0) rectangle (axis cs:508.5,100);
\draw[draw=none,fill=steelblue31119180] (axis cs:508.5,0) rectangle (axis cs:509.5,100);
\draw[draw=none,fill=steelblue31119180] (axis cs:509.5,0) rectangle (axis cs:510.5,100);
\draw[draw=none,fill=steelblue31119180] (axis cs:510.5,0) rectangle (axis cs:511.5,100);
\draw[draw=none,fill=steelblue31119180] (axis cs:511.5,0) rectangle (axis cs:512.5,100);
\draw[draw=none,fill=steelblue31119180] (axis cs:512.5,0) rectangle (axis cs:513.5,100);
\draw[draw=none,fill=steelblue31119180] (axis cs:513.5,0) rectangle (axis cs:514.5,100);
\draw[draw=none,fill=steelblue31119180] (axis cs:514.5,0) rectangle (axis cs:515.5,100);
\draw[draw=none,fill=steelblue31119180] (axis cs:515.5,0) rectangle (axis cs:516.5,100);
\draw[draw=none,fill=steelblue31119180] (axis cs:516.5,0) rectangle (axis cs:517.5,100);
\draw[draw=none,fill=steelblue31119180] (axis cs:517.5,0) rectangle (axis cs:518.5,100);
\draw[draw=none,fill=steelblue31119180] (axis cs:518.5,0) rectangle (axis cs:519.5,100);
\draw[draw=none,fill=steelblue31119180] (axis cs:519.5,0) rectangle (axis cs:520.5,100);
\draw[draw=none,fill=steelblue31119180] (axis cs:520.5,0) rectangle (axis cs:521.5,100);
\draw[draw=none,fill=steelblue31119180] (axis cs:521.5,0) rectangle (axis cs:522.5,100);
\draw[draw=none,fill=steelblue31119180] (axis cs:522.5,0) rectangle (axis cs:523.5,100);
\draw[draw=none,fill=steelblue31119180] (axis cs:523.5,0) rectangle (axis cs:524.5,100);
\draw[draw=none,fill=steelblue31119180] (axis cs:524.5,0) rectangle (axis cs:525.5,100);
\draw[draw=none,fill=steelblue31119180] (axis cs:525.5,0) rectangle (axis cs:526.5,100);
\draw[draw=none,fill=steelblue31119180] (axis cs:526.5,0) rectangle (axis cs:527.5,100);
\draw[draw=none,fill=steelblue31119180] (axis cs:527.5,0) rectangle (axis cs:528.5,100);
\draw[draw=none,fill=steelblue31119180] (axis cs:528.5,0) rectangle (axis cs:529.5,100);
\draw[draw=none,fill=steelblue31119180] (axis cs:529.5,0) rectangle (axis cs:530.5,100);
\draw[draw=none,fill=steelblue31119180] (axis cs:530.5,0) rectangle (axis cs:531.5,100);
\draw[draw=none,fill=steelblue31119180] (axis cs:531.5,0) rectangle (axis cs:532.5,100);
\draw[draw=none,fill=steelblue31119180] (axis cs:532.5,0) rectangle (axis cs:533.5,100);
\draw[draw=none,fill=steelblue31119180] (axis cs:533.5,0) rectangle (axis cs:534.5,100);
\draw[draw=none,fill=steelblue31119180] (axis cs:534.5,0) rectangle (axis cs:535.5,100);
\draw[draw=none,fill=steelblue31119180] (axis cs:535.5,0) rectangle (axis cs:536.5,100);
\draw[draw=none,fill=steelblue31119180] (axis cs:536.5,0) rectangle (axis cs:537.5,100);
\draw[draw=none,fill=steelblue31119180] (axis cs:537.5,0) rectangle (axis cs:538.5,100);
\draw[draw=none,fill=steelblue31119180] (axis cs:538.5,0) rectangle (axis cs:539.5,100);
\draw[draw=none,fill=steelblue31119180] (axis cs:539.5,0) rectangle (axis cs:540.5,100);
\draw[draw=none,fill=steelblue31119180] (axis cs:540.5,0) rectangle (axis cs:541.5,100);
\draw[draw=none,fill=steelblue31119180] (axis cs:541.5,0) rectangle (axis cs:542.5,100);
\draw[draw=none,fill=steelblue31119180] (axis cs:542.5,0) rectangle (axis cs:543.5,100);
\draw[draw=none,fill=steelblue31119180] (axis cs:543.5,0) rectangle (axis cs:544.5,100);
\draw[draw=none,fill=steelblue31119180] (axis cs:544.5,0) rectangle (axis cs:545.5,100);
\draw[draw=none,fill=steelblue31119180] (axis cs:545.5,0) rectangle (axis cs:546.5,100);
\draw[draw=none,fill=steelblue31119180] (axis cs:546.5,0) rectangle (axis cs:547.5,100);
\draw[draw=none,fill=steelblue31119180] (axis cs:547.5,0) rectangle (axis cs:548.5,100);
\draw[draw=none,fill=steelblue31119180] (axis cs:548.5,0) rectangle (axis cs:549.5,100);
\draw[draw=none,fill=steelblue31119180] (axis cs:549.5,0) rectangle (axis cs:550.5,100);
\draw[draw=none,fill=steelblue31119180] (axis cs:550.5,0) rectangle (axis cs:551.5,100);
\draw[draw=none,fill=steelblue31119180] (axis cs:551.5,0) rectangle (axis cs:552.5,100);
\draw[draw=none,fill=steelblue31119180] (axis cs:552.5,0) rectangle (axis cs:553.5,100);
\draw[draw=none,fill=steelblue31119180] (axis cs:553.5,0) rectangle (axis cs:554.5,100);
\draw[draw=none,fill=steelblue31119180] (axis cs:554.5,0) rectangle (axis cs:555.5,100);
\draw[draw=none,fill=steelblue31119180] (axis cs:555.5,0) rectangle (axis cs:556.5,100);
\draw[draw=none,fill=steelblue31119180] (axis cs:556.5,0) rectangle (axis cs:557.5,100);
\draw[draw=none,fill=steelblue31119180] (axis cs:557.5,0) rectangle (axis cs:558.5,100);
\draw[draw=none,fill=steelblue31119180] (axis cs:558.5,0) rectangle (axis cs:559.5,100);
\draw[draw=none,fill=steelblue31119180] (axis cs:559.5,0) rectangle (axis cs:560.5,100);
\draw[draw=none,fill=steelblue31119180] (axis cs:560.5,0) rectangle (axis cs:561.5,100);
\draw[draw=none,fill=steelblue31119180] (axis cs:561.5,0) rectangle (axis cs:562.5,100);
\draw[draw=none,fill=steelblue31119180] (axis cs:562.5,0) rectangle (axis cs:563.5,100);
\draw[draw=none,fill=steelblue31119180] (axis cs:563.5,0) rectangle (axis cs:564.5,100);
\draw[draw=none,fill=steelblue31119180] (axis cs:564.5,0) rectangle (axis cs:565.5,100);
\draw[draw=none,fill=steelblue31119180] (axis cs:565.5,0) rectangle (axis cs:566.5,100);
\draw[draw=none,fill=steelblue31119180] (axis cs:566.5,0) rectangle (axis cs:567.5,100);
\draw[draw=none,fill=steelblue31119180] (axis cs:567.5,0) rectangle (axis cs:568.5,100);
\draw[draw=none,fill=steelblue31119180] (axis cs:568.5,0) rectangle (axis cs:569.5,100);
\draw[draw=none,fill=steelblue31119180] (axis cs:569.5,0) rectangle (axis cs:570.5,100);
\draw[draw=none,fill=steelblue31119180] (axis cs:570.5,0) rectangle (axis cs:571.5,100);
\draw[draw=none,fill=steelblue31119180] (axis cs:571.5,0) rectangle (axis cs:572.5,100);
\draw[draw=none,fill=steelblue31119180] (axis cs:572.5,0) rectangle (axis cs:573.5,100);
\draw[draw=none,fill=steelblue31119180] (axis cs:573.5,0) rectangle (axis cs:574.5,100);
\draw[draw=none,fill=steelblue31119180] (axis cs:574.5,0) rectangle (axis cs:575.5,100);
\draw[draw=none,fill=steelblue31119180] (axis cs:575.5,0) rectangle (axis cs:576.5,100);
\draw[draw=none,fill=steelblue31119180] (axis cs:576.5,0) rectangle (axis cs:577.5,100);
\draw[draw=none,fill=steelblue31119180] (axis cs:577.5,0) rectangle (axis cs:578.5,100);
\draw[draw=none,fill=steelblue31119180] (axis cs:578.5,0) rectangle (axis cs:579.5,100);
\draw[draw=none,fill=steelblue31119180] (axis cs:579.5,0) rectangle (axis cs:580.5,100);
\draw[draw=none,fill=steelblue31119180] (axis cs:580.5,0) rectangle (axis cs:581.5,100);
\draw[draw=none,fill=steelblue31119180] (axis cs:581.5,0) rectangle (axis cs:582.5,100);
\draw[draw=none,fill=steelblue31119180] (axis cs:582.5,0) rectangle (axis cs:583.5,100);
\draw[draw=none,fill=steelblue31119180] (axis cs:583.5,0) rectangle (axis cs:584.5,100);
\draw[draw=none,fill=steelblue31119180] (axis cs:584.5,0) rectangle (axis cs:585.5,100);
\draw[draw=none,fill=steelblue31119180] (axis cs:585.5,0) rectangle (axis cs:586.5,100);
\draw[draw=none,fill=steelblue31119180] (axis cs:586.5,0) rectangle (axis cs:587.5,100);
\draw[draw=none,fill=steelblue31119180] (axis cs:587.5,0) rectangle (axis cs:588.5,100);
\draw[draw=none,fill=steelblue31119180] (axis cs:588.5,0) rectangle (axis cs:589.5,100);
\draw[draw=none,fill=steelblue31119180] (axis cs:589.5,0) rectangle (axis cs:590.5,100);
\draw[draw=none,fill=steelblue31119180] (axis cs:590.5,0) rectangle (axis cs:591.5,100);
\draw[draw=none,fill=steelblue31119180] (axis cs:591.5,0) rectangle (axis cs:592.5,100);
\draw[draw=none,fill=steelblue31119180] (axis cs:592.5,0) rectangle (axis cs:593.5,100);
\draw[draw=none,fill=steelblue31119180] (axis cs:593.5,0) rectangle (axis cs:594.5,100);
\draw[draw=none,fill=steelblue31119180] (axis cs:594.5,0) rectangle (axis cs:595.5,100);
\draw[draw=none,fill=steelblue31119180] (axis cs:595.5,0) rectangle (axis cs:596.5,100);
\draw[draw=none,fill=steelblue31119180] (axis cs:596.5,0) rectangle (axis cs:597.5,100);
\draw[draw=none,fill=steelblue31119180] (axis cs:597.5,0) rectangle (axis cs:598.5,100);
\draw[draw=none,fill=steelblue31119180] (axis cs:598.5,0) rectangle (axis cs:599.5,100);
\draw[draw=none,fill=steelblue31119180] (axis cs:599.5,0) rectangle (axis cs:600.5,100);
\draw[draw=none,fill=steelblue31119180] (axis cs:600.5,0) rectangle (axis cs:601.5,100);
\draw[draw=none,fill=steelblue31119180] (axis cs:601.5,0) rectangle (axis cs:602.5,100);
\draw[draw=none,fill=steelblue31119180] (axis cs:602.5,0) rectangle (axis cs:603.5,100);
\draw[draw=none,fill=steelblue31119180] (axis cs:603.5,0) rectangle (axis cs:604.5,100);
\draw[draw=none,fill=steelblue31119180] (axis cs:604.5,0) rectangle (axis cs:605.5,100);
\draw[draw=none,fill=steelblue31119180] (axis cs:605.5,0) rectangle (axis cs:606.5,100);
\draw[draw=none,fill=steelblue31119180] (axis cs:606.5,0) rectangle (axis cs:607.5,100);
\draw[draw=none,fill=steelblue31119180] (axis cs:607.5,0) rectangle (axis cs:608.5,100);
\draw[draw=none,fill=steelblue31119180] (axis cs:608.5,0) rectangle (axis cs:609.5,100);
\draw[draw=none,fill=steelblue31119180] (axis cs:609.5,0) rectangle (axis cs:610.5,100);
\draw[draw=none,fill=steelblue31119180] (axis cs:610.5,0) rectangle (axis cs:611.5,100);
\draw[draw=none,fill=steelblue31119180] (axis cs:611.5,0) rectangle (axis cs:612.5,100);
\draw[draw=none,fill=steelblue31119180] (axis cs:612.5,0) rectangle (axis cs:613.5,100);
\draw[draw=none,fill=steelblue31119180] (axis cs:613.5,0) rectangle (axis cs:614.5,100);
\draw[draw=none,fill=steelblue31119180] (axis cs:614.5,0) rectangle (axis cs:615.5,100);
\draw[draw=none,fill=steelblue31119180] (axis cs:615.5,0) rectangle (axis cs:616.5,100);
\draw[draw=none,fill=steelblue31119180] (axis cs:616.5,0) rectangle (axis cs:617.5,100);
\draw[draw=none,fill=steelblue31119180] (axis cs:617.5,0) rectangle (axis cs:618.5,100);
\draw[draw=none,fill=steelblue31119180] (axis cs:618.5,0) rectangle (axis cs:619.5,100);
\draw[draw=none,fill=steelblue31119180] (axis cs:619.5,0) rectangle (axis cs:620.5,100);
\draw[draw=none,fill=steelblue31119180] (axis cs:620.5,0) rectangle (axis cs:621.5,100);
\draw[draw=none,fill=steelblue31119180] (axis cs:621.5,0) rectangle (axis cs:622.5,100);
\draw[draw=none,fill=steelblue31119180] (axis cs:622.5,0) rectangle (axis cs:623.5,100);
\draw[draw=none,fill=steelblue31119180] (axis cs:623.5,0) rectangle (axis cs:624.5,100);
\draw[draw=none,fill=steelblue31119180] (axis cs:624.5,0) rectangle (axis cs:625.5,100);
\draw[draw=none,fill=steelblue31119180] (axis cs:625.5,0) rectangle (axis cs:626.5,100);
\draw[draw=none,fill=steelblue31119180] (axis cs:626.5,0) rectangle (axis cs:627.5,100);
\draw[draw=none,fill=steelblue31119180] (axis cs:627.5,0) rectangle (axis cs:628.5,100);
\draw[draw=none,fill=steelblue31119180] (axis cs:628.5,0) rectangle (axis cs:629.5,100);
\draw[draw=none,fill=steelblue31119180] (axis cs:629.5,0) rectangle (axis cs:630.5,100);
\draw[draw=none,fill=steelblue31119180] (axis cs:630.5,0) rectangle (axis cs:631.5,100);
\draw[draw=none,fill=steelblue31119180] (axis cs:631.5,0) rectangle (axis cs:632.5,100);
\draw[draw=none,fill=steelblue31119180] (axis cs:632.5,0) rectangle (axis cs:633.5,100);
\draw[draw=none,fill=steelblue31119180] (axis cs:633.5,0) rectangle (axis cs:634.5,100);
\draw[draw=none,fill=steelblue31119180] (axis cs:634.5,0) rectangle (axis cs:635.5,100);
\draw[draw=none,fill=steelblue31119180] (axis cs:635.5,0) rectangle (axis cs:636.5,100);
\draw[draw=none,fill=steelblue31119180] (axis cs:636.5,0) rectangle (axis cs:637.5,100);
\draw[draw=none,fill=steelblue31119180] (axis cs:637.5,0) rectangle (axis cs:638.5,100);
\draw[draw=none,fill=steelblue31119180] (axis cs:638.5,0) rectangle (axis cs:639.5,100);
\draw[draw=none,fill=steelblue31119180] (axis cs:639.5,0) rectangle (axis cs:640.5,100);
\draw[draw=none,fill=steelblue31119180] (axis cs:640.5,0) rectangle (axis cs:641.5,100);
\draw[draw=none,fill=steelblue31119180] (axis cs:641.5,0) rectangle (axis cs:642.5,100);
\draw[draw=none,fill=steelblue31119180] (axis cs:642.5,0) rectangle (axis cs:643.5,100);
\draw[draw=none,fill=steelblue31119180] (axis cs:643.5,0) rectangle (axis cs:644.5,100);
\draw[draw=none,fill=steelblue31119180] (axis cs:644.5,0) rectangle (axis cs:645.5,100);
\draw[draw=none,fill=steelblue31119180] (axis cs:645.5,0) rectangle (axis cs:646.5,100);
\draw[draw=none,fill=steelblue31119180] (axis cs:646.5,0) rectangle (axis cs:647.5,100);
\draw[draw=none,fill=steelblue31119180] (axis cs:647.5,0) rectangle (axis cs:648.5,100);
\draw[draw=none,fill=steelblue31119180] (axis cs:648.5,0) rectangle (axis cs:649.5,100);
\draw[draw=none,fill=steelblue31119180] (axis cs:649.5,0) rectangle (axis cs:650.5,100);
\draw[draw=none,fill=steelblue31119180] (axis cs:650.5,0) rectangle (axis cs:651.5,100);
\draw[draw=none,fill=steelblue31119180] (axis cs:651.5,0) rectangle (axis cs:652.5,100);
\draw[draw=none,fill=steelblue31119180] (axis cs:652.5,0) rectangle (axis cs:653.5,100);
\draw[draw=none,fill=steelblue31119180] (axis cs:653.5,0) rectangle (axis cs:654.5,100);
\draw[draw=none,fill=steelblue31119180] (axis cs:654.5,0) rectangle (axis cs:655.5,100);
\draw[draw=none,fill=steelblue31119180] (axis cs:655.5,0) rectangle (axis cs:656.5,100);
\draw[draw=none,fill=steelblue31119180] (axis cs:656.5,0) rectangle (axis cs:657.5,100);
\draw[draw=none,fill=steelblue31119180] (axis cs:657.5,0) rectangle (axis cs:658.5,100);
\draw[draw=none,fill=steelblue31119180] (axis cs:658.5,0) rectangle (axis cs:659.5,100);
\draw[draw=none,fill=steelblue31119180] (axis cs:659.5,0) rectangle (axis cs:660.5,100);
\draw[draw=none,fill=steelblue31119180] (axis cs:660.5,0) rectangle (axis cs:661.5,100);
\draw[draw=none,fill=steelblue31119180] (axis cs:661.5,0) rectangle (axis cs:662.5,100);
\draw[draw=none,fill=steelblue31119180] (axis cs:662.5,0) rectangle (axis cs:663.5,100);
\draw[draw=none,fill=steelblue31119180] (axis cs:663.5,0) rectangle (axis cs:664.5,100);
\draw[draw=none,fill=steelblue31119180] (axis cs:664.5,0) rectangle (axis cs:665.5,100);
\draw[draw=none,fill=steelblue31119180] (axis cs:665.5,0) rectangle (axis cs:666.5,100);
\draw[draw=none,fill=steelblue31119180] (axis cs:666.5,0) rectangle (axis cs:667.5,100);
\draw[draw=none,fill=steelblue31119180] (axis cs:667.5,0) rectangle (axis cs:668.5,100);
\draw[draw=none,fill=steelblue31119180] (axis cs:668.5,0) rectangle (axis cs:669.5,100);
\draw[draw=none,fill=steelblue31119180] (axis cs:669.5,0) rectangle (axis cs:670.5,100);
\draw[draw=none,fill=steelblue31119180] (axis cs:670.5,0) rectangle (axis cs:671.5,100);
\draw[draw=none,fill=steelblue31119180] (axis cs:671.5,0) rectangle (axis cs:672.5,100);
\draw[draw=none,fill=steelblue31119180] (axis cs:672.5,0) rectangle (axis cs:673.5,100);
\draw[draw=none,fill=steelblue31119180] (axis cs:673.5,0) rectangle (axis cs:674.5,100);
\draw[draw=none,fill=steelblue31119180] (axis cs:674.5,0) rectangle (axis cs:675.5,100);
\draw[draw=none,fill=steelblue31119180] (axis cs:675.5,0) rectangle (axis cs:676.5,100);
\draw[draw=none,fill=steelblue31119180] (axis cs:676.5,0) rectangle (axis cs:677.5,100);
\draw[draw=none,fill=steelblue31119180] (axis cs:677.5,0) rectangle (axis cs:678.5,100);
\draw[draw=none,fill=steelblue31119180] (axis cs:678.5,0) rectangle (axis cs:679.5,100);
\draw[draw=none,fill=steelblue31119180] (axis cs:679.5,0) rectangle (axis cs:680.5,100);
\draw[draw=none,fill=steelblue31119180] (axis cs:680.5,0) rectangle (axis cs:681.5,100);
\draw[draw=none,fill=steelblue31119180] (axis cs:681.5,0) rectangle (axis cs:682.5,100);
\draw[draw=none,fill=steelblue31119180] (axis cs:682.5,0) rectangle (axis cs:683.5,100);
\draw[draw=none,fill=steelblue31119180] (axis cs:683.5,0) rectangle (axis cs:684.5,100);
\draw[draw=none,fill=steelblue31119180] (axis cs:684.5,0) rectangle (axis cs:685.5,100);
\draw[draw=none,fill=steelblue31119180] (axis cs:685.5,0) rectangle (axis cs:686.5,100);
\draw[draw=none,fill=steelblue31119180] (axis cs:686.5,0) rectangle (axis cs:687.5,100);
\draw[draw=none,fill=steelblue31119180] (axis cs:687.5,0) rectangle (axis cs:688.5,100);
\draw[draw=none,fill=steelblue31119180] (axis cs:688.5,0) rectangle (axis cs:689.5,100);
\draw[draw=none,fill=steelblue31119180] (axis cs:689.5,0) rectangle (axis cs:690.5,100);
\draw[draw=none,fill=steelblue31119180] (axis cs:690.5,0) rectangle (axis cs:691.5,100);
\draw[draw=none,fill=steelblue31119180] (axis cs:691.5,0) rectangle (axis cs:692.5,100);
\draw[draw=none,fill=steelblue31119180] (axis cs:692.5,0) rectangle (axis cs:693.5,100);
\draw[draw=none,fill=steelblue31119180] (axis cs:693.5,0) rectangle (axis cs:694.5,100);
\draw[draw=none,fill=steelblue31119180] (axis cs:694.5,0) rectangle (axis cs:695.5,100);
\draw[draw=none,fill=steelblue31119180] (axis cs:695.5,0) rectangle (axis cs:696.5,100);
\draw[draw=none,fill=steelblue31119180] (axis cs:696.5,0) rectangle (axis cs:697.5,100);
\draw[draw=none,fill=steelblue31119180] (axis cs:697.5,0) rectangle (axis cs:698.5,100);
\draw[draw=none,fill=steelblue31119180] (axis cs:698.5,0) rectangle (axis cs:699.5,100);
\draw[draw=none,fill=steelblue31119180] (axis cs:699.5,0) rectangle (axis cs:700.5,100);
\draw[draw=none,fill=steelblue31119180] (axis cs:700.5,0) rectangle (axis cs:701.5,100);
\draw[draw=none,fill=steelblue31119180] (axis cs:701.5,0) rectangle (axis cs:702.5,100);
\draw[draw=none,fill=steelblue31119180] (axis cs:702.5,0) rectangle (axis cs:703.5,100);
\draw[draw=none,fill=steelblue31119180] (axis cs:703.5,0) rectangle (axis cs:704.5,100);
\draw[draw=none,fill=steelblue31119180] (axis cs:704.5,0) rectangle (axis cs:705.5,100);
\draw[draw=none,fill=steelblue31119180] (axis cs:705.5,0) rectangle (axis cs:706.5,100);
\draw[draw=none,fill=steelblue31119180] (axis cs:706.5,0) rectangle (axis cs:707.5,100);
\draw[draw=none,fill=steelblue31119180] (axis cs:707.5,0) rectangle (axis cs:708.5,100);
\draw[draw=none,fill=steelblue31119180] (axis cs:708.5,0) rectangle (axis cs:709.5,100);
\draw[draw=none,fill=steelblue31119180] (axis cs:709.5,0) rectangle (axis cs:710.5,100);
\draw[draw=none,fill=steelblue31119180] (axis cs:710.5,0) rectangle (axis cs:711.5,100);
\draw[draw=none,fill=steelblue31119180] (axis cs:711.5,0) rectangle (axis cs:712.5,100);
\draw[draw=none,fill=steelblue31119180] (axis cs:712.5,0) rectangle (axis cs:713.5,100);
\draw[draw=none,fill=steelblue31119180] (axis cs:713.5,0) rectangle (axis cs:714.5,100);
\draw[draw=none,fill=steelblue31119180] (axis cs:714.5,0) rectangle (axis cs:715.5,100);
\draw[draw=none,fill=steelblue31119180] (axis cs:715.5,0) rectangle (axis cs:716.5,100);
\draw[draw=none,fill=steelblue31119180] (axis cs:716.5,0) rectangle (axis cs:717.5,100);
\draw[draw=none,fill=steelblue31119180] (axis cs:717.5,0) rectangle (axis cs:718.5,100);
\draw[draw=none,fill=steelblue31119180] (axis cs:718.5,0) rectangle (axis cs:719.5,100);
\draw[draw=none,fill=steelblue31119180] (axis cs:719.5,0) rectangle (axis cs:720.5,100);
\draw[draw=none,fill=steelblue31119180] (axis cs:720.5,0) rectangle (axis cs:721.5,100);
\draw[draw=none,fill=steelblue31119180] (axis cs:721.5,0) rectangle (axis cs:722.5,100);
\draw[draw=none,fill=steelblue31119180] (axis cs:722.5,0) rectangle (axis cs:723.5,100);
\draw[draw=none,fill=steelblue31119180] (axis cs:723.5,0) rectangle (axis cs:724.5,100);
\draw[draw=none,fill=steelblue31119180] (axis cs:724.5,0) rectangle (axis cs:725.5,100);
\draw[draw=none,fill=steelblue31119180] (axis cs:725.5,0) rectangle (axis cs:726.5,100);
\draw[draw=none,fill=steelblue31119180] (axis cs:726.5,0) rectangle (axis cs:727.5,100);
\draw[draw=none,fill=steelblue31119180] (axis cs:727.5,0) rectangle (axis cs:728.5,100);
\draw[draw=none,fill=steelblue31119180] (axis cs:728.5,0) rectangle (axis cs:729.5,100);
\draw[draw=none,fill=steelblue31119180] (axis cs:729.5,0) rectangle (axis cs:730.5,100);
\draw[draw=none,fill=steelblue31119180] (axis cs:730.5,0) rectangle (axis cs:731.5,100);
\draw[draw=none,fill=steelblue31119180] (axis cs:731.5,0) rectangle (axis cs:732.5,100);
\draw[draw=none,fill=steelblue31119180] (axis cs:732.5,0) rectangle (axis cs:733.5,100);
\draw[draw=none,fill=steelblue31119180] (axis cs:733.5,0) rectangle (axis cs:734.5,100);
\draw[draw=none,fill=steelblue31119180] (axis cs:734.5,0) rectangle (axis cs:735.5,100);
\draw[draw=none,fill=steelblue31119180] (axis cs:735.5,0) rectangle (axis cs:736.5,100);
\draw[draw=none,fill=steelblue31119180] (axis cs:736.5,0) rectangle (axis cs:737.5,100);
\draw[draw=none,fill=steelblue31119180] (axis cs:737.5,0) rectangle (axis cs:738.5,100);
\draw[draw=none,fill=steelblue31119180] (axis cs:738.5,0) rectangle (axis cs:739.5,100);
\draw[draw=none,fill=steelblue31119180] (axis cs:739.5,0) rectangle (axis cs:740.5,100);
\draw[draw=none,fill=steelblue31119180] (axis cs:740.5,0) rectangle (axis cs:741.5,100);
\draw[draw=none,fill=steelblue31119180] (axis cs:741.5,0) rectangle (axis cs:742.5,100);
\draw[draw=none,fill=steelblue31119180] (axis cs:742.5,0) rectangle (axis cs:743.5,100);
\draw[draw=none,fill=steelblue31119180] (axis cs:743.5,0) rectangle (axis cs:744.5,100);
\draw[draw=none,fill=steelblue31119180] (axis cs:744.5,0) rectangle (axis cs:745.5,100);
\draw[draw=none,fill=steelblue31119180] (axis cs:745.5,0) rectangle (axis cs:746.5,100);
\draw[draw=none,fill=steelblue31119180] (axis cs:746.5,0) rectangle (axis cs:747.5,100);
\draw[draw=none,fill=steelblue31119180] (axis cs:747.5,0) rectangle (axis cs:748.5,100);
\draw[draw=none,fill=steelblue31119180] (axis cs:748.5,0) rectangle (axis cs:749.5,100);
\draw[draw=none,fill=steelblue31119180] (axis cs:749.5,0) rectangle (axis cs:750.5,100);
\draw[draw=none,fill=steelblue31119180] (axis cs:750.5,0) rectangle (axis cs:751.5,100);
\draw[draw=none,fill=steelblue31119180] (axis cs:751.5,0) rectangle (axis cs:752.5,100);
\draw[draw=none,fill=steelblue31119180] (axis cs:752.5,0) rectangle (axis cs:753.5,100);
\draw[draw=none,fill=steelblue31119180] (axis cs:753.5,0) rectangle (axis cs:754.5,100);
\draw[draw=none,fill=steelblue31119180] (axis cs:754.5,0) rectangle (axis cs:755.5,100);
\draw[draw=none,fill=steelblue31119180] (axis cs:755.5,0) rectangle (axis cs:756.5,100);
\draw[draw=none,fill=steelblue31119180] (axis cs:756.5,0) rectangle (axis cs:757.5,100);
\draw[draw=none,fill=steelblue31119180] (axis cs:757.5,0) rectangle (axis cs:758.5,100);
\draw[draw=none,fill=steelblue31119180] (axis cs:758.5,0) rectangle (axis cs:759.5,100);
\draw[draw=none,fill=steelblue31119180] (axis cs:759.5,0) rectangle (axis cs:760.5,100);
\draw[draw=none,fill=steelblue31119180] (axis cs:760.5,0) rectangle (axis cs:761.5,100);
\draw[draw=none,fill=steelblue31119180] (axis cs:761.5,0) rectangle (axis cs:762.5,100);
\draw[draw=none,fill=steelblue31119180] (axis cs:762.5,0) rectangle (axis cs:763.5,100);
\draw[draw=none,fill=steelblue31119180] (axis cs:763.5,0) rectangle (axis cs:764.5,100);
\draw[draw=none,fill=steelblue31119180] (axis cs:764.5,0) rectangle (axis cs:765.5,100);
\draw[draw=none,fill=steelblue31119180] (axis cs:765.5,0) rectangle (axis cs:766.5,100);
\draw[draw=none,fill=steelblue31119180] (axis cs:766.5,0) rectangle (axis cs:767.5,100);
\draw[draw=none,fill=steelblue31119180] (axis cs:767.5,0) rectangle (axis cs:768.5,100);
\draw[draw=none,fill=steelblue31119180] (axis cs:768.5,0) rectangle (axis cs:769.5,100);
\draw[draw=none,fill=steelblue31119180] (axis cs:769.5,0) rectangle (axis cs:770.5,100);
\draw[draw=none,fill=steelblue31119180] (axis cs:770.5,0) rectangle (axis cs:771.5,100);
\draw[draw=none,fill=steelblue31119180] (axis cs:771.5,0) rectangle (axis cs:772.5,100);
\draw[draw=none,fill=steelblue31119180] (axis cs:772.5,0) rectangle (axis cs:773.5,100);
\draw[draw=none,fill=steelblue31119180] (axis cs:773.5,0) rectangle (axis cs:774.5,100);
\draw[draw=none,fill=steelblue31119180] (axis cs:774.5,0) rectangle (axis cs:775.5,100);
\draw[draw=none,fill=steelblue31119180] (axis cs:775.5,0) rectangle (axis cs:776.5,100);
\draw[draw=none,fill=steelblue31119180] (axis cs:776.5,0) rectangle (axis cs:777.5,100);
\draw[draw=none,fill=steelblue31119180] (axis cs:777.5,0) rectangle (axis cs:778.5,100);
\draw[draw=none,fill=steelblue31119180] (axis cs:778.5,0) rectangle (axis cs:779.5,100);
\draw[draw=none,fill=steelblue31119180] (axis cs:779.5,0) rectangle (axis cs:780.5,100);
\draw[draw=none,fill=steelblue31119180] (axis cs:780.5,0) rectangle (axis cs:781.5,100);
\draw[draw=none,fill=steelblue31119180] (axis cs:781.5,0) rectangle (axis cs:782.5,100);
\draw[draw=none,fill=steelblue31119180] (axis cs:782.5,0) rectangle (axis cs:783.5,100);
\draw[draw=none,fill=steelblue31119180] (axis cs:783.5,0) rectangle (axis cs:784.5,100);
\draw[draw=none,fill=steelblue31119180] (axis cs:784.5,0) rectangle (axis cs:785.5,100);
\draw[draw=none,fill=steelblue31119180] (axis cs:785.5,0) rectangle (axis cs:786.5,100);
\draw[draw=none,fill=darkorange25512714] (axis cs:-0.5,0) rectangle (axis cs:0.5,9.87261146496815);
\addlegendimage{ybar,ybar legend,draw=none,fill=darkorange25512714}
\addlegendentry{$\Delta_L=2$}

\draw[draw=none,fill=darkorange25512714] (axis cs:0.5,0) rectangle (axis cs:1.5,19.7452229299363);
\draw[draw=none,fill=darkorange25512714] (axis cs:1.5,0) rectangle (axis cs:2.5,26.7515923566879);
\draw[draw=none,fill=darkorange25512714] (axis cs:2.5,0) rectangle (axis cs:3.5,31.5286624203822);
\draw[draw=none,fill=darkorange25512714] (axis cs:3.5,0) rectangle (axis cs:4.5,35.3503184713376);
\draw[draw=none,fill=darkorange25512714] (axis cs:4.5,0) rectangle (axis cs:5.5,38.2165605095541);
\draw[draw=none,fill=darkorange25512714] (axis cs:5.5,0) rectangle (axis cs:6.5,40.4458598726115);
\draw[draw=none,fill=darkorange25512714] (axis cs:6.5,0) rectangle (axis cs:7.5,44.9044585987261);
\draw[draw=none,fill=darkorange25512714] (axis cs:7.5,0) rectangle (axis cs:8.5,48.0891719745223);
\draw[draw=none,fill=darkorange25512714] (axis cs:8.5,0) rectangle (axis cs:9.5,51.5923566878981);
\draw[draw=none,fill=darkorange25512714] (axis cs:9.5,0) rectangle (axis cs:10.5,54.7770700636943);
\draw[draw=none,fill=darkorange25512714] (axis cs:10.5,0) rectangle (axis cs:11.5,57.6433121019108);
\draw[draw=none,fill=darkorange25512714] (axis cs:11.5,0) rectangle (axis cs:12.5,58.5987261146497);
\draw[draw=none,fill=darkorange25512714] (axis cs:12.5,0) rectangle (axis cs:13.5,61.1464968152866);
\draw[draw=none,fill=darkorange25512714] (axis cs:13.5,0) rectangle (axis cs:14.5,63.3757961783439);
\draw[draw=none,fill=darkorange25512714] (axis cs:14.5,0) rectangle (axis cs:15.5,66.8789808917197);
\draw[draw=none,fill=darkorange25512714] (axis cs:15.5,0) rectangle (axis cs:16.5,69.7452229299363);
\draw[draw=none,fill=darkorange25512714] (axis cs:16.5,0) rectangle (axis cs:17.5,70.3821656050955);
\draw[draw=none,fill=darkorange25512714] (axis cs:17.5,0) rectangle (axis cs:18.5,72.9299363057325);
\draw[draw=none,fill=darkorange25512714] (axis cs:18.5,0) rectangle (axis cs:19.5,75.1592356687898);
\draw[draw=none,fill=darkorange25512714] (axis cs:19.5,0) rectangle (axis cs:20.5,76.1146496815287);
\draw[draw=none,fill=darkorange25512714] (axis cs:20.5,0) rectangle (axis cs:21.5,78.343949044586);
\draw[draw=none,fill=darkorange25512714] (axis cs:21.5,0) rectangle (axis cs:22.5,79.2993630573248);
\draw[draw=none,fill=darkorange25512714] (axis cs:22.5,0) rectangle (axis cs:23.5,80.2547770700637);
\draw[draw=none,fill=darkorange25512714] (axis cs:23.5,0) rectangle (axis cs:24.5,80.5732484076433);
\draw[draw=none,fill=darkorange25512714] (axis cs:24.5,0) rectangle (axis cs:25.5,82.1656050955414);
\draw[draw=none,fill=darkorange25512714] (axis cs:25.5,0) rectangle (axis cs:26.5,82.484076433121);
\draw[draw=none,fill=darkorange25512714] (axis cs:26.5,0) rectangle (axis cs:27.5,82.8025477707006);
\draw[draw=none,fill=darkorange25512714] (axis cs:27.5,0) rectangle (axis cs:28.5,84.0764331210191);
\draw[draw=none,fill=darkorange25512714] (axis cs:28.5,0) rectangle (axis cs:29.5,85.6687898089172);
\draw[draw=none,fill=darkorange25512714] (axis cs:29.5,0) rectangle (axis cs:30.5,86.3057324840764);
\draw[draw=none,fill=darkorange25512714] (axis cs:30.5,0) rectangle (axis cs:31.5,86.3057324840764);
\draw[draw=none,fill=darkorange25512714] (axis cs:31.5,0) rectangle (axis cs:32.5,87.5796178343949);
\draw[draw=none,fill=darkorange25512714] (axis cs:32.5,0) rectangle (axis cs:33.5,89.171974522293);
\draw[draw=none,fill=darkorange25512714] (axis cs:33.5,0) rectangle (axis cs:34.5,90.1273885350319);
\draw[draw=none,fill=darkorange25512714] (axis cs:34.5,0) rectangle (axis cs:35.5,90.4458598726115);
\draw[draw=none,fill=darkorange25512714] (axis cs:35.5,0) rectangle (axis cs:36.5,90.7643312101911);
\draw[draw=none,fill=darkorange25512714] (axis cs:36.5,0) rectangle (axis cs:37.5,91.7197452229299);
\draw[draw=none,fill=darkorange25512714] (axis cs:37.5,0) rectangle (axis cs:38.5,91.7197452229299);
\draw[draw=none,fill=darkorange25512714] (axis cs:38.5,0) rectangle (axis cs:39.5,92.0382165605096);
\draw[draw=none,fill=darkorange25512714] (axis cs:39.5,0) rectangle (axis cs:40.5,92.9936305732484);
\draw[draw=none,fill=darkorange25512714] (axis cs:40.5,0) rectangle (axis cs:41.5,93.312101910828);
\draw[draw=none,fill=darkorange25512714] (axis cs:41.5,0) rectangle (axis cs:42.5,93.312101910828);
\draw[draw=none,fill=darkorange25512714] (axis cs:42.5,0) rectangle (axis cs:43.5,93.312101910828);
\draw[draw=none,fill=darkorange25512714] (axis cs:43.5,0) rectangle (axis cs:44.5,93.6305732484076);
\draw[draw=none,fill=darkorange25512714] (axis cs:44.5,0) rectangle (axis cs:45.5,95.2229299363057);
\draw[draw=none,fill=darkorange25512714] (axis cs:45.5,0) rectangle (axis cs:46.5,95.2229299363057);
\draw[draw=none,fill=darkorange25512714] (axis cs:46.5,0) rectangle (axis cs:47.5,95.2229299363057);
\draw[draw=none,fill=darkorange25512714] (axis cs:47.5,0) rectangle (axis cs:48.5,95.2229299363057);
\draw[draw=none,fill=darkorange25512714] (axis cs:48.5,0) rectangle (axis cs:49.5,95.2229299363057);
\draw[draw=none,fill=darkorange25512714] (axis cs:49.5,0) rectangle (axis cs:50.5,95.5414012738854);
\draw[draw=none,fill=darkorange25512714] (axis cs:50.5,0) rectangle (axis cs:51.5,95.5414012738854);
\draw[draw=none,fill=darkorange25512714] (axis cs:51.5,0) rectangle (axis cs:52.5,96.1783439490446);
\draw[draw=none,fill=darkorange25512714] (axis cs:52.5,0) rectangle (axis cs:53.5,96.1783439490446);
\draw[draw=none,fill=darkorange25512714] (axis cs:53.5,0) rectangle (axis cs:54.5,96.1783439490446);
\draw[draw=none,fill=darkorange25512714] (axis cs:54.5,0) rectangle (axis cs:55.5,96.1783439490446);
\draw[draw=none,fill=darkorange25512714] (axis cs:55.5,0) rectangle (axis cs:56.5,96.1783439490446);
\draw[draw=none,fill=darkorange25512714] (axis cs:56.5,0) rectangle (axis cs:57.5,96.8152866242038);
\draw[draw=none,fill=darkorange25512714] (axis cs:57.5,0) rectangle (axis cs:58.5,97.1337579617834);
\draw[draw=none,fill=darkorange25512714] (axis cs:58.5,0) rectangle (axis cs:59.5,97.1337579617834);
\draw[draw=none,fill=darkorange25512714] (axis cs:59.5,0) rectangle (axis cs:60.5,97.1337579617834);
\draw[draw=none,fill=darkorange25512714] (axis cs:60.5,0) rectangle (axis cs:61.5,97.1337579617834);
\draw[draw=none,fill=darkorange25512714] (axis cs:61.5,0) rectangle (axis cs:62.5,97.4522292993631);
\draw[draw=none,fill=darkorange25512714] (axis cs:62.5,0) rectangle (axis cs:63.5,98.0891719745223);
\draw[draw=none,fill=darkorange25512714] (axis cs:63.5,0) rectangle (axis cs:64.5,98.0891719745223);
\draw[draw=none,fill=darkorange25512714] (axis cs:64.5,0) rectangle (axis cs:65.5,98.4076433121019);
\draw[draw=none,fill=darkorange25512714] (axis cs:65.5,0) rectangle (axis cs:66.5,98.7261146496815);
\draw[draw=none,fill=darkorange25512714] (axis cs:66.5,0) rectangle (axis cs:67.5,98.7261146496815);
\draw[draw=none,fill=darkorange25512714] (axis cs:67.5,0) rectangle (axis cs:68.5,98.7261146496815);
\draw[draw=none,fill=darkorange25512714] (axis cs:68.5,0) rectangle (axis cs:69.5,99.0445859872611);
\draw[draw=none,fill=darkorange25512714] (axis cs:69.5,0) rectangle (axis cs:70.5,99.0445859872611);
\draw[draw=none,fill=darkorange25512714] (axis cs:70.5,0) rectangle (axis cs:71.5,99.0445859872611);
\draw[draw=none,fill=darkorange25512714] (axis cs:71.5,0) rectangle (axis cs:72.5,99.0445859872611);
\draw[draw=none,fill=darkorange25512714] (axis cs:72.5,0) rectangle (axis cs:73.5,99.3630573248408);
\draw[draw=none,fill=darkorange25512714] (axis cs:73.5,0) rectangle (axis cs:74.5,99.3630573248408);
\draw[draw=none,fill=darkorange25512714] (axis cs:74.5,0) rectangle (axis cs:75.5,99.3630573248408);
\draw[draw=none,fill=darkorange25512714] (axis cs:75.5,0) rectangle (axis cs:76.5,99.3630573248408);
\draw[draw=none,fill=darkorange25512714] (axis cs:76.5,0) rectangle (axis cs:77.5,99.3630573248408);
\draw[draw=none,fill=darkorange25512714] (axis cs:77.5,0) rectangle (axis cs:78.5,99.6815286624204);
\draw[draw=none,fill=darkorange25512714] (axis cs:78.5,0) rectangle (axis cs:79.5,100);
\draw[draw=none,fill=darkorange25512714] (axis cs:79.5,0) rectangle (axis cs:80.5,100);
\draw[draw=none,fill=darkorange25512714] (axis cs:80.5,0) rectangle (axis cs:81.5,100);
\draw[draw=none,fill=darkorange25512714] (axis cs:81.5,0) rectangle (axis cs:82.5,100);
\draw[draw=none,fill=darkorange25512714] (axis cs:82.5,0) rectangle (axis cs:83.5,100);
\draw[draw=none,fill=darkorange25512714] (axis cs:83.5,0) rectangle (axis cs:84.5,100);
\draw[draw=none,fill=darkorange25512714] (axis cs:84.5,0) rectangle (axis cs:85.5,100);
\draw[draw=none,fill=darkorange25512714] (axis cs:85.5,0) rectangle (axis cs:86.5,100);
\draw[draw=none,fill=darkorange25512714] (axis cs:86.5,0) rectangle (axis cs:87.5,100);
\draw[draw=none,fill=darkorange25512714] (axis cs:87.5,0) rectangle (axis cs:88.5,100);
\draw[draw=none,fill=darkorange25512714] (axis cs:88.5,0) rectangle (axis cs:89.5,100);
\draw[draw=none,fill=darkorange25512714] (axis cs:89.5,0) rectangle (axis cs:90.5,100);
\draw[draw=none,fill=darkorange25512714] (axis cs:90.5,0) rectangle (axis cs:91.5,100);
\draw[draw=none,fill=darkorange25512714] (axis cs:91.5,0) rectangle (axis cs:92.5,100);
\draw[draw=none,fill=darkorange25512714] (axis cs:92.5,0) rectangle (axis cs:93.5,100);
\draw[draw=none,fill=darkorange25512714] (axis cs:93.5,0) rectangle (axis cs:94.5,100);
\draw[draw=none,fill=darkorange25512714] (axis cs:94.5,0) rectangle (axis cs:95.5,100);
\draw[draw=none,fill=darkorange25512714] (axis cs:95.5,0) rectangle (axis cs:96.5,100);
\draw[draw=none,fill=darkorange25512714] (axis cs:96.5,0) rectangle (axis cs:97.5,100);
\draw[draw=none,fill=darkorange25512714] (axis cs:97.5,0) rectangle (axis cs:98.5,100);
\draw[draw=none,fill=darkorange25512714] (axis cs:98.5,0) rectangle (axis cs:99.5,100);
\draw[draw=none,fill=darkorange25512714] (axis cs:99.5,0) rectangle (axis cs:100.5,100);
\draw[draw=none,fill=darkorange25512714] (axis cs:100.5,0) rectangle (axis cs:101.5,100);
\draw[draw=none,fill=darkorange25512714] (axis cs:101.5,0) rectangle (axis cs:102.5,100);
\draw[draw=none,fill=darkorange25512714] (axis cs:102.5,0) rectangle (axis cs:103.5,100);
\draw[draw=none,fill=darkorange25512714] (axis cs:103.5,0) rectangle (axis cs:104.5,100);
\draw[draw=none,fill=darkorange25512714] (axis cs:104.5,0) rectangle (axis cs:105.5,100);
\draw[draw=none,fill=darkorange25512714] (axis cs:105.5,0) rectangle (axis cs:106.5,100);
\draw[draw=none,fill=darkorange25512714] (axis cs:106.5,0) rectangle (axis cs:107.5,100);
\draw[draw=none,fill=darkorange25512714] (axis cs:107.5,0) rectangle (axis cs:108.5,100);
\draw[draw=none,fill=darkorange25512714] (axis cs:108.5,0) rectangle (axis cs:109.5,100);
\draw[draw=none,fill=darkorange25512714] (axis cs:109.5,0) rectangle (axis cs:110.5,100);
\draw[draw=none,fill=darkorange25512714] (axis cs:110.5,0) rectangle (axis cs:111.5,100);
\draw[draw=none,fill=darkorange25512714] (axis cs:111.5,0) rectangle (axis cs:112.5,100);
\draw[draw=none,fill=darkorange25512714] (axis cs:112.5,0) rectangle (axis cs:113.5,100);
\draw[draw=none,fill=darkorange25512714] (axis cs:113.5,0) rectangle (axis cs:114.5,100);
\draw[draw=none,fill=darkorange25512714] (axis cs:114.5,0) rectangle (axis cs:115.5,100);
\draw[draw=none,fill=darkorange25512714] (axis cs:115.5,0) rectangle (axis cs:116.5,100);
\draw[draw=none,fill=darkorange25512714] (axis cs:116.5,0) rectangle (axis cs:117.5,100);
\draw[draw=none,fill=darkorange25512714] (axis cs:117.5,0) rectangle (axis cs:118.5,100);
\draw[draw=none,fill=darkorange25512714] (axis cs:118.5,0) rectangle (axis cs:119.5,100);
\draw[draw=none,fill=darkorange25512714] (axis cs:119.5,0) rectangle (axis cs:120.5,100);
\draw[draw=none,fill=darkorange25512714] (axis cs:120.5,0) rectangle (axis cs:121.5,100);
\draw[draw=none,fill=darkorange25512714] (axis cs:121.5,0) rectangle (axis cs:122.5,100);
\draw[draw=none,fill=darkorange25512714] (axis cs:122.5,0) rectangle (axis cs:123.5,100);
\draw[draw=none,fill=darkorange25512714] (axis cs:123.5,0) rectangle (axis cs:124.5,100);
\draw[draw=none,fill=darkorange25512714] (axis cs:124.5,0) rectangle (axis cs:125.5,100);
\draw[draw=none,fill=darkorange25512714] (axis cs:125.5,0) rectangle (axis cs:126.5,100);
\draw[draw=none,fill=darkorange25512714] (axis cs:126.5,0) rectangle (axis cs:127.5,100);
\draw[draw=none,fill=darkorange25512714] (axis cs:127.5,0) rectangle (axis cs:128.5,100);
\draw[draw=none,fill=darkorange25512714] (axis cs:128.5,0) rectangle (axis cs:129.5,100);
\draw[draw=none,fill=darkorange25512714] (axis cs:129.5,0) rectangle (axis cs:130.5,100);
\draw[draw=none,fill=darkorange25512714] (axis cs:130.5,0) rectangle (axis cs:131.5,100);
\draw[draw=none,fill=darkorange25512714] (axis cs:131.5,0) rectangle (axis cs:132.5,100);
\draw[draw=none,fill=darkorange25512714] (axis cs:132.5,0) rectangle (axis cs:133.5,100);
\draw[draw=none,fill=darkorange25512714] (axis cs:133.5,0) rectangle (axis cs:134.5,100);
\draw[draw=none,fill=darkorange25512714] (axis cs:134.5,0) rectangle (axis cs:135.5,100);
\draw[draw=none,fill=darkorange25512714] (axis cs:135.5,0) rectangle (axis cs:136.5,100);
\draw[draw=none,fill=darkorange25512714] (axis cs:136.5,0) rectangle (axis cs:137.5,100);
\draw[draw=none,fill=darkorange25512714] (axis cs:137.5,0) rectangle (axis cs:138.5,100);
\draw[draw=none,fill=darkorange25512714] (axis cs:138.5,0) rectangle (axis cs:139.5,100);
\draw[draw=none,fill=darkorange25512714] (axis cs:139.5,0) rectangle (axis cs:140.5,100);
\draw[draw=none,fill=darkorange25512714] (axis cs:140.5,0) rectangle (axis cs:141.5,100);
\draw[draw=none,fill=darkorange25512714] (axis cs:141.5,0) rectangle (axis cs:142.5,100);
\draw[draw=none,fill=darkorange25512714] (axis cs:142.5,0) rectangle (axis cs:143.5,100);
\draw[draw=none,fill=darkorange25512714] (axis cs:143.5,0) rectangle (axis cs:144.5,100);
\draw[draw=none,fill=darkorange25512714] (axis cs:144.5,0) rectangle (axis cs:145.5,100);
\draw[draw=none,fill=darkorange25512714] (axis cs:145.5,0) rectangle (axis cs:146.5,100);
\draw[draw=none,fill=darkorange25512714] (axis cs:146.5,0) rectangle (axis cs:147.5,100);
\draw[draw=none,fill=darkorange25512714] (axis cs:147.5,0) rectangle (axis cs:148.5,100);
\draw[draw=none,fill=darkorange25512714] (axis cs:148.5,0) rectangle (axis cs:149.5,100);
\draw[draw=none,fill=darkorange25512714] (axis cs:149.5,0) rectangle (axis cs:150.5,100);
\draw[draw=none,fill=darkorange25512714] (axis cs:150.5,0) rectangle (axis cs:151.5,100);
\draw[draw=none,fill=darkorange25512714] (axis cs:151.5,0) rectangle (axis cs:152.5,100);
\draw[draw=none,fill=darkorange25512714] (axis cs:152.5,0) rectangle (axis cs:153.5,100);
\draw[draw=none,fill=darkorange25512714] (axis cs:153.5,0) rectangle (axis cs:154.5,100);
\draw[draw=none,fill=darkorange25512714] (axis cs:154.5,0) rectangle (axis cs:155.5,100);
\draw[draw=none,fill=darkorange25512714] (axis cs:155.5,0) rectangle (axis cs:156.5,100);
\draw[draw=none,fill=darkorange25512714] (axis cs:156.5,0) rectangle (axis cs:157.5,100);
\draw[draw=none,fill=darkorange25512714] (axis cs:157.5,0) rectangle (axis cs:158.5,100);
\draw[draw=none,fill=darkorange25512714] (axis cs:158.5,0) rectangle (axis cs:159.5,100);
\draw[draw=none,fill=darkorange25512714] (axis cs:159.5,0) rectangle (axis cs:160.5,100);
\draw[draw=none,fill=darkorange25512714] (axis cs:160.5,0) rectangle (axis cs:161.5,100);
\draw[draw=none,fill=darkorange25512714] (axis cs:161.5,0) rectangle (axis cs:162.5,100);
\draw[draw=none,fill=darkorange25512714] (axis cs:162.5,0) rectangle (axis cs:163.5,100);
\draw[draw=none,fill=darkorange25512714] (axis cs:163.5,0) rectangle (axis cs:164.5,100);
\draw[draw=none,fill=darkorange25512714] (axis cs:164.5,0) rectangle (axis cs:165.5,100);
\draw[draw=none,fill=darkorange25512714] (axis cs:165.5,0) rectangle (axis cs:166.5,100);
\draw[draw=none,fill=darkorange25512714] (axis cs:166.5,0) rectangle (axis cs:167.5,100);
\draw[draw=none,fill=darkorange25512714] (axis cs:167.5,0) rectangle (axis cs:168.5,100);
\draw[draw=none,fill=darkorange25512714] (axis cs:168.5,0) rectangle (axis cs:169.5,100);
\draw[draw=none,fill=darkorange25512714] (axis cs:169.5,0) rectangle (axis cs:170.5,100);
\draw[draw=none,fill=darkorange25512714] (axis cs:170.5,0) rectangle (axis cs:171.5,100);
\draw[draw=none,fill=darkorange25512714] (axis cs:171.5,0) rectangle (axis cs:172.5,100);
\draw[draw=none,fill=darkorange25512714] (axis cs:172.5,0) rectangle (axis cs:173.5,100);
\draw[draw=none,fill=darkorange25512714] (axis cs:173.5,0) rectangle (axis cs:174.5,100);
\draw[draw=none,fill=darkorange25512714] (axis cs:174.5,0) rectangle (axis cs:175.5,100);
\draw[draw=none,fill=darkorange25512714] (axis cs:175.5,0) rectangle (axis cs:176.5,100);
\draw[draw=none,fill=darkorange25512714] (axis cs:176.5,0) rectangle (axis cs:177.5,100);
\draw[draw=none,fill=darkorange25512714] (axis cs:177.5,0) rectangle (axis cs:178.5,100);
\draw[draw=none,fill=darkorange25512714] (axis cs:178.5,0) rectangle (axis cs:179.5,100);
\draw[draw=none,fill=darkorange25512714] (axis cs:179.5,0) rectangle (axis cs:180.5,100);
\draw[draw=none,fill=darkorange25512714] (axis cs:180.5,0) rectangle (axis cs:181.5,100);
\draw[draw=none,fill=darkorange25512714] (axis cs:181.5,0) rectangle (axis cs:182.5,100);
\draw[draw=none,fill=darkorange25512714] (axis cs:182.5,0) rectangle (axis cs:183.5,100);
\draw[draw=none,fill=darkorange25512714] (axis cs:183.5,0) rectangle (axis cs:184.5,100);
\draw[draw=none,fill=darkorange25512714] (axis cs:184.5,0) rectangle (axis cs:185.5,100);
\draw[draw=none,fill=darkorange25512714] (axis cs:185.5,0) rectangle (axis cs:186.5,100);
\draw[draw=none,fill=darkorange25512714] (axis cs:186.5,0) rectangle (axis cs:187.5,100);
\draw[draw=none,fill=darkorange25512714] (axis cs:187.5,0) rectangle (axis cs:188.5,100);
\draw[draw=none,fill=darkorange25512714] (axis cs:188.5,0) rectangle (axis cs:189.5,100);
\draw[draw=none,fill=darkorange25512714] (axis cs:189.5,0) rectangle (axis cs:190.5,100);
\draw[draw=none,fill=darkorange25512714] (axis cs:190.5,0) rectangle (axis cs:191.5,100);
\draw[draw=none,fill=darkorange25512714] (axis cs:191.5,0) rectangle (axis cs:192.5,100);
\draw[draw=none,fill=darkorange25512714] (axis cs:192.5,0) rectangle (axis cs:193.5,100);
\draw[draw=none,fill=darkorange25512714] (axis cs:193.5,0) rectangle (axis cs:194.5,100);
\draw[draw=none,fill=darkorange25512714] (axis cs:194.5,0) rectangle (axis cs:195.5,100);
\draw[draw=none,fill=darkorange25512714] (axis cs:195.5,0) rectangle (axis cs:196.5,100);
\draw[draw=none,fill=darkorange25512714] (axis cs:196.5,0) rectangle (axis cs:197.5,100);
\draw[draw=none,fill=darkorange25512714] (axis cs:197.5,0) rectangle (axis cs:198.5,100);
\draw[draw=none,fill=darkorange25512714] (axis cs:198.5,0) rectangle (axis cs:199.5,100);
\draw[draw=none,fill=darkorange25512714] (axis cs:199.5,0) rectangle (axis cs:200.5,100);
\draw[draw=none,fill=darkorange25512714] (axis cs:200.5,0) rectangle (axis cs:201.5,100);
\draw[draw=none,fill=darkorange25512714] (axis cs:201.5,0) rectangle (axis cs:202.5,100);
\draw[draw=none,fill=darkorange25512714] (axis cs:202.5,0) rectangle (axis cs:203.5,100);
\draw[draw=none,fill=darkorange25512714] (axis cs:203.5,0) rectangle (axis cs:204.5,100);
\draw[draw=none,fill=darkorange25512714] (axis cs:204.5,0) rectangle (axis cs:205.5,100);
\draw[draw=none,fill=darkorange25512714] (axis cs:205.5,0) rectangle (axis cs:206.5,100);
\draw[draw=none,fill=darkorange25512714] (axis cs:206.5,0) rectangle (axis cs:207.5,100);
\draw[draw=none,fill=darkorange25512714] (axis cs:207.5,0) rectangle (axis cs:208.5,100);
\draw[draw=none,fill=darkorange25512714] (axis cs:208.5,0) rectangle (axis cs:209.5,100);
\draw[draw=none,fill=darkorange25512714] (axis cs:209.5,0) rectangle (axis cs:210.5,100);
\draw[draw=none,fill=darkorange25512714] (axis cs:210.5,0) rectangle (axis cs:211.5,100);
\draw[draw=none,fill=darkorange25512714] (axis cs:211.5,0) rectangle (axis cs:212.5,100);
\draw[draw=none,fill=darkorange25512714] (axis cs:212.5,0) rectangle (axis cs:213.5,100);
\draw[draw=none,fill=darkorange25512714] (axis cs:213.5,0) rectangle (axis cs:214.5,100);
\draw[draw=none,fill=darkorange25512714] (axis cs:214.5,0) rectangle (axis cs:215.5,100);
\draw[draw=none,fill=darkorange25512714] (axis cs:215.5,0) rectangle (axis cs:216.5,100);
\draw[draw=none,fill=darkorange25512714] (axis cs:216.5,0) rectangle (axis cs:217.5,100);
\draw[draw=none,fill=darkorange25512714] (axis cs:217.5,0) rectangle (axis cs:218.5,100);
\draw[draw=none,fill=darkorange25512714] (axis cs:218.5,0) rectangle (axis cs:219.5,100);
\draw[draw=none,fill=darkorange25512714] (axis cs:219.5,0) rectangle (axis cs:220.5,100);
\draw[draw=none,fill=darkorange25512714] (axis cs:220.5,0) rectangle (axis cs:221.5,100);
\draw[draw=none,fill=darkorange25512714] (axis cs:221.5,0) rectangle (axis cs:222.5,100);
\draw[draw=none,fill=darkorange25512714] (axis cs:222.5,0) rectangle (axis cs:223.5,100);
\draw[draw=none,fill=darkorange25512714] (axis cs:223.5,0) rectangle (axis cs:224.5,100);
\draw[draw=none,fill=darkorange25512714] (axis cs:224.5,0) rectangle (axis cs:225.5,100);
\draw[draw=none,fill=darkorange25512714] (axis cs:225.5,0) rectangle (axis cs:226.5,100);
\draw[draw=none,fill=darkorange25512714] (axis cs:226.5,0) rectangle (axis cs:227.5,100);
\draw[draw=none,fill=darkorange25512714] (axis cs:227.5,0) rectangle (axis cs:228.5,100);
\draw[draw=none,fill=darkorange25512714] (axis cs:228.5,0) rectangle (axis cs:229.5,100);
\draw[draw=none,fill=darkorange25512714] (axis cs:229.5,0) rectangle (axis cs:230.5,100);
\draw[draw=none,fill=darkorange25512714] (axis cs:230.5,0) rectangle (axis cs:231.5,100);
\draw[draw=none,fill=darkorange25512714] (axis cs:231.5,0) rectangle (axis cs:232.5,100);
\draw[draw=none,fill=darkorange25512714] (axis cs:232.5,0) rectangle (axis cs:233.5,100);
\draw[draw=none,fill=darkorange25512714] (axis cs:233.5,0) rectangle (axis cs:234.5,100);
\draw[draw=none,fill=darkorange25512714] (axis cs:234.5,0) rectangle (axis cs:235.5,100);
\draw[draw=none,fill=darkorange25512714] (axis cs:235.5,0) rectangle (axis cs:236.5,100);
\draw[draw=none,fill=darkorange25512714] (axis cs:236.5,0) rectangle (axis cs:237.5,100);
\draw[draw=none,fill=darkorange25512714] (axis cs:237.5,0) rectangle (axis cs:238.5,100);
\draw[draw=none,fill=darkorange25512714] (axis cs:238.5,0) rectangle (axis cs:239.5,100);
\draw[draw=none,fill=darkorange25512714] (axis cs:239.5,0) rectangle (axis cs:240.5,100);
\draw[draw=none,fill=darkorange25512714] (axis cs:240.5,0) rectangle (axis cs:241.5,100);
\draw[draw=none,fill=darkorange25512714] (axis cs:241.5,0) rectangle (axis cs:242.5,100);
\draw[draw=none,fill=darkorange25512714] (axis cs:242.5,0) rectangle (axis cs:243.5,100);
\draw[draw=none,fill=darkorange25512714] (axis cs:243.5,0) rectangle (axis cs:244.5,100);
\draw[draw=none,fill=darkorange25512714] (axis cs:244.5,0) rectangle (axis cs:245.5,100);
\draw[draw=none,fill=darkorange25512714] (axis cs:245.5,0) rectangle (axis cs:246.5,100);
\draw[draw=none,fill=darkorange25512714] (axis cs:246.5,0) rectangle (axis cs:247.5,100);
\draw[draw=none,fill=darkorange25512714] (axis cs:247.5,0) rectangle (axis cs:248.5,100);
\draw[draw=none,fill=darkorange25512714] (axis cs:248.5,0) rectangle (axis cs:249.5,100);
\draw[draw=none,fill=darkorange25512714] (axis cs:249.5,0) rectangle (axis cs:250.5,100);
\draw[draw=none,fill=darkorange25512714] (axis cs:250.5,0) rectangle (axis cs:251.5,100);
\draw[draw=none,fill=darkorange25512714] (axis cs:251.5,0) rectangle (axis cs:252.5,100);
\draw[draw=none,fill=darkorange25512714] (axis cs:252.5,0) rectangle (axis cs:253.5,100);
\draw[draw=none,fill=darkorange25512714] (axis cs:253.5,0) rectangle (axis cs:254.5,100);
\draw[draw=none,fill=darkorange25512714] (axis cs:254.5,0) rectangle (axis cs:255.5,100);
\draw[draw=none,fill=darkorange25512714] (axis cs:255.5,0) rectangle (axis cs:256.5,100);
\draw[draw=none,fill=darkorange25512714] (axis cs:256.5,0) rectangle (axis cs:257.5,100);
\draw[draw=none,fill=darkorange25512714] (axis cs:257.5,0) rectangle (axis cs:258.5,100);
\draw[draw=none,fill=darkorange25512714] (axis cs:258.5,0) rectangle (axis cs:259.5,100);
\draw[draw=none,fill=darkorange25512714] (axis cs:259.5,0) rectangle (axis cs:260.5,100);
\draw[draw=none,fill=darkorange25512714] (axis cs:260.5,0) rectangle (axis cs:261.5,100);
\draw[draw=none,fill=darkorange25512714] (axis cs:261.5,0) rectangle (axis cs:262.5,100);
\draw[draw=none,fill=darkorange25512714] (axis cs:262.5,0) rectangle (axis cs:263.5,100);
\draw[draw=none,fill=darkorange25512714] (axis cs:263.5,0) rectangle (axis cs:264.5,100);
\draw[draw=none,fill=darkorange25512714] (axis cs:264.5,0) rectangle (axis cs:265.5,100);
\draw[draw=none,fill=darkorange25512714] (axis cs:265.5,0) rectangle (axis cs:266.5,100);
\draw[draw=none,fill=darkorange25512714] (axis cs:266.5,0) rectangle (axis cs:267.5,100);
\draw[draw=none,fill=darkorange25512714] (axis cs:267.5,0) rectangle (axis cs:268.5,100);
\draw[draw=none,fill=darkorange25512714] (axis cs:268.5,0) rectangle (axis cs:269.5,100);
\draw[draw=none,fill=darkorange25512714] (axis cs:269.5,0) rectangle (axis cs:270.5,100);
\draw[draw=none,fill=darkorange25512714] (axis cs:270.5,0) rectangle (axis cs:271.5,100);
\draw[draw=none,fill=darkorange25512714] (axis cs:271.5,0) rectangle (axis cs:272.5,100);
\draw[draw=none,fill=darkorange25512714] (axis cs:272.5,0) rectangle (axis cs:273.5,100);
\draw[draw=none,fill=darkorange25512714] (axis cs:273.5,0) rectangle (axis cs:274.5,100);
\draw[draw=none,fill=darkorange25512714] (axis cs:274.5,0) rectangle (axis cs:275.5,100);
\draw[draw=none,fill=darkorange25512714] (axis cs:275.5,0) rectangle (axis cs:276.5,100);
\draw[draw=none,fill=darkorange25512714] (axis cs:276.5,0) rectangle (axis cs:277.5,100);
\draw[draw=none,fill=darkorange25512714] (axis cs:277.5,0) rectangle (axis cs:278.5,100);
\draw[draw=none,fill=darkorange25512714] (axis cs:278.5,0) rectangle (axis cs:279.5,100);
\draw[draw=none,fill=darkorange25512714] (axis cs:279.5,0) rectangle (axis cs:280.5,100);
\draw[draw=none,fill=darkorange25512714] (axis cs:280.5,0) rectangle (axis cs:281.5,100);
\draw[draw=none,fill=darkorange25512714] (axis cs:281.5,0) rectangle (axis cs:282.5,100);
\draw[draw=none,fill=darkorange25512714] (axis cs:282.5,0) rectangle (axis cs:283.5,100);
\draw[draw=none,fill=darkorange25512714] (axis cs:283.5,0) rectangle (axis cs:284.5,100);
\draw[draw=none,fill=darkorange25512714] (axis cs:284.5,0) rectangle (axis cs:285.5,100);
\draw[draw=none,fill=darkorange25512714] (axis cs:285.5,0) rectangle (axis cs:286.5,100);
\draw[draw=none,fill=darkorange25512714] (axis cs:286.5,0) rectangle (axis cs:287.5,100);
\draw[draw=none,fill=darkorange25512714] (axis cs:287.5,0) rectangle (axis cs:288.5,100);
\draw[draw=none,fill=darkorange25512714] (axis cs:288.5,0) rectangle (axis cs:289.5,100);
\draw[draw=none,fill=darkorange25512714] (axis cs:289.5,0) rectangle (axis cs:290.5,100);
\draw[draw=none,fill=darkorange25512714] (axis cs:290.5,0) rectangle (axis cs:291.5,100);
\draw[draw=none,fill=darkorange25512714] (axis cs:291.5,0) rectangle (axis cs:292.5,100);
\draw[draw=none,fill=darkorange25512714] (axis cs:292.5,0) rectangle (axis cs:293.5,100);
\draw[draw=none,fill=darkorange25512714] (axis cs:293.5,0) rectangle (axis cs:294.5,100);
\draw[draw=none,fill=darkorange25512714] (axis cs:294.5,0) rectangle (axis cs:295.5,100);
\draw[draw=none,fill=darkorange25512714] (axis cs:295.5,0) rectangle (axis cs:296.5,100);
\draw[draw=none,fill=darkorange25512714] (axis cs:296.5,0) rectangle (axis cs:297.5,100);
\draw[draw=none,fill=darkorange25512714] (axis cs:297.5,0) rectangle (axis cs:298.5,100);
\draw[draw=none,fill=darkorange25512714] (axis cs:298.5,0) rectangle (axis cs:299.5,100);
\draw[draw=none,fill=darkorange25512714] (axis cs:299.5,0) rectangle (axis cs:300.5,100);
\draw[draw=none,fill=darkorange25512714] (axis cs:300.5,0) rectangle (axis cs:301.5,100);
\draw[draw=none,fill=darkorange25512714] (axis cs:301.5,0) rectangle (axis cs:302.5,100);
\draw[draw=none,fill=darkorange25512714] (axis cs:302.5,0) rectangle (axis cs:303.5,100);
\draw[draw=none,fill=darkorange25512714] (axis cs:303.5,0) rectangle (axis cs:304.5,100);
\draw[draw=none,fill=darkorange25512714] (axis cs:304.5,0) rectangle (axis cs:305.5,100);
\draw[draw=none,fill=darkorange25512714] (axis cs:305.5,0) rectangle (axis cs:306.5,100);
\draw[draw=none,fill=darkorange25512714] (axis cs:306.5,0) rectangle (axis cs:307.5,100);
\draw[draw=none,fill=darkorange25512714] (axis cs:307.5,0) rectangle (axis cs:308.5,100);
\draw[draw=none,fill=darkorange25512714] (axis cs:308.5,0) rectangle (axis cs:309.5,100);
\draw[draw=none,fill=darkorange25512714] (axis cs:309.5,0) rectangle (axis cs:310.5,100);
\draw[draw=none,fill=darkorange25512714] (axis cs:310.5,0) rectangle (axis cs:311.5,100);
\draw[draw=none,fill=darkorange25512714] (axis cs:311.5,0) rectangle (axis cs:312.5,100);
\draw[draw=none,fill=darkorange25512714] (axis cs:312.5,0) rectangle (axis cs:313.5,100);
\draw[draw=none,fill=darkorange25512714] (axis cs:313.5,0) rectangle (axis cs:314.5,100);
\draw[draw=none,fill=darkorange25512714] (axis cs:314.5,0) rectangle (axis cs:315.5,100);
\draw[draw=none,fill=darkorange25512714] (axis cs:315.5,0) rectangle (axis cs:316.5,100);
\draw[draw=none,fill=darkorange25512714] (axis cs:316.5,0) rectangle (axis cs:317.5,100);
\draw[draw=none,fill=darkorange25512714] (axis cs:317.5,0) rectangle (axis cs:318.5,100);
\draw[draw=none,fill=darkorange25512714] (axis cs:318.5,0) rectangle (axis cs:319.5,100);
\draw[draw=none,fill=darkorange25512714] (axis cs:319.5,0) rectangle (axis cs:320.5,100);
\draw[draw=none,fill=darkorange25512714] (axis cs:320.5,0) rectangle (axis cs:321.5,100);
\draw[draw=none,fill=darkorange25512714] (axis cs:321.5,0) rectangle (axis cs:322.5,100);
\draw[draw=none,fill=darkorange25512714] (axis cs:322.5,0) rectangle (axis cs:323.5,100);
\draw[draw=none,fill=darkorange25512714] (axis cs:323.5,0) rectangle (axis cs:324.5,100);
\draw[draw=none,fill=darkorange25512714] (axis cs:324.5,0) rectangle (axis cs:325.5,100);
\draw[draw=none,fill=darkorange25512714] (axis cs:325.5,0) rectangle (axis cs:326.5,100);
\draw[draw=none,fill=darkorange25512714] (axis cs:326.5,0) rectangle (axis cs:327.5,100);
\draw[draw=none,fill=darkorange25512714] (axis cs:327.5,0) rectangle (axis cs:328.5,100);
\draw[draw=none,fill=darkorange25512714] (axis cs:328.5,0) rectangle (axis cs:329.5,100);
\draw[draw=none,fill=darkorange25512714] (axis cs:329.5,0) rectangle (axis cs:330.5,100);
\draw[draw=none,fill=darkorange25512714] (axis cs:330.5,0) rectangle (axis cs:331.5,100);
\draw[draw=none,fill=darkorange25512714] (axis cs:331.5,0) rectangle (axis cs:332.5,100);
\draw[draw=none,fill=darkorange25512714] (axis cs:332.5,0) rectangle (axis cs:333.5,100);
\draw[draw=none,fill=darkorange25512714] (axis cs:333.5,0) rectangle (axis cs:334.5,100);
\draw[draw=none,fill=darkorange25512714] (axis cs:334.5,0) rectangle (axis cs:335.5,100);
\draw[draw=none,fill=darkorange25512714] (axis cs:335.5,0) rectangle (axis cs:336.5,100);
\draw[draw=none,fill=darkorange25512714] (axis cs:336.5,0) rectangle (axis cs:337.5,100);
\draw[draw=none,fill=darkorange25512714] (axis cs:337.5,0) rectangle (axis cs:338.5,100);
\draw[draw=none,fill=darkorange25512714] (axis cs:338.5,0) rectangle (axis cs:339.5,100);
\draw[draw=none,fill=darkorange25512714] (axis cs:339.5,0) rectangle (axis cs:340.5,100);
\draw[draw=none,fill=darkorange25512714] (axis cs:340.5,0) rectangle (axis cs:341.5,100);
\draw[draw=none,fill=darkorange25512714] (axis cs:341.5,0) rectangle (axis cs:342.5,100);
\draw[draw=none,fill=darkorange25512714] (axis cs:342.5,0) rectangle (axis cs:343.5,100);
\draw[draw=none,fill=darkorange25512714] (axis cs:343.5,0) rectangle (axis cs:344.5,100);
\draw[draw=none,fill=darkorange25512714] (axis cs:344.5,0) rectangle (axis cs:345.5,100);
\draw[draw=none,fill=darkorange25512714] (axis cs:345.5,0) rectangle (axis cs:346.5,100);
\draw[draw=none,fill=darkorange25512714] (axis cs:346.5,0) rectangle (axis cs:347.5,100);
\draw[draw=none,fill=darkorange25512714] (axis cs:347.5,0) rectangle (axis cs:348.5,100);
\draw[draw=none,fill=darkorange25512714] (axis cs:348.5,0) rectangle (axis cs:349.5,100);
\draw[draw=none,fill=darkorange25512714] (axis cs:349.5,0) rectangle (axis cs:350.5,100);
\draw[draw=none,fill=darkorange25512714] (axis cs:350.5,0) rectangle (axis cs:351.5,100);
\draw[draw=none,fill=darkorange25512714] (axis cs:351.5,0) rectangle (axis cs:352.5,100);
\draw[draw=none,fill=darkorange25512714] (axis cs:352.5,0) rectangle (axis cs:353.5,100);
\draw[draw=none,fill=darkorange25512714] (axis cs:353.5,0) rectangle (axis cs:354.5,100);
\draw[draw=none,fill=darkorange25512714] (axis cs:354.5,0) rectangle (axis cs:355.5,100);
\draw[draw=none,fill=darkorange25512714] (axis cs:355.5,0) rectangle (axis cs:356.5,100);
\draw[draw=none,fill=darkorange25512714] (axis cs:356.5,0) rectangle (axis cs:357.5,100);
\draw[draw=none,fill=darkorange25512714] (axis cs:357.5,0) rectangle (axis cs:358.5,100);
\draw[draw=none,fill=darkorange25512714] (axis cs:358.5,0) rectangle (axis cs:359.5,100);
\draw[draw=none,fill=darkorange25512714] (axis cs:359.5,0) rectangle (axis cs:360.5,100);
\draw[draw=none,fill=darkorange25512714] (axis cs:360.5,0) rectangle (axis cs:361.5,100);
\draw[draw=none,fill=darkorange25512714] (axis cs:361.5,0) rectangle (axis cs:362.5,100);
\draw[draw=none,fill=darkorange25512714] (axis cs:362.5,0) rectangle (axis cs:363.5,100);
\draw[draw=none,fill=darkorange25512714] (axis cs:363.5,0) rectangle (axis cs:364.5,100);
\draw[draw=none,fill=darkorange25512714] (axis cs:364.5,0) rectangle (axis cs:365.5,100);
\draw[draw=none,fill=darkorange25512714] (axis cs:365.5,0) rectangle (axis cs:366.5,100);
\draw[draw=none,fill=darkorange25512714] (axis cs:366.5,0) rectangle (axis cs:367.5,100);
\draw[draw=none,fill=darkorange25512714] (axis cs:367.5,0) rectangle (axis cs:368.5,100);
\draw[draw=none,fill=darkorange25512714] (axis cs:368.5,0) rectangle (axis cs:369.5,100);
\draw[draw=none,fill=darkorange25512714] (axis cs:369.5,0) rectangle (axis cs:370.5,100);
\draw[draw=none,fill=darkorange25512714] (axis cs:370.5,0) rectangle (axis cs:371.5,100);
\draw[draw=none,fill=darkorange25512714] (axis cs:371.5,0) rectangle (axis cs:372.5,100);
\draw[draw=none,fill=darkorange25512714] (axis cs:372.5,0) rectangle (axis cs:373.5,100);
\draw[draw=none,fill=darkorange25512714] (axis cs:373.5,0) rectangle (axis cs:374.5,100);
\draw[draw=none,fill=darkorange25512714] (axis cs:374.5,0) rectangle (axis cs:375.5,100);
\draw[draw=none,fill=darkorange25512714] (axis cs:375.5,0) rectangle (axis cs:376.5,100);
\draw[draw=none,fill=darkorange25512714] (axis cs:376.5,0) rectangle (axis cs:377.5,100);
\draw[draw=none,fill=darkorange25512714] (axis cs:377.5,0) rectangle (axis cs:378.5,100);
\draw[draw=none,fill=darkorange25512714] (axis cs:378.5,0) rectangle (axis cs:379.5,100);
\draw[draw=none,fill=darkorange25512714] (axis cs:379.5,0) rectangle (axis cs:380.5,100);
\draw[draw=none,fill=darkorange25512714] (axis cs:380.5,0) rectangle (axis cs:381.5,100);
\draw[draw=none,fill=darkorange25512714] (axis cs:381.5,0) rectangle (axis cs:382.5,100);
\draw[draw=none,fill=darkorange25512714] (axis cs:382.5,0) rectangle (axis cs:383.5,100);
\draw[draw=none,fill=darkorange25512714] (axis cs:383.5,0) rectangle (axis cs:384.5,100);
\draw[draw=none,fill=darkorange25512714] (axis cs:384.5,0) rectangle (axis cs:385.5,100);
\draw[draw=none,fill=darkorange25512714] (axis cs:385.5,0) rectangle (axis cs:386.5,100);
\draw[draw=none,fill=darkorange25512714] (axis cs:386.5,0) rectangle (axis cs:387.5,100);
\draw[draw=none,fill=darkorange25512714] (axis cs:387.5,0) rectangle (axis cs:388.5,100);
\draw[draw=none,fill=darkorange25512714] (axis cs:388.5,0) rectangle (axis cs:389.5,100);
\draw[draw=none,fill=darkorange25512714] (axis cs:389.5,0) rectangle (axis cs:390.5,100);
\draw[draw=none,fill=darkorange25512714] (axis cs:390.5,0) rectangle (axis cs:391.5,100);
\draw[draw=none,fill=darkorange25512714] (axis cs:391.5,0) rectangle (axis cs:392.5,100);
\draw[draw=none,fill=darkorange25512714] (axis cs:392.5,0) rectangle (axis cs:393.5,100);
\draw[draw=none,fill=darkorange25512714] (axis cs:393.5,0) rectangle (axis cs:394.5,100);
\draw[draw=none,fill=darkorange25512714] (axis cs:394.5,0) rectangle (axis cs:395.5,100);
\draw[draw=none,fill=darkorange25512714] (axis cs:395.5,0) rectangle (axis cs:396.5,100);
\draw[draw=none,fill=darkorange25512714] (axis cs:396.5,0) rectangle (axis cs:397.5,100);
\draw[draw=none,fill=darkorange25512714] (axis cs:397.5,0) rectangle (axis cs:398.5,100);
\draw[draw=none,fill=darkorange25512714] (axis cs:398.5,0) rectangle (axis cs:399.5,100);
\draw[draw=none,fill=darkorange25512714] (axis cs:399.5,0) rectangle (axis cs:400.5,100);
\draw[draw=none,fill=darkorange25512714] (axis cs:400.5,0) rectangle (axis cs:401.5,100);
\draw[draw=none,fill=darkorange25512714] (axis cs:401.5,0) rectangle (axis cs:402.5,100);
\draw[draw=none,fill=darkorange25512714] (axis cs:402.5,0) rectangle (axis cs:403.5,100);
\draw[draw=none,fill=darkorange25512714] (axis cs:403.5,0) rectangle (axis cs:404.5,100);
\draw[draw=none,fill=darkorange25512714] (axis cs:404.5,0) rectangle (axis cs:405.5,100);
\draw[draw=none,fill=darkorange25512714] (axis cs:405.5,0) rectangle (axis cs:406.5,100);
\draw[draw=none,fill=darkorange25512714] (axis cs:406.5,0) rectangle (axis cs:407.5,100);
\draw[draw=none,fill=darkorange25512714] (axis cs:407.5,0) rectangle (axis cs:408.5,100);
\draw[draw=none,fill=darkorange25512714] (axis cs:408.5,0) rectangle (axis cs:409.5,100);
\draw[draw=none,fill=darkorange25512714] (axis cs:409.5,0) rectangle (axis cs:410.5,100);
\draw[draw=none,fill=darkorange25512714] (axis cs:410.5,0) rectangle (axis cs:411.5,100);
\draw[draw=none,fill=darkorange25512714] (axis cs:411.5,0) rectangle (axis cs:412.5,100);
\draw[draw=none,fill=darkorange25512714] (axis cs:412.5,0) rectangle (axis cs:413.5,100);
\draw[draw=none,fill=darkorange25512714] (axis cs:413.5,0) rectangle (axis cs:414.5,100);
\draw[draw=none,fill=darkorange25512714] (axis cs:414.5,0) rectangle (axis cs:415.5,100);
\draw[draw=none,fill=darkorange25512714] (axis cs:415.5,0) rectangle (axis cs:416.5,100);
\draw[draw=none,fill=darkorange25512714] (axis cs:416.5,0) rectangle (axis cs:417.5,100);
\draw[draw=none,fill=darkorange25512714] (axis cs:417.5,0) rectangle (axis cs:418.5,100);
\draw[draw=none,fill=darkorange25512714] (axis cs:418.5,0) rectangle (axis cs:419.5,100);
\draw[draw=none,fill=darkorange25512714] (axis cs:419.5,0) rectangle (axis cs:420.5,100);
\draw[draw=none,fill=darkorange25512714] (axis cs:420.5,0) rectangle (axis cs:421.5,100);
\draw[draw=none,fill=darkorange25512714] (axis cs:421.5,0) rectangle (axis cs:422.5,100);
\draw[draw=none,fill=darkorange25512714] (axis cs:422.5,0) rectangle (axis cs:423.5,100);
\draw[draw=none,fill=darkorange25512714] (axis cs:423.5,0) rectangle (axis cs:424.5,100);
\draw[draw=none,fill=darkorange25512714] (axis cs:424.5,0) rectangle (axis cs:425.5,100);
\draw[draw=none,fill=darkorange25512714] (axis cs:425.5,0) rectangle (axis cs:426.5,100);
\draw[draw=none,fill=darkorange25512714] (axis cs:426.5,0) rectangle (axis cs:427.5,100);
\draw[draw=none,fill=darkorange25512714] (axis cs:427.5,0) rectangle (axis cs:428.5,100);
\draw[draw=none,fill=darkorange25512714] (axis cs:428.5,0) rectangle (axis cs:429.5,100);
\draw[draw=none,fill=darkorange25512714] (axis cs:429.5,0) rectangle (axis cs:430.5,100);
\draw[draw=none,fill=darkorange25512714] (axis cs:430.5,0) rectangle (axis cs:431.5,100);
\draw[draw=none,fill=darkorange25512714] (axis cs:431.5,0) rectangle (axis cs:432.5,100);
\draw[draw=none,fill=darkorange25512714] (axis cs:432.5,0) rectangle (axis cs:433.5,100);
\draw[draw=none,fill=darkorange25512714] (axis cs:433.5,0) rectangle (axis cs:434.5,100);
\draw[draw=none,fill=darkorange25512714] (axis cs:434.5,0) rectangle (axis cs:435.5,100);
\draw[draw=none,fill=darkorange25512714] (axis cs:435.5,0) rectangle (axis cs:436.5,100);
\draw[draw=none,fill=darkorange25512714] (axis cs:436.5,0) rectangle (axis cs:437.5,100);
\draw[draw=none,fill=darkorange25512714] (axis cs:437.5,0) rectangle (axis cs:438.5,100);
\draw[draw=none,fill=darkorange25512714] (axis cs:438.5,0) rectangle (axis cs:439.5,100);
\draw[draw=none,fill=darkorange25512714] (axis cs:439.5,0) rectangle (axis cs:440.5,100);
\draw[draw=none,fill=darkorange25512714] (axis cs:440.5,0) rectangle (axis cs:441.5,100);
\draw[draw=none,fill=darkorange25512714] (axis cs:441.5,0) rectangle (axis cs:442.5,100);
\draw[draw=none,fill=darkorange25512714] (axis cs:442.5,0) rectangle (axis cs:443.5,100);
\draw[draw=none,fill=darkorange25512714] (axis cs:443.5,0) rectangle (axis cs:444.5,100);
\draw[draw=none,fill=darkorange25512714] (axis cs:444.5,0) rectangle (axis cs:445.5,100);
\draw[draw=none,fill=darkorange25512714] (axis cs:445.5,0) rectangle (axis cs:446.5,100);
\draw[draw=none,fill=darkorange25512714] (axis cs:446.5,0) rectangle (axis cs:447.5,100);
\draw[draw=none,fill=darkorange25512714] (axis cs:447.5,0) rectangle (axis cs:448.5,100);
\draw[draw=none,fill=darkorange25512714] (axis cs:448.5,0) rectangle (axis cs:449.5,100);
\draw[draw=none,fill=darkorange25512714] (axis cs:449.5,0) rectangle (axis cs:450.5,100);
\draw[draw=none,fill=darkorange25512714] (axis cs:450.5,0) rectangle (axis cs:451.5,100);
\draw[draw=none,fill=darkorange25512714] (axis cs:451.5,0) rectangle (axis cs:452.5,100);
\draw[draw=none,fill=darkorange25512714] (axis cs:452.5,0) rectangle (axis cs:453.5,100);
\draw[draw=none,fill=darkorange25512714] (axis cs:453.5,0) rectangle (axis cs:454.5,100);
\draw[draw=none,fill=darkorange25512714] (axis cs:454.5,0) rectangle (axis cs:455.5,100);
\draw[draw=none,fill=darkorange25512714] (axis cs:455.5,0) rectangle (axis cs:456.5,100);
\draw[draw=none,fill=darkorange25512714] (axis cs:456.5,0) rectangle (axis cs:457.5,100);
\draw[draw=none,fill=darkorange25512714] (axis cs:457.5,0) rectangle (axis cs:458.5,100);
\draw[draw=none,fill=darkorange25512714] (axis cs:458.5,0) rectangle (axis cs:459.5,100);
\draw[draw=none,fill=darkorange25512714] (axis cs:459.5,0) rectangle (axis cs:460.5,100);
\draw[draw=none,fill=darkorange25512714] (axis cs:460.5,0) rectangle (axis cs:461.5,100);
\draw[draw=none,fill=darkorange25512714] (axis cs:461.5,0) rectangle (axis cs:462.5,100);
\draw[draw=none,fill=darkorange25512714] (axis cs:462.5,0) rectangle (axis cs:463.5,100);
\draw[draw=none,fill=darkorange25512714] (axis cs:463.5,0) rectangle (axis cs:464.5,100);
\draw[draw=none,fill=darkorange25512714] (axis cs:464.5,0) rectangle (axis cs:465.5,100);
\draw[draw=none,fill=darkorange25512714] (axis cs:465.5,0) rectangle (axis cs:466.5,100);
\draw[draw=none,fill=darkorange25512714] (axis cs:466.5,0) rectangle (axis cs:467.5,100);
\draw[draw=none,fill=darkorange25512714] (axis cs:467.5,0) rectangle (axis cs:468.5,100);
\draw[draw=none,fill=darkorange25512714] (axis cs:468.5,0) rectangle (axis cs:469.5,100);
\draw[draw=none,fill=darkorange25512714] (axis cs:469.5,0) rectangle (axis cs:470.5,100);
\draw[draw=none,fill=darkorange25512714] (axis cs:470.5,0) rectangle (axis cs:471.5,100);
\draw[draw=none,fill=darkorange25512714] (axis cs:471.5,0) rectangle (axis cs:472.5,100);
\draw[draw=none,fill=darkorange25512714] (axis cs:472.5,0) rectangle (axis cs:473.5,100);
\draw[draw=none,fill=darkorange25512714] (axis cs:473.5,0) rectangle (axis cs:474.5,100);
\draw[draw=none,fill=darkorange25512714] (axis cs:474.5,0) rectangle (axis cs:475.5,100);
\draw[draw=none,fill=darkorange25512714] (axis cs:475.5,0) rectangle (axis cs:476.5,100);
\draw[draw=none,fill=darkorange25512714] (axis cs:476.5,0) rectangle (axis cs:477.5,100);
\draw[draw=none,fill=darkorange25512714] (axis cs:477.5,0) rectangle (axis cs:478.5,100);
\draw[draw=none,fill=darkorange25512714] (axis cs:478.5,0) rectangle (axis cs:479.5,100);
\draw[draw=none,fill=darkorange25512714] (axis cs:479.5,0) rectangle (axis cs:480.5,100);
\draw[draw=none,fill=darkorange25512714] (axis cs:480.5,0) rectangle (axis cs:481.5,100);
\draw[draw=none,fill=darkorange25512714] (axis cs:481.5,0) rectangle (axis cs:482.5,100);
\draw[draw=none,fill=darkorange25512714] (axis cs:482.5,0) rectangle (axis cs:483.5,100);
\draw[draw=none,fill=darkorange25512714] (axis cs:483.5,0) rectangle (axis cs:484.5,100);
\draw[draw=none,fill=darkorange25512714] (axis cs:484.5,0) rectangle (axis cs:485.5,100);
\draw[draw=none,fill=darkorange25512714] (axis cs:485.5,0) rectangle (axis cs:486.5,100);
\draw[draw=none,fill=darkorange25512714] (axis cs:486.5,0) rectangle (axis cs:487.5,100);
\draw[draw=none,fill=darkorange25512714] (axis cs:487.5,0) rectangle (axis cs:488.5,100);
\draw[draw=none,fill=darkorange25512714] (axis cs:488.5,0) rectangle (axis cs:489.5,100);
\draw[draw=none,fill=darkorange25512714] (axis cs:489.5,0) rectangle (axis cs:490.5,100);
\draw[draw=none,fill=darkorange25512714] (axis cs:490.5,0) rectangle (axis cs:491.5,100);
\draw[draw=none,fill=darkorange25512714] (axis cs:491.5,0) rectangle (axis cs:492.5,100);
\draw[draw=none,fill=darkorange25512714] (axis cs:492.5,0) rectangle (axis cs:493.5,100);
\draw[draw=none,fill=darkorange25512714] (axis cs:493.5,0) rectangle (axis cs:494.5,100);
\draw[draw=none,fill=darkorange25512714] (axis cs:494.5,0) rectangle (axis cs:495.5,100);
\draw[draw=none,fill=darkorange25512714] (axis cs:495.5,0) rectangle (axis cs:496.5,100);
\draw[draw=none,fill=darkorange25512714] (axis cs:496.5,0) rectangle (axis cs:497.5,100);
\draw[draw=none,fill=darkorange25512714] (axis cs:497.5,0) rectangle (axis cs:498.5,100);
\draw[draw=none,fill=darkorange25512714] (axis cs:498.5,0) rectangle (axis cs:499.5,100);
\draw[draw=none,fill=darkorange25512714] (axis cs:499.5,0) rectangle (axis cs:500.5,100);
\draw[draw=none,fill=darkorange25512714] (axis cs:500.5,0) rectangle (axis cs:501.5,100);
\draw[draw=none,fill=darkorange25512714] (axis cs:501.5,0) rectangle (axis cs:502.5,100);
\draw[draw=none,fill=darkorange25512714] (axis cs:502.5,0) rectangle (axis cs:503.5,100);
\draw[draw=none,fill=darkorange25512714] (axis cs:503.5,0) rectangle (axis cs:504.5,100);
\draw[draw=none,fill=darkorange25512714] (axis cs:504.5,0) rectangle (axis cs:505.5,100);
\draw[draw=none,fill=darkorange25512714] (axis cs:505.5,0) rectangle (axis cs:506.5,100);
\draw[draw=none,fill=darkorange25512714] (axis cs:506.5,0) rectangle (axis cs:507.5,100);
\draw[draw=none,fill=darkorange25512714] (axis cs:507.5,0) rectangle (axis cs:508.5,100);
\draw[draw=none,fill=darkorange25512714] (axis cs:508.5,0) rectangle (axis cs:509.5,100);
\draw[draw=none,fill=darkorange25512714] (axis cs:509.5,0) rectangle (axis cs:510.5,100);
\draw[draw=none,fill=darkorange25512714] (axis cs:510.5,0) rectangle (axis cs:511.5,100);
\draw[draw=none,fill=darkorange25512714] (axis cs:511.5,0) rectangle (axis cs:512.5,100);
\draw[draw=none,fill=darkorange25512714] (axis cs:512.5,0) rectangle (axis cs:513.5,100);
\draw[draw=none,fill=darkorange25512714] (axis cs:513.5,0) rectangle (axis cs:514.5,100);
\draw[draw=none,fill=darkorange25512714] (axis cs:514.5,0) rectangle (axis cs:515.5,100);
\draw[draw=none,fill=darkorange25512714] (axis cs:515.5,0) rectangle (axis cs:516.5,100);
\draw[draw=none,fill=darkorange25512714] (axis cs:516.5,0) rectangle (axis cs:517.5,100);
\draw[draw=none,fill=darkorange25512714] (axis cs:517.5,0) rectangle (axis cs:518.5,100);
\draw[draw=none,fill=darkorange25512714] (axis cs:518.5,0) rectangle (axis cs:519.5,100);
\draw[draw=none,fill=darkorange25512714] (axis cs:519.5,0) rectangle (axis cs:520.5,100);
\draw[draw=none,fill=darkorange25512714] (axis cs:520.5,0) rectangle (axis cs:521.5,100);
\draw[draw=none,fill=darkorange25512714] (axis cs:521.5,0) rectangle (axis cs:522.5,100);
\draw[draw=none,fill=darkorange25512714] (axis cs:522.5,0) rectangle (axis cs:523.5,100);
\draw[draw=none,fill=darkorange25512714] (axis cs:523.5,0) rectangle (axis cs:524.5,100);
\draw[draw=none,fill=darkorange25512714] (axis cs:524.5,0) rectangle (axis cs:525.5,100);
\draw[draw=none,fill=darkorange25512714] (axis cs:525.5,0) rectangle (axis cs:526.5,100);
\draw[draw=none,fill=darkorange25512714] (axis cs:526.5,0) rectangle (axis cs:527.5,100);
\draw[draw=none,fill=darkorange25512714] (axis cs:527.5,0) rectangle (axis cs:528.5,100);
\draw[draw=none,fill=darkorange25512714] (axis cs:528.5,0) rectangle (axis cs:529.5,100);
\draw[draw=none,fill=darkorange25512714] (axis cs:529.5,0) rectangle (axis cs:530.5,100);
\draw[draw=none,fill=darkorange25512714] (axis cs:530.5,0) rectangle (axis cs:531.5,100);
\draw[draw=none,fill=darkorange25512714] (axis cs:531.5,0) rectangle (axis cs:532.5,100);
\draw[draw=none,fill=darkorange25512714] (axis cs:532.5,0) rectangle (axis cs:533.5,100);
\draw[draw=none,fill=darkorange25512714] (axis cs:533.5,0) rectangle (axis cs:534.5,100);
\draw[draw=none,fill=darkorange25512714] (axis cs:534.5,0) rectangle (axis cs:535.5,100);
\draw[draw=none,fill=darkorange25512714] (axis cs:535.5,0) rectangle (axis cs:536.5,100);
\draw[draw=none,fill=darkorange25512714] (axis cs:536.5,0) rectangle (axis cs:537.5,100);
\draw[draw=none,fill=darkorange25512714] (axis cs:537.5,0) rectangle (axis cs:538.5,100);
\draw[draw=none,fill=darkorange25512714] (axis cs:538.5,0) rectangle (axis cs:539.5,100);
\draw[draw=none,fill=darkorange25512714] (axis cs:539.5,0) rectangle (axis cs:540.5,100);
\draw[draw=none,fill=darkorange25512714] (axis cs:540.5,0) rectangle (axis cs:541.5,100);
\draw[draw=none,fill=darkorange25512714] (axis cs:541.5,0) rectangle (axis cs:542.5,100);
\draw[draw=none,fill=darkorange25512714] (axis cs:542.5,0) rectangle (axis cs:543.5,100);
\draw[draw=none,fill=darkorange25512714] (axis cs:543.5,0) rectangle (axis cs:544.5,100);
\draw[draw=none,fill=darkorange25512714] (axis cs:544.5,0) rectangle (axis cs:545.5,100);
\draw[draw=none,fill=darkorange25512714] (axis cs:545.5,0) rectangle (axis cs:546.5,100);
\draw[draw=none,fill=darkorange25512714] (axis cs:546.5,0) rectangle (axis cs:547.5,100);
\draw[draw=none,fill=darkorange25512714] (axis cs:547.5,0) rectangle (axis cs:548.5,100);
\draw[draw=none,fill=darkorange25512714] (axis cs:548.5,0) rectangle (axis cs:549.5,100);
\draw[draw=none,fill=darkorange25512714] (axis cs:549.5,0) rectangle (axis cs:550.5,100);
\draw[draw=none,fill=darkorange25512714] (axis cs:550.5,0) rectangle (axis cs:551.5,100);
\draw[draw=none,fill=darkorange25512714] (axis cs:551.5,0) rectangle (axis cs:552.5,100);
\draw[draw=none,fill=darkorange25512714] (axis cs:552.5,0) rectangle (axis cs:553.5,100);
\draw[draw=none,fill=darkorange25512714] (axis cs:553.5,0) rectangle (axis cs:554.5,100);
\draw[draw=none,fill=darkorange25512714] (axis cs:554.5,0) rectangle (axis cs:555.5,100);
\draw[draw=none,fill=darkorange25512714] (axis cs:555.5,0) rectangle (axis cs:556.5,100);
\draw[draw=none,fill=darkorange25512714] (axis cs:556.5,0) rectangle (axis cs:557.5,100);
\draw[draw=none,fill=darkorange25512714] (axis cs:557.5,0) rectangle (axis cs:558.5,100);
\draw[draw=none,fill=darkorange25512714] (axis cs:558.5,0) rectangle (axis cs:559.5,100);
\draw[draw=none,fill=darkorange25512714] (axis cs:559.5,0) rectangle (axis cs:560.5,100);
\draw[draw=none,fill=darkorange25512714] (axis cs:560.5,0) rectangle (axis cs:561.5,100);
\draw[draw=none,fill=darkorange25512714] (axis cs:561.5,0) rectangle (axis cs:562.5,100);
\draw[draw=none,fill=darkorange25512714] (axis cs:562.5,0) rectangle (axis cs:563.5,100);
\draw[draw=none,fill=darkorange25512714] (axis cs:563.5,0) rectangle (axis cs:564.5,100);
\draw[draw=none,fill=darkorange25512714] (axis cs:564.5,0) rectangle (axis cs:565.5,100);
\draw[draw=none,fill=darkorange25512714] (axis cs:565.5,0) rectangle (axis cs:566.5,100);
\draw[draw=none,fill=darkorange25512714] (axis cs:566.5,0) rectangle (axis cs:567.5,100);
\draw[draw=none,fill=darkorange25512714] (axis cs:567.5,0) rectangle (axis cs:568.5,100);
\draw[draw=none,fill=darkorange25512714] (axis cs:568.5,0) rectangle (axis cs:569.5,100);
\draw[draw=none,fill=darkorange25512714] (axis cs:569.5,0) rectangle (axis cs:570.5,100);
\draw[draw=none,fill=darkorange25512714] (axis cs:570.5,0) rectangle (axis cs:571.5,100);
\draw[draw=none,fill=darkorange25512714] (axis cs:571.5,0) rectangle (axis cs:572.5,100);
\draw[draw=none,fill=darkorange25512714] (axis cs:572.5,0) rectangle (axis cs:573.5,100);
\draw[draw=none,fill=darkorange25512714] (axis cs:573.5,0) rectangle (axis cs:574.5,100);
\draw[draw=none,fill=darkorange25512714] (axis cs:574.5,0) rectangle (axis cs:575.5,100);
\draw[draw=none,fill=darkorange25512714] (axis cs:575.5,0) rectangle (axis cs:576.5,100);
\draw[draw=none,fill=darkorange25512714] (axis cs:576.5,0) rectangle (axis cs:577.5,100);
\draw[draw=none,fill=darkorange25512714] (axis cs:577.5,0) rectangle (axis cs:578.5,100);
\draw[draw=none,fill=darkorange25512714] (axis cs:578.5,0) rectangle (axis cs:579.5,100);
\draw[draw=none,fill=darkorange25512714] (axis cs:579.5,0) rectangle (axis cs:580.5,100);
\draw[draw=none,fill=darkorange25512714] (axis cs:580.5,0) rectangle (axis cs:581.5,100);
\draw[draw=none,fill=darkorange25512714] (axis cs:581.5,0) rectangle (axis cs:582.5,100);
\draw[draw=none,fill=darkorange25512714] (axis cs:582.5,0) rectangle (axis cs:583.5,100);
\draw[draw=none,fill=darkorange25512714] (axis cs:583.5,0) rectangle (axis cs:584.5,100);
\draw[draw=none,fill=darkorange25512714] (axis cs:584.5,0) rectangle (axis cs:585.5,100);
\draw[draw=none,fill=darkorange25512714] (axis cs:585.5,0) rectangle (axis cs:586.5,100);
\draw[draw=none,fill=darkorange25512714] (axis cs:586.5,0) rectangle (axis cs:587.5,100);
\draw[draw=none,fill=darkorange25512714] (axis cs:587.5,0) rectangle (axis cs:588.5,100);
\draw[draw=none,fill=darkorange25512714] (axis cs:588.5,0) rectangle (axis cs:589.5,100);
\draw[draw=none,fill=darkorange25512714] (axis cs:589.5,0) rectangle (axis cs:590.5,100);
\draw[draw=none,fill=darkorange25512714] (axis cs:590.5,0) rectangle (axis cs:591.5,100);
\draw[draw=none,fill=darkorange25512714] (axis cs:591.5,0) rectangle (axis cs:592.5,100);
\draw[draw=none,fill=darkorange25512714] (axis cs:592.5,0) rectangle (axis cs:593.5,100);
\draw[draw=none,fill=darkorange25512714] (axis cs:593.5,0) rectangle (axis cs:594.5,100);
\draw[draw=none,fill=darkorange25512714] (axis cs:594.5,0) rectangle (axis cs:595.5,100);
\draw[draw=none,fill=darkorange25512714] (axis cs:595.5,0) rectangle (axis cs:596.5,100);
\draw[draw=none,fill=darkorange25512714] (axis cs:596.5,0) rectangle (axis cs:597.5,100);
\draw[draw=none,fill=darkorange25512714] (axis cs:597.5,0) rectangle (axis cs:598.5,100);
\draw[draw=none,fill=darkorange25512714] (axis cs:598.5,0) rectangle (axis cs:599.5,100);
\draw[draw=none,fill=darkorange25512714] (axis cs:599.5,0) rectangle (axis cs:600.5,100);
\draw[draw=none,fill=darkorange25512714] (axis cs:600.5,0) rectangle (axis cs:601.5,100);
\draw[draw=none,fill=darkorange25512714] (axis cs:601.5,0) rectangle (axis cs:602.5,100);
\draw[draw=none,fill=darkorange25512714] (axis cs:602.5,0) rectangle (axis cs:603.5,100);
\draw[draw=none,fill=darkorange25512714] (axis cs:603.5,0) rectangle (axis cs:604.5,100);
\draw[draw=none,fill=darkorange25512714] (axis cs:604.5,0) rectangle (axis cs:605.5,100);
\draw[draw=none,fill=darkorange25512714] (axis cs:605.5,0) rectangle (axis cs:606.5,100);
\draw[draw=none,fill=darkorange25512714] (axis cs:606.5,0) rectangle (axis cs:607.5,100);
\draw[draw=none,fill=darkorange25512714] (axis cs:607.5,0) rectangle (axis cs:608.5,100);
\draw[draw=none,fill=darkorange25512714] (axis cs:608.5,0) rectangle (axis cs:609.5,100);
\draw[draw=none,fill=darkorange25512714] (axis cs:609.5,0) rectangle (axis cs:610.5,100);
\draw[draw=none,fill=darkorange25512714] (axis cs:610.5,0) rectangle (axis cs:611.5,100);
\draw[draw=none,fill=darkorange25512714] (axis cs:611.5,0) rectangle (axis cs:612.5,100);
\draw[draw=none,fill=darkorange25512714] (axis cs:612.5,0) rectangle (axis cs:613.5,100);
\draw[draw=none,fill=darkorange25512714] (axis cs:613.5,0) rectangle (axis cs:614.5,100);
\draw[draw=none,fill=darkorange25512714] (axis cs:614.5,0) rectangle (axis cs:615.5,100);
\draw[draw=none,fill=darkorange25512714] (axis cs:615.5,0) rectangle (axis cs:616.5,100);
\draw[draw=none,fill=darkorange25512714] (axis cs:616.5,0) rectangle (axis cs:617.5,100);
\draw[draw=none,fill=darkorange25512714] (axis cs:617.5,0) rectangle (axis cs:618.5,100);
\draw[draw=none,fill=darkorange25512714] (axis cs:618.5,0) rectangle (axis cs:619.5,100);
\draw[draw=none,fill=darkorange25512714] (axis cs:619.5,0) rectangle (axis cs:620.5,100);
\draw[draw=none,fill=darkorange25512714] (axis cs:620.5,0) rectangle (axis cs:621.5,100);
\draw[draw=none,fill=darkorange25512714] (axis cs:621.5,0) rectangle (axis cs:622.5,100);
\draw[draw=none,fill=darkorange25512714] (axis cs:622.5,0) rectangle (axis cs:623.5,100);
\draw[draw=none,fill=darkorange25512714] (axis cs:623.5,0) rectangle (axis cs:624.5,100);
\draw[draw=none,fill=darkorange25512714] (axis cs:624.5,0) rectangle (axis cs:625.5,100);
\draw[draw=none,fill=darkorange25512714] (axis cs:625.5,0) rectangle (axis cs:626.5,100);
\draw[draw=none,fill=darkorange25512714] (axis cs:626.5,0) rectangle (axis cs:627.5,100);
\draw[draw=none,fill=darkorange25512714] (axis cs:627.5,0) rectangle (axis cs:628.5,100);
\draw[draw=none,fill=darkorange25512714] (axis cs:628.5,0) rectangle (axis cs:629.5,100);
\draw[draw=none,fill=darkorange25512714] (axis cs:629.5,0) rectangle (axis cs:630.5,100);
\draw[draw=none,fill=darkorange25512714] (axis cs:630.5,0) rectangle (axis cs:631.5,100);
\draw[draw=none,fill=darkorange25512714] (axis cs:631.5,0) rectangle (axis cs:632.5,100);
\draw[draw=none,fill=darkorange25512714] (axis cs:632.5,0) rectangle (axis cs:633.5,100);
\draw[draw=none,fill=darkorange25512714] (axis cs:633.5,0) rectangle (axis cs:634.5,100);
\draw[draw=none,fill=darkorange25512714] (axis cs:634.5,0) rectangle (axis cs:635.5,100);
\draw[draw=none,fill=darkorange25512714] (axis cs:635.5,0) rectangle (axis cs:636.5,100);
\draw[draw=none,fill=darkorange25512714] (axis cs:636.5,0) rectangle (axis cs:637.5,100);
\draw[draw=none,fill=darkorange25512714] (axis cs:637.5,0) rectangle (axis cs:638.5,100);
\draw[draw=none,fill=darkorange25512714] (axis cs:638.5,0) rectangle (axis cs:639.5,100);
\draw[draw=none,fill=darkorange25512714] (axis cs:639.5,0) rectangle (axis cs:640.5,100);
\draw[draw=none,fill=darkorange25512714] (axis cs:640.5,0) rectangle (axis cs:641.5,100);
\draw[draw=none,fill=darkorange25512714] (axis cs:641.5,0) rectangle (axis cs:642.5,100);
\draw[draw=none,fill=darkorange25512714] (axis cs:642.5,0) rectangle (axis cs:643.5,100);
\draw[draw=none,fill=darkorange25512714] (axis cs:643.5,0) rectangle (axis cs:644.5,100);
\draw[draw=none,fill=darkorange25512714] (axis cs:644.5,0) rectangle (axis cs:645.5,100);
\draw[draw=none,fill=darkorange25512714] (axis cs:645.5,0) rectangle (axis cs:646.5,100);
\draw[draw=none,fill=darkorange25512714] (axis cs:646.5,0) rectangle (axis cs:647.5,100);
\draw[draw=none,fill=darkorange25512714] (axis cs:647.5,0) rectangle (axis cs:648.5,100);
\draw[draw=none,fill=darkorange25512714] (axis cs:648.5,0) rectangle (axis cs:649.5,100);
\draw[draw=none,fill=darkorange25512714] (axis cs:649.5,0) rectangle (axis cs:650.5,100);
\draw[draw=none,fill=darkorange25512714] (axis cs:650.5,0) rectangle (axis cs:651.5,100);
\draw[draw=none,fill=darkorange25512714] (axis cs:651.5,0) rectangle (axis cs:652.5,100);
\draw[draw=none,fill=darkorange25512714] (axis cs:652.5,0) rectangle (axis cs:653.5,100);
\draw[draw=none,fill=darkorange25512714] (axis cs:653.5,0) rectangle (axis cs:654.5,100);
\draw[draw=none,fill=darkorange25512714] (axis cs:654.5,0) rectangle (axis cs:655.5,100);
\draw[draw=none,fill=darkorange25512714] (axis cs:655.5,0) rectangle (axis cs:656.5,100);
\draw[draw=none,fill=darkorange25512714] (axis cs:656.5,0) rectangle (axis cs:657.5,100);
\draw[draw=none,fill=darkorange25512714] (axis cs:657.5,0) rectangle (axis cs:658.5,100);
\draw[draw=none,fill=darkorange25512714] (axis cs:658.5,0) rectangle (axis cs:659.5,100);
\draw[draw=none,fill=darkorange25512714] (axis cs:659.5,0) rectangle (axis cs:660.5,100);
\draw[draw=none,fill=darkorange25512714] (axis cs:660.5,0) rectangle (axis cs:661.5,100);
\draw[draw=none,fill=darkorange25512714] (axis cs:661.5,0) rectangle (axis cs:662.5,100);
\draw[draw=none,fill=darkorange25512714] (axis cs:662.5,0) rectangle (axis cs:663.5,100);
\draw[draw=none,fill=darkorange25512714] (axis cs:663.5,0) rectangle (axis cs:664.5,100);
\draw[draw=none,fill=darkorange25512714] (axis cs:664.5,0) rectangle (axis cs:665.5,100);
\draw[draw=none,fill=darkorange25512714] (axis cs:665.5,0) rectangle (axis cs:666.5,100);
\draw[draw=none,fill=darkorange25512714] (axis cs:666.5,0) rectangle (axis cs:667.5,100);
\draw[draw=none,fill=darkorange25512714] (axis cs:667.5,0) rectangle (axis cs:668.5,100);
\draw[draw=none,fill=darkorange25512714] (axis cs:668.5,0) rectangle (axis cs:669.5,100);
\draw[draw=none,fill=darkorange25512714] (axis cs:669.5,0) rectangle (axis cs:670.5,100);
\draw[draw=none,fill=darkorange25512714] (axis cs:670.5,0) rectangle (axis cs:671.5,100);
\draw[draw=none,fill=darkorange25512714] (axis cs:671.5,0) rectangle (axis cs:672.5,100);
\draw[draw=none,fill=darkorange25512714] (axis cs:672.5,0) rectangle (axis cs:673.5,100);
\draw[draw=none,fill=darkorange25512714] (axis cs:673.5,0) rectangle (axis cs:674.5,100);
\draw[draw=none,fill=darkorange25512714] (axis cs:674.5,0) rectangle (axis cs:675.5,100);
\draw[draw=none,fill=darkorange25512714] (axis cs:675.5,0) rectangle (axis cs:676.5,100);
\draw[draw=none,fill=darkorange25512714] (axis cs:676.5,0) rectangle (axis cs:677.5,100);
\draw[draw=none,fill=darkorange25512714] (axis cs:677.5,0) rectangle (axis cs:678.5,100);
\draw[draw=none,fill=darkorange25512714] (axis cs:678.5,0) rectangle (axis cs:679.5,100);
\draw[draw=none,fill=darkorange25512714] (axis cs:679.5,0) rectangle (axis cs:680.5,100);
\draw[draw=none,fill=darkorange25512714] (axis cs:680.5,0) rectangle (axis cs:681.5,100);
\draw[draw=none,fill=darkorange25512714] (axis cs:681.5,0) rectangle (axis cs:682.5,100);
\draw[draw=none,fill=darkorange25512714] (axis cs:682.5,0) rectangle (axis cs:683.5,100);
\draw[draw=none,fill=darkorange25512714] (axis cs:683.5,0) rectangle (axis cs:684.5,100);
\draw[draw=none,fill=darkorange25512714] (axis cs:684.5,0) rectangle (axis cs:685.5,100);
\draw[draw=none,fill=darkorange25512714] (axis cs:685.5,0) rectangle (axis cs:686.5,100);
\draw[draw=none,fill=darkorange25512714] (axis cs:686.5,0) rectangle (axis cs:687.5,100);
\draw[draw=none,fill=darkorange25512714] (axis cs:687.5,0) rectangle (axis cs:688.5,100);
\draw[draw=none,fill=darkorange25512714] (axis cs:688.5,0) rectangle (axis cs:689.5,100);
\draw[draw=none,fill=darkorange25512714] (axis cs:689.5,0) rectangle (axis cs:690.5,100);
\draw[draw=none,fill=darkorange25512714] (axis cs:690.5,0) rectangle (axis cs:691.5,100);
\draw[draw=none,fill=darkorange25512714] (axis cs:691.5,0) rectangle (axis cs:692.5,100);
\draw[draw=none,fill=darkorange25512714] (axis cs:692.5,0) rectangle (axis cs:693.5,100);
\draw[draw=none,fill=darkorange25512714] (axis cs:693.5,0) rectangle (axis cs:694.5,100);
\draw[draw=none,fill=darkorange25512714] (axis cs:694.5,0) rectangle (axis cs:695.5,100);
\draw[draw=none,fill=darkorange25512714] (axis cs:695.5,0) rectangle (axis cs:696.5,100);
\draw[draw=none,fill=darkorange25512714] (axis cs:696.5,0) rectangle (axis cs:697.5,100);
\draw[draw=none,fill=darkorange25512714] (axis cs:697.5,0) rectangle (axis cs:698.5,100);
\draw[draw=none,fill=darkorange25512714] (axis cs:698.5,0) rectangle (axis cs:699.5,100);
\draw[draw=none,fill=darkorange25512714] (axis cs:699.5,0) rectangle (axis cs:700.5,100);
\draw[draw=none,fill=darkorange25512714] (axis cs:700.5,0) rectangle (axis cs:701.5,100);
\draw[draw=none,fill=darkorange25512714] (axis cs:701.5,0) rectangle (axis cs:702.5,100);
\draw[draw=none,fill=darkorange25512714] (axis cs:702.5,0) rectangle (axis cs:703.5,100);
\draw[draw=none,fill=darkorange25512714] (axis cs:703.5,0) rectangle (axis cs:704.5,100);
\draw[draw=none,fill=darkorange25512714] (axis cs:704.5,0) rectangle (axis cs:705.5,100);
\draw[draw=none,fill=darkorange25512714] (axis cs:705.5,0) rectangle (axis cs:706.5,100);
\draw[draw=none,fill=darkorange25512714] (axis cs:706.5,0) rectangle (axis cs:707.5,100);
\draw[draw=none,fill=darkorange25512714] (axis cs:707.5,0) rectangle (axis cs:708.5,100);
\draw[draw=none,fill=darkorange25512714] (axis cs:708.5,0) rectangle (axis cs:709.5,100);
\draw[draw=none,fill=darkorange25512714] (axis cs:709.5,0) rectangle (axis cs:710.5,100);
\draw[draw=none,fill=darkorange25512714] (axis cs:710.5,0) rectangle (axis cs:711.5,100);
\draw[draw=none,fill=darkorange25512714] (axis cs:711.5,0) rectangle (axis cs:712.5,100);
\draw[draw=none,fill=darkorange25512714] (axis cs:712.5,0) rectangle (axis cs:713.5,100);
\draw[draw=none,fill=darkorange25512714] (axis cs:713.5,0) rectangle (axis cs:714.5,100);
\draw[draw=none,fill=darkorange25512714] (axis cs:714.5,0) rectangle (axis cs:715.5,100);
\draw[draw=none,fill=darkorange25512714] (axis cs:715.5,0) rectangle (axis cs:716.5,100);
\draw[draw=none,fill=darkorange25512714] (axis cs:716.5,0) rectangle (axis cs:717.5,100);
\draw[draw=none,fill=darkorange25512714] (axis cs:717.5,0) rectangle (axis cs:718.5,100);
\draw[draw=none,fill=darkorange25512714] (axis cs:718.5,0) rectangle (axis cs:719.5,100);
\draw[draw=none,fill=darkorange25512714] (axis cs:719.5,0) rectangle (axis cs:720.5,100);
\draw[draw=none,fill=darkorange25512714] (axis cs:720.5,0) rectangle (axis cs:721.5,100);
\draw[draw=none,fill=darkorange25512714] (axis cs:721.5,0) rectangle (axis cs:722.5,100);
\draw[draw=none,fill=darkorange25512714] (axis cs:722.5,0) rectangle (axis cs:723.5,100);
\draw[draw=none,fill=darkorange25512714] (axis cs:723.5,0) rectangle (axis cs:724.5,100);
\draw[draw=none,fill=darkorange25512714] (axis cs:724.5,0) rectangle (axis cs:725.5,100);
\draw[draw=none,fill=darkorange25512714] (axis cs:725.5,0) rectangle (axis cs:726.5,100);
\draw[draw=none,fill=darkorange25512714] (axis cs:726.5,0) rectangle (axis cs:727.5,100);
\draw[draw=none,fill=darkorange25512714] (axis cs:727.5,0) rectangle (axis cs:728.5,100);
\draw[draw=none,fill=darkorange25512714] (axis cs:728.5,0) rectangle (axis cs:729.5,100);
\draw[draw=none,fill=darkorange25512714] (axis cs:729.5,0) rectangle (axis cs:730.5,100);
\draw[draw=none,fill=darkorange25512714] (axis cs:730.5,0) rectangle (axis cs:731.5,100);
\draw[draw=none,fill=darkorange25512714] (axis cs:731.5,0) rectangle (axis cs:732.5,100);
\draw[draw=none,fill=darkorange25512714] (axis cs:732.5,0) rectangle (axis cs:733.5,100);
\draw[draw=none,fill=darkorange25512714] (axis cs:733.5,0) rectangle (axis cs:734.5,100);
\draw[draw=none,fill=darkorange25512714] (axis cs:734.5,0) rectangle (axis cs:735.5,100);
\draw[draw=none,fill=darkorange25512714] (axis cs:735.5,0) rectangle (axis cs:736.5,100);
\draw[draw=none,fill=darkorange25512714] (axis cs:736.5,0) rectangle (axis cs:737.5,100);
\draw[draw=none,fill=darkorange25512714] (axis cs:737.5,0) rectangle (axis cs:738.5,100);
\draw[draw=none,fill=darkorange25512714] (axis cs:738.5,0) rectangle (axis cs:739.5,100);
\draw[draw=none,fill=darkorange25512714] (axis cs:739.5,0) rectangle (axis cs:740.5,100);
\draw[draw=none,fill=darkorange25512714] (axis cs:740.5,0) rectangle (axis cs:741.5,100);
\draw[draw=none,fill=darkorange25512714] (axis cs:741.5,0) rectangle (axis cs:742.5,100);
\draw[draw=none,fill=darkorange25512714] (axis cs:742.5,0) rectangle (axis cs:743.5,100);
\draw[draw=none,fill=darkorange25512714] (axis cs:743.5,0) rectangle (axis cs:744.5,100);
\draw[draw=none,fill=darkorange25512714] (axis cs:744.5,0) rectangle (axis cs:745.5,100);
\draw[draw=none,fill=darkorange25512714] (axis cs:745.5,0) rectangle (axis cs:746.5,100);
\draw[draw=none,fill=darkorange25512714] (axis cs:746.5,0) rectangle (axis cs:747.5,100);
\draw[draw=none,fill=darkorange25512714] (axis cs:747.5,0) rectangle (axis cs:748.5,100);
\draw[draw=none,fill=darkorange25512714] (axis cs:748.5,0) rectangle (axis cs:749.5,100);
\draw[draw=none,fill=darkorange25512714] (axis cs:749.5,0) rectangle (axis cs:750.5,100);
\draw[draw=none,fill=darkorange25512714] (axis cs:750.5,0) rectangle (axis cs:751.5,100);
\draw[draw=none,fill=darkorange25512714] (axis cs:751.5,0) rectangle (axis cs:752.5,100);
\draw[draw=none,fill=darkorange25512714] (axis cs:752.5,0) rectangle (axis cs:753.5,100);
\draw[draw=none,fill=darkorange25512714] (axis cs:753.5,0) rectangle (axis cs:754.5,100);
\draw[draw=none,fill=darkorange25512714] (axis cs:754.5,0) rectangle (axis cs:755.5,100);
\draw[draw=none,fill=darkorange25512714] (axis cs:755.5,0) rectangle (axis cs:756.5,100);
\draw[draw=none,fill=darkorange25512714] (axis cs:756.5,0) rectangle (axis cs:757.5,100);
\draw[draw=none,fill=darkorange25512714] (axis cs:757.5,0) rectangle (axis cs:758.5,100);
\draw[draw=none,fill=darkorange25512714] (axis cs:758.5,0) rectangle (axis cs:759.5,100);
\draw[draw=none,fill=darkorange25512714] (axis cs:759.5,0) rectangle (axis cs:760.5,100);
\draw[draw=none,fill=darkorange25512714] (axis cs:760.5,0) rectangle (axis cs:761.5,100);
\draw[draw=none,fill=darkorange25512714] (axis cs:761.5,0) rectangle (axis cs:762.5,100);
\draw[draw=none,fill=darkorange25512714] (axis cs:762.5,0) rectangle (axis cs:763.5,100);
\draw[draw=none,fill=darkorange25512714] (axis cs:763.5,0) rectangle (axis cs:764.5,100);
\draw[draw=none,fill=darkorange25512714] (axis cs:764.5,0) rectangle (axis cs:765.5,100);
\draw[draw=none,fill=darkorange25512714] (axis cs:765.5,0) rectangle (axis cs:766.5,100);
\draw[draw=none,fill=darkorange25512714] (axis cs:766.5,0) rectangle (axis cs:767.5,100);
\draw[draw=none,fill=darkorange25512714] (axis cs:767.5,0) rectangle (axis cs:768.5,100);
\draw[draw=none,fill=darkorange25512714] (axis cs:768.5,0) rectangle (axis cs:769.5,100);
\draw[draw=none,fill=darkorange25512714] (axis cs:769.5,0) rectangle (axis cs:770.5,100);
\draw[draw=none,fill=darkorange25512714] (axis cs:770.5,0) rectangle (axis cs:771.5,100);
\draw[draw=none,fill=darkorange25512714] (axis cs:771.5,0) rectangle (axis cs:772.5,100);
\draw[draw=none,fill=darkorange25512714] (axis cs:772.5,0) rectangle (axis cs:773.5,100);
\draw[draw=none,fill=darkorange25512714] (axis cs:773.5,0) rectangle (axis cs:774.5,100);
\draw[draw=none,fill=darkorange25512714] (axis cs:774.5,0) rectangle (axis cs:775.5,100);
\draw[draw=none,fill=darkorange25512714] (axis cs:775.5,0) rectangle (axis cs:776.5,100);
\draw[draw=none,fill=darkorange25512714] (axis cs:776.5,0) rectangle (axis cs:777.5,100);
\draw[draw=none,fill=darkorange25512714] (axis cs:777.5,0) rectangle (axis cs:778.5,100);
\draw[draw=none,fill=darkorange25512714] (axis cs:778.5,0) rectangle (axis cs:779.5,100);
\draw[draw=none,fill=darkorange25512714] (axis cs:779.5,0) rectangle (axis cs:780.5,100);
\draw[draw=none,fill=darkorange25512714] (axis cs:780.5,0) rectangle (axis cs:781.5,100);
\draw[draw=none,fill=darkorange25512714] (axis cs:781.5,0) rectangle (axis cs:782.5,100);
\draw[draw=none,fill=darkorange25512714] (axis cs:782.5,0) rectangle (axis cs:783.5,100);
\draw[draw=none,fill=darkorange25512714] (axis cs:783.5,0) rectangle (axis cs:784.5,100);
\draw[draw=none,fill=darkorange25512714] (axis cs:784.5,0) rectangle (axis cs:785.5,100);
\draw[draw=none,fill=darkorange25512714] (axis cs:785.5,0) rectangle (axis cs:786.5,100);
\draw[draw=none,fill=forestgreen4416044] (axis cs:-0.5,0) rectangle (axis cs:0.5,0);
\addlegendimage{ybar,ybar legend,draw=none,fill=forestgreen4416044}
\addlegendentry{$\Delta_L=3$}

\draw[draw=none,fill=forestgreen4416044] (axis cs:0.5,0) rectangle (axis cs:1.5,0);
\draw[draw=none,fill=forestgreen4416044] (axis cs:1.5,0) rectangle (axis cs:2.5,0);
\draw[draw=none,fill=forestgreen4416044] (axis cs:2.5,0) rectangle (axis cs:3.5,0);
\draw[draw=none,fill=forestgreen4416044] (axis cs:3.5,0) rectangle (axis cs:4.5,0);
\draw[draw=none,fill=forestgreen4416044] (axis cs:4.5,0) rectangle (axis cs:5.5,0);
\draw[draw=none,fill=forestgreen4416044] (axis cs:5.5,0) rectangle (axis cs:6.5,0);
\draw[draw=none,fill=forestgreen4416044] (axis cs:6.5,0) rectangle (axis cs:7.5,0);
\draw[draw=none,fill=forestgreen4416044] (axis cs:7.5,0) rectangle (axis cs:8.5,0);
\draw[draw=none,fill=forestgreen4416044] (axis cs:8.5,0) rectangle (axis cs:9.5,0);
\draw[draw=none,fill=forestgreen4416044] (axis cs:9.5,0) rectangle (axis cs:10.5,0);
\draw[draw=none,fill=forestgreen4416044] (axis cs:10.5,0) rectangle (axis cs:11.5,0);
\draw[draw=none,fill=forestgreen4416044] (axis cs:11.5,0) rectangle (axis cs:12.5,0);
\draw[draw=none,fill=forestgreen4416044] (axis cs:12.5,0) rectangle (axis cs:13.5,2);
\draw[draw=none,fill=forestgreen4416044] (axis cs:13.5,0) rectangle (axis cs:14.5,2);
\draw[draw=none,fill=forestgreen4416044] (axis cs:14.5,0) rectangle (axis cs:15.5,4);
\draw[draw=none,fill=forestgreen4416044] (axis cs:15.5,0) rectangle (axis cs:16.5,4);
\draw[draw=none,fill=forestgreen4416044] (axis cs:16.5,0) rectangle (axis cs:17.5,6);
\draw[draw=none,fill=forestgreen4416044] (axis cs:17.5,0) rectangle (axis cs:18.5,6);
\draw[draw=none,fill=forestgreen4416044] (axis cs:18.5,0) rectangle (axis cs:19.5,6);
\draw[draw=none,fill=forestgreen4416044] (axis cs:19.5,0) rectangle (axis cs:20.5,6);
\draw[draw=none,fill=forestgreen4416044] (axis cs:20.5,0) rectangle (axis cs:21.5,8);
\draw[draw=none,fill=forestgreen4416044] (axis cs:21.5,0) rectangle (axis cs:22.5,10);
\draw[draw=none,fill=forestgreen4416044] (axis cs:22.5,0) rectangle (axis cs:23.5,10);
\draw[draw=none,fill=forestgreen4416044] (axis cs:23.5,0) rectangle (axis cs:24.5,12);
\draw[draw=none,fill=forestgreen4416044] (axis cs:24.5,0) rectangle (axis cs:25.5,12);
\draw[draw=none,fill=forestgreen4416044] (axis cs:25.5,0) rectangle (axis cs:26.5,12);
\draw[draw=none,fill=forestgreen4416044] (axis cs:26.5,0) rectangle (axis cs:27.5,12);
\draw[draw=none,fill=forestgreen4416044] (axis cs:27.5,0) rectangle (axis cs:28.5,14);
\draw[draw=none,fill=forestgreen4416044] (axis cs:28.5,0) rectangle (axis cs:29.5,14);
\draw[draw=none,fill=forestgreen4416044] (axis cs:29.5,0) rectangle (axis cs:30.5,14);
\draw[draw=none,fill=forestgreen4416044] (axis cs:30.5,0) rectangle (axis cs:31.5,16);
\draw[draw=none,fill=forestgreen4416044] (axis cs:31.5,0) rectangle (axis cs:32.5,16);
\draw[draw=none,fill=forestgreen4416044] (axis cs:32.5,0) rectangle (axis cs:33.5,16);
\draw[draw=none,fill=forestgreen4416044] (axis cs:33.5,0) rectangle (axis cs:34.5,16);
\draw[draw=none,fill=forestgreen4416044] (axis cs:34.5,0) rectangle (axis cs:35.5,18);
\draw[draw=none,fill=forestgreen4416044] (axis cs:35.5,0) rectangle (axis cs:36.5,18);
\draw[draw=none,fill=forestgreen4416044] (axis cs:36.5,0) rectangle (axis cs:37.5,18);
\draw[draw=none,fill=forestgreen4416044] (axis cs:37.5,0) rectangle (axis cs:38.5,18);
\draw[draw=none,fill=forestgreen4416044] (axis cs:38.5,0) rectangle (axis cs:39.5,18);
\draw[draw=none,fill=forestgreen4416044] (axis cs:39.5,0) rectangle (axis cs:40.5,18);
\draw[draw=none,fill=forestgreen4416044] (axis cs:40.5,0) rectangle (axis cs:41.5,18);
\draw[draw=none,fill=forestgreen4416044] (axis cs:41.5,0) rectangle (axis cs:42.5,18);
\draw[draw=none,fill=forestgreen4416044] (axis cs:42.5,0) rectangle (axis cs:43.5,18);
\draw[draw=none,fill=forestgreen4416044] (axis cs:43.5,0) rectangle (axis cs:44.5,20);
\draw[draw=none,fill=forestgreen4416044] (axis cs:44.5,0) rectangle (axis cs:45.5,20);
\draw[draw=none,fill=forestgreen4416044] (axis cs:45.5,0) rectangle (axis cs:46.5,20);
\draw[draw=none,fill=forestgreen4416044] (axis cs:46.5,0) rectangle (axis cs:47.5,22);
\draw[draw=none,fill=forestgreen4416044] (axis cs:47.5,0) rectangle (axis cs:48.5,22);
\draw[draw=none,fill=forestgreen4416044] (axis cs:48.5,0) rectangle (axis cs:49.5,22);
\draw[draw=none,fill=forestgreen4416044] (axis cs:49.5,0) rectangle (axis cs:50.5,22);
\draw[draw=none,fill=forestgreen4416044] (axis cs:50.5,0) rectangle (axis cs:51.5,22);
\draw[draw=none,fill=forestgreen4416044] (axis cs:51.5,0) rectangle (axis cs:52.5,24);
\draw[draw=none,fill=forestgreen4416044] (axis cs:52.5,0) rectangle (axis cs:53.5,26);
\draw[draw=none,fill=forestgreen4416044] (axis cs:53.5,0) rectangle (axis cs:54.5,26);
\draw[draw=none,fill=forestgreen4416044] (axis cs:54.5,0) rectangle (axis cs:55.5,26);
\draw[draw=none,fill=forestgreen4416044] (axis cs:55.5,0) rectangle (axis cs:56.5,26);
\draw[draw=none,fill=forestgreen4416044] (axis cs:56.5,0) rectangle (axis cs:57.5,26);
\draw[draw=none,fill=forestgreen4416044] (axis cs:57.5,0) rectangle (axis cs:58.5,26);
\draw[draw=none,fill=forestgreen4416044] (axis cs:58.5,0) rectangle (axis cs:59.5,26);
\draw[draw=none,fill=forestgreen4416044] (axis cs:59.5,0) rectangle (axis cs:60.5,26);
\draw[draw=none,fill=forestgreen4416044] (axis cs:60.5,0) rectangle (axis cs:61.5,26);
\draw[draw=none,fill=forestgreen4416044] (axis cs:61.5,0) rectangle (axis cs:62.5,26);
\draw[draw=none,fill=forestgreen4416044] (axis cs:62.5,0) rectangle (axis cs:63.5,26);
\draw[draw=none,fill=forestgreen4416044] (axis cs:63.5,0) rectangle (axis cs:64.5,26);
\draw[draw=none,fill=forestgreen4416044] (axis cs:64.5,0) rectangle (axis cs:65.5,26);
\draw[draw=none,fill=forestgreen4416044] (axis cs:65.5,0) rectangle (axis cs:66.5,26);
\draw[draw=none,fill=forestgreen4416044] (axis cs:66.5,0) rectangle (axis cs:67.5,26);
\draw[draw=none,fill=forestgreen4416044] (axis cs:67.5,0) rectangle (axis cs:68.5,26);
\draw[draw=none,fill=forestgreen4416044] (axis cs:68.5,0) rectangle (axis cs:69.5,26);
\draw[draw=none,fill=forestgreen4416044] (axis cs:69.5,0) rectangle (axis cs:70.5,26);
\draw[draw=none,fill=forestgreen4416044] (axis cs:70.5,0) rectangle (axis cs:71.5,26);
\draw[draw=none,fill=forestgreen4416044] (axis cs:71.5,0) rectangle (axis cs:72.5,26);
\draw[draw=none,fill=forestgreen4416044] (axis cs:72.5,0) rectangle (axis cs:73.5,26);
\draw[draw=none,fill=forestgreen4416044] (axis cs:73.5,0) rectangle (axis cs:74.5,26);
\draw[draw=none,fill=forestgreen4416044] (axis cs:74.5,0) rectangle (axis cs:75.5,26);
\draw[draw=none,fill=forestgreen4416044] (axis cs:75.5,0) rectangle (axis cs:76.5,28);
\draw[draw=none,fill=forestgreen4416044] (axis cs:76.5,0) rectangle (axis cs:77.5,28);
\draw[draw=none,fill=forestgreen4416044] (axis cs:77.5,0) rectangle (axis cs:78.5,28);
\draw[draw=none,fill=forestgreen4416044] (axis cs:78.5,0) rectangle (axis cs:79.5,30);
\draw[draw=none,fill=forestgreen4416044] (axis cs:79.5,0) rectangle (axis cs:80.5,30);
\draw[draw=none,fill=forestgreen4416044] (axis cs:80.5,0) rectangle (axis cs:81.5,30);
\draw[draw=none,fill=forestgreen4416044] (axis cs:81.5,0) rectangle (axis cs:82.5,30);
\draw[draw=none,fill=forestgreen4416044] (axis cs:82.5,0) rectangle (axis cs:83.5,30);
\draw[draw=none,fill=forestgreen4416044] (axis cs:83.5,0) rectangle (axis cs:84.5,30);
\draw[draw=none,fill=forestgreen4416044] (axis cs:84.5,0) rectangle (axis cs:85.5,30);
\draw[draw=none,fill=forestgreen4416044] (axis cs:85.5,0) rectangle (axis cs:86.5,32);
\draw[draw=none,fill=forestgreen4416044] (axis cs:86.5,0) rectangle (axis cs:87.5,32);
\draw[draw=none,fill=forestgreen4416044] (axis cs:87.5,0) rectangle (axis cs:88.5,32);
\draw[draw=none,fill=forestgreen4416044] (axis cs:88.5,0) rectangle (axis cs:89.5,32);
\draw[draw=none,fill=forestgreen4416044] (axis cs:89.5,0) rectangle (axis cs:90.5,32);
\draw[draw=none,fill=forestgreen4416044] (axis cs:90.5,0) rectangle (axis cs:91.5,32);
\draw[draw=none,fill=forestgreen4416044] (axis cs:91.5,0) rectangle (axis cs:92.5,36);
\draw[draw=none,fill=forestgreen4416044] (axis cs:92.5,0) rectangle (axis cs:93.5,36);
\draw[draw=none,fill=forestgreen4416044] (axis cs:93.5,0) rectangle (axis cs:94.5,36);
\draw[draw=none,fill=forestgreen4416044] (axis cs:94.5,0) rectangle (axis cs:95.5,36);
\draw[draw=none,fill=forestgreen4416044] (axis cs:95.5,0) rectangle (axis cs:96.5,36);
\draw[draw=none,fill=forestgreen4416044] (axis cs:96.5,0) rectangle (axis cs:97.5,36);
\draw[draw=none,fill=forestgreen4416044] (axis cs:97.5,0) rectangle (axis cs:98.5,36);
\draw[draw=none,fill=forestgreen4416044] (axis cs:98.5,0) rectangle (axis cs:99.5,36);
\draw[draw=none,fill=forestgreen4416044] (axis cs:99.5,0) rectangle (axis cs:100.5,36);
\draw[draw=none,fill=forestgreen4416044] (axis cs:100.5,0) rectangle (axis cs:101.5,36);
\draw[draw=none,fill=forestgreen4416044] (axis cs:101.5,0) rectangle (axis cs:102.5,36);
\draw[draw=none,fill=forestgreen4416044] (axis cs:102.5,0) rectangle (axis cs:103.5,36);
\draw[draw=none,fill=forestgreen4416044] (axis cs:103.5,0) rectangle (axis cs:104.5,36);
\draw[draw=none,fill=forestgreen4416044] (axis cs:104.5,0) rectangle (axis cs:105.5,36);
\draw[draw=none,fill=forestgreen4416044] (axis cs:105.5,0) rectangle (axis cs:106.5,36);
\draw[draw=none,fill=forestgreen4416044] (axis cs:106.5,0) rectangle (axis cs:107.5,36);
\draw[draw=none,fill=forestgreen4416044] (axis cs:107.5,0) rectangle (axis cs:108.5,36);
\draw[draw=none,fill=forestgreen4416044] (axis cs:108.5,0) rectangle (axis cs:109.5,36);
\draw[draw=none,fill=forestgreen4416044] (axis cs:109.5,0) rectangle (axis cs:110.5,36);
\draw[draw=none,fill=forestgreen4416044] (axis cs:110.5,0) rectangle (axis cs:111.5,36);
\draw[draw=none,fill=forestgreen4416044] (axis cs:111.5,0) rectangle (axis cs:112.5,36);
\draw[draw=none,fill=forestgreen4416044] (axis cs:112.5,0) rectangle (axis cs:113.5,36);
\draw[draw=none,fill=forestgreen4416044] (axis cs:113.5,0) rectangle (axis cs:114.5,36);
\draw[draw=none,fill=forestgreen4416044] (axis cs:114.5,0) rectangle (axis cs:115.5,36);
\draw[draw=none,fill=forestgreen4416044] (axis cs:115.5,0) rectangle (axis cs:116.5,36);
\draw[draw=none,fill=forestgreen4416044] (axis cs:116.5,0) rectangle (axis cs:117.5,38);
\draw[draw=none,fill=forestgreen4416044] (axis cs:117.5,0) rectangle (axis cs:118.5,38);
\draw[draw=none,fill=forestgreen4416044] (axis cs:118.5,0) rectangle (axis cs:119.5,38);
\draw[draw=none,fill=forestgreen4416044] (axis cs:119.5,0) rectangle (axis cs:120.5,38);
\draw[draw=none,fill=forestgreen4416044] (axis cs:120.5,0) rectangle (axis cs:121.5,38);
\draw[draw=none,fill=forestgreen4416044] (axis cs:121.5,0) rectangle (axis cs:122.5,38);
\draw[draw=none,fill=forestgreen4416044] (axis cs:122.5,0) rectangle (axis cs:123.5,38);
\draw[draw=none,fill=forestgreen4416044] (axis cs:123.5,0) rectangle (axis cs:124.5,38);
\draw[draw=none,fill=forestgreen4416044] (axis cs:124.5,0) rectangle (axis cs:125.5,38);
\draw[draw=none,fill=forestgreen4416044] (axis cs:125.5,0) rectangle (axis cs:126.5,38);
\draw[draw=none,fill=forestgreen4416044] (axis cs:126.5,0) rectangle (axis cs:127.5,38);
\draw[draw=none,fill=forestgreen4416044] (axis cs:127.5,0) rectangle (axis cs:128.5,38);
\draw[draw=none,fill=forestgreen4416044] (axis cs:128.5,0) rectangle (axis cs:129.5,38);
\draw[draw=none,fill=forestgreen4416044] (axis cs:129.5,0) rectangle (axis cs:130.5,38);
\draw[draw=none,fill=forestgreen4416044] (axis cs:130.5,0) rectangle (axis cs:131.5,40);
\draw[draw=none,fill=forestgreen4416044] (axis cs:131.5,0) rectangle (axis cs:132.5,40);
\draw[draw=none,fill=forestgreen4416044] (axis cs:132.5,0) rectangle (axis cs:133.5,40);
\draw[draw=none,fill=forestgreen4416044] (axis cs:133.5,0) rectangle (axis cs:134.5,40);
\draw[draw=none,fill=forestgreen4416044] (axis cs:134.5,0) rectangle (axis cs:135.5,40);
\draw[draw=none,fill=forestgreen4416044] (axis cs:135.5,0) rectangle (axis cs:136.5,40);
\draw[draw=none,fill=forestgreen4416044] (axis cs:136.5,0) rectangle (axis cs:137.5,40);
\draw[draw=none,fill=forestgreen4416044] (axis cs:137.5,0) rectangle (axis cs:138.5,40);
\draw[draw=none,fill=forestgreen4416044] (axis cs:138.5,0) rectangle (axis cs:139.5,42);
\draw[draw=none,fill=forestgreen4416044] (axis cs:139.5,0) rectangle (axis cs:140.5,42);
\draw[draw=none,fill=forestgreen4416044] (axis cs:140.5,0) rectangle (axis cs:141.5,42);
\draw[draw=none,fill=forestgreen4416044] (axis cs:141.5,0) rectangle (axis cs:142.5,42);
\draw[draw=none,fill=forestgreen4416044] (axis cs:142.5,0) rectangle (axis cs:143.5,42);
\draw[draw=none,fill=forestgreen4416044] (axis cs:143.5,0) rectangle (axis cs:144.5,42);
\draw[draw=none,fill=forestgreen4416044] (axis cs:144.5,0) rectangle (axis cs:145.5,42);
\draw[draw=none,fill=forestgreen4416044] (axis cs:145.5,0) rectangle (axis cs:146.5,44);
\draw[draw=none,fill=forestgreen4416044] (axis cs:146.5,0) rectangle (axis cs:147.5,44);
\draw[draw=none,fill=forestgreen4416044] (axis cs:147.5,0) rectangle (axis cs:148.5,44);
\draw[draw=none,fill=forestgreen4416044] (axis cs:148.5,0) rectangle (axis cs:149.5,44);
\draw[draw=none,fill=forestgreen4416044] (axis cs:149.5,0) rectangle (axis cs:150.5,44);
\draw[draw=none,fill=forestgreen4416044] (axis cs:150.5,0) rectangle (axis cs:151.5,44);
\draw[draw=none,fill=forestgreen4416044] (axis cs:151.5,0) rectangle (axis cs:152.5,44);
\draw[draw=none,fill=forestgreen4416044] (axis cs:152.5,0) rectangle (axis cs:153.5,44);
\draw[draw=none,fill=forestgreen4416044] (axis cs:153.5,0) rectangle (axis cs:154.5,44);
\draw[draw=none,fill=forestgreen4416044] (axis cs:154.5,0) rectangle (axis cs:155.5,44);
\draw[draw=none,fill=forestgreen4416044] (axis cs:155.5,0) rectangle (axis cs:156.5,44);
\draw[draw=none,fill=forestgreen4416044] (axis cs:156.5,0) rectangle (axis cs:157.5,44);
\draw[draw=none,fill=forestgreen4416044] (axis cs:157.5,0) rectangle (axis cs:158.5,44);
\draw[draw=none,fill=forestgreen4416044] (axis cs:158.5,0) rectangle (axis cs:159.5,44);
\draw[draw=none,fill=forestgreen4416044] (axis cs:159.5,0) rectangle (axis cs:160.5,44);
\draw[draw=none,fill=forestgreen4416044] (axis cs:160.5,0) rectangle (axis cs:161.5,44);
\draw[draw=none,fill=forestgreen4416044] (axis cs:161.5,0) rectangle (axis cs:162.5,44);
\draw[draw=none,fill=forestgreen4416044] (axis cs:162.5,0) rectangle (axis cs:163.5,44);
\draw[draw=none,fill=forestgreen4416044] (axis cs:163.5,0) rectangle (axis cs:164.5,44);
\draw[draw=none,fill=forestgreen4416044] (axis cs:164.5,0) rectangle (axis cs:165.5,44);
\draw[draw=none,fill=forestgreen4416044] (axis cs:165.5,0) rectangle (axis cs:166.5,44);
\draw[draw=none,fill=forestgreen4416044] (axis cs:166.5,0) rectangle (axis cs:167.5,44);
\draw[draw=none,fill=forestgreen4416044] (axis cs:167.5,0) rectangle (axis cs:168.5,44);
\draw[draw=none,fill=forestgreen4416044] (axis cs:168.5,0) rectangle (axis cs:169.5,44);
\draw[draw=none,fill=forestgreen4416044] (axis cs:169.5,0) rectangle (axis cs:170.5,44);
\draw[draw=none,fill=forestgreen4416044] (axis cs:170.5,0) rectangle (axis cs:171.5,46);
\draw[draw=none,fill=forestgreen4416044] (axis cs:171.5,0) rectangle (axis cs:172.5,46);
\draw[draw=none,fill=forestgreen4416044] (axis cs:172.5,0) rectangle (axis cs:173.5,46);
\draw[draw=none,fill=forestgreen4416044] (axis cs:173.5,0) rectangle (axis cs:174.5,46);
\draw[draw=none,fill=forestgreen4416044] (axis cs:174.5,0) rectangle (axis cs:175.5,46);
\draw[draw=none,fill=forestgreen4416044] (axis cs:175.5,0) rectangle (axis cs:176.5,46);
\draw[draw=none,fill=forestgreen4416044] (axis cs:176.5,0) rectangle (axis cs:177.5,46);
\draw[draw=none,fill=forestgreen4416044] (axis cs:177.5,0) rectangle (axis cs:178.5,46);
\draw[draw=none,fill=forestgreen4416044] (axis cs:178.5,0) rectangle (axis cs:179.5,46);
\draw[draw=none,fill=forestgreen4416044] (axis cs:179.5,0) rectangle (axis cs:180.5,46);
\draw[draw=none,fill=forestgreen4416044] (axis cs:180.5,0) rectangle (axis cs:181.5,46);
\draw[draw=none,fill=forestgreen4416044] (axis cs:181.5,0) rectangle (axis cs:182.5,46);
\draw[draw=none,fill=forestgreen4416044] (axis cs:182.5,0) rectangle (axis cs:183.5,46);
\draw[draw=none,fill=forestgreen4416044] (axis cs:183.5,0) rectangle (axis cs:184.5,46);
\draw[draw=none,fill=forestgreen4416044] (axis cs:184.5,0) rectangle (axis cs:185.5,46);
\draw[draw=none,fill=forestgreen4416044] (axis cs:185.5,0) rectangle (axis cs:186.5,46);
\draw[draw=none,fill=forestgreen4416044] (axis cs:186.5,0) rectangle (axis cs:187.5,46);
\draw[draw=none,fill=forestgreen4416044] (axis cs:187.5,0) rectangle (axis cs:188.5,46);
\draw[draw=none,fill=forestgreen4416044] (axis cs:188.5,0) rectangle (axis cs:189.5,46);
\draw[draw=none,fill=forestgreen4416044] (axis cs:189.5,0) rectangle (axis cs:190.5,46);
\draw[draw=none,fill=forestgreen4416044] (axis cs:190.5,0) rectangle (axis cs:191.5,46);
\draw[draw=none,fill=forestgreen4416044] (axis cs:191.5,0) rectangle (axis cs:192.5,48);
\draw[draw=none,fill=forestgreen4416044] (axis cs:192.5,0) rectangle (axis cs:193.5,48);
\draw[draw=none,fill=forestgreen4416044] (axis cs:193.5,0) rectangle (axis cs:194.5,48);
\draw[draw=none,fill=forestgreen4416044] (axis cs:194.5,0) rectangle (axis cs:195.5,48);
\draw[draw=none,fill=forestgreen4416044] (axis cs:195.5,0) rectangle (axis cs:196.5,48);
\draw[draw=none,fill=forestgreen4416044] (axis cs:196.5,0) rectangle (axis cs:197.5,48);
\draw[draw=none,fill=forestgreen4416044] (axis cs:197.5,0) rectangle (axis cs:198.5,50);
\draw[draw=none,fill=forestgreen4416044] (axis cs:198.5,0) rectangle (axis cs:199.5,50);
\draw[draw=none,fill=forestgreen4416044] (axis cs:199.5,0) rectangle (axis cs:200.5,50);
\draw[draw=none,fill=forestgreen4416044] (axis cs:200.5,0) rectangle (axis cs:201.5,50);
\draw[draw=none,fill=forestgreen4416044] (axis cs:201.5,0) rectangle (axis cs:202.5,50);
\draw[draw=none,fill=forestgreen4416044] (axis cs:202.5,0) rectangle (axis cs:203.5,50);
\draw[draw=none,fill=forestgreen4416044] (axis cs:203.5,0) rectangle (axis cs:204.5,50);
\draw[draw=none,fill=forestgreen4416044] (axis cs:204.5,0) rectangle (axis cs:205.5,50);
\draw[draw=none,fill=forestgreen4416044] (axis cs:205.5,0) rectangle (axis cs:206.5,50);
\draw[draw=none,fill=forestgreen4416044] (axis cs:206.5,0) rectangle (axis cs:207.5,50);
\draw[draw=none,fill=forestgreen4416044] (axis cs:207.5,0) rectangle (axis cs:208.5,50);
\draw[draw=none,fill=forestgreen4416044] (axis cs:208.5,0) rectangle (axis cs:209.5,50);
\draw[draw=none,fill=forestgreen4416044] (axis cs:209.5,0) rectangle (axis cs:210.5,50);
\draw[draw=none,fill=forestgreen4416044] (axis cs:210.5,0) rectangle (axis cs:211.5,50);
\draw[draw=none,fill=forestgreen4416044] (axis cs:211.5,0) rectangle (axis cs:212.5,52);
\draw[draw=none,fill=forestgreen4416044] (axis cs:212.5,0) rectangle (axis cs:213.5,52);
\draw[draw=none,fill=forestgreen4416044] (axis cs:213.5,0) rectangle (axis cs:214.5,52);
\draw[draw=none,fill=forestgreen4416044] (axis cs:214.5,0) rectangle (axis cs:215.5,56);
\draw[draw=none,fill=forestgreen4416044] (axis cs:215.5,0) rectangle (axis cs:216.5,56);
\draw[draw=none,fill=forestgreen4416044] (axis cs:216.5,0) rectangle (axis cs:217.5,56);
\draw[draw=none,fill=forestgreen4416044] (axis cs:217.5,0) rectangle (axis cs:218.5,56);
\draw[draw=none,fill=forestgreen4416044] (axis cs:218.5,0) rectangle (axis cs:219.5,56);
\draw[draw=none,fill=forestgreen4416044] (axis cs:219.5,0) rectangle (axis cs:220.5,56);
\draw[draw=none,fill=forestgreen4416044] (axis cs:220.5,0) rectangle (axis cs:221.5,56);
\draw[draw=none,fill=forestgreen4416044] (axis cs:221.5,0) rectangle (axis cs:222.5,56);
\draw[draw=none,fill=forestgreen4416044] (axis cs:222.5,0) rectangle (axis cs:223.5,56);
\draw[draw=none,fill=forestgreen4416044] (axis cs:223.5,0) rectangle (axis cs:224.5,56);
\draw[draw=none,fill=forestgreen4416044] (axis cs:224.5,0) rectangle (axis cs:225.5,58);
\draw[draw=none,fill=forestgreen4416044] (axis cs:225.5,0) rectangle (axis cs:226.5,58);
\draw[draw=none,fill=forestgreen4416044] (axis cs:226.5,0) rectangle (axis cs:227.5,58);
\draw[draw=none,fill=forestgreen4416044] (axis cs:227.5,0) rectangle (axis cs:228.5,58);
\draw[draw=none,fill=forestgreen4416044] (axis cs:228.5,0) rectangle (axis cs:229.5,58);
\draw[draw=none,fill=forestgreen4416044] (axis cs:229.5,0) rectangle (axis cs:230.5,58);
\draw[draw=none,fill=forestgreen4416044] (axis cs:230.5,0) rectangle (axis cs:231.5,58);
\draw[draw=none,fill=forestgreen4416044] (axis cs:231.5,0) rectangle (axis cs:232.5,58);
\draw[draw=none,fill=forestgreen4416044] (axis cs:232.5,0) rectangle (axis cs:233.5,58);
\draw[draw=none,fill=forestgreen4416044] (axis cs:233.5,0) rectangle (axis cs:234.5,58);
\draw[draw=none,fill=forestgreen4416044] (axis cs:234.5,0) rectangle (axis cs:235.5,58);
\draw[draw=none,fill=forestgreen4416044] (axis cs:235.5,0) rectangle (axis cs:236.5,58);
\draw[draw=none,fill=forestgreen4416044] (axis cs:236.5,0) rectangle (axis cs:237.5,60);
\draw[draw=none,fill=forestgreen4416044] (axis cs:237.5,0) rectangle (axis cs:238.5,60);
\draw[draw=none,fill=forestgreen4416044] (axis cs:238.5,0) rectangle (axis cs:239.5,60);
\draw[draw=none,fill=forestgreen4416044] (axis cs:239.5,0) rectangle (axis cs:240.5,60);
\draw[draw=none,fill=forestgreen4416044] (axis cs:240.5,0) rectangle (axis cs:241.5,60);
\draw[draw=none,fill=forestgreen4416044] (axis cs:241.5,0) rectangle (axis cs:242.5,60);
\draw[draw=none,fill=forestgreen4416044] (axis cs:242.5,0) rectangle (axis cs:243.5,60);
\draw[draw=none,fill=forestgreen4416044] (axis cs:243.5,0) rectangle (axis cs:244.5,60);
\draw[draw=none,fill=forestgreen4416044] (axis cs:244.5,0) rectangle (axis cs:245.5,60);
\draw[draw=none,fill=forestgreen4416044] (axis cs:245.5,0) rectangle (axis cs:246.5,60);
\draw[draw=none,fill=forestgreen4416044] (axis cs:246.5,0) rectangle (axis cs:247.5,60);
\draw[draw=none,fill=forestgreen4416044] (axis cs:247.5,0) rectangle (axis cs:248.5,64);
\draw[draw=none,fill=forestgreen4416044] (axis cs:248.5,0) rectangle (axis cs:249.5,64);
\draw[draw=none,fill=forestgreen4416044] (axis cs:249.5,0) rectangle (axis cs:250.5,64);
\draw[draw=none,fill=forestgreen4416044] (axis cs:250.5,0) rectangle (axis cs:251.5,64);
\draw[draw=none,fill=forestgreen4416044] (axis cs:251.5,0) rectangle (axis cs:252.5,64);
\draw[draw=none,fill=forestgreen4416044] (axis cs:252.5,0) rectangle (axis cs:253.5,64);
\draw[draw=none,fill=forestgreen4416044] (axis cs:253.5,0) rectangle (axis cs:254.5,64);
\draw[draw=none,fill=forestgreen4416044] (axis cs:254.5,0) rectangle (axis cs:255.5,64);
\draw[draw=none,fill=forestgreen4416044] (axis cs:255.5,0) rectangle (axis cs:256.5,64);
\draw[draw=none,fill=forestgreen4416044] (axis cs:256.5,0) rectangle (axis cs:257.5,64);
\draw[draw=none,fill=forestgreen4416044] (axis cs:257.5,0) rectangle (axis cs:258.5,64);
\draw[draw=none,fill=forestgreen4416044] (axis cs:258.5,0) rectangle (axis cs:259.5,64);
\draw[draw=none,fill=forestgreen4416044] (axis cs:259.5,0) rectangle (axis cs:260.5,64);
\draw[draw=none,fill=forestgreen4416044] (axis cs:260.5,0) rectangle (axis cs:261.5,64);
\draw[draw=none,fill=forestgreen4416044] (axis cs:261.5,0) rectangle (axis cs:262.5,64);
\draw[draw=none,fill=forestgreen4416044] (axis cs:262.5,0) rectangle (axis cs:263.5,64);
\draw[draw=none,fill=forestgreen4416044] (axis cs:263.5,0) rectangle (axis cs:264.5,64);
\draw[draw=none,fill=forestgreen4416044] (axis cs:264.5,0) rectangle (axis cs:265.5,64);
\draw[draw=none,fill=forestgreen4416044] (axis cs:265.5,0) rectangle (axis cs:266.5,64);
\draw[draw=none,fill=forestgreen4416044] (axis cs:266.5,0) rectangle (axis cs:267.5,64);
\draw[draw=none,fill=forestgreen4416044] (axis cs:267.5,0) rectangle (axis cs:268.5,64);
\draw[draw=none,fill=forestgreen4416044] (axis cs:268.5,0) rectangle (axis cs:269.5,64);
\draw[draw=none,fill=forestgreen4416044] (axis cs:269.5,0) rectangle (axis cs:270.5,64);
\draw[draw=none,fill=forestgreen4416044] (axis cs:270.5,0) rectangle (axis cs:271.5,64);
\draw[draw=none,fill=forestgreen4416044] (axis cs:271.5,0) rectangle (axis cs:272.5,64);
\draw[draw=none,fill=forestgreen4416044] (axis cs:272.5,0) rectangle (axis cs:273.5,64);
\draw[draw=none,fill=forestgreen4416044] (axis cs:273.5,0) rectangle (axis cs:274.5,64);
\draw[draw=none,fill=forestgreen4416044] (axis cs:274.5,0) rectangle (axis cs:275.5,64);
\draw[draw=none,fill=forestgreen4416044] (axis cs:275.5,0) rectangle (axis cs:276.5,64);
\draw[draw=none,fill=forestgreen4416044] (axis cs:276.5,0) rectangle (axis cs:277.5,64);
\draw[draw=none,fill=forestgreen4416044] (axis cs:277.5,0) rectangle (axis cs:278.5,64);
\draw[draw=none,fill=forestgreen4416044] (axis cs:278.5,0) rectangle (axis cs:279.5,64);
\draw[draw=none,fill=forestgreen4416044] (axis cs:279.5,0) rectangle (axis cs:280.5,64);
\draw[draw=none,fill=forestgreen4416044] (axis cs:280.5,0) rectangle (axis cs:281.5,64);
\draw[draw=none,fill=forestgreen4416044] (axis cs:281.5,0) rectangle (axis cs:282.5,64);
\draw[draw=none,fill=forestgreen4416044] (axis cs:282.5,0) rectangle (axis cs:283.5,64);
\draw[draw=none,fill=forestgreen4416044] (axis cs:283.5,0) rectangle (axis cs:284.5,64);
\draw[draw=none,fill=forestgreen4416044] (axis cs:284.5,0) rectangle (axis cs:285.5,64);
\draw[draw=none,fill=forestgreen4416044] (axis cs:285.5,0) rectangle (axis cs:286.5,64);
\draw[draw=none,fill=forestgreen4416044] (axis cs:286.5,0) rectangle (axis cs:287.5,64);
\draw[draw=none,fill=forestgreen4416044] (axis cs:287.5,0) rectangle (axis cs:288.5,64);
\draw[draw=none,fill=forestgreen4416044] (axis cs:288.5,0) rectangle (axis cs:289.5,68);
\draw[draw=none,fill=forestgreen4416044] (axis cs:289.5,0) rectangle (axis cs:290.5,68);
\draw[draw=none,fill=forestgreen4416044] (axis cs:290.5,0) rectangle (axis cs:291.5,70);
\draw[draw=none,fill=forestgreen4416044] (axis cs:291.5,0) rectangle (axis cs:292.5,70);
\draw[draw=none,fill=forestgreen4416044] (axis cs:292.5,0) rectangle (axis cs:293.5,70);
\draw[draw=none,fill=forestgreen4416044] (axis cs:293.5,0) rectangle (axis cs:294.5,70);
\draw[draw=none,fill=forestgreen4416044] (axis cs:294.5,0) rectangle (axis cs:295.5,70);
\draw[draw=none,fill=forestgreen4416044] (axis cs:295.5,0) rectangle (axis cs:296.5,70);
\draw[draw=none,fill=forestgreen4416044] (axis cs:296.5,0) rectangle (axis cs:297.5,70);
\draw[draw=none,fill=forestgreen4416044] (axis cs:297.5,0) rectangle (axis cs:298.5,70);
\draw[draw=none,fill=forestgreen4416044] (axis cs:298.5,0) rectangle (axis cs:299.5,70);
\draw[draw=none,fill=forestgreen4416044] (axis cs:299.5,0) rectangle (axis cs:300.5,70);
\draw[draw=none,fill=forestgreen4416044] (axis cs:300.5,0) rectangle (axis cs:301.5,70);
\draw[draw=none,fill=forestgreen4416044] (axis cs:301.5,0) rectangle (axis cs:302.5,70);
\draw[draw=none,fill=forestgreen4416044] (axis cs:302.5,0) rectangle (axis cs:303.5,70);
\draw[draw=none,fill=forestgreen4416044] (axis cs:303.5,0) rectangle (axis cs:304.5,70);
\draw[draw=none,fill=forestgreen4416044] (axis cs:304.5,0) rectangle (axis cs:305.5,72);
\draw[draw=none,fill=forestgreen4416044] (axis cs:305.5,0) rectangle (axis cs:306.5,72);
\draw[draw=none,fill=forestgreen4416044] (axis cs:306.5,0) rectangle (axis cs:307.5,72);
\draw[draw=none,fill=forestgreen4416044] (axis cs:307.5,0) rectangle (axis cs:308.5,72);
\draw[draw=none,fill=forestgreen4416044] (axis cs:308.5,0) rectangle (axis cs:309.5,74);
\draw[draw=none,fill=forestgreen4416044] (axis cs:309.5,0) rectangle (axis cs:310.5,74);
\draw[draw=none,fill=forestgreen4416044] (axis cs:310.5,0) rectangle (axis cs:311.5,74);
\draw[draw=none,fill=forestgreen4416044] (axis cs:311.5,0) rectangle (axis cs:312.5,76);
\draw[draw=none,fill=forestgreen4416044] (axis cs:312.5,0) rectangle (axis cs:313.5,76);
\draw[draw=none,fill=forestgreen4416044] (axis cs:313.5,0) rectangle (axis cs:314.5,76);
\draw[draw=none,fill=forestgreen4416044] (axis cs:314.5,0) rectangle (axis cs:315.5,76);
\draw[draw=none,fill=forestgreen4416044] (axis cs:315.5,0) rectangle (axis cs:316.5,76);
\draw[draw=none,fill=forestgreen4416044] (axis cs:316.5,0) rectangle (axis cs:317.5,76);
\draw[draw=none,fill=forestgreen4416044] (axis cs:317.5,0) rectangle (axis cs:318.5,76);
\draw[draw=none,fill=forestgreen4416044] (axis cs:318.5,0) rectangle (axis cs:319.5,76);
\draw[draw=none,fill=forestgreen4416044] (axis cs:319.5,0) rectangle (axis cs:320.5,76);
\draw[draw=none,fill=forestgreen4416044] (axis cs:320.5,0) rectangle (axis cs:321.5,76);
\draw[draw=none,fill=forestgreen4416044] (axis cs:321.5,0) rectangle (axis cs:322.5,76);
\draw[draw=none,fill=forestgreen4416044] (axis cs:322.5,0) rectangle (axis cs:323.5,76);
\draw[draw=none,fill=forestgreen4416044] (axis cs:323.5,0) rectangle (axis cs:324.5,76);
\draw[draw=none,fill=forestgreen4416044] (axis cs:324.5,0) rectangle (axis cs:325.5,76);
\draw[draw=none,fill=forestgreen4416044] (axis cs:325.5,0) rectangle (axis cs:326.5,76);
\draw[draw=none,fill=forestgreen4416044] (axis cs:326.5,0) rectangle (axis cs:327.5,76);
\draw[draw=none,fill=forestgreen4416044] (axis cs:327.5,0) rectangle (axis cs:328.5,76);
\draw[draw=none,fill=forestgreen4416044] (axis cs:328.5,0) rectangle (axis cs:329.5,76);
\draw[draw=none,fill=forestgreen4416044] (axis cs:329.5,0) rectangle (axis cs:330.5,76);
\draw[draw=none,fill=forestgreen4416044] (axis cs:330.5,0) rectangle (axis cs:331.5,76);
\draw[draw=none,fill=forestgreen4416044] (axis cs:331.5,0) rectangle (axis cs:332.5,76);
\draw[draw=none,fill=forestgreen4416044] (axis cs:332.5,0) rectangle (axis cs:333.5,76);
\draw[draw=none,fill=forestgreen4416044] (axis cs:333.5,0) rectangle (axis cs:334.5,76);
\draw[draw=none,fill=forestgreen4416044] (axis cs:334.5,0) rectangle (axis cs:335.5,76);
\draw[draw=none,fill=forestgreen4416044] (axis cs:335.5,0) rectangle (axis cs:336.5,76);
\draw[draw=none,fill=forestgreen4416044] (axis cs:336.5,0) rectangle (axis cs:337.5,76);
\draw[draw=none,fill=forestgreen4416044] (axis cs:337.5,0) rectangle (axis cs:338.5,76);
\draw[draw=none,fill=forestgreen4416044] (axis cs:338.5,0) rectangle (axis cs:339.5,76);
\draw[draw=none,fill=forestgreen4416044] (axis cs:339.5,0) rectangle (axis cs:340.5,76);
\draw[draw=none,fill=forestgreen4416044] (axis cs:340.5,0) rectangle (axis cs:341.5,76);
\draw[draw=none,fill=forestgreen4416044] (axis cs:341.5,0) rectangle (axis cs:342.5,76);
\draw[draw=none,fill=forestgreen4416044] (axis cs:342.5,0) rectangle (axis cs:343.5,76);
\draw[draw=none,fill=forestgreen4416044] (axis cs:343.5,0) rectangle (axis cs:344.5,76);
\draw[draw=none,fill=forestgreen4416044] (axis cs:344.5,0) rectangle (axis cs:345.5,76);
\draw[draw=none,fill=forestgreen4416044] (axis cs:345.5,0) rectangle (axis cs:346.5,76);
\draw[draw=none,fill=forestgreen4416044] (axis cs:346.5,0) rectangle (axis cs:347.5,76);
\draw[draw=none,fill=forestgreen4416044] (axis cs:347.5,0) rectangle (axis cs:348.5,76);
\draw[draw=none,fill=forestgreen4416044] (axis cs:348.5,0) rectangle (axis cs:349.5,76);
\draw[draw=none,fill=forestgreen4416044] (axis cs:349.5,0) rectangle (axis cs:350.5,76);
\draw[draw=none,fill=forestgreen4416044] (axis cs:350.5,0) rectangle (axis cs:351.5,76);
\draw[draw=none,fill=forestgreen4416044] (axis cs:351.5,0) rectangle (axis cs:352.5,76);
\draw[draw=none,fill=forestgreen4416044] (axis cs:352.5,0) rectangle (axis cs:353.5,76);
\draw[draw=none,fill=forestgreen4416044] (axis cs:353.5,0) rectangle (axis cs:354.5,76);
\draw[draw=none,fill=forestgreen4416044] (axis cs:354.5,0) rectangle (axis cs:355.5,76);
\draw[draw=none,fill=forestgreen4416044] (axis cs:355.5,0) rectangle (axis cs:356.5,76);
\draw[draw=none,fill=forestgreen4416044] (axis cs:356.5,0) rectangle (axis cs:357.5,76);
\draw[draw=none,fill=forestgreen4416044] (axis cs:357.5,0) rectangle (axis cs:358.5,76);
\draw[draw=none,fill=forestgreen4416044] (axis cs:358.5,0) rectangle (axis cs:359.5,76);
\draw[draw=none,fill=forestgreen4416044] (axis cs:359.5,0) rectangle (axis cs:360.5,76);
\draw[draw=none,fill=forestgreen4416044] (axis cs:360.5,0) rectangle (axis cs:361.5,76);
\draw[draw=none,fill=forestgreen4416044] (axis cs:361.5,0) rectangle (axis cs:362.5,76);
\draw[draw=none,fill=forestgreen4416044] (axis cs:362.5,0) rectangle (axis cs:363.5,76);
\draw[draw=none,fill=forestgreen4416044] (axis cs:363.5,0) rectangle (axis cs:364.5,76);
\draw[draw=none,fill=forestgreen4416044] (axis cs:364.5,0) rectangle (axis cs:365.5,76);
\draw[draw=none,fill=forestgreen4416044] (axis cs:365.5,0) rectangle (axis cs:366.5,76);
\draw[draw=none,fill=forestgreen4416044] (axis cs:366.5,0) rectangle (axis cs:367.5,78);
\draw[draw=none,fill=forestgreen4416044] (axis cs:367.5,0) rectangle (axis cs:368.5,78);
\draw[draw=none,fill=forestgreen4416044] (axis cs:368.5,0) rectangle (axis cs:369.5,78);
\draw[draw=none,fill=forestgreen4416044] (axis cs:369.5,0) rectangle (axis cs:370.5,78);
\draw[draw=none,fill=forestgreen4416044] (axis cs:370.5,0) rectangle (axis cs:371.5,78);
\draw[draw=none,fill=forestgreen4416044] (axis cs:371.5,0) rectangle (axis cs:372.5,78);
\draw[draw=none,fill=forestgreen4416044] (axis cs:372.5,0) rectangle (axis cs:373.5,78);
\draw[draw=none,fill=forestgreen4416044] (axis cs:373.5,0) rectangle (axis cs:374.5,78);
\draw[draw=none,fill=forestgreen4416044] (axis cs:374.5,0) rectangle (axis cs:375.5,78);
\draw[draw=none,fill=forestgreen4416044] (axis cs:375.5,0) rectangle (axis cs:376.5,78);
\draw[draw=none,fill=forestgreen4416044] (axis cs:376.5,0) rectangle (axis cs:377.5,78);
\draw[draw=none,fill=forestgreen4416044] (axis cs:377.5,0) rectangle (axis cs:378.5,78);
\draw[draw=none,fill=forestgreen4416044] (axis cs:378.5,0) rectangle (axis cs:379.5,78);
\draw[draw=none,fill=forestgreen4416044] (axis cs:379.5,0) rectangle (axis cs:380.5,78);
\draw[draw=none,fill=forestgreen4416044] (axis cs:380.5,0) rectangle (axis cs:381.5,78);
\draw[draw=none,fill=forestgreen4416044] (axis cs:381.5,0) rectangle (axis cs:382.5,78);
\draw[draw=none,fill=forestgreen4416044] (axis cs:382.5,0) rectangle (axis cs:383.5,78);
\draw[draw=none,fill=forestgreen4416044] (axis cs:383.5,0) rectangle (axis cs:384.5,78);
\draw[draw=none,fill=forestgreen4416044] (axis cs:384.5,0) rectangle (axis cs:385.5,78);
\draw[draw=none,fill=forestgreen4416044] (axis cs:385.5,0) rectangle (axis cs:386.5,78);
\draw[draw=none,fill=forestgreen4416044] (axis cs:386.5,0) rectangle (axis cs:387.5,78);
\draw[draw=none,fill=forestgreen4416044] (axis cs:387.5,0) rectangle (axis cs:388.5,78);
\draw[draw=none,fill=forestgreen4416044] (axis cs:388.5,0) rectangle (axis cs:389.5,78);
\draw[draw=none,fill=forestgreen4416044] (axis cs:389.5,0) rectangle (axis cs:390.5,78);
\draw[draw=none,fill=forestgreen4416044] (axis cs:390.5,0) rectangle (axis cs:391.5,78);
\draw[draw=none,fill=forestgreen4416044] (axis cs:391.5,0) rectangle (axis cs:392.5,78);
\draw[draw=none,fill=forestgreen4416044] (axis cs:392.5,0) rectangle (axis cs:393.5,78);
\draw[draw=none,fill=forestgreen4416044] (axis cs:393.5,0) rectangle (axis cs:394.5,78);
\draw[draw=none,fill=forestgreen4416044] (axis cs:394.5,0) rectangle (axis cs:395.5,80);
\draw[draw=none,fill=forestgreen4416044] (axis cs:395.5,0) rectangle (axis cs:396.5,80);
\draw[draw=none,fill=forestgreen4416044] (axis cs:396.5,0) rectangle (axis cs:397.5,80);
\draw[draw=none,fill=forestgreen4416044] (axis cs:397.5,0) rectangle (axis cs:398.5,80);
\draw[draw=none,fill=forestgreen4416044] (axis cs:398.5,0) rectangle (axis cs:399.5,80);
\draw[draw=none,fill=forestgreen4416044] (axis cs:399.5,0) rectangle (axis cs:400.5,80);
\draw[draw=none,fill=forestgreen4416044] (axis cs:400.5,0) rectangle (axis cs:401.5,80);
\draw[draw=none,fill=forestgreen4416044] (axis cs:401.5,0) rectangle (axis cs:402.5,80);
\draw[draw=none,fill=forestgreen4416044] (axis cs:402.5,0) rectangle (axis cs:403.5,80);
\draw[draw=none,fill=forestgreen4416044] (axis cs:403.5,0) rectangle (axis cs:404.5,80);
\draw[draw=none,fill=forestgreen4416044] (axis cs:404.5,0) rectangle (axis cs:405.5,80);
\draw[draw=none,fill=forestgreen4416044] (axis cs:405.5,0) rectangle (axis cs:406.5,80);
\draw[draw=none,fill=forestgreen4416044] (axis cs:406.5,0) rectangle (axis cs:407.5,80);
\draw[draw=none,fill=forestgreen4416044] (axis cs:407.5,0) rectangle (axis cs:408.5,80);
\draw[draw=none,fill=forestgreen4416044] (axis cs:408.5,0) rectangle (axis cs:409.5,80);
\draw[draw=none,fill=forestgreen4416044] (axis cs:409.5,0) rectangle (axis cs:410.5,80);
\draw[draw=none,fill=forestgreen4416044] (axis cs:410.5,0) rectangle (axis cs:411.5,80);
\draw[draw=none,fill=forestgreen4416044] (axis cs:411.5,0) rectangle (axis cs:412.5,80);
\draw[draw=none,fill=forestgreen4416044] (axis cs:412.5,0) rectangle (axis cs:413.5,80);
\draw[draw=none,fill=forestgreen4416044] (axis cs:413.5,0) rectangle (axis cs:414.5,80);
\draw[draw=none,fill=forestgreen4416044] (axis cs:414.5,0) rectangle (axis cs:415.5,80);
\draw[draw=none,fill=forestgreen4416044] (axis cs:415.5,0) rectangle (axis cs:416.5,82);
\draw[draw=none,fill=forestgreen4416044] (axis cs:416.5,0) rectangle (axis cs:417.5,84);
\draw[draw=none,fill=forestgreen4416044] (axis cs:417.5,0) rectangle (axis cs:418.5,84);
\draw[draw=none,fill=forestgreen4416044] (axis cs:418.5,0) rectangle (axis cs:419.5,84);
\draw[draw=none,fill=forestgreen4416044] (axis cs:419.5,0) rectangle (axis cs:420.5,84);
\draw[draw=none,fill=forestgreen4416044] (axis cs:420.5,0) rectangle (axis cs:421.5,84);
\draw[draw=none,fill=forestgreen4416044] (axis cs:421.5,0) rectangle (axis cs:422.5,84);
\draw[draw=none,fill=forestgreen4416044] (axis cs:422.5,0) rectangle (axis cs:423.5,84);
\draw[draw=none,fill=forestgreen4416044] (axis cs:423.5,0) rectangle (axis cs:424.5,84);
\draw[draw=none,fill=forestgreen4416044] (axis cs:424.5,0) rectangle (axis cs:425.5,84);
\draw[draw=none,fill=forestgreen4416044] (axis cs:425.5,0) rectangle (axis cs:426.5,84);
\draw[draw=none,fill=forestgreen4416044] (axis cs:426.5,0) rectangle (axis cs:427.5,84);
\draw[draw=none,fill=forestgreen4416044] (axis cs:427.5,0) rectangle (axis cs:428.5,84);
\draw[draw=none,fill=forestgreen4416044] (axis cs:428.5,0) rectangle (axis cs:429.5,84);
\draw[draw=none,fill=forestgreen4416044] (axis cs:429.5,0) rectangle (axis cs:430.5,84);
\draw[draw=none,fill=forestgreen4416044] (axis cs:430.5,0) rectangle (axis cs:431.5,84);
\draw[draw=none,fill=forestgreen4416044] (axis cs:431.5,0) rectangle (axis cs:432.5,84);
\draw[draw=none,fill=forestgreen4416044] (axis cs:432.5,0) rectangle (axis cs:433.5,84);
\draw[draw=none,fill=forestgreen4416044] (axis cs:433.5,0) rectangle (axis cs:434.5,84);
\draw[draw=none,fill=forestgreen4416044] (axis cs:434.5,0) rectangle (axis cs:435.5,84);
\draw[draw=none,fill=forestgreen4416044] (axis cs:435.5,0) rectangle (axis cs:436.5,84);
\draw[draw=none,fill=forestgreen4416044] (axis cs:436.5,0) rectangle (axis cs:437.5,84);
\draw[draw=none,fill=forestgreen4416044] (axis cs:437.5,0) rectangle (axis cs:438.5,84);
\draw[draw=none,fill=forestgreen4416044] (axis cs:438.5,0) rectangle (axis cs:439.5,84);
\draw[draw=none,fill=forestgreen4416044] (axis cs:439.5,0) rectangle (axis cs:440.5,84);
\draw[draw=none,fill=forestgreen4416044] (axis cs:440.5,0) rectangle (axis cs:441.5,86);
\draw[draw=none,fill=forestgreen4416044] (axis cs:441.5,0) rectangle (axis cs:442.5,86);
\draw[draw=none,fill=forestgreen4416044] (axis cs:442.5,0) rectangle (axis cs:443.5,86);
\draw[draw=none,fill=forestgreen4416044] (axis cs:443.5,0) rectangle (axis cs:444.5,86);
\draw[draw=none,fill=forestgreen4416044] (axis cs:444.5,0) rectangle (axis cs:445.5,86);
\draw[draw=none,fill=forestgreen4416044] (axis cs:445.5,0) rectangle (axis cs:446.5,86);
\draw[draw=none,fill=forestgreen4416044] (axis cs:446.5,0) rectangle (axis cs:447.5,86);
\draw[draw=none,fill=forestgreen4416044] (axis cs:447.5,0) rectangle (axis cs:448.5,86);
\draw[draw=none,fill=forestgreen4416044] (axis cs:448.5,0) rectangle (axis cs:449.5,86);
\draw[draw=none,fill=forestgreen4416044] (axis cs:449.5,0) rectangle (axis cs:450.5,86);
\draw[draw=none,fill=forestgreen4416044] (axis cs:450.5,0) rectangle (axis cs:451.5,86);
\draw[draw=none,fill=forestgreen4416044] (axis cs:451.5,0) rectangle (axis cs:452.5,86);
\draw[draw=none,fill=forestgreen4416044] (axis cs:452.5,0) rectangle (axis cs:453.5,86);
\draw[draw=none,fill=forestgreen4416044] (axis cs:453.5,0) rectangle (axis cs:454.5,86);
\draw[draw=none,fill=forestgreen4416044] (axis cs:454.5,0) rectangle (axis cs:455.5,86);
\draw[draw=none,fill=forestgreen4416044] (axis cs:455.5,0) rectangle (axis cs:456.5,86);
\draw[draw=none,fill=forestgreen4416044] (axis cs:456.5,0) rectangle (axis cs:457.5,86);
\draw[draw=none,fill=forestgreen4416044] (axis cs:457.5,0) rectangle (axis cs:458.5,86);
\draw[draw=none,fill=forestgreen4416044] (axis cs:458.5,0) rectangle (axis cs:459.5,86);
\draw[draw=none,fill=forestgreen4416044] (axis cs:459.5,0) rectangle (axis cs:460.5,86);
\draw[draw=none,fill=forestgreen4416044] (axis cs:460.5,0) rectangle (axis cs:461.5,86);
\draw[draw=none,fill=forestgreen4416044] (axis cs:461.5,0) rectangle (axis cs:462.5,86);
\draw[draw=none,fill=forestgreen4416044] (axis cs:462.5,0) rectangle (axis cs:463.5,86);
\draw[draw=none,fill=forestgreen4416044] (axis cs:463.5,0) rectangle (axis cs:464.5,86);
\draw[draw=none,fill=forestgreen4416044] (axis cs:464.5,0) rectangle (axis cs:465.5,86);
\draw[draw=none,fill=forestgreen4416044] (axis cs:465.5,0) rectangle (axis cs:466.5,86);
\draw[draw=none,fill=forestgreen4416044] (axis cs:466.5,0) rectangle (axis cs:467.5,86);
\draw[draw=none,fill=forestgreen4416044] (axis cs:467.5,0) rectangle (axis cs:468.5,86);
\draw[draw=none,fill=forestgreen4416044] (axis cs:468.5,0) rectangle (axis cs:469.5,86);
\draw[draw=none,fill=forestgreen4416044] (axis cs:469.5,0) rectangle (axis cs:470.5,86);
\draw[draw=none,fill=forestgreen4416044] (axis cs:470.5,0) rectangle (axis cs:471.5,86);
\draw[draw=none,fill=forestgreen4416044] (axis cs:471.5,0) rectangle (axis cs:472.5,86);
\draw[draw=none,fill=forestgreen4416044] (axis cs:472.5,0) rectangle (axis cs:473.5,86);
\draw[draw=none,fill=forestgreen4416044] (axis cs:473.5,0) rectangle (axis cs:474.5,86);
\draw[draw=none,fill=forestgreen4416044] (axis cs:474.5,0) rectangle (axis cs:475.5,86);
\draw[draw=none,fill=forestgreen4416044] (axis cs:475.5,0) rectangle (axis cs:476.5,86);
\draw[draw=none,fill=forestgreen4416044] (axis cs:476.5,0) rectangle (axis cs:477.5,86);
\draw[draw=none,fill=forestgreen4416044] (axis cs:477.5,0) rectangle (axis cs:478.5,86);
\draw[draw=none,fill=forestgreen4416044] (axis cs:478.5,0) rectangle (axis cs:479.5,86);
\draw[draw=none,fill=forestgreen4416044] (axis cs:479.5,0) rectangle (axis cs:480.5,86);
\draw[draw=none,fill=forestgreen4416044] (axis cs:480.5,0) rectangle (axis cs:481.5,86);
\draw[draw=none,fill=forestgreen4416044] (axis cs:481.5,0) rectangle (axis cs:482.5,86);
\draw[draw=none,fill=forestgreen4416044] (axis cs:482.5,0) rectangle (axis cs:483.5,86);
\draw[draw=none,fill=forestgreen4416044] (axis cs:483.5,0) rectangle (axis cs:484.5,86);
\draw[draw=none,fill=forestgreen4416044] (axis cs:484.5,0) rectangle (axis cs:485.5,86);
\draw[draw=none,fill=forestgreen4416044] (axis cs:485.5,0) rectangle (axis cs:486.5,86);
\draw[draw=none,fill=forestgreen4416044] (axis cs:486.5,0) rectangle (axis cs:487.5,86);
\draw[draw=none,fill=forestgreen4416044] (axis cs:487.5,0) rectangle (axis cs:488.5,86);
\draw[draw=none,fill=forestgreen4416044] (axis cs:488.5,0) rectangle (axis cs:489.5,86);
\draw[draw=none,fill=forestgreen4416044] (axis cs:489.5,0) rectangle (axis cs:490.5,86);
\draw[draw=none,fill=forestgreen4416044] (axis cs:490.5,0) rectangle (axis cs:491.5,88);
\draw[draw=none,fill=forestgreen4416044] (axis cs:491.5,0) rectangle (axis cs:492.5,88);
\draw[draw=none,fill=forestgreen4416044] (axis cs:492.5,0) rectangle (axis cs:493.5,88);
\draw[draw=none,fill=forestgreen4416044] (axis cs:493.5,0) rectangle (axis cs:494.5,88);
\draw[draw=none,fill=forestgreen4416044] (axis cs:494.5,0) rectangle (axis cs:495.5,88);
\draw[draw=none,fill=forestgreen4416044] (axis cs:495.5,0) rectangle (axis cs:496.5,88);
\draw[draw=none,fill=forestgreen4416044] (axis cs:496.5,0) rectangle (axis cs:497.5,88);
\draw[draw=none,fill=forestgreen4416044] (axis cs:497.5,0) rectangle (axis cs:498.5,88);
\draw[draw=none,fill=forestgreen4416044] (axis cs:498.5,0) rectangle (axis cs:499.5,88);
\draw[draw=none,fill=forestgreen4416044] (axis cs:499.5,0) rectangle (axis cs:500.5,88);
\draw[draw=none,fill=forestgreen4416044] (axis cs:500.5,0) rectangle (axis cs:501.5,88);
\draw[draw=none,fill=forestgreen4416044] (axis cs:501.5,0) rectangle (axis cs:502.5,88);
\draw[draw=none,fill=forestgreen4416044] (axis cs:502.5,0) rectangle (axis cs:503.5,88);
\draw[draw=none,fill=forestgreen4416044] (axis cs:503.5,0) rectangle (axis cs:504.5,88);
\draw[draw=none,fill=forestgreen4416044] (axis cs:504.5,0) rectangle (axis cs:505.5,88);
\draw[draw=none,fill=forestgreen4416044] (axis cs:505.5,0) rectangle (axis cs:506.5,88);
\draw[draw=none,fill=forestgreen4416044] (axis cs:506.5,0) rectangle (axis cs:507.5,88);
\draw[draw=none,fill=forestgreen4416044] (axis cs:507.5,0) rectangle (axis cs:508.5,88);
\draw[draw=none,fill=forestgreen4416044] (axis cs:508.5,0) rectangle (axis cs:509.5,88);
\draw[draw=none,fill=forestgreen4416044] (axis cs:509.5,0) rectangle (axis cs:510.5,88);
\draw[draw=none,fill=forestgreen4416044] (axis cs:510.5,0) rectangle (axis cs:511.5,88);
\draw[draw=none,fill=forestgreen4416044] (axis cs:511.5,0) rectangle (axis cs:512.5,88);
\draw[draw=none,fill=forestgreen4416044] (axis cs:512.5,0) rectangle (axis cs:513.5,88);
\draw[draw=none,fill=forestgreen4416044] (axis cs:513.5,0) rectangle (axis cs:514.5,90);
\draw[draw=none,fill=forestgreen4416044] (axis cs:514.5,0) rectangle (axis cs:515.5,90);
\draw[draw=none,fill=forestgreen4416044] (axis cs:515.5,0) rectangle (axis cs:516.5,90);
\draw[draw=none,fill=forestgreen4416044] (axis cs:516.5,0) rectangle (axis cs:517.5,90);
\draw[draw=none,fill=forestgreen4416044] (axis cs:517.5,0) rectangle (axis cs:518.5,90);
\draw[draw=none,fill=forestgreen4416044] (axis cs:518.5,0) rectangle (axis cs:519.5,90);
\draw[draw=none,fill=forestgreen4416044] (axis cs:519.5,0) rectangle (axis cs:520.5,90);
\draw[draw=none,fill=forestgreen4416044] (axis cs:520.5,0) rectangle (axis cs:521.5,90);
\draw[draw=none,fill=forestgreen4416044] (axis cs:521.5,0) rectangle (axis cs:522.5,90);
\draw[draw=none,fill=forestgreen4416044] (axis cs:522.5,0) rectangle (axis cs:523.5,92);
\draw[draw=none,fill=forestgreen4416044] (axis cs:523.5,0) rectangle (axis cs:524.5,92);
\draw[draw=none,fill=forestgreen4416044] (axis cs:524.5,0) rectangle (axis cs:525.5,92);
\draw[draw=none,fill=forestgreen4416044] (axis cs:525.5,0) rectangle (axis cs:526.5,92);
\draw[draw=none,fill=forestgreen4416044] (axis cs:526.5,0) rectangle (axis cs:527.5,92);
\draw[draw=none,fill=forestgreen4416044] (axis cs:527.5,0) rectangle (axis cs:528.5,92);
\draw[draw=none,fill=forestgreen4416044] (axis cs:528.5,0) rectangle (axis cs:529.5,92);
\draw[draw=none,fill=forestgreen4416044] (axis cs:529.5,0) rectangle (axis cs:530.5,92);
\draw[draw=none,fill=forestgreen4416044] (axis cs:530.5,0) rectangle (axis cs:531.5,92);
\draw[draw=none,fill=forestgreen4416044] (axis cs:531.5,0) rectangle (axis cs:532.5,92);
\draw[draw=none,fill=forestgreen4416044] (axis cs:532.5,0) rectangle (axis cs:533.5,92);
\draw[draw=none,fill=forestgreen4416044] (axis cs:533.5,0) rectangle (axis cs:534.5,92);
\draw[draw=none,fill=forestgreen4416044] (axis cs:534.5,0) rectangle (axis cs:535.5,92);
\draw[draw=none,fill=forestgreen4416044] (axis cs:535.5,0) rectangle (axis cs:536.5,92);
\draw[draw=none,fill=forestgreen4416044] (axis cs:536.5,0) rectangle (axis cs:537.5,92);
\draw[draw=none,fill=forestgreen4416044] (axis cs:537.5,0) rectangle (axis cs:538.5,92);
\draw[draw=none,fill=forestgreen4416044] (axis cs:538.5,0) rectangle (axis cs:539.5,92);
\draw[draw=none,fill=forestgreen4416044] (axis cs:539.5,0) rectangle (axis cs:540.5,92);
\draw[draw=none,fill=forestgreen4416044] (axis cs:540.5,0) rectangle (axis cs:541.5,92);
\draw[draw=none,fill=forestgreen4416044] (axis cs:541.5,0) rectangle (axis cs:542.5,92);
\draw[draw=none,fill=forestgreen4416044] (axis cs:542.5,0) rectangle (axis cs:543.5,92);
\draw[draw=none,fill=forestgreen4416044] (axis cs:543.5,0) rectangle (axis cs:544.5,92);
\draw[draw=none,fill=forestgreen4416044] (axis cs:544.5,0) rectangle (axis cs:545.5,92);
\draw[draw=none,fill=forestgreen4416044] (axis cs:545.5,0) rectangle (axis cs:546.5,92);
\draw[draw=none,fill=forestgreen4416044] (axis cs:546.5,0) rectangle (axis cs:547.5,92);
\draw[draw=none,fill=forestgreen4416044] (axis cs:547.5,0) rectangle (axis cs:548.5,92);
\draw[draw=none,fill=forestgreen4416044] (axis cs:548.5,0) rectangle (axis cs:549.5,92);
\draw[draw=none,fill=forestgreen4416044] (axis cs:549.5,0) rectangle (axis cs:550.5,92);
\draw[draw=none,fill=forestgreen4416044] (axis cs:550.5,0) rectangle (axis cs:551.5,92);
\draw[draw=none,fill=forestgreen4416044] (axis cs:551.5,0) rectangle (axis cs:552.5,92);
\draw[draw=none,fill=forestgreen4416044] (axis cs:552.5,0) rectangle (axis cs:553.5,92);
\draw[draw=none,fill=forestgreen4416044] (axis cs:553.5,0) rectangle (axis cs:554.5,92);
\draw[draw=none,fill=forestgreen4416044] (axis cs:554.5,0) rectangle (axis cs:555.5,94);
\draw[draw=none,fill=forestgreen4416044] (axis cs:555.5,0) rectangle (axis cs:556.5,94);
\draw[draw=none,fill=forestgreen4416044] (axis cs:556.5,0) rectangle (axis cs:557.5,94);
\draw[draw=none,fill=forestgreen4416044] (axis cs:557.5,0) rectangle (axis cs:558.5,94);
\draw[draw=none,fill=forestgreen4416044] (axis cs:558.5,0) rectangle (axis cs:559.5,94);
\draw[draw=none,fill=forestgreen4416044] (axis cs:559.5,0) rectangle (axis cs:560.5,94);
\draw[draw=none,fill=forestgreen4416044] (axis cs:560.5,0) rectangle (axis cs:561.5,94);
\draw[draw=none,fill=forestgreen4416044] (axis cs:561.5,0) rectangle (axis cs:562.5,94);
\draw[draw=none,fill=forestgreen4416044] (axis cs:562.5,0) rectangle (axis cs:563.5,94);
\draw[draw=none,fill=forestgreen4416044] (axis cs:563.5,0) rectangle (axis cs:564.5,94);
\draw[draw=none,fill=forestgreen4416044] (axis cs:564.5,0) rectangle (axis cs:565.5,94);
\draw[draw=none,fill=forestgreen4416044] (axis cs:565.5,0) rectangle (axis cs:566.5,94);
\draw[draw=none,fill=forestgreen4416044] (axis cs:566.5,0) rectangle (axis cs:567.5,94);
\draw[draw=none,fill=forestgreen4416044] (axis cs:567.5,0) rectangle (axis cs:568.5,94);
\draw[draw=none,fill=forestgreen4416044] (axis cs:568.5,0) rectangle (axis cs:569.5,94);
\draw[draw=none,fill=forestgreen4416044] (axis cs:569.5,0) rectangle (axis cs:570.5,94);
\draw[draw=none,fill=forestgreen4416044] (axis cs:570.5,0) rectangle (axis cs:571.5,94);
\draw[draw=none,fill=forestgreen4416044] (axis cs:571.5,0) rectangle (axis cs:572.5,94);
\draw[draw=none,fill=forestgreen4416044] (axis cs:572.5,0) rectangle (axis cs:573.5,94);
\draw[draw=none,fill=forestgreen4416044] (axis cs:573.5,0) rectangle (axis cs:574.5,94);
\draw[draw=none,fill=forestgreen4416044] (axis cs:574.5,0) rectangle (axis cs:575.5,94);
\draw[draw=none,fill=forestgreen4416044] (axis cs:575.5,0) rectangle (axis cs:576.5,94);
\draw[draw=none,fill=forestgreen4416044] (axis cs:576.5,0) rectangle (axis cs:577.5,94);
\draw[draw=none,fill=forestgreen4416044] (axis cs:577.5,0) rectangle (axis cs:578.5,94);
\draw[draw=none,fill=forestgreen4416044] (axis cs:578.5,0) rectangle (axis cs:579.5,94);
\draw[draw=none,fill=forestgreen4416044] (axis cs:579.5,0) rectangle (axis cs:580.5,94);
\draw[draw=none,fill=forestgreen4416044] (axis cs:580.5,0) rectangle (axis cs:581.5,94);
\draw[draw=none,fill=forestgreen4416044] (axis cs:581.5,0) rectangle (axis cs:582.5,94);
\draw[draw=none,fill=forestgreen4416044] (axis cs:582.5,0) rectangle (axis cs:583.5,94);
\draw[draw=none,fill=forestgreen4416044] (axis cs:583.5,0) rectangle (axis cs:584.5,96);
\draw[draw=none,fill=forestgreen4416044] (axis cs:584.5,0) rectangle (axis cs:585.5,96);
\draw[draw=none,fill=forestgreen4416044] (axis cs:585.5,0) rectangle (axis cs:586.5,96);
\draw[draw=none,fill=forestgreen4416044] (axis cs:586.5,0) rectangle (axis cs:587.5,96);
\draw[draw=none,fill=forestgreen4416044] (axis cs:587.5,0) rectangle (axis cs:588.5,96);
\draw[draw=none,fill=forestgreen4416044] (axis cs:588.5,0) rectangle (axis cs:589.5,96);
\draw[draw=none,fill=forestgreen4416044] (axis cs:589.5,0) rectangle (axis cs:590.5,96);
\draw[draw=none,fill=forestgreen4416044] (axis cs:590.5,0) rectangle (axis cs:591.5,96);
\draw[draw=none,fill=forestgreen4416044] (axis cs:591.5,0) rectangle (axis cs:592.5,96);
\draw[draw=none,fill=forestgreen4416044] (axis cs:592.5,0) rectangle (axis cs:593.5,96);
\draw[draw=none,fill=forestgreen4416044] (axis cs:593.5,0) rectangle (axis cs:594.5,96);
\draw[draw=none,fill=forestgreen4416044] (axis cs:594.5,0) rectangle (axis cs:595.5,96);
\draw[draw=none,fill=forestgreen4416044] (axis cs:595.5,0) rectangle (axis cs:596.5,96);
\draw[draw=none,fill=forestgreen4416044] (axis cs:596.5,0) rectangle (axis cs:597.5,96);
\draw[draw=none,fill=forestgreen4416044] (axis cs:597.5,0) rectangle (axis cs:598.5,98);
\draw[draw=none,fill=forestgreen4416044] (axis cs:598.5,0) rectangle (axis cs:599.5,98);
\draw[draw=none,fill=forestgreen4416044] (axis cs:599.5,0) rectangle (axis cs:600.5,98);
\draw[draw=none,fill=forestgreen4416044] (axis cs:600.5,0) rectangle (axis cs:601.5,98);
\draw[draw=none,fill=forestgreen4416044] (axis cs:601.5,0) rectangle (axis cs:602.5,98);
\draw[draw=none,fill=forestgreen4416044] (axis cs:602.5,0) rectangle (axis cs:603.5,98);
\draw[draw=none,fill=forestgreen4416044] (axis cs:603.5,0) rectangle (axis cs:604.5,98);
\draw[draw=none,fill=forestgreen4416044] (axis cs:604.5,0) rectangle (axis cs:605.5,98);
\draw[draw=none,fill=forestgreen4416044] (axis cs:605.5,0) rectangle (axis cs:606.5,98);
\draw[draw=none,fill=forestgreen4416044] (axis cs:606.5,0) rectangle (axis cs:607.5,98);
\draw[draw=none,fill=forestgreen4416044] (axis cs:607.5,0) rectangle (axis cs:608.5,98);
\draw[draw=none,fill=forestgreen4416044] (axis cs:608.5,0) rectangle (axis cs:609.5,98);
\draw[draw=none,fill=forestgreen4416044] (axis cs:609.5,0) rectangle (axis cs:610.5,98);
\draw[draw=none,fill=forestgreen4416044] (axis cs:610.5,0) rectangle (axis cs:611.5,98);
\draw[draw=none,fill=forestgreen4416044] (axis cs:611.5,0) rectangle (axis cs:612.5,98);
\draw[draw=none,fill=forestgreen4416044] (axis cs:612.5,0) rectangle (axis cs:613.5,98);
\draw[draw=none,fill=forestgreen4416044] (axis cs:613.5,0) rectangle (axis cs:614.5,98);
\draw[draw=none,fill=forestgreen4416044] (axis cs:614.5,0) rectangle (axis cs:615.5,98);
\draw[draw=none,fill=forestgreen4416044] (axis cs:615.5,0) rectangle (axis cs:616.5,98);
\draw[draw=none,fill=forestgreen4416044] (axis cs:616.5,0) rectangle (axis cs:617.5,98);
\draw[draw=none,fill=forestgreen4416044] (axis cs:617.5,0) rectangle (axis cs:618.5,98);
\draw[draw=none,fill=forestgreen4416044] (axis cs:618.5,0) rectangle (axis cs:619.5,98);
\draw[draw=none,fill=forestgreen4416044] (axis cs:619.5,0) rectangle (axis cs:620.5,98);
\draw[draw=none,fill=forestgreen4416044] (axis cs:620.5,0) rectangle (axis cs:621.5,98);
\draw[draw=none,fill=forestgreen4416044] (axis cs:621.5,0) rectangle (axis cs:622.5,98);
\draw[draw=none,fill=forestgreen4416044] (axis cs:622.5,0) rectangle (axis cs:623.5,98);
\draw[draw=none,fill=forestgreen4416044] (axis cs:623.5,0) rectangle (axis cs:624.5,98);
\draw[draw=none,fill=forestgreen4416044] (axis cs:624.5,0) rectangle (axis cs:625.5,98);
\draw[draw=none,fill=forestgreen4416044] (axis cs:625.5,0) rectangle (axis cs:626.5,98);
\draw[draw=none,fill=forestgreen4416044] (axis cs:626.5,0) rectangle (axis cs:627.5,98);
\draw[draw=none,fill=forestgreen4416044] (axis cs:627.5,0) rectangle (axis cs:628.5,98);
\draw[draw=none,fill=forestgreen4416044] (axis cs:628.5,0) rectangle (axis cs:629.5,98);
\draw[draw=none,fill=forestgreen4416044] (axis cs:629.5,0) rectangle (axis cs:630.5,98);
\draw[draw=none,fill=forestgreen4416044] (axis cs:630.5,0) rectangle (axis cs:631.5,98);
\draw[draw=none,fill=forestgreen4416044] (axis cs:631.5,0) rectangle (axis cs:632.5,98);
\draw[draw=none,fill=forestgreen4416044] (axis cs:632.5,0) rectangle (axis cs:633.5,98);
\draw[draw=none,fill=forestgreen4416044] (axis cs:633.5,0) rectangle (axis cs:634.5,98);
\draw[draw=none,fill=forestgreen4416044] (axis cs:634.5,0) rectangle (axis cs:635.5,98);
\draw[draw=none,fill=forestgreen4416044] (axis cs:635.5,0) rectangle (axis cs:636.5,98);
\draw[draw=none,fill=forestgreen4416044] (axis cs:636.5,0) rectangle (axis cs:637.5,98);
\draw[draw=none,fill=forestgreen4416044] (axis cs:637.5,0) rectangle (axis cs:638.5,98);
\draw[draw=none,fill=forestgreen4416044] (axis cs:638.5,0) rectangle (axis cs:639.5,98);
\draw[draw=none,fill=forestgreen4416044] (axis cs:639.5,0) rectangle (axis cs:640.5,98);
\draw[draw=none,fill=forestgreen4416044] (axis cs:640.5,0) rectangle (axis cs:641.5,98);
\draw[draw=none,fill=forestgreen4416044] (axis cs:641.5,0) rectangle (axis cs:642.5,98);
\draw[draw=none,fill=forestgreen4416044] (axis cs:642.5,0) rectangle (axis cs:643.5,98);
\draw[draw=none,fill=forestgreen4416044] (axis cs:643.5,0) rectangle (axis cs:644.5,98);
\draw[draw=none,fill=forestgreen4416044] (axis cs:644.5,0) rectangle (axis cs:645.5,98);
\draw[draw=none,fill=forestgreen4416044] (axis cs:645.5,0) rectangle (axis cs:646.5,98);
\draw[draw=none,fill=forestgreen4416044] (axis cs:646.5,0) rectangle (axis cs:647.5,98);
\draw[draw=none,fill=forestgreen4416044] (axis cs:647.5,0) rectangle (axis cs:648.5,98);
\draw[draw=none,fill=forestgreen4416044] (axis cs:648.5,0) rectangle (axis cs:649.5,98);
\draw[draw=none,fill=forestgreen4416044] (axis cs:649.5,0) rectangle (axis cs:650.5,98);
\draw[draw=none,fill=forestgreen4416044] (axis cs:650.5,0) rectangle (axis cs:651.5,98);
\draw[draw=none,fill=forestgreen4416044] (axis cs:651.5,0) rectangle (axis cs:652.5,98);
\draw[draw=none,fill=forestgreen4416044] (axis cs:652.5,0) rectangle (axis cs:653.5,98);
\draw[draw=none,fill=forestgreen4416044] (axis cs:653.5,0) rectangle (axis cs:654.5,98);
\draw[draw=none,fill=forestgreen4416044] (axis cs:654.5,0) rectangle (axis cs:655.5,98);
\draw[draw=none,fill=forestgreen4416044] (axis cs:655.5,0) rectangle (axis cs:656.5,98);
\draw[draw=none,fill=forestgreen4416044] (axis cs:656.5,0) rectangle (axis cs:657.5,98);
\draw[draw=none,fill=forestgreen4416044] (axis cs:657.5,0) rectangle (axis cs:658.5,98);
\draw[draw=none,fill=forestgreen4416044] (axis cs:658.5,0) rectangle (axis cs:659.5,98);
\draw[draw=none,fill=forestgreen4416044] (axis cs:659.5,0) rectangle (axis cs:660.5,98);
\draw[draw=none,fill=forestgreen4416044] (axis cs:660.5,0) rectangle (axis cs:661.5,98);
\draw[draw=none,fill=forestgreen4416044] (axis cs:661.5,0) rectangle (axis cs:662.5,98);
\draw[draw=none,fill=forestgreen4416044] (axis cs:662.5,0) rectangle (axis cs:663.5,98);
\draw[draw=none,fill=forestgreen4416044] (axis cs:663.5,0) rectangle (axis cs:664.5,98);
\draw[draw=none,fill=forestgreen4416044] (axis cs:664.5,0) rectangle (axis cs:665.5,98);
\draw[draw=none,fill=forestgreen4416044] (axis cs:665.5,0) rectangle (axis cs:666.5,98);
\draw[draw=none,fill=forestgreen4416044] (axis cs:666.5,0) rectangle (axis cs:667.5,98);
\draw[draw=none,fill=forestgreen4416044] (axis cs:667.5,0) rectangle (axis cs:668.5,98);
\draw[draw=none,fill=forestgreen4416044] (axis cs:668.5,0) rectangle (axis cs:669.5,98);
\draw[draw=none,fill=forestgreen4416044] (axis cs:669.5,0) rectangle (axis cs:670.5,98);
\draw[draw=none,fill=forestgreen4416044] (axis cs:670.5,0) rectangle (axis cs:671.5,98);
\draw[draw=none,fill=forestgreen4416044] (axis cs:671.5,0) rectangle (axis cs:672.5,98);
\draw[draw=none,fill=forestgreen4416044] (axis cs:672.5,0) rectangle (axis cs:673.5,98);
\draw[draw=none,fill=forestgreen4416044] (axis cs:673.5,0) rectangle (axis cs:674.5,98);
\draw[draw=none,fill=forestgreen4416044] (axis cs:674.5,0) rectangle (axis cs:675.5,98);
\draw[draw=none,fill=forestgreen4416044] (axis cs:675.5,0) rectangle (axis cs:676.5,98);
\draw[draw=none,fill=forestgreen4416044] (axis cs:676.5,0) rectangle (axis cs:677.5,98);
\draw[draw=none,fill=forestgreen4416044] (axis cs:677.5,0) rectangle (axis cs:678.5,98);
\draw[draw=none,fill=forestgreen4416044] (axis cs:678.5,0) rectangle (axis cs:679.5,98);
\draw[draw=none,fill=forestgreen4416044] (axis cs:679.5,0) rectangle (axis cs:680.5,98);
\draw[draw=none,fill=forestgreen4416044] (axis cs:680.5,0) rectangle (axis cs:681.5,98);
\draw[draw=none,fill=forestgreen4416044] (axis cs:681.5,0) rectangle (axis cs:682.5,98);
\draw[draw=none,fill=forestgreen4416044] (axis cs:682.5,0) rectangle (axis cs:683.5,98);
\draw[draw=none,fill=forestgreen4416044] (axis cs:683.5,0) rectangle (axis cs:684.5,98);
\draw[draw=none,fill=forestgreen4416044] (axis cs:684.5,0) rectangle (axis cs:685.5,98);
\draw[draw=none,fill=forestgreen4416044] (axis cs:685.5,0) rectangle (axis cs:686.5,98);
\draw[draw=none,fill=forestgreen4416044] (axis cs:686.5,0) rectangle (axis cs:687.5,98);
\draw[draw=none,fill=forestgreen4416044] (axis cs:687.5,0) rectangle (axis cs:688.5,98);
\draw[draw=none,fill=forestgreen4416044] (axis cs:688.5,0) rectangle (axis cs:689.5,98);
\draw[draw=none,fill=forestgreen4416044] (axis cs:689.5,0) rectangle (axis cs:690.5,98);
\draw[draw=none,fill=forestgreen4416044] (axis cs:690.5,0) rectangle (axis cs:691.5,98);
\draw[draw=none,fill=forestgreen4416044] (axis cs:691.5,0) rectangle (axis cs:692.5,98);
\draw[draw=none,fill=forestgreen4416044] (axis cs:692.5,0) rectangle (axis cs:693.5,98);
\draw[draw=none,fill=forestgreen4416044] (axis cs:693.5,0) rectangle (axis cs:694.5,98);
\draw[draw=none,fill=forestgreen4416044] (axis cs:694.5,0) rectangle (axis cs:695.5,98);
\draw[draw=none,fill=forestgreen4416044] (axis cs:695.5,0) rectangle (axis cs:696.5,98);
\draw[draw=none,fill=forestgreen4416044] (axis cs:696.5,0) rectangle (axis cs:697.5,98);
\draw[draw=none,fill=forestgreen4416044] (axis cs:697.5,0) rectangle (axis cs:698.5,98);
\draw[draw=none,fill=forestgreen4416044] (axis cs:698.5,0) rectangle (axis cs:699.5,98);
\draw[draw=none,fill=forestgreen4416044] (axis cs:699.5,0) rectangle (axis cs:700.5,98);
\draw[draw=none,fill=forestgreen4416044] (axis cs:700.5,0) rectangle (axis cs:701.5,98);
\draw[draw=none,fill=forestgreen4416044] (axis cs:701.5,0) rectangle (axis cs:702.5,98);
\draw[draw=none,fill=forestgreen4416044] (axis cs:702.5,0) rectangle (axis cs:703.5,98);
\draw[draw=none,fill=forestgreen4416044] (axis cs:703.5,0) rectangle (axis cs:704.5,98);
\draw[draw=none,fill=forestgreen4416044] (axis cs:704.5,0) rectangle (axis cs:705.5,98);
\draw[draw=none,fill=forestgreen4416044] (axis cs:705.5,0) rectangle (axis cs:706.5,98);
\draw[draw=none,fill=forestgreen4416044] (axis cs:706.5,0) rectangle (axis cs:707.5,98);
\draw[draw=none,fill=forestgreen4416044] (axis cs:707.5,0) rectangle (axis cs:708.5,98);
\draw[draw=none,fill=forestgreen4416044] (axis cs:708.5,0) rectangle (axis cs:709.5,98);
\draw[draw=none,fill=forestgreen4416044] (axis cs:709.5,0) rectangle (axis cs:710.5,98);
\draw[draw=none,fill=forestgreen4416044] (axis cs:710.5,0) rectangle (axis cs:711.5,98);
\draw[draw=none,fill=forestgreen4416044] (axis cs:711.5,0) rectangle (axis cs:712.5,98);
\draw[draw=none,fill=forestgreen4416044] (axis cs:712.5,0) rectangle (axis cs:713.5,98);
\draw[draw=none,fill=forestgreen4416044] (axis cs:713.5,0) rectangle (axis cs:714.5,98);
\draw[draw=none,fill=forestgreen4416044] (axis cs:714.5,0) rectangle (axis cs:715.5,98);
\draw[draw=none,fill=forestgreen4416044] (axis cs:715.5,0) rectangle (axis cs:716.5,98);
\draw[draw=none,fill=forestgreen4416044] (axis cs:716.5,0) rectangle (axis cs:717.5,98);
\draw[draw=none,fill=forestgreen4416044] (axis cs:717.5,0) rectangle (axis cs:718.5,98);
\draw[draw=none,fill=forestgreen4416044] (axis cs:718.5,0) rectangle (axis cs:719.5,98);
\draw[draw=none,fill=forestgreen4416044] (axis cs:719.5,0) rectangle (axis cs:720.5,98);
\draw[draw=none,fill=forestgreen4416044] (axis cs:720.5,0) rectangle (axis cs:721.5,98);
\draw[draw=none,fill=forestgreen4416044] (axis cs:721.5,0) rectangle (axis cs:722.5,98);
\draw[draw=none,fill=forestgreen4416044] (axis cs:722.5,0) rectangle (axis cs:723.5,98);
\draw[draw=none,fill=forestgreen4416044] (axis cs:723.5,0) rectangle (axis cs:724.5,98);
\draw[draw=none,fill=forestgreen4416044] (axis cs:724.5,0) rectangle (axis cs:725.5,98);
\draw[draw=none,fill=forestgreen4416044] (axis cs:725.5,0) rectangle (axis cs:726.5,98);
\draw[draw=none,fill=forestgreen4416044] (axis cs:726.5,0) rectangle (axis cs:727.5,98);
\draw[draw=none,fill=forestgreen4416044] (axis cs:727.5,0) rectangle (axis cs:728.5,98);
\draw[draw=none,fill=forestgreen4416044] (axis cs:728.5,0) rectangle (axis cs:729.5,98);
\draw[draw=none,fill=forestgreen4416044] (axis cs:729.5,0) rectangle (axis cs:730.5,98);
\draw[draw=none,fill=forestgreen4416044] (axis cs:730.5,0) rectangle (axis cs:731.5,98);
\draw[draw=none,fill=forestgreen4416044] (axis cs:731.5,0) rectangle (axis cs:732.5,98);
\draw[draw=none,fill=forestgreen4416044] (axis cs:732.5,0) rectangle (axis cs:733.5,98);
\draw[draw=none,fill=forestgreen4416044] (axis cs:733.5,0) rectangle (axis cs:734.5,98);
\draw[draw=none,fill=forestgreen4416044] (axis cs:734.5,0) rectangle (axis cs:735.5,98);
\draw[draw=none,fill=forestgreen4416044] (axis cs:735.5,0) rectangle (axis cs:736.5,98);
\draw[draw=none,fill=forestgreen4416044] (axis cs:736.5,0) rectangle (axis cs:737.5,98);
\draw[draw=none,fill=forestgreen4416044] (axis cs:737.5,0) rectangle (axis cs:738.5,98);
\draw[draw=none,fill=forestgreen4416044] (axis cs:738.5,0) rectangle (axis cs:739.5,98);
\draw[draw=none,fill=forestgreen4416044] (axis cs:739.5,0) rectangle (axis cs:740.5,98);
\draw[draw=none,fill=forestgreen4416044] (axis cs:740.5,0) rectangle (axis cs:741.5,98);
\draw[draw=none,fill=forestgreen4416044] (axis cs:741.5,0) rectangle (axis cs:742.5,98);
\draw[draw=none,fill=forestgreen4416044] (axis cs:742.5,0) rectangle (axis cs:743.5,98);
\draw[draw=none,fill=forestgreen4416044] (axis cs:743.5,0) rectangle (axis cs:744.5,98);
\draw[draw=none,fill=forestgreen4416044] (axis cs:744.5,0) rectangle (axis cs:745.5,98);
\draw[draw=none,fill=forestgreen4416044] (axis cs:745.5,0) rectangle (axis cs:746.5,98);
\draw[draw=none,fill=forestgreen4416044] (axis cs:746.5,0) rectangle (axis cs:747.5,98);
\draw[draw=none,fill=forestgreen4416044] (axis cs:747.5,0) rectangle (axis cs:748.5,98);
\draw[draw=none,fill=forestgreen4416044] (axis cs:748.5,0) rectangle (axis cs:749.5,98);
\draw[draw=none,fill=forestgreen4416044] (axis cs:749.5,0) rectangle (axis cs:750.5,98);
\draw[draw=none,fill=forestgreen4416044] (axis cs:750.5,0) rectangle (axis cs:751.5,98);
\draw[draw=none,fill=forestgreen4416044] (axis cs:751.5,0) rectangle (axis cs:752.5,98);
\draw[draw=none,fill=forestgreen4416044] (axis cs:752.5,0) rectangle (axis cs:753.5,98);
\draw[draw=none,fill=forestgreen4416044] (axis cs:753.5,0) rectangle (axis cs:754.5,98);
\draw[draw=none,fill=forestgreen4416044] (axis cs:754.5,0) rectangle (axis cs:755.5,98);
\draw[draw=none,fill=forestgreen4416044] (axis cs:755.5,0) rectangle (axis cs:756.5,98);
\draw[draw=none,fill=forestgreen4416044] (axis cs:756.5,0) rectangle (axis cs:757.5,98);
\draw[draw=none,fill=forestgreen4416044] (axis cs:757.5,0) rectangle (axis cs:758.5,98);
\draw[draw=none,fill=forestgreen4416044] (axis cs:758.5,0) rectangle (axis cs:759.5,98);
\draw[draw=none,fill=forestgreen4416044] (axis cs:759.5,0) rectangle (axis cs:760.5,98);
\draw[draw=none,fill=forestgreen4416044] (axis cs:760.5,0) rectangle (axis cs:761.5,98);
\draw[draw=none,fill=forestgreen4416044] (axis cs:761.5,0) rectangle (axis cs:762.5,98);
\draw[draw=none,fill=forestgreen4416044] (axis cs:762.5,0) rectangle (axis cs:763.5,98);
\draw[draw=none,fill=forestgreen4416044] (axis cs:763.5,0) rectangle (axis cs:764.5,98);
\draw[draw=none,fill=forestgreen4416044] (axis cs:764.5,0) rectangle (axis cs:765.5,98);
\draw[draw=none,fill=forestgreen4416044] (axis cs:765.5,0) rectangle (axis cs:766.5,98);
\draw[draw=none,fill=forestgreen4416044] (axis cs:766.5,0) rectangle (axis cs:767.5,98);
\draw[draw=none,fill=forestgreen4416044] (axis cs:767.5,0) rectangle (axis cs:768.5,98);
\draw[draw=none,fill=forestgreen4416044] (axis cs:768.5,0) rectangle (axis cs:769.5,98);
\draw[draw=none,fill=forestgreen4416044] (axis cs:769.5,0) rectangle (axis cs:770.5,98);
\draw[draw=none,fill=forestgreen4416044] (axis cs:770.5,0) rectangle (axis cs:771.5,98);
\draw[draw=none,fill=forestgreen4416044] (axis cs:771.5,0) rectangle (axis cs:772.5,98);
\draw[draw=none,fill=forestgreen4416044] (axis cs:772.5,0) rectangle (axis cs:773.5,98);
\draw[draw=none,fill=forestgreen4416044] (axis cs:773.5,0) rectangle (axis cs:774.5,98);
\draw[draw=none,fill=forestgreen4416044] (axis cs:774.5,0) rectangle (axis cs:775.5,98);
\draw[draw=none,fill=forestgreen4416044] (axis cs:775.5,0) rectangle (axis cs:776.5,98);
\draw[draw=none,fill=forestgreen4416044] (axis cs:776.5,0) rectangle (axis cs:777.5,98);
\draw[draw=none,fill=forestgreen4416044] (axis cs:777.5,0) rectangle (axis cs:778.5,98);
\draw[draw=none,fill=forestgreen4416044] (axis cs:778.5,0) rectangle (axis cs:779.5,98);
\draw[draw=none,fill=forestgreen4416044] (axis cs:779.5,0) rectangle (axis cs:780.5,98);
\draw[draw=none,fill=forestgreen4416044] (axis cs:780.5,0) rectangle (axis cs:781.5,98);
\draw[draw=none,fill=forestgreen4416044] (axis cs:781.5,0) rectangle (axis cs:782.5,98);
\draw[draw=none,fill=forestgreen4416044] (axis cs:782.5,0) rectangle (axis cs:783.5,98);
\draw[draw=none,fill=forestgreen4416044] (axis cs:783.5,0) rectangle (axis cs:784.5,98);
\draw[draw=none,fill=forestgreen4416044] (axis cs:784.5,0) rectangle (axis cs:785.5,98);
\draw[draw=none,fill=forestgreen4416044] (axis cs:785.5,0) rectangle (axis cs:786.5,98);
\end{axis}

\end{tikzpicture}
    \caption{Sample efficiency of Tidyparse at varying Levenshtein radii. After drawing up to $\sim10^5$ samples without replacement we can usually retrieve the human repair for almost all repairs fewer than four edits.}\label{fig:sample_efficiency}
  \end{figure}

%  The advantage of using an enumerative sampler is that it can terminate early, since it knows, after drawing $|\mathbb{T}_2(G_\cap)|$ distinct repairs, it has exhausted the admissible set. This also provides a way to calibrate the threshold for how to sample given a fixed budget: if, for example, we can only afford to sample a small fraction of the admissible set, we should clearly sample with replacement. Otherwise, if we can saturate the admissible set, we should sample without replacement.

  \begin{wrapfigure}{r}{0.35\textwidth}
    \vspace{-0.4cm}
    \resizebox{.35\textwidth}{!}{% This file was created with tikzplotlib v0.10.1.
\begin{tikzpicture}
\begin{axis}[
  legend cell align={left},
  legend style={fill opacity=0.8, draw opacity=1, text opacity=1, legend columns=-1, legend pos=north east},
  tick align=outside,
  tick pos=left,
  axis lines*=left,
  title={\textbf{Average throughput}},
  x grid style={darkgray176},
  xlabel={Lexical tokens},
  xmin=-0.95, xmax=19.95,
  xtick style={color=black},
  xtick={0,1,2,3,4,5,6,7,8,9,10,11,12,13,14,15,16,17,18,19},
  xticklabels={4, , 8, , 12, , 16, , 20, , 24, , 28, , 32, , 36, , 40,},
  y grid style={darkgray176},
  ylabel={Total repairs discovered},
  ymin=0, ymax=200000,
  ymode=log,
  legend pos=north east,
  ytick style={color=black}
]
\addlegendimage{empty legend}
\addlegendentry{$\Delta(\err\sigma, \sigma')=$}
\addlegendentry{\footnotesize{$1,$}}
\addlegendentry{\footnotesize{$2,$}}
\addlegendentry{\footnotesize{$3,$}}
\addlegendentry{\footnotesize{$4$}}
\addplot [semithick, green, mark=*, mark size=3, mark options={solid}]
table {%
  0 7.14285714285714
  1 12.1666666666667
  2 9.27272727272727
  3 10.6041666666667
  4 7.22857142857143
  5 6.4
  6 10.3956043956044
  7 7.86075949367089
  8 9.81818181818182
  9 8.28089887640449
  10 8.08695652173913
  11 10.0430107526882
  12 6.90476190476191
  13 6.94029850746269
  14 8.74698795180723
  15 7.70422535211268
  16 12.1414141414141
  17 4.1
};
\addplot [semithick, blue, mark=*, mark size=3, mark options={solid}]
table {%
0 46.625
1 135.653846153846
2 175.676923076923
3 79.0090909090909
4 87.6589147286822
5 69.9683544303797
6 83.9896373056995
7 85.9777777777778
8 73.3684210526316
9 67.7715517241379
10 92.1878172588832
11 186.568527918782
12 116.195945945946
13 87.8901734104046
14 99.6514285714286
15 154.44
16 123.658064516129
17 36.6111111111111
};
\addplot [semithick, orange, mark=*, mark size=3, mark options={solid}]
table {%
0 1506.8
1 2076.57142857143
2 5194.91304347826
3 413.939393939394
4 3749.34693877551
5 1130.83076923077
6 431.141025641026
7 601.172839506173
8 2135.26666666667
9 1233.85106382979
10 960.263888888889
11 932.225352112676
12 612.833333333333
13 880.2625
14 965.863636363636
15 725.139784946237
16 1199.05882352941
17 114.5
};
\addplot [semithick, red, mark=*, mark size=3, mark options={solid}]
table {%
0 nan
1 73461
2 24965
3 4220.4
4 35145.8181818182
5 2613.28571428571
6 68402.25
7 3300.5
8 10010
9 24536.5217391304
10 17497.6315789474
11 8883.58333333333
12 1948.23076923077
13 4175.27272727273
14 5145.45454545455
15 27785.8095238095
16 2468
17 4027.25
};
\end{axis}

\end{tikzpicture}}
    \caption{Distinct repairs found in 30s.}
    \label{fig:throughput}
    \vspace{-0.4cm}
  \end{wrapfigure}

  End-to-end throughput varies significantly with the edit distance of the repair. Some errors are trivial to fix, while others require a large number of edits to be sampled before the ground truth is discovered. We evaluate throughput by sampling patches across invalid strings $|\err\sigma| \leq 40$ from the StackOverflow dataset balanced across length and distance, and measure the total number of unique valid patches discovered, as a function of string length and edit distance $\Delta\in[1, 4]$. Each trial is terminated after 10 seconds, and the experiment is repeated across 7.3k total repairs. Note the y-axis is log-scaled, as the number of admissible repairs increases sharply with edit distance. Our approach discovers a large number of syntactic repairs in a relatively short amount of time, and is able to quickly saturate the admissible set for $\Delta(\err\sigma, \sigma') \in [1, 4]$ before timeout.

  \begin{wrapfigure}{l}{0.35\textwidth}
    \vspace{-0.4cm}
    \resizebox{.35\textwidth}{!}{% This file was created with tikzplotlib v0.10.1.
\begin{tikzpicture}

\definecolor{green}{RGB}{0,128,0}

\begin{axis}[
log basis y={10},
tick align=outside,
tick pos=left,
x grid style={darkgray176},
xlabel={Snippet length},
xmin=9.5, xmax=42.5,
title={Latency across repair size},
xtick style={color=black},
y grid style={darkgray176},
log basis y={10},
axis lines*=left,
ylabel={Repair Latency (ms)},
ymin=65.2014650905023, ymax=509998.877384793,
ymode=log,
mark size=1.2pt,
ytick style={color=black},
legend style={fill opacity=0.8, draw opacity=1, text opacity=1, legend columns=2, legend pos=south east},
yticklabels={
  \(\displaystyle {10^{0}}\),
  \(\displaystyle {10^{1}}\),
  \(\displaystyle {10^{2}}\),
  \(\displaystyle {10^{3}}\),
  \(\displaystyle {10^{4}}\),
  \(\displaystyle {10^{5}}\),
  \(\displaystyle {10^{6}}\),
  \(\displaystyle {10^{7}}\)
}
]
\addlegendentry{\footnotesize{$\Delta = 1$}}
\addlegendentry{\footnotesize{$\Delta = 4$}}
\addlegendentry{\footnotesize{$\Delta = 2$}}
\addlegendentry{\footnotesize{$\Delta = 3$}}
\addplot [draw=green, fill=green, mark=*, only marks]
table{%
x  y
32 5148
29 1085
38 5383
19 1011
18 763
21 334
39 14789
29 3437
22 2883
17 764
28 2089
22 4131
20 960
21 1591
30 3769
35 4770
15 429
28 1943
25 2659
22 1590
27 851
14 173
36 2007
16 473
18 779
33 1866
39 2838
28 815
31 4644
40 4755
16 338
11 158
29 4790
23 3520
18 750
21 3491
27 2735
40 12157
22 4114
25 988
26 2727
13 193
11 151
25 3897
11 153
15 409
17 523
16 376
33 5443
16 494
34 8625
26 933
20 331
24 1043
12 192
21 3604
36 2593
15 2718
15 252
13 377
28 3815
13 245
12 146
29 2346
22 958
18 641
26 1208
26 3619
39 4494
16 674
30 4110
34 4124
28 3496
21 1387
27 3300
16 398
17 2755
13 282
36 4758
39 7288
35 1486
25 3472
34 1606
15 453
24 1210
19 782
12 173
35 4494
37 6857
12 127
21 1142
16 3631
18 574
39 3727
26 3361
36 1863
39 3152
39 4837
14 416
14 455
12 252
27 2046
11 114
27 5193
26 2286
25 1292
22 607
35 3958
16 263
26 1651
34 7934
31 1462
12 125
20 330
26 1767
23 479
17 539
37 3038
31 1084
18 880
32 2256
31 4098
37 1607
16 227
17 1050
24 3343
28 5643
40 2478
28 3858
27 1994
22 1267
24 1994
24 617
27 2141
19 414
11 2761
13 169
39 9225
28 3785
21 1130
34 1299
17 1016
20 3013
34 2118
21 4000
34 6822
21 303
21 400
20 857
37 1936
20 1063
18 2657
16 357
31 1229
25 1999
26 748
15 404
18 438
26 583
27 1157
31 6206
26 2629
16 229
17 612
33 10128
37 5953
11 196
14 304
12 303
41 4299
25 1908
20 3028
33 4834
40 2165
27 4406
18 396
38 2757
25 569
30 2030
29 3317
15 544
21 731
15 3022
31 1594
31 1047
27 2908
25 5094
12 173
11 150
30 3979
16 339
33 1593
31 3451
36 4199
28 6294
35 1726
16 637
11 165
15 750
25 2259
17 329
15 242
12 202
22 384
29 4869
39 5517
23 765
33 5001
32 1396
34 8403
20 1191
27 2758
33 1190
16 3353
13 186
26 5396
34 2910
33 2022
13 2877
28 1016
18 499
25 1246
33 3780
13 163
23 3873
19 806
39 8507
19 3459
34 3809
25 1802
20 585
26 2232
23 3450
32 4602
26 5929
18 347
19 3242
37 2207
17 746
22 6797
19 704
15 353
12 2762
23 580
22 1373
22 3589
25 1160
31 1695
20 1950
29 3053
39 2198
30 3956
32 2285
13 294
32 3911
13 237
27 4485
24 2690
19 934
29 4071
27 2122
24 995
24 507
18 625
24 3284
37 8257
12 3040
35 9141
14 2982
28 734
28 1306
17 2695
24 930
30 950
19 2784
32 6357
18 659
20 348
15 308
14 244
16 408
13 2797
27 3366
19 960
31 3625
11 193
30 3366
30 4727
33 5697
33 6494
29 1069
36 5357
37 3787
27 643
41 6979
32 1640
17 474
23 402
19 398
20 1209
25 1904
12 206
12 350
36 2470
35 1505
14 478
12 221
15 2811
28 3637
19 1122
35 10465
14 316
11 98
15 2795
13 3284
26 1998
12 177
36 7469
20 1307
36 3463
19 643
37 2210
36 2022
33 5903
35 4781
13 260
20 290
40 12774
35 4327
23 1640
14 318
13 163
14 276
39 7636
16 176
18 3151
32 3490
20 773
39 2417
11 143
14 326
11 135
37 9387
23 618
15 532
};
\addplot [draw=red, fill=red, mark=*, only marks]
table{%
x  y
20 58108
12 51666
15 52872
27 107841
22 45114
31 195597
29 107688
28 173325
39 307961
20 19806
41 191265
28 119322
38 116176
34 142491
23 67212
25 72841
23 72985
22 62959
36 136040
36 178035
14 48610
28 170679
36 111571
33 123007
22 61457
26 96451
41 175655
26 77194
29 35609
28 169819
14 46105
39 195626
24 50751
21 48576
37 231823
11 47882
21 84741
22 53728
33 90374
35 208976
33 162398
21 90295
39 322474
29 202614
35 168255
31 140663
28 148913
37 269034
26 115719
33 216191
25 71448
38 173519
13 43689
34 102807
13 57551
31 112047
21 61169
32 163450
22 64456
18 60343
33 109713
36 225779
19 57508
33 166085
12 51689
16 50373
19 59679
25 264645
22 78703
14 78482
33 87321
31 110919
20 82528
26 100040
37 339313
40 102371
32 145489
24 99148
18 55190
25 110437
24 80558
36 184598
18 53084
20 71368
33 313482
12 40595
13 45971
24 81069
39 137454
34 171818
20 111020
32 111393
15 51481
22 64615
20 63536
30 179273
16 15228
19 61004
31 125079
36 93285
34 142710
20 64707
19 46655
24 157899
34 112240
16 52176
34 172346
14 18390
13 19427
13 46764
31 134228
36 240347
27 120224
20 98838
21 34748
15 69259
13 50058
31 100419
20 47676
14 30778
27 69207
37 156018
40 172073
18 41322
36 84168
19 171901
17 71764
33 186628
26 127454
24 97074
23 68230
18 55273
32 154745
24 72679
16 60711
22 67811
35 112980
35 109433
14 49207
23 85428
21 67406
38 190668
34 311037
39 124408
};
\addplot [draw=blue, fill=blue, mark=*, only marks]
table{%
x  y
15 1546
36 31862
38 31248
26 16119
33 29070
20 7825
19 12970
23 4255
32 11597
35 21227
15 5376
13 2141
30 39009
19 15202
38 15296
25 13377
32 9555
24 5876
20 15019
38 22100
33 35495
11 719
28 20635
27 9206
31 6117
23 20579
36 42335
28 15167
25 4452
20 9944
11 934
19 8080
12 1903
24 5029
18 2931
37 44732
34 17401
11 1400
30 17034
33 18639
34 34009
22 19118
12 5107
18 12158
19 6940
30 15976
12 5113
18 4296
21 6390
28 21456
34 13034
25 5045
22 4333
34 22718
23 2505
30 39479
16 4844
33 35755
15 2477
22 8883
22 13650
23 5820
22 4991
14 3418
35 7532
31 11363
25 29496
20 5739
12 3618
25 18211
17 16099
18 9777
28 8551
18 8232
17 1239
15 2497
39 9304
32 25118
14 1484
26 10561
21 11388
29 7176
30 17223
21 8634
37 7977
17 5433
21 8490
19 17680
18 4292
25 5967
34 13023
28 14019
24 8835
16 3006
24 25420
25 12155
13 1941
29 11108
31 18410
26 8491
19 7119
17 3180
17 3805
28 15120
32 32599
12 5333
14 2206
13 8839
23 8906
33 30905
38 14948
31 22032
38 9392
27 9183
36 13507
37 18594
26 22120
17 4754
34 35382
21 10508
24 22208
36 40758
36 31016
37 28948
14 4118
14 5684
20 3662
29 14922
12 1274
13 4091
22 7839
20 15066
23 12974
14 973
21 8622
37 46911
14 3107
32 32976
37 9228
38 26957
22 3555
37 27104
34 9138
13 3353
15 8690
26 16640
34 23305
23 4885
22 10846
28 7864
15 7957
15 8747
30 6519
12 4219
20 4469
35 47446
19 4889
27 8550
39 57886
23 14245
23 21774
20 2934
26 14781
32 37106
35 11013
20 8515
15 2539
29 24498
31 25400
22 13575
28 32976
33 37348
30 17818
39 11397
29 20198
12 4619
29 17768
28 18812
14 9168
26 21072
29 28933
16 11018
12 2620
17 8626
23 19402
41 61109
33 25388
16 7368
33 6903
20 7721
29 9658
15 1256
11 4263
34 9959
18 2161
37 47273
30 21254
31 10892
26 11905
11 1763
26 13949
19 7328
29 28416
16 12111
20 6234
17 4271
27 11313
29 12137
21 7872
35 21597
27 21491
39 10804
19 15818
19 3881
38 47207
19 3066
28 25870
30 12676
17 5394
23 20199
28 15303
22 29638
39 56108
38 7822
30 13293
31 9281
26 37151
30 19684
17 3761
39 20861
32 14503
22 13901
16 7406
17 13836
28 11129
19 2229
28 33492
38 16182
15 4529
33 31173
23 12037
22 9964
20 5659
32 24562
12 680
24 14414
30 15023
23 12023
38 17828
18 2136
34 13575
28 22060
21 16814
25 13082
38 47841
35 7347
26 18942
32 14484
34 10717
25 15950
34 32785
27 5235
23 19239
12 1326
30 20466
27 22080
26 16954
23 13937
30 18121
32 18321
15 5283
15 10042
29 19455
32 22029
12 3696
25 14987
38 15618
15 2085
38 13985
37 23038
18 4363
17 3734
21 6877
36 30820
36 40020
23 11704
36 28883
33 47148
41 13381
34 11999
19 4825
29 35034
15 5870
20 8343
23 9744
36 8368
31 37536
37 67404
38 14815
25 17157
31 8009
22 11478
23 17884
27 16520
11 1158
16 7454
28 9036
26 7471
14 1098
15 6672
14 4076
};
\addplot [draw=orange, fill=orange, mark=*, only marks]
table{%
x  y
40 42016
20 24189
33 33848
21 10526
19 29639
30 36650
20 40015
23 8429
38 32819
27 69823
29 43136
26 63140
17 42342
32 17391
12 40540
28 54691
26 43793
19 9754
34 173208
36 80096
12 16827
33 33888
16 39156
29 64485
23 14735
37 126065
36 118322
39 101483
25 57385
16 36604
31 66360
11 10512
38 126586
38 137665
41 32922
25 32976
37 65246
40 65888
22 66806
27 51204
29 24094
32 39285
37 72613
26 65788
28 44251
16 5413
35 63518
32 110601
31 36090
29 47350
19 12132
31 52986
36 149418
11 37909
41 32579
26 72372
29 87247
22 50063
29 56161
23 9463
11 39435
32 48779
35 73904
28 67568
29 71226
29 30317
25 26198
35 47708
37 132547
39 174702
17 13890
30 38364
15 41787
28 83550
29 55774
14 15458
30 46008
39 60620
22 45732
37 102710
14 42288
30 87541
38 102615
30 319242
23 56070
33 58231
20 31861
15 29942
14 22163
15 15176
19 48163
19 16105
29 53572
37 159305
39 62454
28 18109
20 30293
15 11821
33 89003
27 30070
31 63481
30 55899
29 61285
19 50656
32 16841
19 16676
21 42955
15 36952
19 11367
15 14400
29 64709
30 48081
29 53699
21 57294
12 38724
31 63965
23 49107
12 37097
12 37806
34 19660
17 17929
31 64083
33 88215
21 58401
15 39393
12 36692
17 39872
32 77042
35 127043
35 126005
12 18445
20 22042
26 44232
12 17708
11 38234
36 94269
35 120430
33 75175
21 39828
31 81582
36 96767
11 5778
19 43704
13 8415
23 44485
32 79172
29 60849
12 3793
27 41624
20 26682
17 14129
33 126568
24 46496
29 48634
16 42893
27 77985
24 65710
27 85586
31 62940
41 68300
33 102270
11 5412
41 191452
31 59553
17 52058
15 38902
32 28359
17 45407
34 135818
21 5610
31 112023
13 37517
19 14117
18 33770
29 58688
38 110486
21 44776
23 8391
19 23909
25 60406
30 28910
36 138052
36 135891
39 80921
30 108163
21 52436
20 21960
37 106582
35 78295
31 69280
22 27743
37 79928
12 15280
30 116392
14 7227
37 35707
26 49205
35 58959
35 149408
23 12075
37 101615
33 60528
18 59862
37 177210
33 77206
28 57257
13 7439
12 17730
24 68492
39 112291
22 47825
24 76251
13 34728
23 42931
39 145463
37 98906
33 72393
23 44026
28 23606
16 27186
22 40794
18 46860
26 34273
21 9723
26 74168
24 44922
38 118479
30 70306
27 67938
22 61961
35 95075
26 111028
17 42197
13 41792
13 10854
34 61113
15 39520
34 62196
29 50751
41 133897
18 22246
27 21455
16 9835
14 11967
37 35868
21 45493
37 69815
35 27949
28 44323
32 47729
32 83133
22 47076
};
\end{axis}

\end{tikzpicture}
}
    \caption{End-to-end repair timings.}
    \label{fig:timings}
    \vspace{-0.4cm}
  \end{wrapfigure}

  In Fig.~\ref{fig:timings}, we plot the end-to-end repair timings by collecting 1000 samples balanced across length and edit distance, then measure the wallclock time until the sampler retrieves the human repair and report the log latency. While short repairs finish quickly, latency is positively correlated with length and edit distance. Our method is typically able to saturate the admissible set for 1- and 2-edit repairs before timeout, while 4+-edit throughput starts becoming constrained by compute around 30s, when Python's admissible set approaches a volume of $10^5$ valid edits. This bottleneck can be relaxed with a longer timeout or additional CPU cores. We anticipate that a much longer delay will begin to tax the patience of most users, and so we consider 30s a reasonable upper bound for repair latency. As we will now show, end-to-end latency can be improved by doing rejection sampling, albeit at the cost of naturalness and sample efficiency.

%  In the following benchmark, we measure the Precision@k of our repair procedure against human repairs of varying edit distances and latency cutoffs, using an adaptive resampling procedure described in \S\ref{sec:adaptive}. This sampler maintains a buffer of successful repairs ranked by perplexity and uses stochastic local search to resample edits within a neighborhood. Initially, edits are sampled uniformly at random. Over time and as the admissible set grows, it prioritizes edits nearby low-perplexity repairs. This technique offers a significant advantage in the low-latency setting.

  \clearpage\subsection{Subcomponent ablation}\label{sec:rq3}

  Originally, we used an adaptive rejection-based sampler, which did not sample directly from the admissible set, but the entire Levenshtein ball, and then rejected invalid samples. Although rejection sampling has a much lower minimum latency threshold to return admissible repairs, i.e., a few seconds at most, the average time required to attain a desired precision on human repairs is much higher. We present the results from the rejection-based evaluation for comparison below.

  \begin{figure}[H]
    \resizebox{.24\textwidth}{!}{% This file was created with tikzplotlib v0.10.1.
\begin{tikzpicture}

\definecolor{darkgray176}{RGB}{176,176,176}
\definecolor{darkviolet1270255}{RGB}{127,0,255}
\definecolor{deepskyblue3176236}{RGB}{3,176,236}
\definecolor{dodgerblue45123246}{RGB}{45,123,246}
\definecolor{lightgray204}{RGB}{204,204,204}
\definecolor{royalblue8762253}{RGB}{87,62,253}

\begin{axis}[
legend cell align={left},
legend style={fill opacity=0.8, draw opacity=1, text opacity=1, draw=lightgray204, legend columns=-1, legend pos=north west},
tick align=outside,
tick pos=left,
title={\(\displaystyle \Delta\in[1,3]\) Repair Precision},
x grid style={darkgray176},
xlabel={Seconds},
xmin=-0.5925, xmax=8.5925,
ymin=0.0, ymax=0.7,
xtick style={color=black},
xtick={0,1,2,3,4,5,6,7,8},
xticklabels={0.1,0.3,0.5,0.7,0.9,1.1,1.3,1.5,1.7},
y grid style={darkgray176},
ylabel={Precision@k},
ytick style={color=black},
width=0.5\textwidth,
height=0.48\textwidth,
]
\addlegendimage{empty legend}
\addlegendentry{P@}
\draw[draw=none,fill=darkviolet1270255] (axis cs:-0.175,0) rectangle (axis cs:0.175,0.245);
\addlegendimage{ybar,ybar legend,draw=none,fill=darkviolet1270255}
\addlegendentry{All}

\draw[draw=none,fill=darkviolet1270255] (axis cs:0.825,0) rectangle (axis cs:1.175,0.417);
\draw[draw=none,fill=darkviolet1270255] (axis cs:1.825,0) rectangle (axis cs:2.175,0.496);
\draw[draw=none,fill=darkviolet1270255] (axis cs:2.825,0) rectangle (axis cs:3.175,0.535);
\draw[draw=none,fill=darkviolet1270255] (axis cs:3.825,0) rectangle (axis cs:4.175,0.562);
\draw[draw=none,fill=darkviolet1270255] (axis cs:4.825,0) rectangle (axis cs:5.175,0.584);
\draw[draw=none,fill=darkviolet1270255] (axis cs:5.825,0) rectangle (axis cs:6.175,0.603);
\draw[draw=none,fill=darkviolet1270255] (axis cs:6.825,0) rectangle (axis cs:7.175,0.619);
\draw[draw=none,fill=darkviolet1270255] (axis cs:7.825,0) rectangle (axis cs:8.175,0.629);
\draw[draw=none,fill=royalblue8762253] (axis cs:-0.175,0) rectangle (axis cs:0.175,0.234);
\addlegendimage{ybar,ybar legend,draw=none,fill=royalblue8762253}
\addlegendentry{10}

\draw[draw=none,fill=royalblue8762253] (axis cs:0.825,0) rectangle (axis cs:1.175,0.387);
\draw[draw=none,fill=royalblue8762253] (axis cs:1.825,0) rectangle (axis cs:2.175,0.451);
\draw[draw=none,fill=royalblue8762253] (axis cs:2.825,0) rectangle (axis cs:3.175,0.481);
\draw[draw=none,fill=royalblue8762253] (axis cs:3.825,0) rectangle (axis cs:4.175,0.502);
\draw[draw=none,fill=royalblue8762253] (axis cs:4.825,0) rectangle (axis cs:5.175,0.517);
\draw[draw=none,fill=royalblue8762253] (axis cs:5.825,0) rectangle (axis cs:6.175,0.53);
\draw[draw=none,fill=royalblue8762253] (axis cs:6.825,0) rectangle (axis cs:7.175,0.54);
\draw[draw=none,fill=royalblue8762253] (axis cs:7.825,0) rectangle (axis cs:8.175,0.546);
\draw[draw=none,fill=dodgerblue45123246] (axis cs:-0.175,0) rectangle (axis cs:0.175,0.218);
\addlegendimage{ybar,ybar legend,draw=none,fill=dodgerblue45123246}
\addlegendentry{5}

\draw[draw=none,fill=dodgerblue45123246] (axis cs:0.825,0) rectangle (axis cs:1.175,0.353);
\draw[draw=none,fill=dodgerblue45123246] (axis cs:1.825,0) rectangle (axis cs:2.175,0.406);
\draw[draw=none,fill=dodgerblue45123246] (axis cs:2.825,0) rectangle (axis cs:3.175,0.43);
\draw[draw=none,fill=dodgerblue45123246] (axis cs:3.825,0) rectangle (axis cs:4.175,0.444);
\draw[draw=none,fill=dodgerblue45123246] (axis cs:4.825,0) rectangle (axis cs:5.175,0.455);
\draw[draw=none,fill=dodgerblue45123246] (axis cs:5.825,0) rectangle (axis cs:6.175,0.467);
\draw[draw=none,fill=dodgerblue45123246] (axis cs:6.825,0) rectangle (axis cs:7.175,0.474);
\draw[draw=none,fill=dodgerblue45123246] (axis cs:7.825,0) rectangle (axis cs:8.175,0.479);
\draw[draw=none,fill=deepskyblue3176236] (axis cs:-0.175,0) rectangle (axis cs:0.175,0.117);
\addlegendimage{ybar,ybar legend,draw=none,fill=deepskyblue3176236}
\addlegendentry{1}

\draw[draw=none,fill=deepskyblue3176236] (axis cs:0.825,0) rectangle (axis cs:1.175,0.187);
\draw[draw=none,fill=deepskyblue3176236] (axis cs:1.825,0) rectangle (axis cs:2.175,0.219);
\draw[draw=none,fill=deepskyblue3176236] (axis cs:2.825,0) rectangle (axis cs:3.175,0.233);
\draw[draw=none,fill=deepskyblue3176236] (axis cs:3.825,0) rectangle (axis cs:4.175,0.244);
\draw[draw=none,fill=deepskyblue3176236] (axis cs:4.825,0) rectangle (axis cs:5.175,0.25);
\draw[draw=none,fill=deepskyblue3176236] (axis cs:5.825,0) rectangle (axis cs:6.175,0.258);
\draw[draw=none,fill=deepskyblue3176236] (axis cs:6.825,0) rectangle (axis cs:7.175,0.262);
\draw[draw=none,fill=deepskyblue3176236] (axis cs:7.825,0) rectangle (axis cs:8.175,0.265);
\addplot [red, thick] coordinates {(-0.8925,0.35) (14.8925,0.35)};
\end{axis}

\end{tikzpicture}
}
    \resizebox{.25\textwidth}{!}{% This file was created with tikzplotlib v0.10.1.
\begin{tikzpicture}

\definecolor{darkgray176}{RGB}{176,176,176}
\definecolor{darkviolet1270255}{RGB}{127,0,255}
\definecolor{deepskyblue3176236}{RGB}{3,176,236}
\definecolor{dodgerblue45123246}{RGB}{45,123,246}
\definecolor{lightgray204}{RGB}{204,204,204}
\definecolor{royalblue8762253}{RGB}{87,62,253}

\begin{axis}[
legend cell align={left},
legend style={fill opacity=0.8, draw opacity=1, text opacity=1, draw=lightgray204, legend columns=-1, legend pos=north west},
tick align=outside,
tick pos=left,
title={$\Delta=1$ Repair Precision},
x grid style={darkgray176},
xlabel={Seconds},
xmin=-0.5925, xmax=8.5925,
xtick style={color=black},
xtick={0,1,2,3,4,5,6,7,8},
xticklabels={0.1,0.3,0.5,0.7,0.9,1.1,1.3,1.5,1.7},
y grid style={darkgray176},
ymin=0, ymax=1.0,
ytick style={color=black}
width=0.43\textwidth,
height=0.49\textwidth,
]
\addlegendimage{empty legend}
\addlegendentry{P@}
\draw[draw=none,fill=darkviolet1270255] (axis cs:-0.175,0) rectangle (axis cs:0.175,0.416);
\addlegendimage{ybar,ybar legend,draw=none,fill=darkviolet1270255}
\addlegendentry{All}

\draw[draw=none,fill=darkviolet1270255] (axis cs:0.825,0) rectangle (axis cs:1.175,0.672);
\draw[draw=none,fill=darkviolet1270255] (axis cs:1.825,0) rectangle (axis cs:2.175,0.773);
\draw[draw=none,fill=darkviolet1270255] (axis cs:2.825,0) rectangle (axis cs:3.175,0.822);
\draw[draw=none,fill=darkviolet1270255] (axis cs:3.825,0) rectangle (axis cs:4.175,0.849);
\draw[draw=none,fill=darkviolet1270255] (axis cs:4.825,0) rectangle (axis cs:5.175,0.873);
\draw[draw=none,fill=darkviolet1270255] (axis cs:5.825,0) rectangle (axis cs:6.175,0.891);
\draw[draw=none,fill=darkviolet1270255] (axis cs:6.825,0) rectangle (axis cs:7.175,0.904);
\draw[draw=none,fill=darkviolet1270255] (axis cs:7.825,0) rectangle (axis cs:8.175,0.912);
\draw[draw=none,fill=royalblue8762253] (axis cs:-0.175,0) rectangle (axis cs:0.175,0.402);
\addlegendimage{ybar,ybar legend,draw=none,fill=royalblue8762253}
\addlegendentry{10}

\draw[draw=none,fill=royalblue8762253] (axis cs:0.825,0) rectangle (axis cs:1.175,0.645);
\draw[draw=none,fill=royalblue8762253] (axis cs:1.825,0) rectangle (axis cs:2.175,0.741);
\draw[draw=none,fill=royalblue8762253] (axis cs:2.825,0) rectangle (axis cs:3.175,0.786);
\draw[draw=none,fill=royalblue8762253] (axis cs:3.825,0) rectangle (axis cs:4.175,0.812);
\draw[draw=none,fill=royalblue8762253] (axis cs:4.825,0) rectangle (axis cs:5.175,0.834);
\draw[draw=none,fill=royalblue8762253] (axis cs:5.825,0) rectangle (axis cs:6.175,0.851);
\draw[draw=none,fill=royalblue8762253] (axis cs:6.825,0) rectangle (axis cs:7.175,0.862);
\draw[draw=none,fill=royalblue8762253] (axis cs:7.825,0) rectangle (axis cs:8.175,0.87);
\draw[draw=none,fill=dodgerblue45123246] (axis cs:-0.175,0) rectangle (axis cs:0.175,0.379);
\addlegendimage{ybar,ybar legend,draw=none,fill=dodgerblue45123246}
\addlegendentry{5}

\draw[draw=none,fill=dodgerblue45123246] (axis cs:0.825,0) rectangle (axis cs:1.175,0.601);
\draw[draw=none,fill=dodgerblue45123246] (axis cs:1.825,0) rectangle (axis cs:2.175,0.686);
\draw[draw=none,fill=dodgerblue45123246] (axis cs:2.825,0) rectangle (axis cs:3.175,0.724);
\draw[draw=none,fill=dodgerblue45123246] (axis cs:3.825,0) rectangle (axis cs:4.175,0.743);
\draw[draw=none,fill=dodgerblue45123246] (axis cs:4.825,0) rectangle (axis cs:5.175,0.762);
\draw[draw=none,fill=dodgerblue45123246] (axis cs:5.825,0) rectangle (axis cs:6.175,0.777);
\draw[draw=none,fill=dodgerblue45123246] (axis cs:6.825,0) rectangle (axis cs:7.175,0.785);
\draw[draw=none,fill=dodgerblue45123246] (axis cs:7.825,0) rectangle (axis cs:8.175,0.791);
\draw[draw=none,fill=deepskyblue3176236] (axis cs:-0.175,0) rectangle (axis cs:0.175,0.211);
\addlegendimage{ybar,ybar legend,draw=none,fill=deepskyblue3176236}
\addlegendentry{1}

\draw[draw=none,fill=deepskyblue3176236] (axis cs:0.825,0) rectangle (axis cs:1.175,0.327);
\draw[draw=none,fill=deepskyblue3176236] (axis cs:1.825,0) rectangle (axis cs:2.175,0.378);
\draw[draw=none,fill=deepskyblue3176236] (axis cs:2.825,0) rectangle (axis cs:3.175,0.401);
\draw[draw=none,fill=deepskyblue3176236] (axis cs:3.825,0) rectangle (axis cs:4.175,0.418);
\draw[draw=none,fill=deepskyblue3176236] (axis cs:4.825,0) rectangle (axis cs:5.175,0.428);
\draw[draw=none,fill=deepskyblue3176236] (axis cs:5.825,0) rectangle (axis cs:6.175,0.437);
\draw[draw=none,fill=deepskyblue3176236] (axis cs:6.825,0) rectangle (axis cs:7.175,0.442);
\draw[draw=none,fill=deepskyblue3176236] (axis cs:7.825,0) rectangle (axis cs:8.175,0.446);
\addplot [red, thick] coordinates {(-0.8925,0.45) (14.8925,0.45)};
\end{axis}

\end{tikzpicture}
}
    \resizebox{.24\textwidth}{!}{% This file was created with tikzplotlib v0.10.1.
\begin{tikzpicture}

\definecolor{darkgray176}{RGB}{176,176,176}
\definecolor{darkviolet1270255}{RGB}{127,0,255}
\definecolor{deepskyblue3176236}{RGB}{3,176,236}
\definecolor{dodgerblue45123246}{RGB}{45,123,246}
\definecolor{lightgray204}{RGB}{204,204,204}
\definecolor{royalblue8762253}{RGB}{87,62,253}

\begin{axis}[
legend cell align={left},
legend style={fill opacity=0.8, draw opacity=1, text opacity=1, draw=lightgray204, legend columns=-1, legend pos=north west},
tick align=outside,
tick pos=left,
axis lines*=left,
title={$\Delta=2$ Repair Precision},
x grid style={darkgray176},
xlabel={Seconds},
xmin=-0.5925, xmax=8.5925,
xtick style={color=black},
xtick={0,1,2,3,4,5,6,7,8},
xticklabels={0.1,0.3,0.5,0.7,0.9,1.1,1.3,1.5,1.7},
y grid style={darkgray176},
ymin=0, ymax=0.3843,
ytick style={color=black},
width=0.55\textwidth,
height=0.49\textwidth,
]
\addlegendimage{empty legend}
\addlegendentry{P@}
\draw[draw=none,fill=darkviolet1270255] (axis cs:-0.175,0) rectangle (axis cs:0.175,0.067);
\addlegendimage{ybar,ybar legend,draw=none,fill=darkviolet1270255}
\addlegendentry{All}

\draw[draw=none,fill=darkviolet1270255] (axis cs:0.825,0) rectangle (axis cs:1.175,0.165);
\draw[draw=none,fill=darkviolet1270255] (axis cs:1.825,0) rectangle (axis cs:2.175,0.225);
\draw[draw=none,fill=darkviolet1270255] (axis cs:2.825,0) rectangle (axis cs:3.175,0.261);
\draw[draw=none,fill=darkviolet1270255] (axis cs:3.825,0) rectangle (axis cs:4.175,0.292);
\draw[draw=none,fill=darkviolet1270255] (axis cs:4.825,0) rectangle (axis cs:5.175,0.309);
\draw[draw=none,fill=darkviolet1270255] (axis cs:5.825,0) rectangle (axis cs:6.175,0.335);
\draw[draw=none,fill=darkviolet1270255] (axis cs:6.825,0) rectangle (axis cs:7.175,0.353);
\draw[draw=none,fill=darkviolet1270255] (axis cs:7.825,0) rectangle (axis cs:8.175,0.366);
\draw[draw=none,fill=royalblue8762253] (axis cs:-0.175,0) rectangle (axis cs:0.175,0.06);
\addlegendimage{ybar,ybar legend,draw=none,fill=royalblue8762253}
\addlegendentry{10}

\draw[draw=none,fill=royalblue8762253] (axis cs:0.825,0) rectangle (axis cs:1.175,0.128);
\draw[draw=none,fill=royalblue8762253] (axis cs:1.825,0) rectangle (axis cs:2.175,0.164);
\draw[draw=none,fill=royalblue8762253] (axis cs:2.825,0) rectangle (axis cs:3.175,0.182);
\draw[draw=none,fill=royalblue8762253] (axis cs:3.825,0) rectangle (axis cs:4.175,0.2);
\draw[draw=none,fill=royalblue8762253] (axis cs:4.825,0) rectangle (axis cs:5.175,0.208);
\draw[draw=none,fill=royalblue8762253] (axis cs:5.825,0) rectangle (axis cs:6.175,0.221);
\draw[draw=none,fill=royalblue8762253] (axis cs:6.825,0) rectangle (axis cs:7.175,0.23);
\draw[draw=none,fill=royalblue8762253] (axis cs:7.825,0) rectangle (axis cs:8.175,0.238);
\draw[draw=none,fill=dodgerblue45123246] (axis cs:-0.175,0) rectangle (axis cs:0.175,0.05);
\addlegendimage{ybar,ybar legend,draw=none,fill=dodgerblue45123246}
\addlegendentry{5}

\draw[draw=none,fill=dodgerblue45123246] (axis cs:0.825,0) rectangle (axis cs:1.175,0.1);
\draw[draw=none,fill=dodgerblue45123246] (axis cs:1.825,0) rectangle (axis cs:2.175,0.123);
\draw[draw=none,fill=dodgerblue45123246] (axis cs:2.825,0) rectangle (axis cs:3.175,0.132);
\draw[draw=none,fill=dodgerblue45123246] (axis cs:3.825,0) rectangle (axis cs:4.175,0.142);
\draw[draw=none,fill=dodgerblue45123246] (axis cs:4.825,0) rectangle (axis cs:5.175,0.144);
\draw[draw=none,fill=dodgerblue45123246] (axis cs:5.825,0) rectangle (axis cs:6.175,0.155);
\draw[draw=none,fill=dodgerblue45123246] (axis cs:6.825,0) rectangle (axis cs:7.175,0.163);
\draw[draw=none,fill=dodgerblue45123246] (axis cs:7.825,0) rectangle (axis cs:8.175,0.167);
\draw[draw=none,fill=deepskyblue3176236] (axis cs:-0.175,0) rectangle (axis cs:0.175,0.016);
\addlegendimage{ybar,ybar legend,draw=none,fill=deepskyblue3176236}
\addlegendentry{1}

\draw[draw=none,fill=deepskyblue3176236] (axis cs:0.825,0) rectangle (axis cs:1.175,0.038);
\draw[draw=none,fill=deepskyblue3176236] (axis cs:1.825,0) rectangle (axis cs:2.175,0.052);
\draw[draw=none,fill=deepskyblue3176236] (axis cs:2.825,0) rectangle (axis cs:3.175,0.057);
\draw[draw=none,fill=deepskyblue3176236] (axis cs:3.825,0) rectangle (axis cs:4.175,0.061);
\draw[draw=none,fill=deepskyblue3176236] (axis cs:4.825,0) rectangle (axis cs:5.175,0.064);
\draw[draw=none,fill=deepskyblue3176236] (axis cs:5.825,0) rectangle (axis cs:6.175,0.072);
\draw[draw=none,fill=deepskyblue3176236] (axis cs:6.825,0) rectangle (axis cs:7.175,0.076);
\draw[draw=none,fill=deepskyblue3176236] (axis cs:7.825,0) rectangle (axis cs:8.175,0.078);
\addplot [red, thick] coordinates {(-0.8925,0.16) (14.8925,0.16)};
\end{axis}

\end{tikzpicture}
}
    \resizebox{.24\textwidth}{!}{% This file was created with tikzplotlib v0.10.1.
\begin{tikzpicture}

\definecolor{darkgray176}{RGB}{176,176,176}
\definecolor{darkviolet1270255}{RGB}{127,0,255}
\definecolor{deepskyblue3176236}{RGB}{3,176,236}
\definecolor{dodgerblue45123246}{RGB}{45,123,246}
\definecolor{lightgray204}{RGB}{204,204,204}
\definecolor{royalblue8762253}{RGB}{87,62,253}

\begin{axis}[
legend cell align={left},
legend style={fill opacity=0.8, draw opacity=1, text opacity=1, draw=lightgray204, legend columns=-1, legend pos=north west},
tick align=outside,
tick pos=left,
axis lines*=left,
title={$\Delta=3$ Repair Precision},
x grid style={darkgray176},
xlabel={Seconds},
xmin=-0.5925, xmax=8.5925,
xtick style={color=black},
xtick={0,1,2,3,4,5,6,7,8},
xticklabels={0.1,0.3,0.5,0.7,0.9,1.1,1.3,1.5,1.7},
yticklabels={0, 0.05, 0.10, 0.15, 0.20},
y grid style={darkgray176},
ymin=0, ymax=0.15225,
ytick style={color=black},
width=0.53\textwidth,
height=0.48\textwidth,
]
\addlegendimage{empty legend}
\addlegendentry{P@}
\draw[draw=none,fill=darkviolet1270255] (axis cs:-0.175,0) rectangle (axis cs:0.175,0.01);
\addlegendimage{ybar,ybar legend,draw=none,fill=darkviolet1270255}
\addlegendentry{All}

\draw[draw=none,fill=darkviolet1270255] (axis cs:0.825,0) rectangle (axis cs:1.175,0.035);
\draw[draw=none,fill=darkviolet1270255] (axis cs:1.825,0) rectangle (axis cs:2.175,0.064);
\draw[draw=none,fill=darkviolet1270255] (axis cs:2.825,0) rectangle (axis cs:3.175,0.072);
\draw[draw=none,fill=darkviolet1270255] (axis cs:3.825,0) rectangle (axis cs:4.175,0.09);
\draw[draw=none,fill=darkviolet1270255] (axis cs:4.825,0) rectangle (axis cs:5.175,0.109);
\draw[draw=none,fill=darkviolet1270255] (axis cs:5.825,0) rectangle (axis cs:6.175,0.116);
\draw[draw=none,fill=darkviolet1270255] (axis cs:6.825,0) rectangle (axis cs:7.175,0.133);
\draw[draw=none,fill=darkviolet1270255] (axis cs:7.825,0) rectangle (axis cs:8.175,0.145);
\draw[draw=none,fill=royalblue8762253] (axis cs:-0.175,0) rectangle (axis cs:0.175,0.003);
\addlegendimage{ybar,ybar legend,draw=none,fill=royalblue8762253}
\addlegendentry{10}

\draw[draw=none,fill=royalblue8762253] (axis cs:0.825,0) rectangle (axis cs:1.175,0.01);
\draw[draw=none,fill=royalblue8762253] (axis cs:1.825,0) rectangle (axis cs:2.175,0.01);
\draw[draw=none,fill=royalblue8762253] (axis cs:2.825,0) rectangle (axis cs:3.175,0.014);
\draw[draw=none,fill=royalblue8762253] (axis cs:3.825,0) rectangle (axis cs:4.175,0.018);
\draw[draw=none,fill=royalblue8762253] (axis cs:4.825,0) rectangle (axis cs:5.175,0.019);
\draw[draw=none,fill=royalblue8762253] (axis cs:5.825,0) rectangle (axis cs:6.175,0.018);
\draw[draw=none,fill=royalblue8762253] (axis cs:6.825,0) rectangle (axis cs:7.175,0.023);
\draw[draw=none,fill=royalblue8762253] (axis cs:7.825,0) rectangle (axis cs:8.175,0.023);
\draw[draw=none,fill=dodgerblue45123246] (axis cs:-0.175,0) rectangle (axis cs:0.175,0.002);
\addlegendimage{ybar,ybar legend,draw=none,fill=dodgerblue45123246}
\addlegendentry{5}

\draw[draw=none,fill=dodgerblue45123246] (axis cs:0.825,0) rectangle (axis cs:1.175,0.005);
\draw[draw=none,fill=dodgerblue45123246] (axis cs:1.825,0) rectangle (axis cs:2.175,0.008);
\draw[draw=none,fill=dodgerblue45123246] (axis cs:2.825,0) rectangle (axis cs:3.175,0.01);
\draw[draw=none,fill=dodgerblue45123246] (axis cs:3.825,0) rectangle (axis cs:4.175,0.011);
\draw[draw=none,fill=dodgerblue45123246] (axis cs:4.825,0) rectangle (axis cs:5.175,0.011);
\draw[draw=none,fill=dodgerblue45123246] (axis cs:5.825,0) rectangle (axis cs:6.175,0.011);
\draw[draw=none,fill=dodgerblue45123246] (axis cs:6.825,0) rectangle (axis cs:7.175,0.014);
\draw[draw=none,fill=dodgerblue45123246] (axis cs:7.825,0) rectangle (axis cs:8.175,0.018);
\draw[draw=none,fill=deepskyblue3176236] (axis cs:-0.175,0) rectangle (axis cs:0.175,0);
\addlegendimage{ybar,ybar legend,draw=none,fill=deepskyblue3176236}
\addlegendentry{1}

\draw[draw=none,fill=deepskyblue3176236] (axis cs:0.825,0) rectangle (axis cs:1.175,0.003);
\draw[draw=none,fill=deepskyblue3176236] (axis cs:1.825,0) rectangle (axis cs:2.175,0.008);
\draw[draw=none,fill=deepskyblue3176236] (axis cs:2.825,0) rectangle (axis cs:3.175,0.008);
\draw[draw=none,fill=deepskyblue3176236] (axis cs:3.825,0) rectangle (axis cs:4.175,0.011);
\draw[draw=none,fill=deepskyblue3176236] (axis cs:4.825,0) rectangle (axis cs:5.175,0.011);
\draw[draw=none,fill=deepskyblue3176236] (axis cs:5.825,0) rectangle (axis cs:6.175,0.011);
\draw[draw=none,fill=deepskyblue3176236] (axis cs:6.825,0) rectangle (axis cs:7.175,0.013);
\draw[draw=none,fill=deepskyblue3176236] (axis cs:7.825,0) rectangle (axis cs:8.175,0.013);
\addplot [red, thick] coordinates {(-0.8925,0.096) (14.8925,0.096)};
\addplot [orange, thick] coordinates {(-0.8925,0.02) (14.8925,0.02)};
\end{axis}

\end{tikzpicture}
}
    \caption{Adaptive sampling repairs. The red line indicates Seq2Parse Precision@1, and the orange indicates BIFI's precision at single-shot repair, all three of which were evaluated on the exact same repairs.}\label{fig:adaptive}
  \end{figure}

  We also evaluate Seq2Parse on the same dataset. Seq2Parse only supports Precision@1 repairs, and so we only report Seq2parse Precision@1 from the StackOverflow benchmark for comparison. Unlike our approach which only produces syntactically correct repairs, Seq2Parse and BIFI also produce syntactically incorrect repairs in practice. The overall latency of Seq2Parse varies depending on the length of the repair, averaging 1.5s for $\Delta=1$ to 2.7s for $\Delta=3$, across the entire StackOverflow dataset, while BIFI consistently achieves subsecond latency across all repairs and distances.

  Next, we conduct an ablation study to compare the effectiveness of PCFG sampling versus enumeration. In both experiments, we balance the StackOverflow dataset across edit distances and run the repair sampler for up to 30 seconds (in either enumerative or PCFG mode), then rerank all repairs by n-gram perplexity and measure the Precision@1 for sampling, enumeration, and the default hybrid approach, which uses enumeration if the intersection grammar size $|G_\cap'|$ contains fewer than 10,000 total productions, and PCFG sampling otherwise.

  \begin{figure}[h]
    \begin{tikzpicture}[scale=0.9]
  \begin{axis}[
  ybar,
  xlabel={$\Delta(\err\sigma, \sigma)$},
  ylabel={Precision@1},
  title={Enumeration Only},
  axis x line*=bottom,
  axis y line*=left,
  ymin=0,
  ymax=1,
  xtick=data,
  xticklabels={1, 2, 3, 4},
  width=5cm,
  height=4cm,
  enlarge x limits=0.15,
  legend style={at={(0.5,-0.15)},
  anchor=north,legend columns=-1},
  ]
  \addplot[fill=black!30] table[x=Lev, y=P@1] {
    Lev P@1
    1 0.98
    2 0.61
    3 0.35
    4 0.29
  };
  \end{axis}
\end{tikzpicture}
    \begin{tikzpicture}[scale=0.9]
  \begin{axis}[
  ybar,
  xlabel={$\Delta(\err\sigma, \sigma)$},
  ylabel={Precision@1},
  title={PCFG sampling (Alg. 1)},
  axis x line*=bottom,
  axis y line*=left,
  ymin=0,
  ymax=1,
  xtick=data,
  width=5cm,
  height=4cm,
  xticklabels={1, 2, 3, 4},
  enlarge x limits=0.15,
  legend style={at={(0.5,-0.15)},
  anchor=north,legend columns=-1},
  ]
  \addplot[fill=black!30] table[x=Lev, y=P@1] {
    Lev P@1
    1 0.33
    2 0.15
    3 0.08
    4 0.04
  };
  \end{axis}
\end{tikzpicture}
    \begin{tikzpicture}[scale=0.9]
  \begin{axis}[
  ybar,
  xlabel={$\Delta(\err\sigma, \sigma)$},
  ylabel={Precision@1},
  title={Hybrid},
  axis x line*=bottom,
  axis y line*=left,
  ymin=0,
  ymax=1,
  xtick=data,
  xticklabels={1, 2, 3, 4},
  width=5cm,
  height=4cm,
  enlarge x limits=0.15,
  legend style={at={(0.5,-0.15)},
  anchor=north,legend columns=-1},
  ]
  \addplot table[x=Lev, y=P@1] {
    Lev P@1
    1 0.98
    2 0.66
    3 0.42
    4 0.35
  };
  \end{axis}
\end{tikzpicture}
  \end{figure}

  While the overall precision is notably lower for PCFG sampling than enumeration, the average number of samples drawn is also significantly lower, indicating a higher sample efficiency. This illustrates the tradeoff between sample efficiency and diversity, and suggests a hybrid approach is most effective at retrieval across density regimes. When the CFL is very large, uniform sampling is unlikely to retrieve the most natural samples, and top-down PCFG sampling offers a more informed prior, albeit at the cost of lower coverage on small languages.

  In general, enumeration has an advantage when the CFL is very small. For example, if the CFL contains 2,000 sentences, enumeration is likely to recover all 2,000 before timeout, whereas PCFG sampling may only recover 100 of the most probable samples depending on the transition probabilities. However, if the CFL has 200,000 sentences, enumeration may only recover 10,000 uniform random samples in the time allotted and the PCFG sampler only 5,000, but due to its higher sample efficiency, is likely to contain a much higher proportion of natural repairs.

  \begin{wrapfigure}{r}{0.4\textwidth}
    \begin{tikzpicture}[scale=0.47]
  \begin{axis}[
    legend cell align={left},
    legend style={fill opacity=0.8, draw opacity=1, text opacity=1, legend columns=-1, legend pos=north west},
    ybar,
    enlargelimits=0.15,
    legend style={at={(0.03,0.97)},anchor=north west},
    title={\textbf{Relative Total Distinct Solutions Found vs. Single Core}},
    x=1cm,
    axis lines*=left,
    ylabel={Relative improvement},
    xlabel={Number of assigned cores},
    symbolic x coords={2,3,4,5,6,7,8,9,10},
    xtick=data,
    bar width=4pt,
  ]
    \addlegendimage{empty legend}
    \addlegendentry{$\Delta(\err\sigma, \sigma')=$}
    \addlegendentry{\footnotesize{$1,$}}
    \addlegendentry{\footnotesize{$2,$}}
    \addlegendentry{\footnotesize{$3,$}}
    \addlegendentry{\footnotesize{$4$}}
    \addplot[green, fill=green!50] coordinates { (2,0.1343) (3,0.1955) (4,0.2249) (5,0.2475) (6,0.2760) (7,0.2994) (8,0.3073) (9,0.3151) (10,0.3229) };
    \addplot[blue, fill=blue!50] coordinates { (2,0.2655) (3,0.4353) (4,0.5614) (5,0.6644) (6,0.7462) (7,0.7783) (8,0.8347) (9,0.8005) (10,0.7798) };
    \addplot[orange, fill=orange!50] coordinates { (2,0.3972) (3,0.6928) (4,0.9327) (5,1.1834) (6,1.3138) (7,1.3988) (8,1.6039) (9,1.5500) (10,1.5691) };
    \addplot[red, fill=red!50] coordinates { (2,0.4863) (3,0.8326) (4,1.1368) (5,1.3879) (6,1.5873) (7,1.7494) (8,1.8802) (9,1.9059) (10,1.9625) };
  \end{axis}
\end{tikzpicture}
    \caption{Observed improvement in throughput relative to total CPU cores assigned.}
    \label{fig:speedup}
  \end{wrapfigure}

  We also measure the relative improvement in throughput (measured by the number of distinct repairs found after 30s) as a function of the number of additional CPU cores, averaged across 1000 trials. We observe from Fig.\ref{fig:speedup} the relative throughput increases logarithmically with the number of additional CPU cores, with at least four CPU cores needed to offset the parallelization overhead. Generally, increasing parallelism only helps when the size of the admissible set is large enough to absorb the additional computation, which is seldom the case for small-radii Levenshtein balls. Further speedups may be possible to realize by rewriting the sampler in CUDA, which we leave for future work.

  \section{Discussion}\label{sec:discussion}

  The main lesson we draw from our experiments is that it is possible to leverage compute to compete with large language models on practical program repair tasks. Though sample-efficient, their size comes at the cost of expensive training, and domain adaptation requires fine-tuning or retraining on pairwise repairs. Our approach uses a small grammar and a relatively cheap ranking metric to achieve significantly higher precision. This allows us to repair errors in languages with little to no training data and provides far more flexibility and controllability.

  Our primary insight leading to state-of-the-art precision is that repairs are typically concentrated near the center of a small Levenshtein ball, and by enumerating or sampling it carefully, then reranking all repairs found we can achieve a significant improvement over one-shot neural repair. This is especially true for small-radii Levenshtein balls, where the admissible set is small enough to be enumerated completely and ranked. For larger radii, we can still achieve competitive precision by using a PCFG to sample from the admissible set and reranking by perplexity.

   Unexpectedly, we find that Precision@1 of our method is competitive with BIFI's Precision@20k, while requiring only a fraction of the data and compute. This is likely because BIFI's training set does not cover the full space of syntactically valid repairs. As Tidyparse uses its own grammar, it can sample from the language directly, and does not require training distribution to suggest valid repairs, only to rank them by naturalness. The emphasis on completeness is especially useful for discovering small repairs, which may be overlooked by neural models.

   Although latency and precision are ultimately the deciding usability factors, repair throughput is a crucial intermediate factor to consider when evaluating the performance of a repair system. Even with a perfectly accurate scoring function, if the correct repair is never retrieved, it will be for naught. By maximizing the total number of unique valid repairs, we increase the likelihood of retrieving natural repairs to give the scoring function the best chance of ranking them successfully. For this reason, we prioritize throughput heavily in our design (Def.~\ref{def:linear-convergence}) and evaluation (Fig.~\ref{fig:throughput}).

   Rejection sampling can be a useful technique for quickly retrieving a subset of valid repairs, but has the disadvantage of converging very slowly, requiring far too many samples to achieve competitive precision on natural repairs. One avenue may be to use rejection sampling to find probable edit locations, then switch to an exhaustive method to retrieve all repairs in that region. This approach however, would not offer the same completeness guarantees as language intersection.

  \subsection{Limitations and future work}

%  We identify three broad categories of limitations in evaluating Tidyparse and suggest directions for future work: naturalness, complexity, and semantics.

  \subsubsection{Naturalness}

  Firstly, Tidyparse does not currently support intersections between weighted CFGs and weighted finite automata, a la Pasti et al.~\cite{pasti2023intersection}. This feature would allow us to put transition probabilities on the Levenshtein automaton corresponding to more likely edits and construct a weighted intersection grammar. It would be interesting to train a weighted intersection grammar and compare the precision versus vanilla PCFG sampling. We also hope to explore more incremental sampling strategies such as sequential Monte-Carlo~\cite{lew2023sequential}.

  The scoring function is currently computed over lexical tokens. We expect that a more precise scoring function could be constructed by splicing candidate repairs back into the original source code and then scoring plaintext, however this would require special handling for insertions and substitutions of names, numbers and identifiers that were absent from the original source code. For this reason, we currently perform the scoring in lexical space, which discards a useful signal, but even this coarse approximation is sufficient to achieve state-of-the-art precision.

  Furthermore, the scoring function only considers each candidate repair $P_\theta(\sigma')$ in isolation, returning the most plausible candidate independent of the original error. One way to improve this would be to incorporate the broken sequence ($\err\sigma$), parser error message ($m$), original source ($s$), and possibly other contextual priors to inform the scoring function. This would require a more expressive probabilistic language model to faithfully model the joint distribution $P_\theta(\sigma' \mid \err\sigma, m, s, \ldots)$, but would significantly improve the precision of the generated repairs.

  \subsubsection{Complexity}

  Latency can vary depending on several factors including string length, grammar size, and critically the Levenshtein edit distance. This can be an advantage because in the absence of any contextual or statistical information, syntax and Levenshtein edits are often sufficiently constrained to identify a small number of valid repairs. It is also a limitation, because as the number of edits grows, the admissible set expands rapidly and the number of valid repairs quickly becomes too large to be useful without a very precise naturalness metric to distinguish equidistant repairs.

  Space complexity increases sharply with edit distance and to a lesser extent with length. This can be partly alleviated with more precise criteria to avoid creating superfluous productions, but the memory overhead is still considerable. Memory pressure can be attributed to engineering factors such as the grammar encoding, but is also an inherent challenge of grammar intersection. Managing the size of the intersection grammar during construction is a critical factor in scaling up this technique. Weighted intersections would also help to limit the size of the intersection grammar, as it would enable us to prune improbable productions without diminishing precision.

  \subsubsection{Toolchain integration}

  Lastly and perhaps most significantly, Tidyparse does not incorporate any semantic constraints, so its repairs whilst syntactically admissible, are not guaranteed to be semantically valid. This can partly be alleviated by filtering the results through an incremental compiler or linter, however, the latency necessary to check every repair may be non-negligible. It may also be possible to add a type-based semantic refinement to our language intersection, however this would require a more expressive grammatical formalism than CFGs can naturally provide.

  Slicing is an important preprocessing consideration which has so far gone unmentioned. The current implementation expects pre-sliced code fragments, however in a more practical scenario, it would be necessary to leverage information from the programming environment to identify the boundaries of the code fragment to be repaired. This could be done with careful integration with the development environment or via ad hoc slicing. We leave this aspect for future work.

%  We envision a few primary use cases for our tool: (1) helping novice programmers become more quickly familiar with a new programming language, (2) autocorrecting common typos among proficient but forgetful programmers, (3) as a prototyping tool for PL designers and educators, and (4) as a pluggable library or service for parser-generators and language servers.

  \clearpage\section{Related Work}\label{sec:related}

  Three important questions arise when repairing syntax errors: (1) is the program broken in the first place? (2) if so, where are the errors located? (3) how should those locations then be altered? Those questions are addressed by three theoretical areas, (1)~parsing, (2)~language equations and (3)~syntax repair. We survey each of those areas, then turn our attention to more engineering-oriented research, including (4) string solving, (5) error-correction, (6) decoding and finally (7) neural program repair.

  \subsection{Parsing}

  Context-free language (CFL) parsing is the well-studied problem of how to turn a string into a unique tree, with many different algorithms and implementations (e.g., shift-reduce, recursive-descent, LR). Many of those algorithms expect grammars to be expressed in a certain form (e.g., left- or right- recursive) or are optimized for a narrow class of grammars (e.g., regular, linear).

  General CFL parsing allows ambiguity (non-unique trees) and can be formulated as a dynamic programming problem, as shown by Cocke-Younger-Kasami (CYK)~\cite{sakai1961syntax}, Earley~\cite{earley1970efficient} and others. These parsers have roughly cubic complexity with respect to the length of the input string.

  As shown by Valiant~\cite{valiant1975general}, Lee~\cite{lee2002fast} and others, general CFL recognition is in some sense equivalent to binary matrix multiplication, another well-studied combinatorial problem with broad applications, known to be at worst subcubic. This reduction opens the door to a range of complexity-theoretic speedups to CFL recognition, however large constants tend to limit their practical utility.

%  Okhotin~\cite{okhotin2001conjunctive} extends CFGs with language conjunction in \textit{conjunctive grammars}, followed by Zhang \& Su~\cite{zhang2017context} who apply conjunctive language reachability to dataflow analysis.

  \subsection{Language equations}

  Language equations are a powerful tool for reasoning about formal languages and their inhabitants. First proposed by Ginsburg et al.~\cite{ginsburg1962two} for the ALGOL language, language equations are essentially systems of inequalities with variables representing \textit{holes}, i.e., unknown values, in the language or grammar. Solutions to these equations can be obtained using various fixpoint techniques, yielding members of the language. This insight reveals the true algebraic nature of CFLs and their cousins.

  Being an algebraic formalism, language equations naturally give rise to a kind of calculus, vaguely reminiscent of Leibniz' and Newton's. First studied by Brzozowski~\cite{brzozowski1964derivatives, brzozowski1980equations} and Antimirov~\cite{antimirov1996partial}, one can take the derivative of a language equation, yielding another equation. This can be interpreted as a kind of continuation or language quotient, revealing the suffixes that complete a given prefix. This technique leads to an elegant family of algorithms for incremental parsing~\cite{might2011parsing, adams2016complexity} and automata minimization~\cite{brzozowski1962canonical}. In our setting, differentiation corresponds to code completion.

  Bar-Hillel~\cite{bar1961formal} establishes the closure of CFLs under intersection with regular languages, but does not elaborate on how to construct the corresponding grammar in order to recognize it. Beigel~\cite{beigelproof} and Pasti et al.~\cite{pasti2023intersection} provide helpful insights into the construction of the intersection grammar, and Nederhof and Satta~\cite{nederhof2004language} specifically consider finite CFL intersections, but neither considers Levenshtein intersections. Our work specializes Bar-Hillel intersections to Levenshtein automata.

  More concretely, we restrict our attention to language equations over CFLs, whose variables coincide with edit locations in the source code of a computer program, and solutions correspond to syntax repairs. While prior work has studied the use of language equations for parsing~\cite{might2011parsing}, to our knowledge they were never specifically applied to code completion or syntax error correction.

  \subsection{Syntax repair}

  In finite languages, syntax repair corresponds to spelling correction, a more restrictive and largely solved problem. Schulz and Stoyan~\cite{schulz2002fast} construct a finite automaton that returns the nearest dictionary entry by Levenshtein edit distance. Though considerably simpler than syntax correction, their work shares similar challenges and offers insights for handling more general repair scenarios.

  When a sentence is grammatically invalid, parsing grows more challenging. Like spelling, the problem is to find the minimum number of edits required to transform an arbitrary string into a syntactically valid one, where validity is defined as containment in a (typically) context-free language. Early work, including Irons~\cite{irons1963error} and Aho~\cite{aho1972minimum} propose a dynamic programming algorithm to compute the minimum number of edits required to fix an invalid string. Prior work on error correcting parsing only considers the shortest edit(s), and does not study multiple edits over the Levenshtein ball. Furthermore, the problem of actually generating the repairs is not well-posed, as there are usually many valid strings that can be obtained within a given number of edits. We instead focus on bounded Levenshtein reachability, which is the problem of finding useful repairs within a fixed Levenshtein distance of the broken string, which requires language intersection.

%Bar-Hillel~\cite{bar1961formal} establishes the closure of CFLs under intersection with regular languages, and can be used to construct the corresponding minimum-cost CFG in a straightforward manner. However while generally quite expressive, CFLs are themselves not closed under intersection and have other practical limitations, i.e., are unable to express indentation or variable binding. These limitations motivate us towards more expressive yet still efficiently parsable formalisms.

%Conveniently, Okhotin~\cite{okhotin2001conjunctive} offers exactly the formalism needed: language equations augmented with logical operators like conjunction or disjunction. These operators afford us the flexibility to encode both language union and intersection, which are difficult or impossible to express using a single grammar, as well as incorporate extra-syntactic side-constraints during solving.

%Naively, one might think to enumerate all edits within a certain Levenshtein radius and reject invalid strings. However, this approach is intractable for long strings. Another approach is to compute the \textit{edit distance} between the string and the language as proposed by Aho. This approach is also intractable depending on the size of the grammar. A third approach is to use a language equation to compute the intersection between the Levenshtein hypersphere and the language.

%The literature on parsers is vast and deep, covering far more ground than we could possibly hope to survey. We take inspiration from their findings, but restrict ourselves to a dozen most closely related papers, in four major research areas: (1) formal language theory, (2) constraint satisfaction (3) program synthesis, and (4) error correction. We survey each of these areas in turn.
%
%It was Noam Chomsky who first developed the algebraic theory of \textit{context-free grammars} (CFGs) in 1959~\cite{chomsky1959algebraic}, and since then, CFGs have been the subject of extensive research in formal language theory and program analysis. In particular, we take inspiration from Leslie Valiant~\cite{valiant1975general} who first discovered the connection to matrix multiplication in 1975 and Alexander Okhotin~\cite{okhotin2001conjunctive} who later introduced the idea of \textit{conjunctive grammars} in 2001. More recently, Azimov \& Grigorev~\cite{azimov2018cfpq} and Qirun Zhang shed light into the practical applicability of these ideas and made the theory of CFG reachability more accessible to the general public. In our work, we show how to compile their ideas onto a SAT solver to fix syntax errors, which seems like a perfectly natural extension, but was heretofore previously never considered to the best of our knowledge.

  \subsection{String solving}

  There is related work on string constraints in the constraint programing literature, featuring solvers like CFGAnalyzer and HAMPI~\cite{kiezun2009hampi}, which consider bounded context free grammars and intersections thereof. Boja{\'n}czyk et al. (2014)~\cite{bojanczyk2014automata} introduce the theory of nominal automata. Around the same time, D'Antoni et al. (2014) introduce \textit{symbolic automata}~\cite{dantoni2014minimization}, a generalization of finite automata which allow infinite alphabets and symbolic expressions over them. Hague et al. (2024)~\cite{hague2024parikh} use Parikh's theorem in the context of symbolic automata to speed up string constraint solving, from which we draw partial inspiration for the Levenshtein-Bar-Hillel construction in \S~\ref{sec:lev_bh}. In none of the constraint programming literature we surveyed do any of the approaches specifically consider the problem of syntax error correction, which is the main focus of our work.

  \subsection{Error correcting codes}

  Our work focuses on errors arising from human factors in computer programming, in particular \textit{syntax error correction}, which is the problem of fixing partially corrupted programs. Modern research on error correction, however, can be traced back to the early days of coding theory when researchers designed \textit{error-correcting codes} (ECCs) to denoise transmission errors induced by external interference, e.g., collision with a high-energy proton, manipulation by an adversary or even typographical mistake. In this context, \textit{code} can be any logical representation for communicating information between two parties (such as a human and a computer), and an ECC is a carefully-designed scheme which ensures that even if some portion of the message should become corrupted, one can still recover the original message by solving a linear system of equations. When designing ECCs, one typically assumes a noise model over a certain sample space, such as the Hamming~\cite{titsias2017hamming, dong2023number} or Levenshtein~\cite{levenshtein1966binary, becerra2008learning, barlev2021levenshtein} balls, from which we draw inspiration for this work.

  \subsection{Decoding}

Decoding is a key problem in machine translation, speech recognition, and other sequence-to-sequence tasks. Given a compressed encoding of some finite distribution, its goal is find the maximum likelihood samples. A classic example is Viterbi decoding, which is used to find the most likely sequence of hidden states in a hidden Markov model and is closely related to the CYK algorithm for parsing. For PCFGs, the problem is more challenging, as the total support can be exponentially larger relative to the number of transitions.

In particular, we care about the problem of \textit{top-k decoding}, which attempts to find the exact or approximate $k$-most likely samples in order of decreasing likelihood. This is closely related to the $k$-best enumeration~\cite{eppstein2014k} problem, a carefully studied problem in graph theory and combinatorial optimization. An exact solution to this problem for large acyclic PCFGs is often intractable, but we can approximate it using a beam search or cube-pruning technique.

A popular solution to k-best decoding in the NLP literature is a technique called cube-pruning~\cite{huang2005better, huang2007forest, chiang2007hierarchical}, which samples maximum likelihood paths through a hypergraph. We take inspiration from this technique, and adapt it with a probabilistic higher-order grammar (PHOG)~\cite{bielik2016phog}. Our approach is similar to work by Zhang and McDonald~\cite{zhang2012generalized}, but specialized to language intersections.

  An alternate approach would be to use MCMC or sequential Monte Carlo method to steer a transformer-based large langauge model (LLM), as proposed by Lew et al.~\cite{lew2023sequential}. This technique is particularly useful for constrained sampling from LLMs, and could be adapted to our setting to improve sample efficiency. The downside is that it introduces a dependency on an LLM, which requires a very large dataset to train and is more computationally expensive than cube-pruning. Furthermore, distinct sampling is unclear how to do properly, as LLMs are not trained to generate unique samples, and sampling without replacement is a fundamentally non-Markovian process. One potential solution proposed by Shi and Bieber~\cite{shi2020incremental} assumes trace injectivity and constructs a trie, however their solution is not stateless and can require a significant latency overhead.

  Our approach is complementary to existing work in constrained decoding. The bijection proposed in Eq.~\ref{eq:pairing} guarantees that all repairs are well-formed and converge linearly to the exact top-k maximum likelihood samples. This method is completely stateless and can be used to enumerate a bounded Levenshtein ball with linear parallelization speedup. Alternately, in the case of approximate ranked repair over a very large sample space, this technique can be adapted to sample with high probability a representative subset of the most likely sentences in a finite but large PCFG.

  \subsection{Neural program repair}

  The recent success of deep learning has lead to a variety of work on neural program repair~\cite{allamanis2021self, chirkova2021empirical, drain2021generating}. These approaches typically employ Transformer-based large language models (LLMs) and model the problem as a sequence-to-sequence transformation. In particular, two papers stand out being closely related to our own: Break-It-Fix-It (BIFI)~\cite{yasunaga2021break} and Seq2Parse~\cite{sakkas2022seq2parse}. BIFI adapts techniques from semi-supervised learning to generate synthetic errors in clean code and fixes them. This reduces the amount of pairwise training data, but tends to generalize poorly to length and out-of-distribution repairs. Seq2Parse combines a transformer-based model with an augmented version of the Early parser to suggest error rules, but only suggests a single repair. Our work differs from both in that we suggest multiple repairs at much higher precision, do not require a pairwise repair dataset, and can fix syntax errors in any language with a well-defined grammar. We note our approach is complementary to existing work in neural program repair, and may be used to generate synthetic repairs for training or employ an LLM for ranking.

  \clearpage\section{Conclusion}\label{sec:conclusion}

  Our work, while a case study on syntax repair, is part of a broader line of inquiry in program synthesis that investigates how to weave formal language theory and machine learning into helpful programming tools for everyday developers. In some ways, syntax repair serves as a test bench for integrating learning and language theory, as it lacks the intricacies of type-checking and semantic analysis, but is still rich enough to be an interesting challenge. By starting with syntax repair, we hope to lay the foundation for more organic hybrid approaches to program synthesis.

  Two high level codesign patterns have emerged to combine the naturalness of neural language models with the precision of formal methods. One seeks to filter the outputs of a generative language model to satisfy a formal specification, typically by some form of rejection sampling. Alternatively, some attempt to use language models to steer an incremental search for valid programs via a reinforcement learning or hybrid neurosymbolic strategy. However, implementing these strategies is often painstaking and their generalization behavior can be difficult to analyze.

  In our work, we take a more pragmatic tack - by incorporating the distance metric into a formal language, we attempt to exhaustively enumerate repairs by increasing distance, then use the stochastic language model to sort the resulting solutions by naturalness. The more constraints we can incorporate into formal language, the more efficient sampling becomes, and the more precise control we have over the output. This reduces the need for training a large, expensive language model to relearn syntax, and allows us to leverage compute for more efficient search and ranking.

  %  The great compromise in program synthesis is that of efficiency versus expressiveness. The more expressive a language, the more concise and varied the programs it can represent, but the harder those programs are to synthesize without resorting to domain-specific heuristics. Likewise, the simpler a language is to synthesize, the weaker its concision and expressive power. A large body of work focuses on general $\lambda$-calculi, or narrow languages such as finite sets or regular expressions. The former are too expressive to be efficiently synthesized or verified, whilst the latter are too restrictive to be useful for syntax repair. %Our work, we focus on context-free languages, which are expressive enough to capture a variety of practical synthesis tasks while retaining the benefits of compositionality and reusability across domains.

  There is a delicate balance in formal methods between soundness and completeness. Often these two seem at odds because the target language is too expressive to achieve them both simultaneously. In syntax repair, we also care about \textit{naturalness}. Fortunately, syntax repair is tractable enough to achieve all three by modeling the problem using language intersection. Completeness helps us to avoid missing simple repairs that might be easily overlooked, soundness guarantees all repairs will be valid, and naturalness ensures the most likely repairs receive the highest priority.

%  The second great compromise in program synthesis is that of reusability versus specialization. In programming, as in human communications, there is a vast constellation of languages, each requiring specialized generators and interpreters. Are these languages truly irreconcilable? Or, as Noam Chomsky argues, are these merely dialects of a universal language? \textit{Synthesis} then, might be a misnomer, and more aptly called \textit{recognition}, in the analytic tradition.

%  In our work, we argue these two compromises are not mutually exclusive, but complementary and reciprocal. Programs and the languages they inhabit are indeed synthetic, but can be analyzed and reused in the metalanguage of algebraic language theory. Not only does this admit an efficient synthesis algorithm, but allows users to introduce additional constraints without breaking compositionality, one of the most sacred tenets in programming language design.

%  Over the last few years, there has been a surge of progress in applying language models to write programs. That work is primarily based on methods from differential calculus and continuous optimization, leading to the so-called \textit{naturalness hypothesis}, which suggests programming languages are not so different from natural ones. In contrast, programming language theory takes the view that languages are essentially discrete and finitely-generated sets governed by logical calculi. Programming, thus viewed, is more like a mathematical exercise in constraint satisfaction. These constraints naturally arise at various stages of syntax validation, type-checking and runtime verification, and help to ensure programs fulfill their intended purpose.

%  As our work shows, not only is linear algebra over finite fields an expressive language for probabilistic inference, but also an efficient framework for inference on languages themselves. Borrowing analysis techniques from multilinear algebra and tensor completion in the machine learning setting, we develop an equational theory that allows us to translate various decision problems on formal languages into a system of inequalities over finite fields. We demonstrate the effectiveness of our approach for syntax repair in context-free languages, and show that our approach is competitive with state-of-the-art methods in terms of both accuracy and efficiency. In future work, we hope to extend our method to more natural grammars like conjunctive languages, TAG, LCFRS and other mildly context-sensitive languages.

  From a usability standpoint, syntax repair tools should be as user-friendly and widely accessible as autocorrection tools in word processors. We argue it is possible to reduce disruption from manual syntax repair and improve the efficiency of working programmers by driving down the latency needed to synthesize an acceptable repair. In contrast with program synthesizers that require intermediate editor states to be well-formed, our synthesizer does not impose any constraints on the code itself being written and is possible to use in an interactive programming setting.

%  The design of the tool itself is relatively simple. Tidyparse accepts a context-free language and a string. If the string is valid, it returns the parse forest, otherwise, it returns a set of repairs, ordered by likelihood. This approach has many advantages, enabling us to repair broken syntax, correct typos and recover from small errors, while being provably sound and complete with respect to the grammatical specification and a Levenshtein bound. It is also compatible with neural program synthesis and repair techniques, which can be used to score and rank the generated repairs.

  We have implemented our approach and demonstrated its viability as a tool for syntax assistance in real-world programming languages. Tidyparse is capable of generating repairs for invalid source code in a range of practical languages with little to no data required. We plan to continue expanding the prototype's autocorrection functionality to cover a broader range of languages and hope to conduct a more thorough user study to validate its effectiveness in practical programming scenarios.

  \section*{Data-Availability Statement}

  An artifact for Tidyparse is currently available as a browser application.~\footnote{\url{https://tidyparse.github.io}} While the browser demo is single-threaded and does not support ranking synthetic repairs by naturalness, it is capable of automatically repairing syntax errors in arbitrary context-free languages. The data and source code for the experiments contained in this paper will be made available upon publication.

%\subsection{Ranking}
%
%Since the number of solutions can be very large, we can use a language model to rank the results maximizing likelihood, or minimizing perplexity, subject to the constraints. This ranking can be used to guide the propagation, sample the choice function, sample hole locations or as a post-processing step after a fixed timeout has expired.

%Alternatively, this expression can be rewritten as a polynomial over GF(2):
%
%\[
%  (v_1 \times w_2 + y_3 + 1) \Leftrightarrow [S \in Y] \Leftrightarrow [Q R \in L(G)]
%\]

%
% ---- Bibliography ----
%
% BibTeX users should specify bibliography style 'splncs04'.
% References will then be sorted and formatted in the correct style.
%
%  \bibliographystyle{splncs04}
  \clearpage\bibliography{../bib/acmart}

  \pagebreak\appendix

  \section{Example Repairs}\label{sec:exaple_repairs}

  Below, we provide a few representative examples of broken code snippets and the corresponding human repairs that were successfully ranked first by our method. On the left is a complete snippet fed to the model and on the right, the corresponding human repair that was correctly predicted.

  \begin{figure}[H]
      \begin{tabular}{|m{6.6cm}|m{6.6cm}|}
        \hline \rule{0pt}{2.5ex}\textbf{Original broken code}\rule[-1ex]{0pt}{2ex} &  \rule{0pt}{2.5ex}\textbf{First predicted repair}\rule[-1ex]{0pt}{2ex} \\\hline
        \begin{lstlisting}[basicstyle=\ttfamily\lst@ifdisplaystyle\footnotesize\fi, language=python]

  (*@\hlorange{form}@*) sympy import *
  x = Symbol('x', real=True)
  x, re(x), im(x)

        \end{lstlisting} & \begin{lstlisting}[basicstyle=\ttfamily\lst@ifdisplaystyle\footnotesize\fi, language=python]

  (*@\hlorange{\textbf{from}}@*) sympy import *
  x = Symbol('x', real=True)
  x, re(x), im(x)

        \end{lstlisting} \\\hline
        \begin{lstlisting}[basicstyle=\ttfamily\lst@ifdisplaystyle\footnotesize\fi, language=python]

  result = (*@\hlorange{yeald}@*) From(item.create())
  raise Return(result)

        \end{lstlisting} & \begin{lstlisting}[basicstyle=\ttfamily\lst@ifdisplaystyle\footnotesize\fi, language=python]

  result = (*@\hlorange{\textbf{yield}}@*) From(item.create())
  raise Return(result)

        \end{lstlisting} \\\hline
        \begin{lstlisting}[basicstyle=\ttfamily\lst@ifdisplaystyle\footnotesize\fi, language=python]

  df.apply(lambda row: list(set(row['ids'(*@\hlorange{)}@*))))

        \end{lstlisting} & \begin{lstlisting}[basicstyle=\ttfamily\lst@ifdisplaystyle\footnotesize\fi, language=python]

  df.apply(lambda row: list(set(row['ids'(*@\hlorange{]}@*))))

        \end{lstlisting} \\\hline
%        \begin{lstlisting}[basicstyle=\ttfamily\lst@ifdisplaystyle\footnotesize\fi, language=python]
%
%  import numpy (*@\hlorange{ad}@*) np
%  A_concate = np.array([a_0, a_1, a_2,..., a_n])
%
%        \end{lstlisting} & \begin{lstlisting}[basicstyle=\ttfamily\lst@ifdisplaystyle\footnotesize\fi, language=python]
%
%  import numpy (*@\hlorange{\textbf{as}}@*) np
%  A_concate = np.array([a_0, a_1, a_2,..., a_n])
%
%        \end{lstlisting} \\\hline
        \begin{lstlisting}[basicstyle=\ttfamily\lst@ifdisplaystyle\footnotesize\fi, language=python]

  sum(len(v) for v items.values())(*@\hlred{)}@*)

        \end{lstlisting} & \begin{lstlisting}[basicstyle=\ttfamily\lst@ifdisplaystyle\footnotesize\fi, language=python]

  sum(len(v) for v (*@\hlgreen{\textbf{in}}@*) items.values())

        \end{lstlisting} \\\hline
        \begin{lstlisting}[basicstyle=\ttfamily\lst@ifdisplaystyle\footnotesize\fi, language=python]

  def average(values):
    if values == (1,2,3):
      return (1+2+3)/3
    (*@\hlorange{else}@*) (*@\hlred{if}@*) values == (-3,2,8,-1):
      return (-3+2+8-1)/4

        \end{lstlisting} & \begin{lstlisting}[basicstyle=\ttfamily\lst@ifdisplaystyle\footnotesize\fi, language=python]

  def average(values):
    if values == (1,2,3):
      return (1+2+3)/3
    (*@\hlorange{elif}@*) values == (-3,2,8,-1):
      return (-3+2+8-1)/4

        \end{lstlisting} \\\hline
        \begin{lstlisting}[basicstyle=\ttfamily\lst@ifdisplaystyle\footnotesize\fi, language=python]

  dict = {
     "Jan": 1
     "January": 1
     "Feb": 2 # and so on
  }

        \end{lstlisting} & \begin{lstlisting}[basicstyle=\ttfamily\lst@ifdisplaystyle\footnotesize\fi, language=python]

  dict = {
     "Jan": 1(*@\hlgreen{,}@*)
     "January": 1(*@\hlgreen{,}@*)
     "Feb": 2 # and so on
  }

        \end{lstlisting} \\\hline
        \begin{lstlisting}[basicstyle=\ttfamily\lst@ifdisplaystyle\footnotesize\fi, language=python]

  class MixIn(object)
    def m():
      pass

  class classA(MixIn):

  class classB(MixIn):

        \end{lstlisting} & \begin{lstlisting}[basicstyle=\ttfamily\lst@ifdisplaystyle\footnotesize\fi, language=python]

  class MixIn(object)(*@\hlgreen{:}@*)
    def m():
      pass

  class classA(MixIn): (*@\hlgreen{\textbf{pass}}@*)

  class classB(MixIn): (*@\hlgreen{\textbf{pass}}@*)

        \end{lstlisting} \\\hline
      \end{tabular}
  \end{figure}

  \clearpage\section{Raw data}\label{sec:raw_prec_data}

  Raw data from Precision@k experiments across snippet length and Levenshtein distance from \S~\ref{sec:rq2}. $|\err\sigma|$ indicates the snippet length and $\Delta$ indicates the Levenshtein distance between the broken and code and human fix computed over lexical tokens. For Tidyparse, we sample until exhausting the admissible set or a timeout of 30s is reached, whichever happens first, then rank the results. For the other models Precision@1, we sample one repair and report the percentage of repairs matching the human repair. For Precision@All, we report the percentage of repairs matching the human repair within the top 20000 samples.

  \begin{table}[!h]
    \centering
    \begin{tabular}{c|c|cccccccc}
      \hline\hline
      & $\Delta$ & \multicolumn{8}{c}{Precision@1} \\ \hline
      $|\err\sigma|$ &  & $(0,10)$ & $[10,20)$ & $[20,30)$ & $[30, 40)$ & $[40,50)$ & $[50, 60)$ & $[60,70)$ & $[70, 80)$ \\ \hline
      Tidyparse
      & 1 & 1.00 & 1.00 & 0.98 & 0.98 & 1.00 & 1.00 & 0.95 & 0.90 \\
      & 2 & 0.51 & 0.36 & 0.24 & 0.26 & 0.24 & 0.23 & 0.12 & 0.10 \\
      & 3 & 0.38 & 0.26 & 0.37 & 0.25 & 0.22 & 0.16 & 0.14 & 0.08 \\ \hline
      Seq2Parse
      & 1 & 0.35 & 0.41 & 0.40 & 0.37 & 0.31 & 0.29 & 0.27 & 0.21 \\
      & 2 & 0.12 & 0.13 & 0.14 & 0.12 & 0.11 & 0.11 & 0.10 & 0.12 \\
      & 3 & 0.03 & 0.07 & 0.08 & 0.09 & 0.09 & 0.02 & 0.07 & 0.06 \\ \hline
      BIFI
      & 1 & 0.20 & 0.33 & 0.32 & 0.27 & 0.21 & 0.21 & 0.25 & 0.18 \\
      & 2 & 0.18 & 0.18 & 0.21 & 0.19 & 0.19 & 0.18 & 0.11 & 0.11 \\
      & 3 & 0.02 & 0.02 & 0.03 & 0.02 & 0.03 & 0.05 & 0.03 & 0.02 \\ \hline
      & & \multicolumn{8}{c}{Precision@All} \\ \hline
      Tidyparse
      & 1 & 1.00 & 1.00 & 0.98 & 0.98 & 1.00 & 1.00 & 1.00 & 0.91 \\
      & 2 & 0.91 & 0.89 & 0.85 & 0.82 & 0.68 & 0.82 & 0.58 & 0.50 \\
      & 3 & 0.50 & 0.37 & 0.53 & 0.40 & 0.44 & 0.27 & 0.34 & 0.22 \\ \hline
      BIFI
      & 1 & 0.65 & 0.67 & 0.70 & 0.65 & 0.60 & 0.62 & 0.60 & 0.64 \\
      & 2 & 0.52 & 0.41 & 0.37 & 0.32 & 0.27 & 0.27 & 0.21 & 0.24 \\
      & 3 & 0.20 & 0.13 & 0.08 & 0.17 & 0.15 & 0.18 & 0.17 & 0.07 \\ \hline\hline
    \end{tabular}
  \end{table}

  Synthetic evaluation

  \begin{table}[!h]
    \centering
    \begin{tabular}{c|c|cccccccc}
      \hline\hline
      & $\Delta$ & \multicolumn{8}{c}{Precision@1} \\ \hline
      $|\err\sigma|$ &  & $(0,10)$ & $[10,20)$ & $[20,30)$ & $[30, 40)$ & $[40,50)$ & $[50, 60)$ & $[60,70)$ & $[70, 80)$ \\ \hline
      Tidyparse
      & 1 & 1.00 & 1.00 & 1.00 & 0.99 & 1.00 & 1.00 & 1.00 & 0.98 \\
      & 2 & 0.45 & 0.63 & 0.66 & 0.68 & 0.65 & 0.81 & 0.64 & 0.62 \\
      & 3 & 0.06 & 0.20 & 0.29 & 0.36 & 0.29 & 0.39 & 0.38 & 0.32 \\ \hline
      & & \multicolumn{8}{c}{Precision@All} \\ \hline
      Tidyparse
      & 1 & 1.00 & 1.00 & 1.00 & 0.99 & 1.00 & 1.00 & 0.98 & 1.00 \\
      & 2 & 0.98 & 0.98 & 0.94 & 0.94 & 0.98 & 0.97 & 0.89 & 0.90 \\
      & 3 & 1.00 & 0.97 & 0.92 & 0.84 & 0.87 & 0.90 & 0.84 & 0.72 \\ \hline\hline
    \end{tabular}
  \end{table}

\clearpage\section{Supplemental Proofs}

  The problem of syntax error correction under a finite number of typographic errors is reducible to the bounded Levenshtein-CFL reachability problem, which can be formally stated as follows:

  \begin{definition}
    The language edit distance (LED) is the minimum number of edits required to transform an invalid string into a valid one, where validity is defined as containment in a context-free language, $\ell: \mathcal{L}$, i.e., $\Delta^*(\err{\sigma}, \ell) \coloneqq \min_{\sigma \in \ell}\Delta(\err{\sigma}, \sigma)$, and $\Delta$ is the Levenshtein distance.
  \end{definition}

  We seek to find the set of strings $S$ such that $\forall \sigma'\in S, \Delta(\err{\sigma}, \sigma') \leq q$, where $q$ is greater than or equal to the language edit distance. We call this set the \textit{Levenshtein ball} of $\err{\sigma}$ and denote it $\Delta_q(\err{\sigma})$. Since $1 \leq \Delta^*(\err{\sigma}, \ell) \leq q$, we have $1 \leq q$. We now consider an upper bound on $\Delta^*(\err{\sigma}, \ell)$, i.e., the greatest lower bound on $q$ such that $\Delta_q(\err{\sigma}) \cap \ell \neq \varnothing$.

  \begin{lemma}\label{lemma:upper-bound}
  For any nonempty language $\ell: \mathcal{L}$ and invalid string $\err{\sigma}: \Sigma^n \cap \bar\ell$, there exists an $(\sigma', m)$ such that $\sigma' \in \ell\cap\Sigma^m$ and $0 < \Delta(\err{\sigma}, \ell) \leq \max(m, n) < \infty$.
  \end{lemma}

  \begin{proof}
    Since $\ell$ is nonempty, it must have at least one inhabitant $\sigma \in \ell$. Let $\sigma'$ be the smallest such member. Since $\sigma'$ is a valid sentence in $\ell$, by definition it must be that $|\sigma'|<\infty$. Let $m\coloneqq|\sigma'|$. Since we know $\err{\sigma} \notin \ell$, it follows that $0 < \Delta(\err{\sigma}, \ell)$. Let us consider two cases, either $\sigma' = \varepsilon$, or $0 < |\sigma'|$:

    \begin{itemize}
      \item If $\sigma' = \varepsilon$, then $\Delta(\err{\sigma}, \sigma') = n$ by full erasure of $\err{\sigma}$, or
      \item If $0 < m$, then $\Delta(\err{\sigma}, \sigma') \leq \max(m, n)$ by overwriting.
    \end{itemize}

    In either case, it follows $\Delta(\err{\sigma}, \ell) \leq \max(m, n)$ and $\ell$ is always reachable via a finite nonempty set of Levenshtein edits, i.e., $0 < \Delta(\err{\sigma}, \ell) < \infty$.
  \end{proof}

  Let us now consider the maximum growth rate of the \textit{admissible set}, $A \coloneqq \Delta_q(\err{\sigma}) \cap \ell$, as a function of $q$ and $n$. Let $\bar\ell \coloneqq \{\err{\sigma}\}$. Since $\bar\ell$ is finite and thus regular, $\ell = \Sigma^* \setminus \{\err{\sigma}\}$ is regular by the closure of regular languages under complementation, and thus context-free a fortiori. Since $\ell$ accepts every string except $\err{\sigma}$, it represents the worst CFL in terms of asymptotic growth of $A$.

  \begin{lemma}\label{lemma:interleaving}
  The complexity of enumerating $A$ is upper bounded by $\mathcal{O}\left(\sum_{c=1}^q{{cn + n + c} \choose c}(|\Sigma| + 1)^c\right)$.
  \end{lemma}

  \begin{proof}
    We can overestimate the size of $A$ by considering the number of unique ways to insert, delete, or substitute $c$ terminals into a string $\err{\sigma}$ of length $n$. This can be overaproximated by interleaving $\varepsilon^c$ around every token, i.e., $\err{\sigma}_\varepsilon\coloneqq \left(\varepsilon^c\err{\sigma}_i\right)_{i=1}^n\varepsilon^c$, where $|\err{\sigma}_\varepsilon| = cn + n + c$, and only considering substitution. We augment $\Sigma_\varepsilon \coloneqq \Sigma \cup \{\varepsilon\}$ so that deletions and insertions may be treated as special cases of substitution. Thus, we have $cn + n + c$ positions to substitute $(|\Sigma_\varepsilon|)$ tokens, i.e., ${{cn + n + c} \choose c}|\Sigma_\varepsilon|^c$ ways to edit $\err{\sigma}_\varepsilon$ for each $c \in [1, q]$. This upper bound is not tight, as overcounts many identical edits w.r.t. $\err{\sigma}$. Nonetheless, it is sufficient to show $|A| < \sum_{c=1}^q{{cn + n + c} \choose c}|\Sigma_\varepsilon|^c$.
  \end{proof}

  We note that the above bound applies to all strings and languages, and relates to the Hamming bound on $H_q(\err{\sigma}_\varepsilon)$, which only considers substitutions. In practice, much tighter bounds may be obtained by considering the structure of $\ell$ and $\err{\sigma}$. For example, based on an empirical evaluation from a dataset of human errors and repairs in Python code snippets ($|\Sigma| = 50, |\err{\sigma}| < 40, \Delta(\err{\sigma}, \ell) \in [1, 3]$), we estimate the \textit{filtration rate}, i.e., the density of the admissible set relative to the Levenshtein ball, $D=|A|/|\Delta_q(\err{\sigma})|$ to have empirical mean $E_\sigma[D] \approx 2.6\times 10^{-4}$, and variance $\mathrm{Var}_\sigma[D] \approx 3.8\times10^{-7}$.

%  In practice, this problem is ill-posed even when $q = \Delta^*(\err{\sigma}, \ell) \approx 1$. For example, consider the language of ursine dietary preferences. Although $\err{\sigma}\coloneqq$ ``Bears like to eat plastic'' is not a valid sentence, e.g., $\sigma'\coloneqq$``Bears like to eat'' is $(\Delta^*=1)$, however there are many others with roughly the same edit distance, e.g., ``Bears like to eat \{\hlorange{berries}, \hlorange{honey}, \hlorange{fish}\}'', or ``\{\hlgreen{Polar}, \hlgreen{Panda}\} bears like to eat \{\hlgreen{seals}, \hlgreen{bamboo}\}''. In general, there are usually many strings nearby $\err{\sigma}$, and we seek to find those among them which are both syntactically valid and semantically plausible as quickly as possible.

\end{document}